\chapter{Poincar\'e--Lelong and local psh functions}
\index{Poincare--Lelong formula@Poincar\'e--Lelong formula!CLD}\index{plurisubharmonic function!CLD}
\label{chap:cld-poincare-lelong}

The Poincar\'e--Lelong formula is the first point at which the analytic
divisor of a function and the real tropical differential calculus become the
same object.  It also fixes a sign that propagates through the entire book:
with $v(f)=-\log|f|$,
\[
       \dprimeddprime[\log|f|]=\delta_{\operatorname{div}(f)},
       \qquad
       \dprimeddprime[v(f)]=-\delta_{\operatorname{div}(f)}.
\]
We prove the formal reductions, explain the analytic heart of the theorem,
and then compare two notions of non-Archimedean plurisubharmonicity that have
different definitions and presently different scopes.

Throughout, $X$ is a topologically separated, pure $n$-dimensional
$K$-analytic space without boundary.

It is worth saying at once what kind of statement the displayed formula is,
because its shape is so familiar from complex geometry that one may fail to
notice how much has to be arranged before it can even be written down.  On a
complex manifold, $\log|f|$ is a locally integrable function because $|f|$
vanishes to finite order along a hypersurface and $\log$ has an integrable
singularity there; the current $\ddc\log|f|$ is then computed by a residue
argument, and the coefficient of the resulting divisor is an order of
vanishing.  In the non-Archimedean setting none of these three steps is
available in the same form.  There is no residue calculus; there is no
integration against a background volume form on $X$; and $\log|f|$ is not a
function with an integrable logarithmic pole in any classical sense, since
the locus $\{f=0\}$ is a Zariski-closed subset which is \emph{not} charged by
the measures at issue but which also carries no obvious transversal along
which to integrate.

What replaces all of this is the calibrated tropical geometry of
\Cref{chap:tropical-spaces,chap:tropical-integration,chap:analytic-tropicalization}.
Adjoining $f$ to a moment map turns a neighbourhood of $\{f=0\}$ into a
tropical product $\Trop(V\cap\{f=0\})\times[\log s,\log t]$ over a short
logarithmic interval; the extra coordinate is $\log|f|$ itself; and the whole
theorem becomes a computation of one boundary face in Stokes' formula.  The
multiplicity that appears is not read off from a Taylor expansion but from
the canonical calibration of \Cref{def:trop-canonical-calibration}, and it is
this that makes the right-hand side the current of the \emph{Cartier}
divisor, with its lengths and residue degrees, rather than the current of the
reduced zero set.  The reader who wants a one-line summary of the mechanism
may take this: \emph{Poincar\'e--Lelong is the Stokes formula of
\Cref{thm:ch11-balancing-iff-stokes} applied to the cylinder over the zero
locus, and the divisor multiplicity is the calibration of the base of that
cylinder}.

\section{Orientation and prerequisites}
\label{sec:ch15-orientation}

The reader should have available: the superform calculus and the Hessian sign
convention \eqref{eq:appJ-real-hessian} of \Cref{chap:superforms}; the
calibrated integration theory and Stokes' formula of
\Cref{chap:tropical-integration}; canonical calibrations and tropical
multiplicities from \Cref{chap:analytic-tropicalization}; the sheaves
$A_X^{p,q}$ and the analytic Stokes and Green formulae of
\Cref{chap:cld-forms}; and the definition, positivity cones and integration
currents of \Cref{chap:cld-currents}.  From \Cref{part:foundations} we use only the
graph Laplacian \eqref{eq:graph-laplacian} and the slope formula
\Cref{thm:ch08-slope-formula}, and these only in the worked examples, where
they serve as an independent check on every sign.

By the end of the chapter the reader should be able to:
\begin{enumerate}
\item say precisely which meromorphic functions have a Poincar\'e--Lelong
  formula, and exhibit the failure when the hypothesis is dropped;
\item derive the constant and the sign in
  $\dprimeddprime[\log|f|]=\delta_{\divis(f)}$ from the Hessian convention,
  rather than remembering them;
\item verify the formula by hand on the Berkovich affine line for a
  polynomial with several roots, and reconcile the answer with the graph
  Laplacian of \Cref{chap:model-calculations};
\item state the Bochner--Martinelli and Levine formulae with their exact
  hypotheses, and explain why the left-hand side of the first cannot be
  abbreviated as a $d$-fold Monge--Amp\`ere product;
\item work with the CLD sheaf of psh functions, including the transfinite
  presentation and the class of \emph{locally approachable} functions, and
  prove the elementary stability properties;
\item recognize which maximum principles are theorems, which are theorems for
  a different psh class, and which are open.
\end{enumerate}

\subsection*{Roadmap}

\Cref{sec:ch15-meromorphic,sec:ch15-current-log} set up the two objects in
the formula: the Cartier divisor $\divis(f)$ of a regular meromorphic
function, and the current $[\log|f|]$.  Neither is automatic.  Regularity is
what makes $\divis(f)$ a Cartier divisor of pure codimension one, and it is
also what makes the local product description at the heart of the proof
correct; the construction of $[\log|f|]$ is an extension theorem across a
Zariski-closed set of small local tropical dimension, and it is
\emph{constructed}, not characterized by a uniqueness principle.

\Cref{sec:ch15-pl} proves the formula.  The proof is organized in four
movements: a formal reduction to a holomorphic function defining a morphism
to $\mathbf A^1$ of pure relative dimension $n-1$; a cutoff that splits the
test form into a part supported near $\{f=0\}$ and a part supported away from
it; Green's formula, which moves $\dprimeddprime$ onto the test form and
leaves a boundary integral over the level set $\{|f|=s\}$; and a tropical
computation of that boundary integral.  We also record the sign derivation
and the normalization checks against \Cref{app:normalizations}.

\Cref{sec:ch15-tests} is a battery of examples.  The centrepiece is a
polynomial of degree four on the Berkovich affine line, carried to explicit
numbers on a compact skeleton and cross-checked against
\eqref{eq:graph-laplacian} and \eqref{eq:ch08-slope-formula}.

\Cref{sec:ch15-bm,sec:ch15-levine} treat the higher-codimension formulae.
The Bochner--Martinelli formula of Mihatsch replaces the divisor by a regular
complete intersection and the potential $\log|f|$ by $\log\max_i|f_i|$; the
Levine formula computes the difference between a power of the curvature of
$\cO(1)$ and the class of a linear subspace of $\mathbf P^n$.  Both are
proved, and both are stated with the extension conventions that make their
left-hand sides meaningful.

\Cref{sec:ch15-psh,sec:ch15-hierarchy,sec:ch15-convexity} develop the CLD
class of psh functions: the generated sheaf, its transfinite presentation,
the first-stage class of locally approachable functions, the stability
calculus, and the model case of convex functions on $\mathbf R^n$, where the
whole theory collapses to classical convex analysis.

\Cref{sec:ch15-polar,sec:ch15-maximum,sec:ch15-pullback} treat polar loci,
maximum principles, and functoriality.  \Cref{sec:ch15-classical} compares
the CLD class with the Gubler--Rabinoff curve-test class, and
\Cref{sec:ch15-dictionary} collects the dictionary.

\subsection*{What is deliberately absent}

Four things a reader might expect are not here.

There are no Bedford--Taylor products of unbounded potentials.  Every product
of currents occurring in this chapter is either a product of a current with a
smooth form, or a Bedford--Taylor product of \emph{continuous} psh functions
constructed in \Cref{chap:cld-bedford-taylor}, or a canonical extension whose
construction is part of the statement being proved.  We never write
$(\dprimeddprime u)^d$ when $u$ is $-\infty$ somewhere.

There is no global theory.  Poincar\'e--Lelong here is a local statement
about an analytic space; the passage to metrized line bundles, curvature
currents, and Chambert--Loir measures is \Cref{chap:cld-metrized-bundles}.
We state the section form of the formula once, as a corollary, because it is
the form that \Cref{app:normalizations} uses to compare with complex
geometry, but we do not develop the metric formalism here.

There is no general regularization theorem proved in this chapter.  The class of psh functions is
\emph{defined} by an approximation property precisely because no analogue of
convolution exists on a Berkovich space; the only place where genuine
regularization occurs in this chapter is \Cref{sec:ch15-convexity}, on
$\mathbf R^n$, where translations are available.  The continuous
smooth-proper-curve theorem is proved later in
\Cref{thm:p063-current-psh-ordinary-regularization}.

There is no delta-form calculus.  The Poincar\'e--Lelong formula of
Gubler--K\"unnemann, in which $\delta_{\divis(f)}$ becomes a delta-form and
can therefore be multiplied, is \Cref{thm:gk-poincare-lelong}.

\section{Regular meromorphic functions and their divisors}
\label{sec:ch15-meromorphic}

On a reducible or nonreduced analytic space, a meromorphic function can be a
zero divisor.  Divisors and Poincar\'e--Lelong then cease to behave as in the
integral case.

\subsection{The sheaf of meromorphic functions}
\label{sec:ch15-meromorphic-sheaf}
\index{meromorphic function!on an analytic space}

Let $X$ be a $K$-analytic space.  For an affinoid domain $V=\cM(\mathcal A)$ of $X$
let $S_V\subset\mathcal A$ be the multiplicative set of non-zero-divisors and put
$\cM_X(V)=S_V^{-1}\mathcal A$.  Sheafifying the presheaf $V\mapsto\cM_X(V)$ on the
$G$-topology gives the sheaf $\cM_X$ of \emph{meromorphic functions}, exactly
as in algebraic geometry \cite[20.1.8]{EGAIV4}; the corresponding statement
in the analytic category is \cite[\S18.1.5]{CLD25}.  A section of $\cM_X$ is
thus locally of the form $g/h$ with $g,h\in\cO_X(V)$ and $h$ a
non-zero-divisor.  We write $\cM_X^\times$ for the sheaf of invertible
sections of $\cM_X$.

Three degeneracies must be kept apart, and it costs nothing to name them.

First, $X$ may be reducible.  If $X=X_1\cup X_2$ with $X_1\cap X_2$ of
dimension $<n$, a function may vanish identically on $X_1$ and be invertible
on $X_2\setminus X_1$.  Such a function is a zero divisor in $\cO_X$: it is
not a section of $\cM_X^\times$, and it has no divisor.

Second, $X$ may be nonreduced.  A nilpotent function is a zero divisor.  More
importantly, even a non-zero-divisor on a nonreduced space has a divisor
whose coefficients are lengths rather than orders of vanishing of the reduced
structure; this is where \Cref{thm:current-fundamental-cycle} enters the
picture.

Third, and least visible, $X$ may fail Serre's condition $(S_1)$: it may have
embedded components.  This is not detected by looking at $X_{\mathrm{red}}$
and it is exactly the hypothesis that has to be created by hand in
\Cref{lem:pl-formal-reductions}.

\subsection{Regularity}

\begin{definition}
A meromorphic function $f$ on $X$ is \emph{regular} if, on every affinoid
domain on which it is represented in the total meromorphic ring, it is
invertible there; equivalently its numerator and denominator can locally be
chosen to be non-zero-divisors.  Its Cartier divisor is denoted
$\operatorname{div}(f)$.
\end{definition}

In sheaf-theoretic language, a regular meromorphic function is a section of
$\cM_X^\times$.  Its divisor is the associated Cartier divisor: if
$f=g/h$ on $V$ with $g$ and $h$ non-zero-divisors, then
\[
 \divis(f)|_V=\divis(g)|_V-\divis(h)|_V,
\]
and $\divis(g)$ is the Cartier divisor cut out by the single equation $g=0$.
The next lemma collects the properties of $\divis$ used below; all of them
are formal, but the reader who has met divisors only on integral schemes
should notice which of them survive.

\begin{lemma}
\label{lem:ch15-divisor-basic}
Let $f,f_1,f_2$ be regular meromorphic functions on $X$.
\begin{enumerate}
\item\label{itm:ch15-div-supp} The support of $\divis(f)$ is a Zariski-closed
  subset, purely of codimension one, on whose complement $f$ is invertible.
\item\label{itm:ch15-div-add} $\divis(f_1f_2)=\divis(f_1)+\divis(f_2)$ and
  $\divis(f^{-1})=-\divis(f)$.
\item\label{itm:ch15-div-unit} $\divis(f)=0$ if and only if $f$ is a section
  of $\cO_X^\times$.
\item\label{itm:ch15-div-cycle} The associated cycle $[\divis(f)]$ is a
  locally finite $\mathbf Z$-linear combination of irreducible closed
  analytic subsets of codimension one, and the multiplicity along a component
  $Y$ is the length of the local ring of $\divis(f)$ at the generic point of
  $Y$.
\item\label{itm:ch15-div-pullback} If $g:X'\to X$ is a morphism such that
  $g^*f$ is again a regular meromorphic function, then
  $\divis(g^*f)=g^*\divis(f)$ as Cartier divisors.
\end{enumerate}
\end{lemma}

\begin{proof}
\cref{itm:ch15-div-supp}
The assertion is local, so assume $f=g/h$ with $g,h$ non-zero-divisors on an
affinoid $V=\cM(\mathcal A)$.  The locus where $f$ is not invertible is
$V(g)\cup V(h)$, a Zariski-closed subset.  If $g$ is a non-zero-divisor of
$\mathcal A$, then $g$ lies in no associated prime of $\mathcal A$; hence for every minimal
prime $\mathfrak p$ of $\mathcal A$ the image of $g$ in $\mathcal A/\mathfrak p$ is
nonzero, and Krull's Hauptidealsatz gives that every minimal prime over
$(g)$ has height one.  Therefore $V(g)$ is purely of codimension one in $V$;
the same for $h$, and the union of two pure codimension-one sets is again
pure of codimension one.  The Hauptidealsatz in the analytic category is
\cite[Lemma~11.5.2]{CLD25}.

\cref{itm:ch15-div-add} and \cref{itm:ch15-div-unit} are immediate from
the definition of the Cartier divisor of a section of $\cM_X^\times$: the map
$f\mapsto\divis(f)$ is the connecting map of the exact sequence
$1\to\cO_X^\times\to\cM_X^\times\to\cM_X^\times/\cO_X^\times\to1$, hence a
homomorphism of sheaves of abelian groups with kernel $\cO_X^\times$.

\cref{itm:ch15-div-cycle} is the general passage from a Cartier divisor to
its associated Weil cycle, together with \Cref{thm:current-fundamental-cycle}
for the identification of the multiplicity with a length.

\cref{itm:ch15-div-pullback} follows because $\divis$ is defined by local
equations and pullback of local equations computes the pullback of a Cartier
divisor; the hypothesis that $g^*f$ is again regular is exactly what makes
the pullback of the local equations again a non-zero-divisor.  Without it,
the pullback Cartier divisor is undefined, and this is not a technicality:
see \Cref{rem:ch15-nonregular} below.
\end{proof}

\begin{remark}[Regularity is not implied by codimension]
\label{rem:ch15-nonregular}
It is tempting to believe that if the non-invertibility locus of a
meromorphic function has codimension $\ge1$, then the function is regular and
the Poincar\'e--Lelong formula holds.  This is false.  CLD exhibit a
meromorphic function whose non-invertibility locus has dimension $<n$ and for
which the formula fails \cite[Example~16.1.4]{CLD25}.  The mechanism is the
one described above: on a space with an embedded component, or on a
reducible space, a function can vanish on a component of the structure sheaf
without vanishing on any component of the reduction, so that its
``divisor'' is not a Cartier divisor at all.  The correct statement is not
about dimensions but about associated primes.
\end{remark}

\begin{warning}{Three separate hypotheses}
The proof of \Cref{thm:cld-poincare-lelong} uses three conditions that are
easy to conflate: that $f$ be a section of $\cM_X^\times$; that $X$ satisfy
$(S_1)$; and that the morphism $f:X\to\mathbf A^1$ be purely of relative
dimension $n-1$.  The first is a hypothesis of the theorem.  The second and
third are \emph{produced} locally from the first in
\Cref{lem:pl-formal-reductions}, and neither is a formal consequence of the
mere existence of a non-zero-divisor: the deduction uses the Hauptidealsatz
and a semicontinuity theorem for relative dimension.
\end{warning}

\section{The current of a regular meromorphic function}
\label{sec:ch15-current-log}

Let $U\subset X$ be the locus where $f$ is represented by an invertible
analytic function.  The function $\log|f|$ is smooth on $U$.  For a regular
meromorphic $f$, CLD prove that $\log|f|$ is locally integrable and construct
its canonical $L^1_{\mathrm{loc}}$ current on $X$; see
\cite[\S16.7.2, using Proposition~7.1.12]{CLD25}.  We denote this current by
$[\log|f|]$.  Its construction is theorem-specific, not a consequence of a
general uniqueness-by-codimension principle, and the notation does not say
that $\log|f|$ is everywhere finite.

We spell out both halves of this construction, because the two are used
separately later and because the phrase ``canonical extension'' is used in
this chapter in three different technical senses.

\subsection{Why \texorpdfstring{$\log|f|$}{log|f|} is smooth on the
invertibility locus}

\begin{lemma}
\label{lem:ch15-log-smooth}
Let $f$ be an invertible analytic function on an open set $U\subset X$.  Then
$\log|f|$ is a smooth function on $U$ in the sense of
\Cref{chap:cld-forms}, and
\[
 \dprime\log|f|=-\dprime t,\qquad
 \ddprime\log|f|=-\ddprime t,\qquad
 \dprimeddprime\log|f|=0 ,
\]
where $t$ is the affine coordinate on $\mathbf R$ pulled back by the moment
map $f$, that is, $t=-\log|f|$.
\end{lemma}

\begin{proof}
An invertible analytic function is a moment map $f:U\to\mathbf G_{m,K}$ of
rank one, and $-\log|f|=f_{\mathrm{trop}}$ by the convention of
\Cref{chap:cld-forms}.  Hence $\log|f|=\varphi\circ f_{\mathrm{trop}}$ with
$\varphi(t)=-t$, a smooth --- indeed affine --- function on $\mathbf R$, so
$\log|f|$ is smooth by \Cref{prop:cld-local-normal-form}.  The differentials
are computed on the chart representative: $\dprime\varphi=-\dprime t$,
$\ddprime\varphi=-\ddprime t$ and $\dprimeddprime\varphi=\varphi''\,\dprime
t\wedge\ddprime t=0$ since $\varphi''=0$.
\end{proof}

\Cref{lem:ch15-log-smooth} already contains the first half of
Poincar\'e--Lelong: \emph{on the invertibility locus both sides of the formula
vanish}.  It also contains a warning.  The function $\log|f|$ is smooth and
$\dprimeddprime$-closed on $U$, but it is not the pullback of a bounded
function, and its extension across $X\setminus U$ is where all the content
lies.

\subsection{Local integrability and the extension across the polar locus}

\begin{proposition}
\label{prop:ch15-log-current}
Let $f$ be a regular meromorphic function on $X$ and let $U$ be its
invertibility locus, a dense Zariski-open subset whose complement $Z$
satisfies $d_K(x)<\dim X$ for every $x\in Z$.  Then:
\begin{enumerate}
\item\label{itm:ch15-logcur-int} $\log|f|$ is locally integrable on $X$ for
  every measure of the form $|\omega|$ with $\omega$ a smooth $(n,n)$-form;
\item\label{itm:ch15-logcur-ext} the current $[\log|f|]$ on $U$ extends
  uniquely to a current of bidegree $(0,0)$ on $X$, again denoted
  $[\log|f|]$;
\item\label{itm:ch15-logcur-add} the extension is additive:
  $[\log|f_1f_2|]=[\log|f_1|]+[\log|f_2|]$ and
  $[\log|f^{-1}|]=-[\log|f|]$.
\end{enumerate}
\end{proposition}

\begin{proof}
That $U$ is Zariski-dense and that $d_K(x)<\dim X$ on its complement is
\cite[\S16.7.2]{CLD25}: $Z$ is a Zariski-closed subset of codimension one by
\Cref{lem:ch15-divisor-basic}\,\cref{itm:ch15-div-supp}, and a point of a
Zariski-closed subset of dimension $<n$ has $d_K(x)<n$ by
\Cref{prop:analytic-tropical-dimension}.

\cref{itm:ch15-logcur-int}
The key point is that a top-degree test form has no support at all near $Z$.
Indeed, a smooth form of bidegree $(p,q)$ vanishes as a germ at every point
$x$ whose local tropical dimension is smaller than $\max\{p,q\}$
\cite[Proposition~15.2.1]{CLD25}; for $(p,q)=(n,n)$ this applies at every
point of $Z$, since $d_K(x)<n$ there.  Vanishing as a germ means vanishing on
a neighbourhood, so
\[
 \Supp\omega\cap Z=\varnothing .
\]
Now $\Supp\omega$ is compact and $Z$ is closed, so $|f|$ attains a positive
minimum on $\Supp\omega$ and $\log|f|$ is bounded there.  A bounded function
is integrable against the Radon measure $|\omega|$ of
\Cref{thm:cld-analytic-integration}, whose support is contained in
$\Supp\omega$.  Hence $\int_X\log|f|\,\omega$ is a well-defined finite
number.

The same remark shows why no growth estimate on $\log|f|$ is needed here, in
contrast with the complex case: in the complex setting the analogous integral
really does have a logarithmic singularity along the divisor and converges
only because $\log$ is locally integrable, whereas here the test form
disappears before the singularity is reached.

\cref{itm:ch15-logcur-ext}
Existence follows from \cref{itm:ch15-logcur-int}: the assignment
$\omega\mapsto\int_X\log|f|\,\omega$ is a well-defined linear functional on
$A_{X,c}^{n,n}$, and it is continuous, since on fixed compact chart data the
integral is bounded by
$\bigl(\sup_{\Supp\omega}|\log|f||\bigr)\,\|\omega\|_{C^0}$ times a fixed
finite mass, by \Cref{prop:test-form-continuity-criterion}.  For uniqueness,
suppose $T$ and $T'$ are two currents on $X$ of bidegree $(0,0)$ restricting
to $[\log|f|]$ on $U$.  Their difference is supported on $Z$.  A current of
bidegree $(0,0)$ supported on a closed set all of whose points have local
tropical dimension $<n$ vanishes: by the germ-vanishing statement just
recalled, every $(n,n)$-test form vanishes near $Z$, hence is supported in
$U$, where the difference is zero.  This is \cite[Lemma~7.1.13]{CLD25}.

\cref{itm:ch15-logcur-add}
Both sides are currents on $X$ restricting on the common invertibility locus
to the same current, by multiplicativity of $|\cdot|$; apply the uniqueness
just proved.
\end{proof}

\begin{warning}{Three meanings of ``canonical extension''}
This chapter uses the phrase three times, with three different contents, and
they must not be interchanged.
\begin{enumerate}
\item In \Cref{prop:ch15-log-current} the extension is by \emph{local
  integrability}: the extended current is integration against a function that
  happens to be integrable.  Uniqueness holds because the difference of two
  extensions is a $(0,0)$-current supported on a small set.
\item In \Cref{thm:cld-local-bochner-martinelli} the extension is
  \emph{by zero} across a set of codimension $d$, applied to currents of
  bidegree $(r,r)$ with $r<d$.  Uniqueness there is
  \cite[Lemma~7.1.13]{CLD25} and uses the bidegree bound; it fails for
  $r=d$, which is precisely why the theorem is stated as it is.
\item In \Cref{thm:cld-levine} the extension is of $h\Omega$ across
  $Y=\mathbf P^{n-p}$, and it is proved by exhibiting the extension as a
  pushforward from a blowup.  Canonicity there means independence of the
  choices, not a uniqueness-by-codimension statement.
\end{enumerate}
None of the three follows from the codimension of the removed set alone.  A
current supported on a Zariski-closed set need not vanish once its bidegree
is large enough to survive restriction; the integration current
$\delta_{\divis(f)}$ is itself an example.
\end{warning}

\section{The Poincar\'e--Lelong formula}
\label{sec:ch15-pl}

\begin{theorem}\label{thm:cld-poincare-lelong}
Let $f$ be a regular meromorphic function on $X$.  Then
\[
                    \boxed{\quad
                    \dprimeddprime[\log|f|]
                    =\delta_{\operatorname{div}(f)}.
                    \quad}
\]
\end{theorem}
\begin{proof}
Use the local reductions in \Cref{lem:pl-formal-reductions} below.  Test
against a compactly supported symmetric $(n-1,n-1)$-form $\omega$ and choose
finitely many compact analytic domains tropicalizing it.  Adjoin $f$ to each
moment map and truncate at $s\le|f|\le t$.  Green's formula moves
$\dprimeddprime$ from $\log|f|$ to $\omega$ and reduces the pairing to the
boundary at $|f|=s$.  For sufficiently small $s$, the calibrated local
tropicalization has the product form
\[
          \operatorname{Trop}(V\cap\{f=0\})
                    \times[\log s,\log t].
\]
The outer face has zero contribution by the cutoff and the lateral faces
cancel by balancing.  With the boundary orientation fixed in
\Cref{thm:polyhedral-stokes}, the inner face contributes exactly the
calibrated integral over $f=0$.  Its calibration is the order of vanishing,
including the generic length and residue degree, so the contribution is the
Cartier divisor current rather than the current of the reduced zero set.
Inclusion--exclusion over the finite chart cover now gives
\[
 \langle\dprimeddprime[\log|f|],\omega\rangle
       =\int_{\operatorname{div}(f)}\omega.
\]
This is the product-tropicalization proof of
\cite[Theorem~16.7.3 and \S16.7]{CLD25}.  Regularity supplies the purity and
product description used at the inner face.

The argument is written out in full in
\Cref{sec:ch15-pl-proof-detail}, in six steps, each of which is stated so
that it can be checked separately.
\end{proof}

This is \cite[Theorem~16.7.3]{CLD25}.  The regularity hypothesis is
essential; CLD exhibit the failure for a nonregular meromorphic function in
\cite[Example~16.1.4]{CLD25}.

\begin{lemma}\label{lem:pl-formal-reductions}
To prove \Cref{thm:cld-poincare-lelong}, it suffices locally to treat a
regular holomorphic function $f$ such that $X$ satisfies Serre's condition
$(S_1)$ and $f:X\to\mathbf A^1$ is pure of relative dimension $n-1$.
On the invertibility locus both sides vanish.
\end{lemma}
\begin{proof}
Locally a regular meromorphic function is a quotient $g/h$ of regular
holomorphic functions.  Both divisor and $\log|\cdot|$ are additive under
products and quotients, so it suffices to treat a holomorphic function.  The
remaining reduction is precisely \cite[Lemma~16.7.5]{CLD25}: the ambient
$(S_1)$ condition and purity of the morphism to $\mathbf A^1$ are separate
parts of that reduction.  In particular, the mere existence of the
non-zero-divisor $f$ does not imply that $X$ is $(S_1)$.  If $f$ is
invertible, its divisor is zero and
$\log|f|$ is the pullback of an affine coordinate on $\mathbf R$, whose
$\dprimeddprime$ is zero.

We record the geometric content of that reduction, since it is invoked
repeatedly.  Both sides of the asserted identity are currents, hence local,
so it suffices to prove the equality on a neighbourhood of each $x\in X$.
Near $x$ a section of $\cM_X^\times$ is a quotient $g/h$ of non-zero-divisors
in $\cO_{X,x}$; by \Cref{lem:ch15-divisor-basic}\,\cref{itm:ch15-div-add} and
\Cref{prop:ch15-log-current}\,\cref{itm:ch15-logcur-add} both sides are
additive under this decomposition, so we may assume $f$ holomorphic.  If
$f(x)\neq0$ then $f$ is invertible near $x$ and both sides vanish there, by
\Cref{lem:ch15-log-smooth} and
\Cref{lem:ch15-divisor-basic}\,\cref{itm:ch15-div-unit}.  Assume then
$f(x)=0$.  Because $f$ is a non-zero-divisor, the fibre of
$f:X\to\mathbf A^1$ through $x$ has dimension $n-1$ at $x$, and by upper
semicontinuity of fibre dimension the morphism has relative dimension
$\le n-1$ near $x$; the Hauptidealsatz \cite[Lemma~11.5.2]{CLD25} gives the
reverse inequality, so $f$ is purely of relative dimension $n-1$ on a
neighbourhood of $x$.  Finally, the germ of $f$ at $x$ lies in the maximal
ideal of $\cO_{V,x}$ for an affinoid neighbourhood $V$ and is a
non-zero-divisor there; the existence of a non-zero-divisor in the maximal
ideal is exactly the statement that $\cO_{V,x}$ has depth $\ge1$, that is,
that $X$ satisfies $(S_1)$ at $x$.  Shrinking once more makes $(S_1)$ hold
on a whole neighbourhood.  This is why the two extra hypotheses are
legitimate, and it also shows in what order they are obtained: purity first,
$(S_1)$ afterwards.
\end{proof}

\subsection{The proof in six steps}
\label{sec:ch15-pl-proof-detail}

We now expand the proof of \Cref{thm:cld-poincare-lelong}.  Nothing new is
proved; the point is that each step can be isolated and tested, and that the
place where each hypothesis is used becomes visible.  Throughout, $f$ is a
regular holomorphic function on a boundaryless pure $n$-dimensional $X$
satisfying $(S_1)$, with $f:X\to\mathbf A^1$ purely of relative dimension
$n-1$, as supplied by \Cref{lem:pl-formal-reductions}.  Write
$X_0=\{f=0\}$.

\paragraph{Step 1: reduction to symmetric test forms.}
Both $\delta_{\divis(f)}$ and $\dprimeddprime[\log|f|]$ are symmetric
currents, in the sense that they are invariant under the involution $J$ of
\Cref{chap:superforms}: the first because integration currents of analytic
cycles are symmetric \cite[Proposition~7.7.9]{CLD25}, the second because
$\dprimeddprime$ commutes with $J$ up to the sign that $J$-invariance
absorbs.  Two symmetric currents agreeing on symmetric test forms are equal,
because every test form decomposes as a symmetric part plus a part killed by
pairing with a symmetric current.  We may therefore fix a compactly
supported \emph{symmetric} $\omega\in A_{X,c}^{n-1,n-1}$ and prove
\[
 \langle\dprimeddprime[\log|f|],\omega\rangle=\int_{X_0}\omega .
\]

\paragraph{Step 2: a finite tropicalizing cover and inclusion--exclusion.}
Since $\Supp\omega$ is compact, there is a finite family $(V_i)_{1\le i\le m}$
of compact analytic domains of $X$ tropicalizing $\omega$, such that every
point of $\Supp\omega$ lies in the topological interior of some $V_i$.  Put
$V=\bigcup_iV_i$ and, for $\varnothing\neq E\subset\{1,\dots,m\}$,
$V_E=\bigcap_{i\in E}V_i$.  Each $V_E$ is a finite intersection of compact
analytic domains, hence again a compact analytic domain, and it inherits
$(S_1)$.  On $V_E$ the form $\omega$ is tropical: there are a moment map
$g_E:V_E\to\mathbf G_m^{N}$ and a superform $\alpha_E$ on
$g_{E,\mathrm{trop}}(V_E)$ with $\omega|_{V_E}=g_{E,\mathrm{trop}}^*\alpha_E$.
Adjoin $f$: set $h_E=(g_E,f)$.  Since $(V_{i,0})_i$ is a finite $G$-cover of
$V_0=V\cap X_0$ and $(V_{i,\ge s})_i$ is a finite $G$-cover of
$V_{\ge s}=V\cap\{|f|\ge s\}$, inclusion--exclusion converts every integral
over $V_0$ or $V_{\ge s}$ into an alternating sum of integrals over the
$V_E$.  All the estimates below are therefore made on a single $V_E$.

\paragraph{Step 3: the cutoff.}
Choose $r>0$ so small that the product conclusion of
\Cref{thm:trop-local-germ} holds for every pair $(V_E,g_E)$ on
$\{|f|\le r\}$, and let $(u_1,u_2)$ be a smooth partition of unity on
$\mathbf R$ subordinate to
$(\,]-\infty,\log r[\,,\,]\log r-1,+\infty[\,)$.  Let $v_j$ be the smooth
function on $X$ equal to $u_j(\log|f|)$ on the invertibility locus and to
$u_j(-\infty)$ near $X_0$ --- that is, $v_1\equiv1$ and $v_2\equiv0$ near
$X_0$ --- and put $\omega_j=v_j\omega$.  Then $\omega=\omega_1+\omega_2$,
the support of $\omega_2$ misses $X_0$, and $\omega_1=\omega$ near $X_0$.
Since $\dprimeddprime\log|f|=0$ on $\{f\neq0\}$ by
\Cref{lem:ch15-log-smooth}, both sides of the desired identity vanish on
$\omega_2$.  It remains to treat $\omega_1$.

\paragraph{Step 4: Green's formula and the reduction to one face.}
The form $\ddprime\omega_1$ has type $(n-1,n)$, so its support meets no
Zariski-closed subset of dimension $<n$ \cite[Proposition~3.3.6]{CLD25}; in
particular it misses $X_0$.  Its support being compact and $|f|$ continuous,
there is $s>0$ with $|f|>s$ on $\Supp\ddprime\omega_1$; shrink $s$ further so
that $u_1\equiv1$ on $]-\infty,\log s]$.  By the definition of $[\log|f|]$
and the fact that $V_{\ge s}$ contains $\Supp\dprimeddprime\omega_1$,
\[
 \langle\dprimeddprime[\log|f|],\omega_1\rangle
 =\int_{V_{\ge s}}\log|f|\ \dprimeddprime\omega_1 .
\]
Green's formula on the compact analytic domain $V_{\ge s}$
(\Cref{cor:analytic-stokes-green}, in the version with boundary of
\cite[Proposition~5.11.1]{CLD25}) gives
\[
 \begin{aligned}
  \int_{V_{\ge s}}\log|f|\,\dprimeddprime\omega_1
   ={}&\int_{V_{\ge s}}(\dprimeddprime\log|f|)\wedge\omega_1\\
     &+\int_{V_{\ge s}}^{\partial}\log|f|\ \ddprime\omega_1
      -\int_{V_{\ge s}}^{\partial}
       (\ddprime\log|f|)\wedge\omega_1 .
 \end{aligned}
\]
The first term vanishes because $f$ is invertible on $V_{\ge s}$; the second
vanishes because $\Supp\ddprime\omega_1$ does not meet
$\partial(V_{\ge s})$.  Hence everything reduces to the single boundary
integral
\[
 \langle\dprimeddprime[\log|f|],\omega_1\rangle
 =-\int^{\partial}_{V_{\ge s}}(\ddprime\log|f|)\wedge\omega_1 ,
\]
and, since $\omega_1=\omega$ near $X_0$, the theorem becomes the assertion
\begin{equation}
\label{eq:ch15-pl-face}
 \int^{\partial}_{V_{\ge s}}\ddprime(\log|f|)\wedge\omega_1
 =-\int_{V_0}\omega .
\end{equation}

\paragraph{Step 5: the tropical product.}
Fix $E$ and put $\sigma=\log s$, $\tau=\log\rho$ for a fixed
$\rho\in\,]s,r[$ beyond which $u_1$ vanishes; write
$V_{E,[s,\rho]}=V_E\cap\{s\le|f|\le\rho\}$.  On $V_{E,[s,\rho]}$ the map
$h_E=(g_E,f)$ is a moment map.  Write $w=\log|f|$ for the parameter of the
interval factor; the tropical coordinate of \Cref{lem:ch15-log-smooth} is
$t=-\log|f|=-w$, so that $\ddprime\log|f|=\ddprime w=-\ddprime t$, and
\[
 \omega_1|_{V_{E,[s,\rho]}}=h_{E,\mathrm{trop}}^*(u_1\alpha_E),
 \qquad
 \ddprime\log|f|\wedge\omega_1
 =h_{E,\mathrm{trop}}^*\bigl(\ddprime w\wedge u_1\alpha_E\bigr).
\]
Let $(P,\mu)$ be the calibrated polyhedron
$g_{E,\mathrm{trop}}(V_{E,0})^{(n-1)}$ and $(Q,\nu)$ the calibrated
polyhedron $h_{E,\mathrm{trop}}(V_{E,[s,\rho]})^{(n)}$.  The choice of $r$ in
Step 3 is exactly the choice for which
\[
 (Q,\nu)=(P\times[\sigma,\tau],\ \mu\times\lambda),
\]
$\lambda$ being the standard calibration of the interval.  This is the
content of \Cref{thm:trop-local-germ}; it is the only place where the purity
of the relative dimension is used, and it is what converts an analytic
statement into a polyhedral one.

\paragraph{Step 6: the three faces.}
Choose a cellular decomposition $\mathscr C$ of $P$ adapted to $\mu$.  The
cells $C\times\{\sigma\}$, $C\times\{\tau\}$ and $C\times[\sigma,\tau]$ for
$C\in\mathscr C$ decompose $P\times[\sigma,\tau]$ compatibly with its
calibration, and the boundary of $Q$ has three parts.
\begin{enumerate}
\item \emph{The outer face $P\times\{\tau\}$ contributes zero}, because
    $u_1$ vanishes near $\tau$ by the choice of $\rho$.
\item \emph{The lateral faces $\partial P\times[\sigma,\tau]$ contribute
  zero}.  Each such contribution is an integral over $F\times[\sigma,\tau]$,
  with $F$ a cell of $\mathscr C$ of dimension $n-2$, of an $(n-1)$-form that
  has already had the differential $\ddprime w$ contracted away and therefore
  comes from a form on $F$ alone; such a form vanishes on a product with an
  interval for degree reasons.
\item \emph{The inner face $P\times\{\sigma\}$ contributes
  $-\int_{V_{E,0}}\omega$}.  For an $(n-1)$-cell $F$ of $P$, the discordance
  $\partial\nu_{F\times\{\sigma\}}$ equals $-\mu_F$ times the volume vector
  of the interval direction; since $\ddprime t\wedge\alpha=\alpha\wedge
  \ddprime t$ and $u_1(\sigma)=1$, the contribution of $F\times\{\sigma\}$
  is $-\int_F\langle\alpha,\mu_F\rangle$.  Summing over $F$ gives
  $-\int_P\langle\alpha,\mu\rangle=-\int_{V_{E,0}}\omega$.
\end{enumerate}
This proves \eqref{eq:ch15-pl-face} on each $V_E$; inclusion--exclusion from
Step 2 assembles the pieces, and the theorem follows.

Three comments on the proof are worth recording.

First, \emph{the multiplicity is never computed}.  Nowhere does the argument
differentiate $f$ or expand it in a local parameter.  The coefficient
attached to a component of $\{f=0\}$ enters only through the calibration
$\mu$ of $\Trop(V_{E,0})$, that is, through
\Cref{def:trop-canonical-calibration}, and by
\Cref{thm:trop-multiplicity-formula} that calibration is a product of a
length, a residue degree, and a lattice index.  This is why the right-hand
side is the Cartier divisor current and not the current of the reduced zero
set: the length factor is built into the calibration.  The reader who wants
to see the length appear concretely should look at
\Cref{ex:trop-multiplicity-two,ex:trop-length-factor}.

Second, \emph{the boundary orientation is where the sign lives}.  Two minus
signs occur: the one in front of the last term of Green's formula, and the
one in the discordance of the inner face.  They multiply to $+1$, and the
result is the stated sign.  Reversing either one --- for instance by
orienting the truncation as $\{|f|\le s\}$ rather than $\{|f|\ge s\}$ ---
reverses the sign of the whole formula.  \Cref{sec:ch15-tests} checks the
outcome against the graph Laplacian, where the sign is fixed independently.

Third, \emph{boundarylessness of $X$ is used twice}, once through the version
of Green's formula in \Cref{cor:analytic-stokes-green} and once in the
identification of $\partial V_{\ge s}$ with the level set $\{|f|=s\}$.  On a
space with boundary the formula acquires an extra boundary term; we do not
state it here, because the only later use of a boundary version is on
compact analytic domains inside a boundaryless space, where the term has
already been computed above.

\subsection{Deriving the constant and the sign}
\label{sec:ch15-pl-constant}

The normalization in \Cref{thm:cld-poincare-lelong} is not a convention that
one is free to choose once $\dprimeddprime$ and $\delta_{\divis(f)}$ have
been fixed: given the Hessian convention \eqref{eq:appJ-real-hessian} and the
calibrated integration of \Cref{chap:tropical-integration}, the constant
$+1$ and the sign are forced.  The following computation is the shortest
derivation, and it is worth doing once so that the formula need never be
remembered.

\begin{lemma}
\label{lem:ch15-pl-constant}
Let $K$ be complete and nontrivially valued, let
$D^\circ=\{x\in\mathbf A^{1,\An}_K:|T(x)|<1\}$ be the open unit disc, and let
$o$ be the rigid point $T=0$.  Then
\[
 \dprimeddprime[\log|T|]=\delta_o
\]
as currents on the boundaryless one-dimensional space $D^\circ$.
\end{lemma}

\begin{proof}
Write $\Gamma=\{\zeta_{0,\rme^{-\theta}}:0<\theta\le+\infty\}$ for the
skeleton of $D^\circ$, with $\theta=v(T)$ and $\theta=+\infty$ corresponding
to $o$, and let $r:D^\circ\to\Gamma$ be the retraction.

The current $\dprimeddprime[\log|T|]$ has bidegree $(1,1)$ on a
one-dimensional space, hence bidimension $(0,0)$: it is a signed Radon
measure, by the identification in \Cref{chap:cld-currents}.  On
$D^\circ\setminus\{o\}$ the function $T$ is invertible, so
\Cref{lem:ch15-log-smooth} gives $\dprimeddprime\log|T|=0$ there.  Therefore
the measure is supported on the single point $o$, that is,
$\dprimeddprime[\log|T|]=c\,\delta_o$ for some $c\in\mathbf R$.  It remains to
compute $c$, and for that one test function suffices.

Choose $0<\theta_0<\theta_1<\infty$ and a smooth $\psi:\mathbf R\to[0,1]$
with $\psi\equiv0$ on $]-\infty,\theta_0]$ and $\psi\equiv1$ on
$[\theta_1,+\infty[$, and set $\varphi=\psi\circ r$, which is a smooth
compactly supported function on $D^\circ$ with $\varphi(o)=1$.  All the data
are pulled back from $\Gamma$, so the pairing may be computed there.  The
calibrated Lebesgue measure of $\Gamma$ is $\rd\theta$ for the lattice
$\mathbf Z\subset\mathbf R$ (\Cref{def:ch11-lattice-measure}), and on
$\Gamma$ one has $\log|T|=-\theta$; the Hessian convention
\eqref{eq:appJ-real-hessian} in one variable reads
$\dprimeddprime\varphi=\psi''(\theta)\,\dprime\theta\wedge\ddprime\theta$.
Hence, by \Cref{prop:current-differential-calculus},
\[
 \langle\dprimeddprime[\log|T|],\varphi\rangle
 =\langle[\log|T|],\dprimeddprime\varphi\rangle
 =\int_0^{+\infty}(-\theta)\,\psi''(\theta)\,\rd\theta .
\]
Two integrations by parts on $[0,+\infty)$, legitimate because $\psi''$ is
supported in the compact interval $[\theta_0,\theta_1]$, give
\[
 \int_0^{+\infty}(-\theta)\psi''(\theta)\rd\theta
 =\bigl[-\theta\psi'(\theta)\bigr]_0^{+\infty}
   +\int_0^{+\infty}\psi'(\theta)\rd\theta
 =0+\bigl[\psi\bigr]_0^{+\infty}=1-0=1 .
\]
Therefore $c=\langle c\delta_o,\varphi\rangle=1$.
\end{proof}

\begin{remark}[Why one test function is enough]
The step that makes \Cref{lem:ch15-pl-constant} short is the observation that
the current is a measure supported at a point, so that a \emph{single}
normalized test function determines it.  One should resist the temptation to
run the same computation on the closed disc $\cM(K\{T\})$: there the Gauss
point $\zeta_{0,1}$ is the analytic boundary, the space is not boundaryless,
and the two integrations by parts produce an extra endpoint term
$-\psi(0)$ which is precisely the boundary contribution excluded by the
hypotheses of \Cref{thm:cld-poincare-lelong}.  The same term reappears, with
the same sign, in \Cref{ex:ch15-boundary-max} as a maximum-principle failure.
\end{remark}

\begin{remark}[Reading the constant off the Hessian convention]
\label{rem:ch15-constant-reading}
The mechanism behind \Cref{lem:ch15-pl-constant} is worth stating without
the analytic dressing.  On $\mathbf R$ with lattice $\mathbf Z$, the Hessian
convention \eqref{eq:appJ-real-hessian} reads
$\dprimeddprime u=u''\,\dprime\theta\wedge\ddprime\theta$ and the positive
orientation \eqref{eq:appJ-lagerberg-orientation} is
$\tau_1=\dprime\theta\wedge\ddprime\theta$.  Therefore the current
$\dprimeddprime[u]$ is the distributional second derivative $u''$ with
\emph{no} extra constant: no factor $2$, no factor $2\pi$, no factor
$\rmi$.  The function $\theta\mapsto\max\{0,\theta\}$ has
$\dprimeddprime=\delta_0$, as computed in \Cref{chap:cld-currents}, and
$\theta\mapsto|\theta|$ has $\dprimeddprime=2\delta_0$, matching the graph
Laplacian convention below \eqref{eq:graph-laplacian}.  Since
$-\log|T|=\theta$ has a slope break of size $+1$ at the end of the skeleton
where $T=0$, the mass is $1$.  It is the \emph{absence} of the complex
normalizing constant that distinguishes
$\dprimeddprime\log|f|=\delta_{\divis(f)}$ from
$\ddc\log|f|=\tfrac12[\divis(f)]$: the complex convention
\eqref{eq:appJ-ddc} divides by $2\pi$ and works with $\log|f|^2$, so it
carries the factor $\tfrac12$ visible in \eqref{eq:appJ-complex-pl}, whereas
the real Hessian convention carries none.  The two are compared once and for
all in \Cref{sec:appJ-cld}; they are never identified.
\end{remark}

\begin{corollary}
\label{cor:ch15-appJ-consistency}
\Cref{thm:cld-poincare-lelong} is \eqref{eq:appJ-cld-pl}, and it implies
\eqref{eq:appJ-valuation-pl}:
\[
 \dprimeddprime[v(f)]=-\delta_{\divis(f)},
 \qquad v(f)=-\log|f| .
\]
Moreover the normalization passes the principal-divisor test and the
special-fibre test of \Cref{sec:appJ-unit-tests}.
\end{corollary}

\begin{proof}
The first assertion is \Cref{thm:cld-poincare-lelong} together with linearity
of $\dprimeddprime$ and $v(f)=-\log|f|$, which is
\eqref{eq:appJ-valuation-sign}.  For the principal-divisor test, note that
the two displayed formulas differ by a sign, so any identity in which
$\log|f|$ has been silently replaced by $v(f)$ changes sign and is detected.
For the special-fibre test, let $K$ be discretely valued with $v_K(\pi)=1$
and let $D=\cX_s$ be the special fibre of a model, locally cut out by $\pi$.
By \eqref{eq:appJ-model-divisor} the associated model function is
$\phi_D=v(\pi)=1$, a constant, whose $\dprimeddprime$ vanishes; and indeed
$\divis(\pi)$ is trivial on the generic fibre $X$, since $\pi\in K^\times$ is
invertible there.  Both sides of the formula are zero, as they must be.  The
apparent paradox that an effective vertical divisor gives a nonzero,
nonnegative model function while its Poincar\'e--Lelong current vanishes is
resolved in \Cref{sec:appJ-cld}: the curvature current of a model metric is
not the bare current of a principal function on a chart, and the two are only
compared through \eqref{eq:appJ-cld-section-pl}.
\end{proof}

\subsection{The section form of the formula}
\label{sec:ch15-pl-section}

For later comparison with complex geometry it is convenient to record the
version of \Cref{thm:cld-poincare-lelong} for a meromorphic section of a
metrized line bundle.  The metric formalism itself belongs to
\Cref{chap:cld-metrized-bundles}; we use only the definition of the curvature
current there and the fact that it is local.

\begin{corollary}
\label{cor:ch15-pl-section}
Let $L$ be a line bundle on $X$ equipped with a locally integrable
semimetric, with curvature current $c_1(L)$ as in
\Cref{chap:cld-metrized-bundles}, and let $s$ be a regular meromorphic
section of $L$.  Then
\[
 \dprimeddprime\bigl[-\log\|s\|\bigr]+\delta_{\divis(s)}=c_1(L).
\]
\end{corollary}

\begin{proof}
The assertion is local, so work on an open $U$ on which $L$ has an invertible
section $\varepsilon$, and write $s|_U=f\varepsilon$ with $f$ the unique
regular meromorphic function determined by this equation.  Put
$\lambda_\varepsilon=-\log\|\varepsilon\|$.  By definition of the curvature
current, \eqref{eq:appJ-cld-c1}, one has
$c_1(L)|_U=\dprimeddprime[\lambda_\varepsilon]$, and
$-\log\|s\|=\lambda_\varepsilon-\log|f|$ on the locus where $f$ is
invertible.  \Cref{prop:ch15-log-current} extends both terms to $U$, and
\Cref{thm:cld-poincare-lelong} gives
\[
 \begin{aligned}
  \dprimeddprime[-\log\|s\|]
   &=\dprimeddprime[\lambda_\varepsilon]-\dprimeddprime[\log|f|]\\
   &=c_1(L)|_U-\delta_{\divis(f)}\\
   &=c_1(L)|_U-\delta_{\divis(s)}|_U,
 \end{aligned}
\]
the last equality because $\varepsilon$ is invertible, so
$\divis(s)|_U=\divis(f)$.  Since the two sides agree on a cover, they agree
on $X$.  This is \cite[Theorem~18.1.6]{CLD25}.
\end{proof}

\Cref{cor:ch15-pl-section} is \eqref{eq:appJ-cld-section-pl}.  Comparing it
with the complex identity \eqref{eq:appJ-complex-section-pl} shows that the
two have the same formal shape, provided one remembers that the complex
potential is $-\log|s|_h^2$ and the non-Archimedean one is $-\log\|s\|$; the
factor $2$ between them is exactly the factor recorded in
\eqref{eq:appJ-complex-weight}.  This is the safest of all the comparisons in
\Cref{app:normalizations}, and the reader who is unsure of a sign in a
mixed formula should reduce it to this one.

\begin{corollary}
\label{cor:ch15-degree-zero}
Let $C$ be a proper $K$-analytic curve and $f\in\cM_C^\times(C)$ a global
regular meromorphic function.  Then the signed measure
$\dprimeddprime[\log|f|]$ has total mass zero, and
\[
 \deg\divis(f)=\sum_Pm_P\,[\kappa(P):K]=0 .
\]
\end{corollary}

\begin{proof}
Pair \Cref{thm:cld-poincare-lelong} with the constant function $1$, which is
a legitimate test function because $C$ is compact and boundaryless:
\[
 \langle\delta_{\divis(f)},1\rangle
 =\langle\dprimeddprime[\log|f|],1\rangle
 =\langle[\log|f|],\dprimeddprime1\rangle=0
\]
by \Cref{prop:current-differential-calculus}, since $\dprimeddprime1=0$.  The
left-hand side is $\sum_Pm_P\langle\delta_P,1\rangle$ and
$\langle\delta_P,1\rangle=[\kappa(P):K]$ by the definition of the integration
current of a rigid point.  The argument is \cite[Corollary~18.1.7]{CLD25}
with the metric taken trivial.  Note that it does \emph{not} use algebraicity
of $C$; that is needed only for the version with a metrized line bundle,
where the total mass is $\deg L$.
\end{proof}

\Cref{cor:ch15-degree-zero} is the current-theoretic form of the identity
\eqref{eq:lap-mass-zero} that a graph Laplacian has total mass zero.  On a
curve with a semistable model the two statements are literally the same, by
the slope formula \Cref{thm:ch08-slope-formula}; the next section makes this
explicit.

\section{Sign and normalization tests}
\label{sec:ch15-tests}

A formula whose sign is wrong is worse than no formula, and a formula whose
normalization is wrong is harder to detect.  This section is a battery of
checks.  Each of them computes both sides of \Cref{thm:cld-poincare-lelong}
independently --- the left-hand side by the differential calculus of
\Cref{chap:cld-forms,chap:cld-currents}, the right-hand side by counting
zeros --- and compares.  Two of the checks are also checks on
\Cref{chap:model-calculations}: they recompute a graph Laplacian that was
obtained there from the slope formula, without using
\Cref{thm:cld-poincare-lelong}, so that agreement is evidence and not
tautology.

\begin{example}[A zero on an analytic curve]
\label{ex:ch15-curve-zero}
Let $C$ be a smooth analytic curve and let $z$ be a parameter at a rigid
point $a$.  Then
\[
                    \dprimeddprime[\log|z|]=\delta_a.
\]
On a skeleton emanating from $a$, the function $\log|z|$ has outgoing slope
$+1$ when distance is oriented away from the zero using the logarithmic
radius convention.  The graph Laplacian therefore has mass $+1$ at $a$.
For $v(z)=-\log|z|$, the mass is $-1$.
\end{example}

\begin{example}[The unit disc, in coordinates]
\label{ex:ch15-disc-coordinate}
Take $K$ nontrivially valued and algebraically closed, and let
$D^\circ=\{|T|<1\}$ be the open unit disc.  Its skeleton is the half-open
segment $\Gamma=\{\zeta_{0,\rme^{-\theta}}:0<\theta\le+\infty\}$, with
$\theta=v(T)$ increasing towards the rigid point $T=0$ at $\theta=+\infty$;
see \eqref{eq:gauss-disc}.  \Cref{lem:ch15-pl-constant} computes the current
directly.  In graph language the same computation reads: the function
$G(\theta)=\log|T|=-\theta$ is affine of slope $-1$ in the direction of
increasing $\theta$, so its only slope break in $\Gamma$ is at the endpoint
$\theta=+\infty$, where the single outgoing slope of $G$ --- taken in the
direction of \emph{decreasing} $\theta$, the only tangent direction at an
endpoint --- equals $+1$.  Hence $\Delta_\Gamma G=\delta_{T=0}$ by
\eqref{eq:graph-laplacian}, the same answer.

Two traps are worth naming.  First, the point $T=0$ is an endpoint of
$\Gamma$, and an endpoint has exactly one tangent direction
(\Cref{ex:ch08-valence}); a computation that counted two would produce mass
$2$.  Second, the sign of the slope depends on which end of $\Gamma$ one
calls the origin: $\theta=v(T)$ increases towards the zero of $T$, whereas
$\log|T|$ decreases towards it.  This is the sign
\eqref{eq:appJ-valuation-sign}, and it is why
$\dprimeddprime[v(T)]=-\delta_{T=0}$.
\end{example}

\begin{example}[Rational functions on $\mathbf P^1$]
\label{ex:ch15-rational-p1}
For
\[
                    f(T)=c\frac{\prod_i(T-a_i)^{m_i}}
                                  {\prod_j(T-b_j)^{n_j}},
\]
Poincar\'e--Lelong gives, on the affine line,
\[
 \dprimeddprime[\log|f|]
       =\sum_i m_i\delta_{a_i}-\sum_jn_j\delta_{b_j}.
\]
On $\mathbf P^{1,\An}$ the divisor of $f$ carries the further coefficient
$\sum_jn_j-\sum_im_i$ at $\infty$, so there
\[
 \dprimeddprime[\log|f|]
       =\sum_i m_i\delta_{a_i}-\sum_jn_j\delta_{b_j}
        +\Bigl(\sum_jn_j-\sum_im_i\Bigr)\delta_\infty ,
\]
and only now is the total mass zero, because a principal divisor on
$\mathbf P^1$ has degree zero.  Dropping the term at infinity leaves mass
$\sum_im_i-\sum_jn_j$, which is nonzero whenever the numerator and
denominator degrees differ; the smallest instance is $f(T)=T$, computed in
\Cref{ex:ch15-toric-compactified} as $\delta_0-\delta_\infty$.  Retracting to a finite skeleton recovers the slope formula:
the sum of outgoing slopes at a vertex is the coefficient of the retracted
divisor there.
\end{example}

\subsection{A quartic on the Berkovich affine line, computed}
\label{sec:ch15-quartic}

We now verify \Cref{thm:cld-poincare-lelong} by hand, in explicit numbers, for
a polynomial with four distinct roots.  The point of the computation is that
neither side is quoted: the divisor is read off from a Newton polygon, the
current is computed as a tree Laplacian from the outgoing slopes of
$\log|f|$, and the two are then compared.  All the branching facts used are
from \Cref{sec:ch08-p1-tree,sec:ch08-p1-newton}.

Let $K=\mathbf C_p$, so that $K$ is complete and algebraically closed with
$v(p)=1$ and $v(K^\times)=\mathbf Q$.  Consider
\begin{equation}
\label{eq:ch15-quartic}
 f(T)=T^4+p\,T^3+p^3T^2+p^6T+p^{10}
 \qquad\text{on }\mathbf A^{1,\An}_K .
\end{equation}

\paragraph{The divisor, from the Newton polygon.}
Along the skeleton of $\mathbf G_m^{\An}$, parametrized by
$\theta=v_\theta(T)$ so that the point at parameter $\theta$ is the Gauss
point $\zeta_{0,\rme^{-\theta}}$ of \eqref{eq:gauss-disc}, one has
\[
 F(\theta):=v_\theta(f)
 =\min\{4\theta,\ 1+3\theta,\ 3+2\theta,\ 6+\theta,\ 10\}.
\]
The five affine functions have slopes $4,3,2,1,0$ and cross consecutively:
$4\theta=1+3\theta$ at $\theta=1$; $1+3\theta=3+2\theta$ at $\theta=2$;
$3+2\theta=6+\theta$ at $\theta=3$; $6+\theta=10$ at $\theta=4$.  A direct
check at the four candidate breakpoints confirms that the minimum is attained
by consecutive terms only; for instance at $\theta=2$ the five values are
$8,7,7,8,10$, so the minimum $7$ is attained by the terms of slope $3$ and
$2$ and by no other.  Hence
\begin{equation}
\label{eq:ch15-quartic-F}
 F(\theta)=
 \begin{cases}
  4\theta, & \theta\le1,\\
  1+3\theta, & 1\le\theta\le2,\\
  3+2\theta, & 2\le\theta\le3,\\
  6+\theta, & 3\le\theta\le4,\\
  10, & \theta\ge4,
 \end{cases}
 \qquad
 F(1)=4,\ F(2)=7,\ F(3)=9,\ F(4)=10 .
\end{equation}
So $F$ is concave, piecewise affine with integral slopes $4,3,2,1,0$, and
drops slope by exactly $1$ at each of $\theta=1,2,3,4$.  Each segment of the
Newton polygon therefore has horizontal length one, and $f$ has four distinct
roots $P_1,P_2,P_3,P_4$ with
\[
 v(P_j)=j,\qquad j=1,2,3,4 ,
\]
each simple.  As a divisor on $\mathbf A^{1,\An}$,
\begin{equation}
\label{eq:ch15-quartic-divisor}
 \divis(f)=P_1+P_2+P_3+P_4 ,
\end{equation}
and on $\mathbf P^{1,\An}$ one adds $-4\,(\infty)$.

\paragraph{The tree, and where the branches separate.}
Since $v(P_i)=i\neq j=v(P_j)$ for $i\neq j$, the ultrametric inequality is an
equality:
\[
 v(P_i-P_j)=\min\{v(P_i),v(P_j)\}=\min\{i,j\}.
\]
By \eqref{eq:ch08-path-metric} the paths $[\zeta_{0,1},P_i]$ and
$[\zeta_{0,1},P_j]$ therefore separate exactly at the point of parameter
$\theta=\min\{i,j\}$.  Consequently the convex hull of
$\{\zeta_{0,1},P_1,P_2,P_3,P_4\}$ in $\mathbf A^{1,\An}$ is the segment
$\{\zeta_{0,\rme^{-\theta}}:0\le\theta\le4\}$ with a pendant edge attached at
each $\theta=j$ running out to $P_j$.  Note in particular that the branch
from $\theta=4$ towards $P_4$ is \emph{not} the continuation of the segment:
the discs $\zeta_{0,\rme^{-\theta}}$ with $\theta>4$ have radius smaller than
$|P_4|$ and do not contain $P_4$.  Exactly three of the tangent directions at $\theta=j$ carry a nonzero
outgoing slope for $j=1,2,3$, and two do at $\theta=4$; on every other
direction $\log|f|$ is constant.  There are of course infinitely many
tangent directions at each of these type-II points, one for each element of
$\mathbf P^1(\widetilde K)$, by \Cref{prop:ch08-branching}.

\paragraph{The function $\log|f|$ on each branch.}
Because $K$ is algebraically closed and $f$ is monic,
\begin{equation}
\label{eq:ch15-quartic-product}
 |f|(\zeta_{a,r})=\prod_{j=1}^4\max\{|a-P_j|,\,r\},
\end{equation}
which is the multiplicativity of the Gauss seminorm
(\Cref{lem:trop-gauss-multiplicative}) applied to the factorization
$f=\prod_j(T-P_j)$, together with
$|T-P_j|(\zeta_{a,r})=\max\{|a-P_j|,r\}$.  We compute the outgoing slopes of
$\log|f|$ in every tangent direction of the tree; distances are logarithmic
radii, as in \eqref{eq:ch08-path-metric}.

\emph{On the main segment.}  There $\log|f|=-F$, so by
\eqref{eq:ch15-quartic-F} its slope in the direction of increasing $\theta$ is
$-4,-3,-2,-1,0$ on the successive intervals.  For $\theta<0$, that is, on
discs of radius $>1$, the minimum in $F$ is attained by the term $4\theta$,
so $\log|f|=-4\theta$ there as well, with slope $-4$.

\emph{On the pendant edge towards $P_j$.}  Parametrize it by
$\zeta_{P_j,\rme^{-\rho}}$ with $\rho\ge j$, so $\rho=j$ is the branch point
and $\rho=+\infty$ is $P_j$.  For $i\neq j$ we have
$|P_j-P_i|=\rme^{-\min\{i,j\}}\ge\rme^{-j}\ge\rme^{-\rho}$, so the $i$-th
factor of \eqref{eq:ch15-quartic-product} is the constant
$\rme^{-\min\{i,j\}}$; and the $j$-th factor is $\rme^{-\rho}$.  Hence
\[
 \log|f|\bigl(\zeta_{P_j,\rme^{-\rho}}\bigr)
 =-\rho-\sum_{i\neq j}\min\{i,j\},
\]
an affine function of $\rho$ with slope $-1$.

\emph{On the other branches at the Gauss point.}  If $|a|=1$ and $r<1$, then
$|a-P_j|=1$ for every $j$ because $|P_j|<1$, so
\eqref{eq:ch15-quartic-product} gives $|f|=1$: the function $\log|f|$ is
constant on every residue-class branch at $\zeta_{0,1}$ other than the two
lying on the main segment.

\paragraph{The Laplacian, vertex by vertex.}
We now sum outgoing slopes, using \eqref{eq:graph-laplacian}.  Recall that an
outgoing slope is computed with the tangent direction pointing \emph{away}
from the vertex.
\begin{center}
\small
\begin{tabular}{@{}llll@{}}
\toprule
Point & outgoing slopes & sum & mass \\
\midrule
$\zeta_{0,1}$ ($\theta=0$)
 & $+4$ (towards $\theta<0$), $-4$ (towards $\theta>0$), $0$ (others)
 & $0$ & $0$\\
$\theta=1$
 & $+4$, $-3$, $-1$ (towards $P_1$) & $0$ & $0$\\
$\theta=2$
 & $+3$, $-2$, $-1$ (towards $P_2$) & $0$ & $0$\\
$\theta=3$
 & $+2$, $-1$, $-1$ (towards $P_3$) & $0$ & $0$\\
$\theta=4$
 & $+1$, $\ \ 0$, $-1$ (towards $P_4$) & $0$ & $0$\\
$P_j$ ($j=1,2,3,4$)
 & $+1$ (the unique direction) & $+1$ & $+1$\\
\bottomrule
\end{tabular}
\end{center}
At an interior vertex of the main segment, the first entry is the outgoing
slope in the direction of decreasing $\theta$, which is the negative of the
slope of $\log|f|$ on the interval to the left; the second is the outgoing
slope in the direction of increasing $\theta$, which is that slope on the
interval to the right; and the third is the slope $-1$ computed on the
pendant edge.  Thus at $\theta=2$, for instance, the three numbers are
$-(-3)=+3$, $-2$ and $-1$, whose sum is zero.  At the rigid point $P_j$ there
is a single tangent direction (\Cref{ex:ch08-valence}), pointing back towards
$\rho$ decreasing, along which $\log|f|$ increases at unit rate, so the
outgoing slope is $+1$.

Every vertex mass vanishes except at the four roots, where it is $+1$.
Therefore
\[
 \dprimeddprime[\log|f|]
 =\delta_{P_1}+\delta_{P_2}+\delta_{P_3}+\delta_{P_4}
 =\delta_{\divis(f)},
\]
which is \Cref{thm:cld-poincare-lelong} in this case, with the divisor
\eqref{eq:ch15-quartic-divisor} obtained independently from the Newton
polygon.  The total mass is $4=\deg f$.

\paragraph{The check on a compact skeleton.}
It is instructive to repeat the computation after retracting to the finite
skeleton $\Gamma=\{\zeta_{0,\rme^{-\theta}}:0\le\theta\le5\}$, that is, on
the skeleton of the model of $\mathbf P^1$ obtained from $\mathbf P^1_R$ by
five successive blowups of the point $T=0$ of the special fibre.  Write
$F=v_\bullet(f)|_\Gamma$ as in \eqref{eq:ch15-quartic-F}.  The outgoing-slope
sums are
\begin{align*}
 \theta=0:&\ d_\nu F=+4 &&\Longrightarrow\ \text{mass}\ +4;\\
 \theta=1:&\ d_\nu F=-4,\ +3 &&\Longrightarrow\ \text{mass}\ -1;\\
 \theta=2:&\ d_\nu F=-3,\ +2 &&\Longrightarrow\ \text{mass}\ -1;\\
 \theta=3:&\ d_\nu F=-2,\ +1 &&\Longrightarrow\ \text{mass}\ -1;\\
 \theta=4:&\ d_\nu F=-1,\ \ \,0 &&\Longrightarrow\ \text{mass}\ -1;\\
 \theta=5:&\ d_\nu F=\ \ \,0 &&\Longrightarrow\ \text{mass}\ \ \ 0,
\end{align*}
so that
\[
 \Delta_\Gamma F
 =4\delta_{\theta=0}-\delta_{\theta=1}-\delta_{\theta=2}
   -\delta_{\theta=3}-\delta_{\theta=4},
\]
of total mass $4-4=0$, as \eqref{eq:lap-mass-zero} demands.  On the other
side, $\divis(f)=P_1+P_2+P_3+P_4-4\,(\infty)$ on $\mathbf P^1$, and the
retraction sends $P_j$ to $\theta=j$ and sends $\infty$, which lies outside
the closed unit disc, to $\theta=0$.  Hence
\[
 r_*\divis(f)
 =\delta_{\theta=1}+\delta_{\theta=2}+\delta_{\theta=3}+\delta_{\theta=4}
   -4\delta_{\theta=0}
 =-\Delta_\Gamma F ,
\]
which is \eqref{eq:ch08-slope-formula}, and equivalently
$\Delta_\Gamma(\log|f|\,|_\Gamma)=r_*\divis(f)$, which is
\eqref{eq:p1-poincare-lelong}.  The two computations agree, and the second is
the first pushed forward along the retraction; the pendant-edge slopes $-1$
of the first computation are exactly what the retraction collapses into the
vertex masses $-1$ of the second.

\begin{warning}{The retracted formula is not the formula}
The identity $\Delta_\Gamma(\log|f|\,|_\Gamma)=r_*\divis(f)$ is the
\emph{image} of \Cref{thm:cld-poincare-lelong} under $r_*$, not an equivalent
statement.  Two different divisors with the same retraction give the same
right-hand side.  In the quartic above, replacing $P_3$ by a root of the same
valuation lying in a different residue class of the disc of radius
$\rme^{-2}$ changes the tree but not $\Gamma$ or the retracted formula.  This
is the reason the slope formula \Cref{thm:ch08-slope-formula} is stated with
a retraction and the current formula is not.
\end{warning}

\begin{example}[A double root, and where the multiplicity comes from]
\label{ex:ch15-double-root}
Modify \eqref{eq:ch15-quartic} to
\[
 g(T)=(T-P_1)^2(T-P_3),\qquad v(P_1)=1,\ v(P_3)=3 .
\]
Along the main segment $\log|g|$ has slopes $-3,-1,0$ on $\theta\le1$,
$1\le\theta\le3$, $\theta\ge3$; on the pendant edge towards $P_1$ the
computation of \eqref{eq:ch15-quartic-product} gives
$\log|g|=-2\rho+\text{const}$, of slope $-2$, and on the edge towards $P_3$
it gives slope $-1$.  The vertex sums are then $3-1-2=0$ at $\theta=1$ and
$1+0-1=0$ at $\theta=3$; the endpoint masses are $+2$ at $P_1$ and $+1$ at
$P_3$.  Thus
\[
 \dprimeddprime[\log|g|]=2\delta_{P_1}+\delta_{P_3}
 =\delta_{\divis(g)} .
\]
The multiplicity $2$ has entered through the \emph{slope} of $\log|g|$ on the
pendant edge, not through any local expansion.  In the general proof this
slope is the calibration of the base of the tropical cylinder in Step 5 of
\Cref{sec:ch15-pl-proof-detail}, and \Cref{thm:trop-multiplicity-formula}
identifies it with the length of the local ring.  Over a field that is not
algebraically closed the same slope would additionally record the residue
degree $[\kappa(P):K]$, as in \Cref{cor:ch15-degree-zero}.
\end{example}

\begin{example}[A nonmonic leading coefficient does not move the divisor]
\label{ex:ch15-constant-shift}
Replacing $f$ by $cf$ with $c\in K^\times$ adds the constant $\log|c|$ to
$\log|f|$ and does not change $\divis(f)$.  Both sides of
\Cref{thm:cld-poincare-lelong} are therefore unchanged, since
$\dprimeddprime$ kills constants.  This is the constant test of
\Cref{sec:appJ-unit-tests} in its simplest form, and it is the reason the
formula involves $\log|f|$ rather than, say, $\log(|f|/\|f\|)$ for some
normalization $\|f\|$: any such normalization would have to be constant to
preserve the identity.
\end{example}

\begin{example}[A toric character]
\label{ex:ch15-toric-character}
On the torus, $\log|\chi^m|=-\langle m,u\rangle$ for
$u=-\log|\chi|$, so its tropical Hessian is zero.  This does not contradict
Poincar\'e--Lelong: the character has no zero or pole on the torus.  After a
toric compactification, the missing boundary divisor is exactly where the
piecewise affine extension changes slope.
\end{example}

The last example deserves to be pushed one step further, because it is the
cleanest illustration of the difference between an identity of currents and
an identity of piecewise affine functions.

\begin{example}[Where the toric divisor went]
\label{ex:ch15-toric-compactified}
Take $N=\mathbf Z$, $T=\mathbf G_{m,K}$, and the character $\chi(T)=T$.  On
$T^{\An}$ the tropical coordinate is $u=-\log|T|=v(T)$, by
\eqref{eq:appJ-toric-coordinate}, and $\log|\chi|=-u$ is affine, so
$\dprimeddprime[\log|\chi|]=0$, in agreement with $\divis(\chi)=0$ on the
torus.  Compactify to $\mathbf P^{1,\An}$.  The skeleton becomes the closed
interval $[-\infty,+\infty]$ in the coordinate $u$, with $u=+\infty$ the
point $T=0$ and $u=-\infty$ the point $T=\infty$.  The function $-u$ is still
affine on the open interval, but the two endpoints are now points of the
space, and the current acquires the two atoms
\[
 \dprimeddprime[\log|\chi|]=\delta_{0}-\delta_{\infty} ,
\]
which is \eqref{eq:p1-poincare-lelong}.  Nothing about the function on the
skeleton has changed; what changed is the space over which it is a current.
This is the tropical shadow of the fact that a character of a torus becomes a
rational function with a zero and a pole on any equivariant
compactification, and it is also the reason that toric potentials are always
described by a recession function \eqref{eq:appJ-toric-recession}: the
recession data is exactly the slope information at the two ends, hence
exactly the boundary divisor.
\end{example}

\subsection{The unit tests, applied}
\label{sec:ch15-unit-tests}

We close the section by running \Cref{thm:cld-poincare-lelong} through the
first three tests of \Cref{sec:appJ-unit-tests}, since these are the tests
that a mistaken normalization fails.

\emph{Constant test.}  Replacing $\log|f|$ by $\log|f|+c$ changes neither
side; see \Cref{ex:ch15-constant-shift}.  A formula in which the divisor
depended on an additive constant would be detected at once.

\emph{Principal-divisor test.}  The formula distinguishes $\log|f|$ from
$v(f)=-\log|f|$, by \Cref{cor:ch15-appJ-consistency}.  If a later formula
mixes a metric weight, which by \eqref{eq:appJ-na-weight} is a
$-\log$ of a norm, with a Poincar\'e--Lelong term written in $\log|f|$, the
two signs must be reconciled; the safe route is
\Cref{cor:ch15-pl-section}, in which every term is written with the weight
convention.

\emph{Special-fibre test.}  Over a discretely valued $K$ with $v_K(\pi)=1$,
the model function of the special fibre is the constant $+1$ by
\eqref{eq:appJ-special-fibre-test}, and $\dprimeddprime$ of a constant is
zero, consistently with $\divis(\pi)=0$ on the generic fibre.  A convention
in which the special fibre received $-1$ would still pass this particular
test, since it too would give a constant; the special-fibre test therefore
discriminates metric conventions, not the Poincar\'e--Lelong sign, and the
principal-divisor test must be run as well.  This is the reason
\Cref{sec:appJ-unit-tests} lists them separately and insists that they be run
in order.

\section{The Bochner--Martinelli formula}
\label{sec:ch15-bm}

Poincar\'e--Lelong is the codimension-one case of a higher-codimension
phenomenon.  Its exact CLD form requires regular sequences.

\subsection{The setting, and the currents that make sense}
\label{sec:ch15-bm-setting}

Fix a boundaryless pure $n$-dimensional analytic space $X$ and a
\emph{regular sequence} $f_1,\dots,f_d$ of holomorphic functions on $X$: for
each $i$, the germ of $f_i$ is a non-zero-divisor in
$\cO_X/(f_1,\dots,f_{i-1})$.  Let
\[
 Y=V(f_1,\dots,f_d),\qquad U=X\setminus Y,
 \qquad u=\log\max\{|f_1|,\dots,|f_d|\} .
\]
The regularity hypothesis has two consequences that are used constantly and
that fail without it.

\begin{lemma}
\label{lem:ch15-bm-regular}
With the notation above:
\begin{enumerate}
\item\label{itm:ch15-bm-codim} $Y$ is purely of codimension $d$ in $X$, and
  $d_K(x)\le n-d$ for every $x\in Y$;
\item\label{itm:ch15-bm-Sd} the locus where $X$ satisfies Serre's condition
  $(S_d)$ is a Zariski-open subset containing $Y$;
\item\label{itm:ch15-bm-usc} $u$ is a continuous psh function on $U$ with
  values in $\mathbf R$, and $u(x)\to-\infty$ as $x\to Y$; more precisely, if
  $u$ is extended by $-\infty$ on $Y$ then it is upper semicontinuous on $X$.
\end{enumerate}
\end{lemma}

\begin{proof}
\cref{itm:ch15-bm-codim}
That the zero locus of a regular sequence of length $d$ has pure codimension
$d$ is the analytic Hauptidealsatz applied $d$ times: each successive
$f_i$ is a non-zero-divisor on the previously cut-out space, so it drops the
dimension by exactly one and no component is lost.  See
\cite[Lemma~11.5.2]{CLD25} for the analytic statement.  A point of a
Zariski-closed subset of dimension $\le n-d$ satisfies $d_K(x)\le n-d$ by
\Cref{prop:analytic-tropical-dimension}.

\cref{itm:ch15-bm-Sd}
The existence of a regular sequence of length $d$ in the maximal ideal of
$\cO_{X,x}$ means $\operatorname{depth}\cO_{X,x}\ge d$, which is $(S_d)$ at
$x$; and the $(S_d)$ locus is Zariski-open, which is the analytic form of
\cite[Theorem~4.4]{Duc21} quoted as such in \cite[\S18.7.3]{CLD25}.

\cref{itm:ch15-bm-usc}
On $U$ at least one $|f_i|$ is nonzero, so $u$ is finite; it is continuous
because each $|f_i|$ is.  Each $\log|f_i|$ is psh --- indeed smooth where
$f_i$ is invertible, and locally approachable in general by
\Cref{prop:ch15-log-approachable} below --- and a finite maximum of psh
functions is psh by \Cref{prop:ch15-max-stability}; hence $u$ is psh on $U$.
Upper semicontinuity of the extension is the continuity of
$x\mapsto\max_i|f_i(x)|$ together with the fact that $\log$ is increasing and
continuous on $[0,\infty)$ with value $-\infty$ at $0$.
\end{proof}

Because $u$ is continuous and psh on $U$, the Bedford--Taylor theory of
\Cref{chap:cld-bedford-taylor} defines the currents
\[
 (\dprimeddprime u)^k
 \qquad\text{and}\qquad
 u\,(\dprimeddprime u)^k
 \qquad\text{on }U,
\]
of bidegree $(k,k)$ and $(k,k)$ respectively, for every $k\ge0$.  The
question is how far they extend across $Y$.

\begin{lemma}
\label{lem:ch15-bm-extension}
Let $T$ be a current of bidegree $(r,r)$ on $U$, with $r<d$.  If
$d_K(x)\le n-d$ for every $x\in Y$, then $T$ extends by zero to a unique
current of bidegree $(r,r)$ on $X$; and if $r+1<d$, the extensions satisfy
\[
 \dprimeddprime\bigl(u(\dprimeddprime u)^r\bigr)=(\dprimeddprime u)^{r+1}
\]
on $X$.
\end{lemma}

\begin{proof}
A test form for a current of bidegree $(r,r)$ has bidegree $(n-r,n-r)$ with
$n-r>n-d$.  By \cite[Proposition~15.2.1]{CLD25} such a form vanishes as a
germ at every point whose local tropical dimension is $<n-r$; since every
point of $Y$ has $d_K(x)\le n-d<n-r$, every test form vanishes on a
neighbourhood of $Y$ and is therefore supported in $U$.  Consequently the
formula $\langle T,\omega\rangle$ makes sense for every test form on $X$, and
it is the only possible extension: two extensions differ by a current
supported on $Y$, which pairs to zero with every test form of the relevant
bidegree.  This is \cite[Lemma~7.1.13]{CLD25}.

For the last assertion, the identity holds on $U$ by the Bedford--Taylor
calculus of \Cref{chap:cld-bedford-taylor}.  Both sides extend by zero,
because both have bidegree $\le(r+1,r+1)$ with $r+1<d$, so the identity
persists.  For $r+1=d$ the right-hand side has bidegree $(d,d)$ and the
extension-by-zero argument is unavailable; this is exactly the borderline
that \Cref{thm:cld-local-bochner-martinelli} occupies.
\end{proof}

\begin{warning}{Why $(\dprimeddprime u)^d$ is not written}
The theorem below computes
$\dprimeddprime\bigl(u(\dprimeddprime u)^{d-1}\bigr)$, where the inner
current is the extension by zero supplied by \Cref{lem:ch15-bm-extension} for
$r=d-1<d$.  One would like to call this $(\dprimeddprime u)^d$, since
$(\dprimeddprime u)^{d-1}$ is closed.  That abbreviation is illegitimate:
$u=-\infty$ on $Y$, so the Bedford--Taylor product of $d$ factors is not
defined at points of $Y$, and no extension-by-zero argument covers bidegree
$(d,d)$.  The theorem is precisely the statement that the well-defined
left-hand side equals $\delta_Y$, which is what one would have wanted the
undefined right-hand side to be.
\end{warning}

\subsection{Statement and proof}

\begin{theorem}\label{thm:cld-local-bochner-martinelli}
Let $f_1,\ldots,f_d$ be a regular sequence of holomorphic functions on $X$,
let $Y=V(f_1,\ldots,f_d)$, and put on $X\setminus Y$
\[
                      u=\log\max_i|f_i|.
\]
For $0\leq k<d$, let the current $u(\dprimeddprime u)^k$ denote the
canonical extension constructed in \cite[\S18.7.1]{CLD25}.  Then
\[
 \boxed{\quad
 \dprimeddprime\!\left(u(\dprimeddprime u)^{d-1}\right)=\delta_Y.
 \quad}
\]
\end{theorem}
\begin{proof}
For $r>0$, truncate
\[
                   u_r=\log\max\{r,|f_1|,\ldots,|f_d|\}.
\]
This is a continuous PL psh function and its $d$-fold Monge--Amp\`ere current
is the integration current over the fibre above the Gauss point of radius
$r$.  Against a fixed test form, $u_r=u$ away from $Y$ for small $r$, so
integration by parts identifies the desired current with
$(\dprimeddprime u_r)^d$.  The calibrated fibre integral stabilizes as
$r\downarrow0$ and equals integration over the regular complete intersection
$Y$; this is \cite[Proposition~18.7.3]{CLD25}.  Regularity supplies purity of
codimension and the correct intersection multiplicity.

In detail, let $\omega\in A_{X,c}^{n-d,n-d}$ be a test form.  By definition
of the differential of a current and by
\Cref{prop:current-differential-calculus},
\[
 \bigl\langle\dprimeddprime\bigl(u(\dprimeddprime u)^{d-1}\bigr),
   \omega\bigr\rangle
 =\bigl\langle u(\dprimeddprime u)^{d-1},\dprimeddprime\omega\bigr\rangle
 =\bigl\langle u(\dprimeddprime u)^{d-1},
      (\dprimeddprime\omega)|_U\bigr\rangle ,
\]
the second equality because the current is the extension by zero and
$(\dprimeddprime\omega)|_U$ has compact support in $U$.  That support is
compact and disjoint from $Y$, so $\max_i|f_i|$ is bounded below there by
some $r_0>0$; for $0<r\le r_0$ one therefore has $u_r=u$ on it.  Hence
\[
 \begin{aligned}
  &\bigl\langle\dprimeddprime
     \bigl(u(\dprimeddprime u)^{d-1}\bigr),\omega\bigr\rangle\\
  &\qquad=\bigl\langle u_r(\dprimeddprime u_r)^{d-1},
      (\dprimeddprime\omega)|_U\bigr\rangle\\
  &\qquad=\bigl\langle\dprimeddprime
     \bigl(u_r(\dprimeddprime u_r)^{d-1}\bigr),\omega\bigr\rangle ,
 \end{aligned}
\]
running the same two identities backwards, now with the globally continuous
$u_r$.  Since $u_r$ is continuous and psh on all of $X$, the
Bedford--Taylor product of \Cref{chap:cld-bedford-taylor} applies without
restriction and
$\dprimeddprime\bigl(u_r(\dprimeddprime u_r)^{d-1}\bigr)
=(\dprimeddprime u_r)^d$.  By \Cref{prop:ch15-bm-fibre} below,
$(\dprimeddprime u_r)^d=\delta_{f^{-1}(\eta_r)}$, where
$f=(f_1,\dots,f_d):X\to\mathbf A^d$ and $\eta_r$ is the Gauss point of the
polydisc of polyradius $(r,\dots,r)$; and by
\Cref{prop:ch15-bm-fibre-stabilizes}, for $r$ small enough
$\langle\delta_{f^{-1}(\eta_r)},\omega\rangle=\langle\delta_Y,\omega\rangle$.
Combining, the pairing equals $\langle\delta_Y,\omega\rangle$ for every
$\omega$, which is the assertion.
\end{proof}

The two inputs used at the end are recorded separately, because each is a
theorem in its own right and each has its own hypotheses.

\begin{proposition}
\label{prop:ch15-bm-fibre}
Let $f=(f_1,\dots,f_d):X\to\mathbf A^d$ be a morphism given by $d$
holomorphic functions, let $r>0$, and put
$u_r=\log\max\{r,|f_1|,\dots,|f_d|\}$.  Then $u_r$ is a continuous
piecewise linear psh function on $X$ and
\[
 (\dprimeddprime u_r)^d=\delta_{f^{-1}(\eta_r)} ,
\]
where $\eta_r\in\mathbf A^{d,\An}$ is the Gauss point of the closed polydisc
of polyradius $(r,\dots,r)$.
\end{proposition}

\begin{proof}
That $u_r$ is continuous and PL is clear, and it is psh because it is a
finite maximum of the psh functions $\log r$ and $\log|f_i|$
(\Cref{prop:ch15-max-stability}), each of which is continuous after
truncation.  The identity is a pullback statement.  On $\mathbf A^{d,\An}$
with coordinates $T_1,\dots,T_d$, the function
$w_r=\log\max\{r,|T_1|,\dots,|T_d|\}$ is the weight of a formal metric, and
its $d$-fold Monge--Amp\`ere measure is the Dirac mass at $\eta_r$: this is
the model computation for $\mathbf A^d$, generalizing
\eqref{eq:p1-standard-measure}, and it is
\cite[Corollary~18.6.2]{CLD25}.  Pulling back along $f$ and using
compatibility of Bedford--Taylor products with pullback for continuous psh
potentials \Cref{prop:bt-current-stability} converts the Dirac mass into the
integration current of the fibre $f^{-1}(\eta_r)$, with its analytic
multiplicities.  The total mass bookkeeping is the calibration of
\Cref{def:trop-canonical-calibration}, which is why the fibre appears with
its scheme structure rather than reduced.
\end{proof}

\begin{proposition}
\label{prop:ch15-bm-fibre-stabilizes}
Let $X$ be an $n$-dimensional analytic space and let
$f=(f_1,\dots,f_d):X\to\mathbf A^d$ be purely of relative dimension $n-d$,
with $(f_1,\dots,f_d)$ a regular sequence and $Y=f^{-1}(0)$.  Let $\omega$ be
a compactly supported $G$-form of bidegree $(n-d,n-d)$ with locally bounded
tropically Borel coefficients.  Then, for $t\in\mathbf A^{d,\An}$ close
enough to the origin,
\[
 \langle\delta_{f^{-1}(t)},\omega\rangle=\langle\delta_Y,\omega\rangle .
\]
\end{proposition}

\begin{proof}
By \Cref{lem:ch15-bm-regular}\,\cref{itm:ch15-bm-Sd} the $(S_d)$ locus is
Zariski-open and contains $Y$, so after shrinking $X$ we may assume $X$ is
$(S_d)$.  The fibre $Y$ is purely of dimension $n-d$ by
\Cref{lem:ch15-bm-regular}\,\cref{itm:ch15-bm-codim}; upper semicontinuity of
fibre dimension and the Hauptidealsatz then make $f$ purely of relative
dimension $n-d$ after shrinking to a neighbourhood of $Y$.

Cover $\Supp\omega$ by finitely many compact analytic domains tropicalizing
$\omega$ and argue by additivity, so that we may assume $X$ compact and
$\omega=g_{\mathrm{trop}}^*\alpha$ tropical for a moment map
$g:X\to\mathbf G_m^N$.  Put $h=(g,f)$.  The local structure theorem for
tropicalizations, \Cref{thm:trop-local-germ}, supplies $r>0$ such that for
every $t$ in the polydisc of polyradius $r$ the calibrated polyhedra
$g_{\mathrm{trop}}(X_t)^{(n-d)}$ coincide, where $X_t=f^{-1}(t)$.  Since the
pairing $\langle\delta_{X_t},\omega\rangle$ is by definition the integral of
$\alpha$ over that calibrated polyhedron, it is independent of $t$, and
taking $t=0$ gives the claim.  This is
\cite[Proposition~18.7.3]{CLD25}.
\end{proof}

This is Mihatsch's formula
\cite[Theorem~18.7.2]{CLD25}.  The left side should not be abbreviated as
$(\dprimeddprime u)^d$ on all of $X$: $u=-\infty$ on $Y$, so the ordinary
continuous-potential Bedford--Taylor product is not defined there.  The
canonical extensions are part of the theorem's construction; their
uniqueness does not follow merely from the codimension of $Y$.

\begin{example}[$d=1$ recovers Poincar\'e--Lelong]
\label{ex:ch15-bm-d1}
For $d=1$ the current $u(\dprimeddprime u)^0$ is $[\log|f_1|]$, the extension
of \Cref{prop:ch15-log-current}, and $Y=\divis(f_1)$ as a cycle, because a
single non-zero-divisor generates a Cartier divisor whose associated cycle is
$[V(f_1)]$ by \Cref{thm:current-fundamental-cycle}.  The formula becomes
$\dprimeddprime[\log|f_1|]=\delta_{\divis(f_1)}$, that is,
\Cref{thm:cld-poincare-lelong}.  Notice that the two extension mechanisms
agree here: extension by local integrability and extension by zero in
bidegree $(0,0)$ both give the same current, because a $(0,0)$-current
supported on a set of local tropical dimension $<n$ vanishes.
\end{example}

\begin{example}[$d=n$: a finite fibre on the affine plane]
\label{ex:ch15-bm-dn}
Let $K$ be algebraically closed and nontrivially valued and let
$X=\mathbf A^{2,\An}_K$ with coordinates $T_1,T_2$; take $f_1=T_1$,
$f_2=T_2$, which is a regular sequence with $Y=\{(0,0)\}$, a single rigid
point, and $d=n=2$.  Then
$u_r=\log\max\{r,|T_1|,|T_2|\}$ and, by \Cref{prop:ch15-bm-fibre},
$(\dprimeddprime u_r)^2=\delta_{\eta_r}$, the Dirac mass at the Gauss point
of the polydisc of polyradius $(r,r)$: a single point of mass one.  As
$r\downarrow0$ these Dirac masses converge weakly to
$\delta_{(0,0)}=\delta_Y$, since $\eta_r\to(0,0)$ in the Berkovich topology.
The theorem asserts the corresponding statement for the extended currents,
and the limit is visible: all the mass is at one point for every $r$, and
that point moves to the origin.

It is worth contrasting this with the naive guess.  The function
$u=\log\max\{|T_1|,|T_2|\}$ is $-\infty$ only at the origin, and one might
hope to define $(\dprimeddprime u)^2$ as a limit of $(\dprimeddprime u_r)^2$.
That limit exists here and is $\delta_Y$, but the definition would be
circular: it is precisely what the theorem proves, and there is no
independent Bedford--Taylor product to compare it with, because $u$ is not
locally bounded.
\end{example}

\begin{example}[Regularity is not automatic]
\label{ex:ch15-bm-notregular}
On $X=\mathbf A^{2,\An}$, the pair $(T_1,T_1T_2)$ has the same zero locus as
$(T_1,T_2)$ set-theoretically only along $T_1=0$; in fact
$V(T_1,T_1T_2)=V(T_1)$ has codimension one, not two, so $(T_1,T_1T_2)$ is not
a regular sequence: $T_1T_2$ is a zero divisor modulo $(T_1)$.  Here
$u=\log\max\{|T_1|,|T_1T_2|\}=\log|T_1|+\log\max\{1,|T_2|\}$, and
$\dprimeddprime u$ is not the current of a codimension-two cycle at all.  The
hypothesis in \Cref{thm:cld-local-bochner-martinelli} is therefore not
cosmetic: it is what guarantees both the codimension in
\Cref{lem:ch15-bm-regular}\,\cref{itm:ch15-bm-codim} and the applicability of
\Cref{lem:ch15-bm-extension}.
\end{example}

\section{The Levine formula}
\label{sec:ch15-levine}

For completeness we also record the projective Levine formula.  On
$X=\mathbf P^n$, let $1\leq p\leq n$, let
$Y=\{x_0=\cdots=x_{p-1}=0\}\simeq\mathbf P^{n-p}$, and let
$\alpha$ be the first Chern current of the standard formal psh metric on
$\mathcal O(1)$.  On $U=X\setminus Y$, let $\beta$ be the pullback of the
corresponding current on $\mathbf P^{p-1}$ and put
\[
 h=\log\frac{\max_{0\leq i\leq n}|x_i|}
                   {\max_{0\leq i<p}|x_i|},
 \qquad
 \Omega=\sum_{j=0}^{p-1}\alpha^{p-1-j}\wedge\beta^j.
\]

Before the statement, three remarks on the objects involved.

The current $\alpha=c_1(\cO(1))$ for the standard formal metric is
\emph{parlinear} and positive of bidegree $(1,1)$; it is the higher-dimensional
analogue of \eqref{eq:p1-standard-measure} and is computed in
\Cref{chap:cld-metrized-bundles}.  The rational map
$f:[x_0:\cdots:x_n]\mapsto[x_0:\cdots:x_{p-1}]$ is defined exactly on $U$,
and $\beta=c_1(f^*\cO(1))$ there.  Both are curvature currents of continuous
psh metrics, so the Bedford--Taylor theory of
\Cref{chap:cld-bedford-taylor} defines all the products $\alpha^i\beta^j$ on
$U$, and they are parlinear.

The function $h$ is the logarithm of the ratio of the two metrics on the line
bundle $\cO(1)|_U\simeq f^*\cO(1)$, so it is continuous and parlinear on $U$;
by \eqref{eq:appJ-na-weight} it is the relative weight of the two metrics.
It is bounded below by $0$ and tends to $+\infty$ towards $Y$.

Finally, $\Omega$ is a current of bidimension $(n-p+1,n-p+1)$ on $U$, and
$h\Omega$ is defined there because $h$ is continuous.  The content of the
theorem is that $h\Omega$ extends, and what its $\dprimeddprime$ is.

\begin{theorem}\label{thm:cld-levine}
The current $h\Omega$ on $U$ extends canonically to a current $\Lambda$ on
$X$, and
\[
                       \dprimeddprime\Lambda+\delta_Y=\alpha^p.
\]
\end{theorem}
\begin{proof}
Let $p:\widetilde X=\operatorname{Bl}_Y X\to X$ be the blowup and let $E$ be
its exceptional divisor.  The rational map
$[x_0:\cdots:x_{p-1}]$ extends to a morphism
$q:\widetilde X\to\mathbf P^{p-1}$.  Write
$\widetilde\alpha=p^*\alpha$ and let
$\widetilde\beta=q^*\beta_{\mathbf P^{p-1}}$.  The quotient of the two
tautological model norms is represented by a regular meromorphic section
whose divisor is $E$.  Poincar\'e--Lelong, with the sign convention at the
beginning of the chapter, gives
\begin{equation}\label{eq:levine-blowup-ddc}
 \dprimeddprime[\widetilde h]
       =\widetilde\alpha-\widetilde\beta-\delta_E,
 \qquad \widetilde h=p^*h\quad\text{on }\widetilde X\setminus E.
\end{equation}
The locally integrable extension in
\eqref{eq:levine-blowup-ddc} defines the canonical extension of
$\widetilde h\widetilde\Omega$, where
\[
 \widetilde\Omega=
 \sum_{j=0}^{p-1}\widetilde\alpha^{p-1-j}\wedge
                         \widetilde\beta^j.
\]

Both curvature forms are closed.  Multiplying
\eqref{eq:levine-blowup-ddc} by $\widetilde\Omega$ and using the finite
geometric-series identity gives
\[
 \dprimeddprime[\widetilde h\widetilde\Omega]
 =\widetilde\alpha^p-\widetilde\beta^p
       -\delta_E\wedge\widetilde\Omega.
\]
Since $q$ has target dimension $p-1$, one has
$\widetilde\beta^p=0$.  Push the identity forward by $p$.  The projection
formula gives $p_*\widetilde\alpha^p=\alpha^p$.  On each fibre
$E_y\simeq\mathbf P^{p-1}$, only the term
$\widetilde\beta^{p-1}$ has vertical top degree, and its integral is one;
the terms containing $\widetilde\alpha$ have insufficient vertical degree.
Consequently
\[
                  p_*(\delta_E\wedge\widetilde\Omega)=\delta_Y.
\]
Finally put
$\Lambda=p_*[\widetilde h\widetilde\Omega]$.  It agrees with $h\Omega$ on
$X\setminus Y$, is canonical because the blowup and tautological metrics are
canonical, and satisfies
\(
                  \dprimeddprime\Lambda+\delta_Y=\alpha^p.
\)
This is the proof of \cite[Theorem~18.7.5]{CLD25}.

We expand three of the steps.

\emph{The Poincar\'e--Lelong input \eqref{eq:levine-blowup-ddc}.}  On
$\widetilde X$ the two line bundles $L_a=p^*\cO(1)$ and $L_b=q^*\cO(1)$ carry
parlinear psh metrics with curvature currents $\widetilde\alpha$ and
$\widetilde\beta$.  Over $U'=\widetilde X\setminus E$ they are canonically
isomorphic, and this isomorphism extends to a morphism
$\rho:L_b\to L_a$ of line bundles on $\widetilde X$ whose zero locus is
exactly $E$; equivalently, $\rho$ is a regular meromorphic section of
$L_b^{-1}\otimes L_a$ with $\divis(\rho)=E$.  The direction matters: on the
blowup one has $q^*\cO(1)=p^*\cO(1)(-E)$, that is
$L_b=L_a(-E)$, so it is $L_b^{-1}\otimes L_a$ that is $\cO(E)$ and carries a
section with divisor $E$, while $L_a^{-1}\otimes L_b=\cO(-E)$ does not.  Its
logarithmic weight is $\log\|\rho\|^{-1}=\widetilde h$.
\Cref{cor:ch15-pl-section} applied to $\rho$ gives
$\dprimeddprime[\widetilde h]+\delta_E=c_1(L_b^{-1}\otimes L_a)
=\widetilde\alpha-\widetilde\beta$, which is
\eqref{eq:levine-blowup-ddc} after moving terms; the sign is the one fixed in
\Cref{sec:ch15-pl-constant} and it is the only place in the proof where a
sign can go wrong.  The case $p=1$ is the check: there
$\widetilde\beta=0$ and \Cref{ex:ch15-levine-p1} reads
$\dprimeddprime[h]+\delta_Y=\alpha$, which the displayed orientation
reproduces and the opposite one does not.

\emph{The geometric-series identity.}  Writing $A=\widetilde\alpha$ and
$B=\widetilde\beta$ for brevity, and using that both are closed so that
$\dprimeddprime$ passes through the product,
\[
 (A-B)\wedge\Bigl(\sum_{j=0}^{p-1}A^{p-1-j}B^{j}\Bigr)
 =\sum_{j=0}^{p-1}A^{p-j}B^j-\sum_{j=0}^{p-1}A^{p-1-j}B^{j+1}
 =A^p-B^p ,
\]
the sum telescoping.  The multiplication of \eqref{eq:levine-blowup-ddc} by
$\widetilde\Omega$ is legitimate because $\widetilde\Omega$ is a
Bedford--Taylor product of curvature currents of continuous psh metrics, and
the relevant product rule is \cite[Proposition~18.3.9]{CLD25}; it is not the
multiplication of two arbitrary currents.

\emph{The vertical degree count.}  Since $q$ factors through
$\mathbf P^{p-1}$, which has dimension $p-1$, the pullback $B=q^*(\cdot)$
satisfies $B^p=0$ by \Cref{prop:form-restriction-dimension-gap}.  On the
exceptional divisor $E$, which is a $\mathbf P^{p-1}$-bundle over
$Y\simeq\mathbf P^{n-p}$, the restriction $B|_E$ is the curvature of the
relative $\cO(1)$ and $A|_E$ is pulled back from $Y$; a term
$A^{p-1-j}B^j$ with $j<p-1$ has vertical degree $j<p-1$ and therefore
integrates to zero along the fibre, while $j=p-1$ contributes $\int_{E_y}
B^{p-1}=1$, the degree of $\cO(1)$ on $\mathbf P^{p-1}$ computed with
\eqref{eq:appJ-cl-total-mass}.  Fibre integration
\Cref{thm:cld-fibre-integration} then yields
$p_*(\delta_E\wedge\widetilde\Omega)=\delta_Y$.
\end{proof}

\begin{example}[$p=1$ recovers Poincar\'e--Lelong again]
\label{ex:ch15-levine-p1}
For $p=1$ one has $Y=\{x_0=0\}$, a hyperplane; $\beta$ is the curvature of a
metric on the pullback of $\cO(1)$ from $\mathbf P^0$, hence zero; and
$\Omega=\alpha^0=1$.  Thus $h=\log\bigl(\max_i|x_i|/|x_0|\bigr)$ and the
formula reads $\dprimeddprime[h]+\delta_Y=\alpha$.  This is exactly
\Cref{cor:ch15-pl-section} for the section $x_0$ of $\cO(1)$ with the
standard metric: indeed $-\log\|x_0\|=h$ by the definition of that metric,
generalizing \eqref{eq:p1-standard-weight}.
\end{example}

\begin{example}[$p=n$ and the top power]
\label{ex:ch15-levine-pn}
For $p=n$, $Y$ is a single rational point of $\mathbf P^n$ and $\alpha^n$ is
the Chambert--Loir measure of the standard formal metric, a Dirac mass at the
Gauss point of the standard torus with total mass $1=(\cO(1)^n)$ by
\eqref{eq:appJ-cl-total-mass}.  The formula
$\dprimeddprime\Lambda+\delta_Y=\alpha^n$ says that these two point masses,
one at a rigid point and one at the Gauss point, differ by a
$\dprimeddprime$-exact current, which is the non-Archimedean analogue of the
classical statement that a point and the Fubini--Study volume are homologous.
Integrating both sides against the constant $1$ on the compact boundaryless
$\mathbf P^{n,\An}$ gives $0+1=1$, so the identity passes the total-mass test
of \Cref{sec:appJ-unit-tests}.
\end{example}

\section{The CLD approximation definition of psh}
\label{sec:ch15-psh}

In complex geometry the plurisubharmonic functions on a domain admit several
equivalent descriptions: by the sign of $\rmi\partial\bar\partial u$ in the
sense of currents, by the sub-mean-value property on complex lines, and by
decreasing approximation by smooth psh functions.  The last of these is a
theorem --- on a Stein space it is the smoothing theorem of Fornaess and
Narasimhan \cite{FN80} --- and it is the one that CLD elevate to a
definition.  The reason is structural.  On a Berkovich space there is no
convolution, hence no way to produce a smooth approximant of a given
function; there is a notion of smooth function and a notion of positive
$(1,1)$-form, hence a notion of \emph{smooth} psh function, but nothing
forces the resulting class to be closed under the limits one needs.  Defining
the psh class as the closure therefore guarantees, by construction, the one
property that every later argument uses: that psh functions can be
approximated from above.

The price is that the class is defined by generation and not by a pointwise
test.  A section of the resulting sheaf need not be a single decreasing limit
of smooth psh functions; it is only guaranteed to be reachable from the
smooth ones by iterated decreasing limits.  How many iterations are needed is
the subject of \Cref{sec:ch15-hierarchy}.

Let $S$ be an affine $G$-tropical space, in particular an analytic space with
its CLD structure.

\begin{definition}[CLD-psh]
The sheaf $\PSH_{\mathrm{CLD}}$ is the smallest sheaf of functions with
values in $\mathbf R\cup\{-\infty\}$ that
\begin{enumerate}
\item contains every smooth $u$ with $\dprimeddprime u\geq0$;
\item is closed under decreasing limits of sequences.
\end{enumerate}
A continuous CLD-psh function will be abbreviated \emph{cpsh}.
\end{definition}

We write $\operatorname{Psh}_S$ for $\PSH_{\mathrm{CLD}}$ when the ambient space is to be
displayed, and $\operatorname{CPsh}_S$ for the analogous sheaf generated by smooth psh
functions under local \emph{uniform} limits.  Three points of the definition
deserve emphasis before anything is proved.

\emph{The limits are of sequences, not of nets.}  A Berkovich space need not
be first countable, and a decreasing net of psh functions is not
automatically the limit of a decreasing sequence.  Every statement below is
therefore about sequences, and where a net occurs --- as in the
Gubler--Rabinoff class of \Cref{sec:ch15-classical} --- we say so.

\emph{The value $-\infty$ is allowed, but $+\infty$ is not.}  The sheaf takes
values in $\mathbf R\cup\{-\infty\}$.  A decreasing sequence of real-valued
functions has a limit in $\mathbf R\cup\{-\infty\}$, so the class is closed
under the operation used to define it; no truncation convention is needed.

\emph{Local integrability is not part of the definition.}  The constant
$-\infty$ is psh, being the decreasing limit of the smooth psh constants
$-k$, and it is not locally integrable.  The sheaf of psh functions that are
also locally integrable is a separate object, and
\Cref{prop:ch15-L1-hierarchy} shows that it is compatible with the hierarchy
level by level.  This is why the positivity statement
\Cref{prop:cld-psh-positive-current} carries a local integrability
hypothesis.

\subsection{Elementary consequences of the definition}
\label{sec:ch15-psh-elementary}

Everything in this subsection is proved by the same device: to show that
every psh function has a property $P$, exhibit the subsheaf of functions with
property $P$, check that it contains the smooth psh functions and is stable
under decreasing sequential limits, and invoke minimality.  We write the
device out once, in full, and use it thereafter without comment.

\begin{lemma}
\label{lem:ch15-minimality}
Let $\mathscr F$ be a subsheaf of the sheaf of
$\mathbf R\cup\{-\infty\}$-valued functions on $S$ such that
\begin{enumerate}
\item $\mathscr F(U)$ contains every smooth $u$ on $U$ with
  $\dprimeddprime u\ge0$, for every open $U$;
\item if $(u_k)$ is a decreasing sequence in $\mathscr F(U)$ converging
  pointwise to $u$, then $u\in\mathscr F(U)$.
\end{enumerate}
Then $\operatorname{Psh}_S\subseteq\mathscr F$.
\end{lemma}

\begin{proof}
The intersection $\mathscr F\cap\operatorname{Psh}_S$ is again a subsheaf
satisfying (1) and (2), and it is contained in $\operatorname{Psh}_S$.  By
the minimality in the definition of $\operatorname{Psh}_S$ it equals
$\operatorname{Psh}_S$, which is the assertion.
\end{proof}

\begin{proposition}
\label{prop:ch15-psh-cone}
For every open $U\subseteq S$, the set $\operatorname{Psh}_S(U)$ consists of
upper semicontinuous functions, contains the constants, and is a convex cone:
if $u,v$ are psh on $U$ and $\lambda>0$, then $u+v$ and $\lambda u$ are
psh on $U$; the zero element is the constant zero function.  Thus no
product of $0$ with the value $-\infty$ is part of the assertion.
\end{proposition}

\begin{proof}
\emph{Upper semicontinuity and scaling.}  Let $\mathscr F(U)$ be the set of
upper semicontinuous $u$ on $U$ such that $\lambda u$ is psh for every
$\lambda>0$.  This is a subsheaf, because upper semicontinuity and
membership in a sheaf are both local.  It contains the smooth psh functions:
a smooth function is continuous, and $\dprimeddprime(\lambda u)
=\lambda\,\dprimeddprime u\ge0$.  It is stable under decreasing sequential
limits: a decreasing limit of upper semicontinuous functions is upper
semicontinuous, since $\{u<c\}=\bigcup_k\{u_k<c\}$ is open; and if
$u_k\downarrow u$ then $\lambda u_k\downarrow\lambda u$ for $\lambda>0$.
By \Cref{lem:ch15-minimality}, $\mathscr F=\operatorname{Psh}_S$.  Constants are
smooth with vanishing $\dprimeddprime$, hence psh.

\emph{Sums.}  We argue in two steps, because the naive one-step argument
fails: the set of $u$ such that $u+v$ is psh, for a \emph{fixed} psh $v$, is
not obviously stable under the operations, since we do not yet know that it
contains the smooth psh functions.

First fix a \emph{smooth} psh $w$ on $U$ and let $\mathscr F_w$ be the
subsheaf of functions $u$ with $u+w$ psh.  It contains the smooth psh
functions, because a sum of two smooth psh functions is smooth with
$\dprimeddprime(u+w)=\dprimeddprime u+\dprimeddprime w\ge0$; and it is stable
under decreasing limits, since $u_k\downarrow u$ implies
$u_k+w\downarrow u+w$ (no indeterminate form arises, $w$ being real-valued).
By \Cref{lem:ch15-minimality}, $u+w$ is psh for every psh $u$ and every
smooth psh $w$.

Now fix an arbitrary psh $v$ and let $\mathscr G_v$ be the subsheaf of
functions $u$ with $u+v$ psh.  By the previous paragraph $\mathscr G_v$
contains the smooth psh functions.  It is stable under decreasing limits by
the same argument as before, with the caveat that if $v=-\infty$ at a point
then $u_k+v=-\infty=u+v$ there, so the convergence still holds in
$\mathbf R\cup\{-\infty\}$.  By \Cref{lem:ch15-minimality} again,
$\mathscr G_v=\operatorname{Psh}_S$, that is, $u+v$ is psh for all psh
$u,v$.  This is \cite[Lemma~8.1.2]{CLD25}.
\end{proof}

\begin{remark}[Why the two-step argument is necessary]
The order of quantifiers in the proof of \Cref{prop:ch15-psh-cone} is not an
artefact.  If one tries to prove stability under sums in one step by fixing
$v$ psh and applying \Cref{lem:ch15-minimality} to $\mathscr G_v$, the
verification of hypothesis (1) --- that $w+v$ is psh for $w$ smooth psh ---
is precisely the statement being proved.  The first step breaks the circle by
putting the smooth function in the fixed slot.  The same two-step pattern
recurs in \Cref{prop:ch15-composition} and is worth recognizing.
\end{remark}

\section{The hierarchy and locally approachable functions}
\label{sec:ch15-hierarchy}

\subsection{The transfinite presentation}

\begin{definition}[The CLD hierarchy]
\label{def:ch15-hierarchy}
\index{plurisubharmonic function!CLD hierarchy}
For an ordinal $\alpha$ define a subsheaf
$\operatorname{Psh}_{S,\alpha}$ of the sheaf of
$\mathbf R\cup\{-\infty\}$-valued functions on $S$ by transfinite recursion:
\begin{enumerate}
\item $\operatorname{Psh}_{S,0}$ is the sheaf of smooth functions $u$ with
  $\dprimeddprime u\ge0$;
\item $\operatorname{Psh}_{S,\alpha+1}$ is the subsheaf of functions that are
  \emph{locally} pointwise decreasing limits of sequences in
  $\operatorname{Psh}_{S,\alpha}$;
\item for a limit ordinal $\alpha$, $\operatorname{Psh}_{S,\alpha}$ is the
  subsheaf of functions belonging locally to some
  $\operatorname{Psh}_{S,\beta}$ with $\beta<\alpha$.
\end{enumerate}
The sheaves $\operatorname{CPsh}_{S,\alpha}$ are defined identically with
``decreasing limit'' replaced by ``uniform limit'' and with values in
$\mathbf R$.
\end{definition}

\begin{proposition}
\label{prop:ch15-hierarchy-stabilizes}
The family $(\operatorname{Psh}_{S,\alpha})_\alpha$ is an increasing family
of subsheaves of $\operatorname{Psh}_S$, it is eventually constant, and its
eventual value is $\operatorname{Psh}_S$.  The same holds for
$(\operatorname{CPsh}_{S,\alpha})_\alpha$ with limit
$\operatorname{CPsh}_S$.
\end{proposition}

\begin{proof}
That each $\operatorname{Psh}_{S,\alpha}$ is contained in
$\operatorname{Psh}_S$ follows by transfinite induction directly from
\Cref{def:ch15-hierarchy} and the two defining properties of
$\operatorname{Psh}_S$: the base case is property (1) of the definition, the
successor case property (2) together with the sheaf property, and the limit
case the sheaf property alone.  Monotonicity in $\alpha$ is clear from the
construction, since a constant sequence is decreasing.

\emph{Eventual constancy.}  Fix an open $U$.  The sets
$\operatorname{Psh}_{S,\alpha}(U)$ form an increasing ordinal-indexed chain
of subsets of the fixed set $\mathcal G(U)$ of all functions
$U\to\mathbf R\cup\{-\infty\}$.  An increasing ordinal-indexed chain of
subsets of a fixed set is eventually constant: otherwise one could choose,
for every ordinal $\alpha$ below the cardinal successor of
$\#\mathcal G(U)$, an element of
$\operatorname{Psh}_{S,\alpha+1}(U)\setminus\operatorname{Psh}_{S,\alpha}(U)$,
producing an injection of that cardinal successor into $\mathcal G(U)$, which
is absurd.  Let $\alpha_U$ be a stabilization ordinal for $U$.  The opens of
$S$ form a set, so $\alpha_0=\sup_U\alpha_U$ is an ordinal, and the whole
family is constant from $\alpha_0$ on.

\emph{The eventual value.}  Let $\alpha\ge\alpha_0$.  We check that
$\operatorname{Psh}_{S,\alpha}$ satisfies the two conditions of
\Cref{lem:ch15-minimality}.  It contains $\operatorname{Psh}_{S,0}$, hence
every smooth psh function.  If $(u_k)$ is a decreasing sequence in
$\operatorname{Psh}_{S,\alpha}(U)$ with pointwise limit $u$, then
$u\in\operatorname{Psh}_{S,\alpha+1}(U)=\operatorname{Psh}_{S,\alpha}(U)$ by
the choice of $\alpha$.  Hence
$\operatorname{Psh}_S\subseteq\operatorname{Psh}_{S,\alpha}$, and the reverse
inclusion was proved above.  This is \cite[Lemma~8.1.4]{CLD25}; the argument
for $\operatorname{CPsh}$ is identical, replacing decreasing by uniform
limits, and is \cite[Lemma~8.1.8]{CLD25}.
\end{proof}

\begin{definition}[Locally approachable]
\label{def:ch15-approachable}
\index{plurisubharmonic function!locally approachable}
A psh function on an open subset of $S$ is \emph{locally approachable} if it
is a section of $\operatorname{Psh}_{S,1}$, that is, if every point has a
neighbourhood on which the function is the pointwise limit of a decreasing
sequence of smooth psh functions.
\end{definition}

This is CLD's \emph{localement approchable}
\cite[Definition~8.1.13]{CLD25}.  We keep their terminology, in the English
rendering \emph{locally approachable}, because it names exactly the property
that the theorems of \Cref{chap:cld-bedford-taylor} require and that the
general psh class is not known to have.

\begin{warning}{What is known and what is open}
\index{CLD hierarchy!early collapse}
Two distinct questions must be kept apart.

\emph{Eventual stabilization is a theorem.}
\Cref{prop:ch15-hierarchy-stabilizes} proves that the hierarchy is eventually
constant, by the cardinality argument of \cite[Lemma~8.1.4]{CLD25}, and CLD
themselves use that stabilization repeatedly, for instance in transporting a
statement proved for every $\alpha$ to a statement about
$\operatorname{Psh}_S$.  Read literally, the sentence in
\cite[Remark~8.1.6]{CLD25} saying that stationarity of the hierarchy is not
known contradicts that lemma; the remark's own second sentence identifies the
intended content, namely that every known example lies at the first stage.
We therefore use the corrected formulation throughout and never assert that
eventual stabilization is unknown.

\emph{Collapse at stage one is open in general.}  Positive cases include
real affine spaces, tropical toric partial compactifications, finite
ambient-induced real-affine stars, their products with open convex Euclidean
domains or with one convex polyhedron for jointly convex ambient product
generators, the two constant-tail interval models, and the exact CLD stars
$X_B$ of \cite[\S6.7.3]{CLD25} for arbitrary nonempty branch sets
(\Cref{thm:ch15-convexity-model,prop:ch16-collapse-tropical-toric,%
thm:ch16-collapse-finite-tropical-curves,%
thm:p071-finite-star-euclidean-collapse,%
thm:p071-finite-star-polyhedral-collapse,%
thm:p071-cld-constant-collar-endpoint,thm:p071-cld-infinite-collar-endpoint,%
thm:p071-infinite-star-collapse-candidate}).  Separately,
\Cref{thm:p063-current-psh-ordinary-regularization} settles continuous
current-psh functions on smooth proper algebraic curves in its stated field
scope.  The unrestricted question is
\Cref{prob:frontier-cld-hierarchy}.  In complex geometry the corresponding
statement is the smoothing theorem \cite{FN80}, cited by CLD as motivation
for their definition; the unrestricted non-Archimedean analogue has neither
a known proof nor a known counterexample.
\end{warning}

\begin{proposition}
\label{prop:ch15-collapse-reduction}
Let $\alpha$ be an ordinal.  The following are equivalent:
\begin{enumerate}
\item $\operatorname{Psh}_{S,\alpha}=\operatorname{Psh}_S$;
\item for every open $U$ and every decreasing sequence $(u_k)$ in
  $\operatorname{Psh}_{S,\alpha}(U)$ with pointwise limit $u$, one has
  $u\in\operatorname{Psh}_{S,\alpha}(U)$.
\end{enumerate}
Equivalently, $\operatorname{Psh}_{S,\alpha}=\operatorname{Psh}_{S,\alpha+1}$
if and only if $\operatorname{Psh}_{S,\alpha}=\operatorname{Psh}_S$.
\end{proposition}

\begin{proof}
(1)$\Rightarrow$(2) is property (2) of the definition of
$\operatorname{Psh}_S$.  For (2)$\Rightarrow$(1), the subsheaf
$\operatorname{Psh}_{S,\alpha}$ contains $\operatorname{Psh}_{S,0}$, hence
every smooth psh function, and by hypothesis it is stable under decreasing
sequential limits; \Cref{lem:ch15-minimality} gives
$\operatorname{Psh}_S\subseteq\operatorname{Psh}_{S,\alpha}$, and the reverse
inclusion is \Cref{prop:ch15-hierarchy-stabilizes}.

For the reformulation, condition (2) says exactly that
$\operatorname{Psh}_{S,\alpha+1}\subseteq\operatorname{Psh}_{S,\alpha}$ ---
membership in $\operatorname{Psh}_{S,\alpha+1}$ being a local condition and
$\operatorname{Psh}_{S,\alpha}$ a sheaf --- and the reverse inclusion always
holds.
\end{proof}

\Cref{prop:ch15-collapse-reduction} localizes the whole question.  The
transfinite apparatus is not where the difficulty lies: the entire content of
early collapse is the single question \emph{is a decreasing limit of locally
approachable functions locally approachable?}  Everything else in
\Cref{def:ch15-hierarchy} is bookkeeping.  The next subsections show that the
approachable class is closed under every other operation the theory performs,
which is why, as CLD observe, no known example escapes the first stage.

\subsection{The regularized maximum}
\label{sec:ch15-regmax}

The single most useful construction in the psh calculus is a smooth
substitute for $\max$.  We fix it once, with all the properties that will be
needed, because two of them --- monotonicity in the parameter and the
truncation property --- are not standard and are exactly what makes the
diagonal arguments below work.

\begin{lemma}
\label{lem:ch15-regmax}
\index{regularized maximum}
Let $\theta:\mathbf R\to[0,\infty)$ be smooth, even, supported in $[-1,1]$,
with $\int_{\mathbf R}\theta=1$.  For $\varepsilon>0$ and
$s=(s_1,\dots,s_m)\in\mathbf R^m$ put
\[
 M_\varepsilon(s)
 =\int_{\mathbf R^m}\max\{s_1+t_1,\dots,s_m+t_m\}
     \prod_{i=1}^m\varepsilon^{-1}\theta(t_i/\varepsilon)\,\rd t .
\]
Then:
\begin{enumerate}
\item\label{itm:ch15-regmax-smooth} $M_\varepsilon$ is smooth on
  $\mathbf R^m$;
\item\label{itm:ch15-regmax-convex} $M_\varepsilon$ is convex and
  nondecreasing in each variable, and
  $M_\varepsilon(s+c\mathbf 1)=M_\varepsilon(s)+c$ for $c\in\mathbf R$;
\item\label{itm:ch15-regmax-sandwich}
  $\max_is_i\le M_\varepsilon(s)\le\max_is_i+\varepsilon$;
\item\label{itm:ch15-regmax-trunc} if $s_j\le\max_is_i-2\varepsilon$, then
  $M_\varepsilon(s)$ does not depend on $s_j$;
\item\label{itm:ch15-regmax-mono} for $0<\varepsilon\le\varepsilon'$ one has
  $M_\varepsilon\le M_{\varepsilon'}$ pointwise.
\end{enumerate}
\end{lemma}

\begin{proof}
Write $\rho_\varepsilon(t)=\prod_i\varepsilon^{-1}\theta(t_i/\varepsilon)$, a
nonnegative smooth probability density supported in
$[-\varepsilon,\varepsilon]^m$ and invariant under $t\mapsto-t$, and note
$M_\varepsilon=\max*\rho_\varepsilon$.

\cref{itm:ch15-regmax-smooth}
Convolution of a continuous function of at most linear growth with a smooth
compactly supported kernel is smooth, by differentiation under the integral
sign.

\cref{itm:ch15-regmax-convex}
The function $\max$ is convex and nondecreasing in each variable; a positive
average of translates of such a function has the same two properties.  For
the translation law, $\max\{(s_i+c)+t_i\}=\max\{s_i+t_i\}+c$ and
$\int\rho_\varepsilon=1$.

\cref{itm:ch15-regmax-sandwich}
The kernel is even, so its barycentre is the origin; Jensen's inequality for
the convex function $\max$ gives
$M_\varepsilon(s)\ge\max\{s_i+0\}=\max_is_i$.  For the upper bound, on the
support of $\rho_\varepsilon$ one has $t_i\le\varepsilon$, so
$\max_i\{s_i+t_i\}\le\max_is_i+\varepsilon$; integrate.

\cref{itm:ch15-regmax-trunc}
Let $i_0$ realize $\max_is_i$ and suppose
$s_j\le s_{i_0}-2\varepsilon$.  For $t$ in the support of
$\rho_\varepsilon$,
\[
 s_j+t_j\le s_{i_0}-2\varepsilon+\varepsilon
 =s_{i_0}-\varepsilon\le s_{i_0}+t_{i_0}
 \le\max_i\{s_i+t_i\},
\]
so the $j$-th entry never attains the maximum, and lowering $s_j$ further
does not change the integrand.  Formally, $M_\varepsilon(s)$ agrees with
$M_\varepsilon(s')$ where $s'$ is obtained from $s$ by replacing $s_j$ by any
value $\le s_{i_0}-2\varepsilon$.

\cref{itm:ch15-regmax-mono}
Fix $s$ and consider $g(\lambda)=\int\max\{s_i+\lambda t_i\}\rho_1(t)\rd t$
for $\lambda\in\mathbf R$, so that $M_\varepsilon(s)=g(\varepsilon)$ after
the change of variables $t\mapsto t/\varepsilon$.  The function $g$ is convex,
being an average of the convex functions
$\lambda\mapsto\max_i\{s_i+\lambda t_i\}$ with the nonnegative weights
$\rho_1(t)\rd t$, and it is even, because $\rho_1(-t)=\rho_1(t)$.  A convex
even function on $\mathbf R$ is nondecreasing on $[0,\infty)$: for
$0\le a<b$ write $a=\frac{b-a}{2b}(-b)+\frac{b+a}{2b}b$, a convex
combination, and use convexity and evenness to get $g(a)\le g(b)$.  Hence
$M_\varepsilon(s)=g(\varepsilon)\le g(\varepsilon')=M_{\varepsilon'}(s)$.
\end{proof}

\begin{remark}[Two regularized maxima]
\Cref{lem:ch15-regmax} is the kernel-based construction used by CLD in the
proof of \cite[Proposition~16.9.1]{CLD25}, following
\cite[Lemma~5.18]{DemBook}.  A second, more familiar smooth substitute for
$\max$ is the log-sum-exp function
$L_\eta(s)=\eta\log\sum_i\rme^{s_i/\eta}$ of
\Cref{cor:superform-convex-composition}; it satisfies
\cref{itm:ch15-regmax-smooth}--\cref{itm:ch15-regmax-sandwich} with
$\varepsilon$ replaced by $\eta\log m$, and
\cref{itm:ch15-regmax-mono} as well, since
$L_\eta(s)=\log\|(\rme^{s_1},\dots,\rme^{s_m})\|_{\ell^{1/\eta}}$ and
$p\mapsto\|x\|_{\ell^p}$ is nonincreasing.  What log-sum-exp does \emph{not}
satisfy is the truncation property \cref{itm:ch15-regmax-trunc}: every
argument contributes, however small.  That property is indispensable in
\Cref{prop:ch15-log-approachable}, where one argument is $-\infty$ on a
Zariski-closed set and the approximant must nevertheless be smooth there.
Both constructions are used below, and the choice is dictated by whether
truncation is needed.
\end{remark}

\subsection{Composition and stability}

\begin{proposition}
\label{prop:ch15-composition}
Let $u_1,\dots,u_m$ be psh on $S$, let $\Omega\subseteq\mathbf R^m$ be open
and convex, and let $g:\Omega\to\mathbf R$ be smooth, convex, and
nondecreasing in each variable.  Assume $(u_1(x),\dots,u_m(x))\in\Omega$ for
every $x\in S$; in particular the $u_i$ are real-valued.  Then
$g(u_1,\dots,u_m)$ is psh.  More precisely, if $\alpha$ is an ordinal with
$u_i\in\operatorname{Psh}_{S,\alpha}$ for every $i$, then
$g(u_1,\dots,u_m)\in\operatorname{Psh}_{S,\alpha}$.
\end{proposition}

\begin{proof}
Transfinite induction on $\alpha$.

\emph{Base case $\alpha=0$.}  All $u_i$ are smooth, hence so is
$u=g(u_1,\dots,u_m)$, and the chain rule gives
\[
 \dprimeddprime u
 =\sum_{i,j}\frac{\partial^2g}{\partial y_i\partial y_j}(u)\,
     \dprime u_i\wedge\ddprime u_j
 +\sum_j\frac{\partial g}{\partial y_j}(u)\,\dprimeddprime u_j .
\]
The first sum is a strongly positive $(1,1)$-form.  Indeed, write the
positive semidefinite matrix $C=(\partial_i\partial_jg(u))$ --- positive
semidefinite by convexity of $g$ --- in spectral form
$C=\sum_k\lambda_ka^{(k)}(a^{(k)})^{\mathsf T}$ with $\lambda_k\ge0$ and
$a^{(k)}\in\mathbf R^m$.  Setting $w_k=\sum_ia^{(k)}_iu_i$, a smooth
function, one gets
\[
 \sum_{i,j}C_{ij}\,\dprime u_i\wedge\ddprime u_j
 =\sum_k\lambda_k\,\dprime w_k\wedge\ddprime w_k ,
\]
and each $\dprime w\wedge\ddprime w=\eta\wedge J\eta$ with $\eta=\dprime w$
is a generator of the strongly positive cone, by
\Cref{lem:ch09-generators} and the dictionary
\Cref{prop:ch09-positivity-dictionary}; a nonnegative combination of
generators stays in the cone.  This is the mechanism recorded after
\Cref{lem:superform-composition}.  In the second sum each coefficient
$\partial_jg(u)$ is $\ge0$ because $g$ is nondecreasing in each variable, and
each $\dprimeddprime u_j$ is positive by hypothesis; the positive cone is
stable under multiplication by nonnegative functions and under addition.
Hence $\dprimeddprime u\ge0$ and $u\in\operatorname{Psh}_{S,0}$.

\emph{Limit ordinals.}  If $\alpha$ is a limit ordinal, each point has a
neighbourhood $U$ and an ordinal $\beta<\alpha$ with
$u_i|_U\in\operatorname{Psh}_{S,\beta}(U)$ for all $i$; the inductive
hypothesis gives $g(u_1,\dots,u_m)|_U\in\operatorname{Psh}_{S,\beta}(U)
\subseteq\operatorname{Psh}_{S,\alpha}(U)$, and the sheaf property concludes.

\emph{Successor ordinals.}  Let $\alpha=\beta+1$.  Near a given point choose,
for each $i$, a decreasing sequence $(v_{i,k})_k$ in
$\operatorname{Psh}_{S,\beta}$ converging pointwise to $u_i$.  Since $g$ is
continuous and nondecreasing in each variable, the sequence
$\bigl(g(v_{1,k},\dots,v_{m,k})\bigr)_k$ is decreasing and converges
pointwise to $g(u_1,\dots,u_m)$.  By the inductive hypothesis its terms lie
in $\operatorname{Psh}_{S,\beta}$, whence
$g(u_1,\dots,u_m)\in\operatorname{Psh}_{S,\beta+1}$.  This is
\cite[Proposition~8.2.1]{CLD25}.
\end{proof}

\begin{proposition}
\label{prop:ch15-max-stability}
Let $u_1,\dots,u_m$ be psh on $S$ and let $\alpha$ be an ordinal with
$u_i\in\operatorname{Psh}_{S,\alpha}$ for all $i$.  Then
\[
 \max\{u_1,\dots,u_m\}\in\operatorname{Psh}_{S,\sup\{\alpha,1\}} .
\]
In particular a finite maximum of psh functions is psh, and a finite maximum
of locally approachable functions is locally approachable.
\end{proposition}

\begin{proof}
The statement for general $\alpha$ is \cite[Corollary~8.2.3]{CLD25}: the
function $\max$ on $\mathbf R^m$ is the uniform limit of the smooth convex
functions $L_\eta$ that are nondecreasing in each variable, so
\Cref{prop:ch15-composition} together with the uniform-limit refinement
\cite[Corollary~8.2.2]{CLD25} gives the assertion.  We write out the case
$\alpha=1$, which is the one used repeatedly below and which also covers
$-\infty$-valued arguments, by an explicit diagonal.

Let $u_1,\dots,u_m\in\operatorname{Psh}_{S,1}$.  The assertion is local, so
fix a point and a neighbourhood $U$ on which each $u_i$ is the pointwise
limit of a decreasing sequence $(u_{i,k})_k$ of \emph{smooth} psh functions
on $U$.  Choose $\varepsilon_k\downarrow0$, for instance
$\varepsilon_k=1/k$, and set
\[
 w_k=M_{\varepsilon_k}(u_{1,k},\dots,u_{m,k})
\]
with $M_\varepsilon$ as in \Cref{lem:ch15-regmax}.  Three verifications
finish the proof.

\emph{$w_k$ is smooth psh.}  Each $u_{i,k}$ is smooth and real-valued, and
$M_{\varepsilon_k}$ is smooth, convex and nondecreasing in each variable on
$\Omega=\mathbf R^m$ by
\Cref{lem:ch15-regmax}\,\cref{itm:ch15-regmax-smooth}--\cref{itm:ch15-regmax-convex};
\Cref{prop:ch15-composition} with $\alpha=0$ applies.

\emph{$(w_k)$ is decreasing.}  Using
\Cref{lem:ch15-regmax}\,\cref{itm:ch15-regmax-mono} and then monotonicity in
each argument together with $u_{i,k+1}\le u_{i,k}$,
\[
 \begin{aligned}
 w_{k+1}=M_{\varepsilon_{k+1}}(u_{1,k+1},\dots,u_{m,k+1})
 &\le M_{\varepsilon_k}(u_{1,k+1},\dots,u_{m,k+1})\\
 &\le M_{\varepsilon_k}(u_{1,k},\dots,u_{m,k})=w_k .
 \end{aligned}
\]

\emph{$w_k\to\max_iu_i$ pointwise.}  By
\Cref{lem:ch15-regmax}\,\cref{itm:ch15-regmax-sandwich},
\[
 \max_iu_{i,k}\le w_k\le\max_iu_{i,k}+\varepsilon_k .
\]
Since each $u_{i,k}\downarrow u_i$ and $\max$ is continuous and monotone on
$(\mathbf R\cup\{-\infty\})^m$, one has $\max_iu_{i,k}\downarrow\max_iu_i$
pointwise, including at points where some or all $u_i$ equal $-\infty$; and
$\varepsilon_k\to0$.  Hence $w_k\downarrow\max_iu_i$.

Thus $\max_iu_i$ is a local decreasing limit of smooth psh functions, that
is, a section of $\operatorname{Psh}_{S,1}$.
\end{proof}

\begin{theorem}
\label{thm:ch15-approachable-stability}
The sheaf $\operatorname{Psh}_{S,1}$ of locally approachable psh functions
contains the constants and is stable under
\begin{enumerate}
\item\label{itm:ch15-app-sum} sums;
\item\label{itm:ch15-app-scal} multiplication by positive real scalars;
\item\label{itm:ch15-app-max} finite maxima;
\item\label{itm:ch15-app-comp} composition $g(u_1,\dots,u_m)$ with a smooth
  convex $g$ nondecreasing in each variable, defined on a convex open subset
  of $\mathbf R^m$ containing the image, and more generally with a uniform
  decreasing limit of such functions;
\item\label{itm:ch15-app-pull} pullback along a morphism of affine
  $G$-tropical spaces.
\end{enumerate}
It is not known to be stable under decreasing sequential limits; by
\Cref{prop:ch15-collapse-reduction}, that single further stability is
equivalent to $\operatorname{Psh}_{S,1}=\operatorname{Psh}_S$.
\end{theorem}

\begin{proof}
Constants are smooth psh, hence in $\operatorname{Psh}_{S,0}\subseteq
\operatorname{Psh}_{S,1}$.

\cref{itm:ch15-app-sum} and \cref{itm:ch15-app-scal} are immediate: if
$u=\lim\downarrow u_k$ and $v=\lim\downarrow v_k$ locally with $u_k,v_k$
smooth psh, then $u_k+v_k$ is smooth psh and decreases to $u+v$, and
$\lambda u_k$ is smooth psh and decreases to $\lambda u$ for $\lambda>0$.
The constant $0$ was handled separately, so no indeterminate form
arises, since the $u_k,v_k$ are real-valued and decreasing limits in
$\mathbf R\cup\{-\infty\}$ add without ambiguity.

\cref{itm:ch15-app-max} is \Cref{prop:ch15-max-stability} with $\alpha=1$.

\cref{itm:ch15-app-comp} is \Cref{prop:ch15-composition} with $\alpha=1$,
together with \cite[Corollary~8.2.2]{CLD25} for the case of a uniform
decreasing limit of such $g$: in that case the same diagonal as in the proof
of \Cref{prop:ch15-max-stability} applies, with $M_{\varepsilon_k}$ replaced
by the $k$-th approximant $g_k$ of $g$.

\cref{itm:ch15-app-pull} is \Cref{prop:ch15-pullback-stagewise} below with
$\alpha=1$.

The final sentence is \Cref{prop:ch15-collapse-reduction} with $\alpha=1$.
\end{proof}

\Cref{thm:ch15-approachable-stability} explains the empirical observation
recorded by CLD, that every known psh function lies at the first stage.  The
approachable class is closed under every operation that the local theory
actually performs on psh functions --- sums, scalars, maxima, smooth convex
compositions, pullbacks --- with exactly one exception: passing to a
decreasing limit.  A candidate function of higher closure rank would
therefore have to be produced by a decreasing limit that is not achievable in
any other way, and no mechanism for producing one is known.  It is worth
noting which naive mechanism does \emph{not} work: an abstract diagonal
argument would require a psh-preserving smooth regularized \emph{minimum},
and no such operator exists, since a minimum of psh functions is not psh.
This is also the reason the complex statement is proved by convolution rather
than by a formal argument.

\begin{proposition}\label{prop:cld-psh-calculus}
CLD-psh functions are upper semicontinuous and form a convex cone containing
the constants.  They are local, stable under finite maxima, and preserved by
pullback of affine $G$-tropical spaces.  A local uniform limit of smooth psh
functions is continuous and CLD-psh.  More precisely, at every ordinal stage
$\alpha$ of the CLD hierarchy one has the stagewise identity
\[
       \operatorname{CPsh}_{S,\alpha}
       =C^0_S\cap\operatorname{Psh}_{S,\alpha}.
\]
This identity does not assert that a continuous function appearing at a
higher stage is a direct first-stage limit of smooth psh functions.
\end{proposition}
\begin{proof}
Upper semicontinuity, addition, and nonnegative scaling are preserved both
at the smooth stage and under decreasing limits.  The sheaf construction
gives locality.  Smooth regularized maxima give the maximum at the initial
stage; passage through the hierarchy gives the general case.  Pullback
preserves smooth Hessian positivity and commutes with decreasing limits.
For the sufficient approximation statement, work on a relatively compact
open set and pass to a subsequence for which
$\epsilon_j:=\|u_j-u\|_\infty$ satisfies
$3\epsilon_{j+1}\leq\epsilon_j$.  Then
$v_j=u_j+2\epsilon_j$ is smooth psh,
$v_{j+1}\leq v_j$, and $v_j\downarrow u$.  The stagewise identity follows by
the transfinite induction of \cite[\S8.1]{CLD25}.  At a successor stage the
predecessor functions may already be nonsmooth, so Dini's theorem does not
collapse the hierarchy.  Indeed, CLD leave stabilization at stage one open;
locally approachable psh functions are, by definition, the stage-one class
\cite[Definition~8.1.13]{CLD25}.
\end{proof}

\subsection{Uniform limits and the continuous class}
\label{sec:ch15-uniform}

The stagewise identity in \Cref{prop:cld-psh-calculus} rests on an elementary
but indispensable device: a uniformly convergent sequence can be shifted so
as to become decreasing.  We isolate it, since it is used again in
\Cref{chap:cld-bedford-taylor}.

\begin{lemma}
\label{lem:ch15-shift}
Let $X$ be a set and let $(f_k)$ be a sequence of real-valued functions on
$X$ converging uniformly to $f$.  Then the sequence
$\bigl(f_k+2\|f-f_k\|_\infty\bigr)_k$ converges uniformly to $f$ and admits a
decreasing subsequence.
\end{lemma}

\begin{proof}
Uniform convergence of the shifted sequence is clear, since
$\|f-f_k\|_\infty\to0$.  If $\|f-f_k\|_\infty=0$ for infinitely many $k$, the
corresponding constant subsequence works.  Otherwise assume
$\|f-f_k\|_\infty>0$ for all $k$.  Given $k$, choose $l>k$ with
$\|f_l-f\|_\infty<\|f-f_k\|_\infty/3$.  From $f_l\le f+\|f_l-f\|_\infty$ and
$f-\|f_k-f\|_\infty\le f_k$ one gets
\[
 f_l+2\|f_l-f\|_\infty
 \le f+3\|f_l-f\|_\infty
 \le f_k+3\|f_l-f\|_\infty+\|f_k-f\|_\infty
 \le f_k+2\|f_k-f\|_\infty ,
\]
the last step because $3\|f_l-f\|_\infty<\|f_k-f\|_\infty$.  Iterating this
choice produces the required decreasing subsequence.  This is
\cite[Lemma~8.1.10]{CLD25}.
\end{proof}

\begin{proposition}
\label{prop:ch15-uniform-limit}
Let $(f_k)$ be a sequence of psh functions on $S$ with values in $\mathbf R$
converging locally uniformly to $f$.  Then $f$ is psh; and if all $f_k$ lie
in $\operatorname{Psh}_{S,\alpha}$, then $f$ lies in
$\operatorname{Psh}_{S,\alpha+1}$.  In particular a locally uniform limit of
smooth psh functions is locally approachable.
\end{proposition}

\begin{proof}
The assertion is local, so work on a relatively compact open set on which the
convergence is uniform.  Adding a real constant $c$ preserves the stage: the
function $y\mapsto y+c$ is smooth, convex and nondecreasing on $\mathbf R$,
so \Cref{prop:ch15-composition} gives
$f_k+c\in\operatorname{Psh}_{S,\alpha}$.  Hence the shifted functions
$f_k+2\|f-f_k\|_\infty$ lie in $\operatorname{Psh}_{S,\alpha}$, and
\Cref{lem:ch15-shift} extracts from them a decreasing subsequence converging
uniformly, in particular pointwise, to $f$.  Therefore
$f\in\operatorname{Psh}_{S,\alpha+1}$.  Taking $\alpha=0$ gives the last
assertion.  The unrefined statement is
\cite[Proposition~8.1.9]{CLD25}.
\end{proof}

\begin{corollary}
\label{cor:ch15-uniform-approachable}
If $u$ is continuous on $S$ and every point of $S$ has one fixed
neighbourhood on which $u$ is a uniform limit of smooth psh functions, then
$u$ is locally approachable.
\end{corollary}

\begin{proof}
Apply \Cref{prop:ch15-uniform-limit} with $\alpha=0$ on the given
neighbourhood, and use that $\operatorname{Psh}_{S,1}$ is a sheaf.
\end{proof}

\begin{warning}{A fixed neighbourhood, not a shrinking family}
\Cref{cor:ch15-uniform-approachable} requires one neighbourhood of each point
serving all accuracies.  If one is only given, for every $\varepsilon>0$, a
neighbourhood $W_\varepsilon$ and a smooth psh $v$ on $W_\varepsilon$ with
$\sup_{W_\varepsilon}|v-u|<\varepsilon$, the neighbourhoods may shrink to the
point and the conclusion does not follow.  This distinction is not pedantry:
it is exactly the gap between the two formulations, and closing it in general
would require a patching construction --- a regularized maximum together with
a strictly positive smooth reserve --- for which no analogue is available on
an arbitrary Berkovich germ.  See \Cref{prob:frontier-cld-hierarchy}.
\end{warning}

\begin{proposition}
\label{prop:ch15-L1-hierarchy}
Let $L^1\!\operatorname{Psh}_S$ denote the sheaf of psh functions on $S$ that
are locally integrable.  Then $L^1\!\operatorname{Psh}_S$ is the smallest
subsheaf of the sheaf of locally integrable functions that contains the
smooth psh functions and is stable under decreasing sequential limits with
locally integrable limit; and, for every ordinal $\alpha$,
\[
 L^1\!\operatorname{Psh}_{S,\alpha}
 =L^1_{\loc}\cap\operatorname{Psh}_{S,\alpha} .
\]
Consequently,
\[
 \operatorname{Psh}_{S,\alpha}=\operatorname{Psh}_S
 \quad\Longrightarrow\quad
 L^1\!\operatorname{Psh}_{S,\alpha}
   =L^1\!\operatorname{Psh}_S .
\]
Conversely,
\[
 L^1\!\operatorname{Psh}_{S,\alpha}
   =L^1\!\operatorname{Psh}_S
 \quad\Longrightarrow\quad
 \operatorname{Psh}_{S,\alpha+1}=\operatorname{Psh}_S .
\]
Write $\beta$ and $\gamma$ for the least collapse ordinals of the plain and
$L^1$ hierarchies, respectively.  They exist by
\Cref{prop:ch15-hierarchy-stabilizes} and the forward implication above, and
they satisfy
\[
                         \gamma\leq\beta\leq\gamma+1.
\]
Equality does not follow from the stagewise intersection identity.
\end{proposition}

\begin{proof}
Define $L^1\!\operatorname{Psh}_{S,\alpha}$ by the recursion of
\Cref{def:ch15-hierarchy} inside the sheaf of locally integrable functions,
and prove the displayed identity by transfinite induction.  The cases
$\alpha=0$ and $\alpha$ limit are immediate.  For a successor, the inclusion
$L^1\!\operatorname{Psh}_{S,\alpha+1}\subseteq
L^1_{\loc}\cap\operatorname{Psh}_{S,\alpha+1}$ is clear.  Conversely, let $u$
be locally integrable and lie in $\operatorname{Psh}_{S,\alpha+1}(U)$; the
question is local, so assume $u=\lim\downarrow u_k$ with
$u_k\in\operatorname{Psh}_{S,\alpha}(U)$.  Each $u_k$ is upper
semicontinuous, hence locally bounded above, and $u_k\ge u$ with $u$ locally
integrable; therefore $u_k$ is locally integrable, so
$u_k\in L^1_{\loc}\cap\operatorname{Psh}_{S,\alpha}
=L^1\!\operatorname{Psh}_{S,\alpha}$ by induction.  Hence
$u\in L^1\!\operatorname{Psh}_{S,\alpha+1}$.  The first assertion follows by
the same minimality argument as in \Cref{lem:ch15-minimality}.  This is
\cite[Proposition~8.1.12]{CLD25}.

The forward implication follows immediately from the stagewise
identity.  For the converse, let $u\in\operatorname{Psh}_S(U)$ and set
$u_m:=\max\{u,-m\}$.  Stability under finite maxima shows that $u_m$ is psh;
it is locally bounded and hence locally integrable.  Under the hypothesis of
the second implication, $u_m\in\operatorname{Psh}_{S,\alpha}(U)$ for every
$m$, while $u_m\downarrow u$.  Hence
$u\in\operatorname{Psh}_{S,\alpha+1}(U)$.
\end{proof}

The practical consequence of \Cref{prop:ch15-L1-hierarchy} is that no
separate integrability bookkeeping is needed: a statement proved for
approachable psh functions restricts to locally integrable approachable psh
functions without loss.  Plain collapse at a specified stage implies
$L^1$-collapse at that stage; in the reverse direction the possible loss of
one successor stage must be retained.

\subsection{Psh functions and positive currents}

\begin{proposition}\label{prop:cld-psh-positive-current}
On a pure-dimensional boundaryless tropical or analytic space, a locally
integrable CLD-psh function $u$ satisfies
\[
                          \dprimeddprime[u]\geq0.
\]
For a smooth $u$, the converse holds.
\end{proposition}
\begin{proof}
For a decreasing smooth approximation, test against a compactly supported
positive complementary form and use local integrability to pass to the
distributional limit.  The transfinite assertion follows by induction.
For smooth $u$, the current pairing is integration against the smooth Hessian;
local positive test forms detect pointwise positivity.

In more detail, let $\mathscr P(U)$ be the set of locally integrable
functions $u$ on $U$ such that $\dprimeddprime[u]\ge0$.  This is a subsheaf
of the sheaf of locally integrable functions, since positivity of a current
is local by \Cref{prop:current-positivity-calculus}.  It contains the smooth
psh functions, because for such a $u$ the current $\dprimeddprime[u]$ is the
current of the smooth positive form $\dprimeddprime u$.  It is stable under
decreasing limits with locally integrable limit: if $u_k\downarrow u$ with
all terms locally integrable and $\omega$ is a compactly supported positive
$(n,n)$-form, then $\int u_k\omega\to\int u\omega$ by monotone convergence,
and an arbitrary test form is a difference of two positive ones
\cite[Proposition~6.5.2]{CLD25}, so $[u_k]\to[u]$ weakly; the operator
$\dprimeddprime$ is weakly continuous by
\Cref{prop:current-differential-calculus}, and the positive cone is weakly
closed, so $\dprimeddprime[u]\ge0$.  By the minimality statement of
\Cref{prop:ch15-L1-hierarchy}, $\mathscr P$ contains every locally
integrable psh function.

For the converse in the smooth case, $\dprimeddprime[u]$ is the current
attached to the smooth form $\dprimeddprime u$, and a smooth form whose
current is positive is itself positive, because positivity of a smooth form
can be tested pointwise and pointwise positivity is detected by test forms
concentrated near a point.  This is \cite[Proposition~8.3.1]{CLD25}.
\end{proof}

\begin{warning}{The converse is not automatic outside the smooth case}
For a merely locally integrable $u$ with $\dprimeddprime[u]\ge0$, one may
\emph{not} conclude that $u$ is psh on a general tropical or analytic space.
On $\mathbf R^n$ the converse does hold, and it is the classical
distributional characterization of convexity, proved as
\Cref{thm:ch09-distributional-convexity} and restated in the present
language as \Cref{thm:ch15-convexity-model}.  The reason it holds there is
that convolution is available: one regularizes, uses the smooth case, and
passes to the limit.  On a general Berkovich space there is no analogous
convolution regularization theorem, and the implication ``positive current
$\Rightarrow$ psh'' is not available in this generality;
what is available instead is the curve test of \Cref{sec:ch15-classical},
which is a different condition with a different proof.  This is the single
most frequently misquoted point in the local theory.  In dimension one,
Wanner identifies continuous functions with positive current with
Thuillier-subharmonic functions; this does not by itself identify them with
the approximation-generated CLD hierarchy.  The separate theorem
\Cref{thm:p063-current-psh-ordinary-regularization} supplies that
identification on smooth proper curves over algebraically closed complete
nontrivially valued fields.
\end{warning}

\subsection{The curve case, after Thuillier}
\label{sec:ch15-thuillier-curves}
\index{subharmonic function!on a Berkovich curve}
\index{Thuillier, Amaury!potential theory on curves}

Everything in this section so far has been stated in a generality where the
converse implication fails or is unknown.  On a smooth strictly analytic
\emph{curve} there is a complete and unconditional potential theory, due to
Thuillier \cite{Thu05}, in which the analogue of the converse is a theorem
and in which subharmonic functions are regularized from above by smooth ones.
We record its two relevant statements here, with their hypotheses, because
they are the primary source for the curve theory used in
\Cref{sec:ch15-classical} and in \Cref{chap:cld-examples}.  We also record
precisely why these imported statements alone do not settle the ordinary
CLD-smooth stage-one question of \Cref{def:ch15-approachable}; the additional
bridge is proved in Chapter~19.

\paragraph{The setting.}
Let $k$ be a complete non-Archimedean field with \emph{nontrivial} absolute
value and let $X$ be a smooth strictly $k$-analytic curve.  For a strictly
$k$-affinoid domain $Y\subseteq X$, write $H(Y)$ for the finite-dimensional
real vector space of continuous real functions on $Y$ that are harmonic on
its interior, in the sense of \Cref{sec:ch08-harmonic} transported by
retraction to a skeleton \cite[\S2.3]{Thu05}.  An upper semicontinuous
function $u:\Omega\to\mathbf R\cup\{-\infty\}$ on an open
$\Omega\subseteq X$ is \emph{subharmonic} if it is not identically $-\infty$
on any connected component of $\Omega$ and, for every strictly $k$-affinoid
domain $Y\subseteq\Omega$ and every $h\in H(Y)$,
\begin{equation}
\label{eq:ch15-thuillier-sh}
 u|_{\partial Y}\le h|_{\partial Y}
 \quad\Longrightarrow\quad
 u|_Y\le h .
\end{equation}
This is \cite[D\'efinition~3.1.5]{Thu05}, and it is the notion appealed to in
the definition of a classically psh function in
\Cref{sec:ch15-classical}.  The cone of subharmonic functions on $\Omega$ is
written $\operatorname{SH}(\Omega)$.

Two auxiliary classes are needed.  Thuillier's \emph{smooth} functions form
the sheaf $\mathcal A^0_X$ associated to the presheaf
$V\mapsto\mathbf R\log|\mathcal O_X(V)^\times|$
\cite[Proposition~3.2.5]{Thu05}; concretely, $\varphi$ is smooth exactly when
it is locally of the form $\Phi\circ\tau_{\mathcal Y}$ for the retraction
$\tau_{\mathcal Y}$ onto the skeleton of a simply semistable formal model,
with $\Phi$ piecewise affine on that skeleton
\cite[Proposition~3.2.4]{Thu05}.  Currents of degree $0$ and $1$ are the
topological duals $\mathrm D^0(X)=\mathcal A^1_c(X)'$ and
$\mathrm D^1(X)=\mathcal A^0_c(X)'$, and $\dprimeddprime:\mathrm D^0(X)\to\mathrm
D^1(X)$ is the transpose of $\dprimeddprime$ on smooth functions
\cite[D\'efinition~3.3.1 and Proposition~3.3.4]{Thu05}.  A degree-one current
$S$ is \emph{positive} if $\langle S,\varphi\rangle\ge0$ for every
nonnegative $\varphi\in\mathcal A^0_c(X)$.  Finally $I(X)\subseteq X$ denotes
the set of points of type $(2)$ or $(3)$, equivalently the set of points that
occur in the Shilov boundary of some $k$-affinoid domain of $X$; a
degree-zero current is the same thing as an arbitrary real function on
$I(X)$ \cite[Proposition~3.3.3]{Thu05}.

\begin{remark}[Normalizations]
\label{rem:ch15-thuillier-normalization}
Three conventions must be matched before the statements below can be used
with the rest of this book, and all three match without a conversion factor.
First, $\log$ denotes the natural logarithm, and \cite{Thu05} points out
that its $\ddc$ depends on that choice through the integral affine structure
it puts on skeleta \cite[Remarque~3.3.16]{Thu05}; the same choice is in force
here, by \Cref{app:normalizations}.  Second, the Poincar\'e--Lelong
normalization of
\cite[Proposition~3.3.15]{Thu05} is
$\dprimeddprime\log|f|=\delta_{\divis(f)}$, which is exactly
\eqref{eq:appJ-cld-pl}, that is, \Cref{thm:cld-poincare-lelong} in dimension
one; no sign or factor of $2$ intervenes.  Third, positivity of a degree-one
current is tested against nonnegative smooth functions, which is the
convention of \Cref{sec:ch14-positivity} in bidegree $(1,1)$ on a curve.
Consequently the operator written $\ddc$ in \cite{Thu05} is the
operator written $\dprimeddprime$ here, with no intervening constant, and we
use the latter notation throughout.
\end{remark}

\paragraph{The two foundational inputs.}

\refstepcounter{theorem}\label{thm:ch15-thuillier-currents}
\begin{quote}
\textbf{Foundational input \thetheorem\ (Thuillier's positive-Laplacian
characterization; source boundary).}
Let $k$ be complete and nontrivially valued and let $X$ be a smooth strictly
$k$-analytic curve.  Then:
\begin{enumerate}
\item\label{itm:ch15-thu-measures} the positive currents of degree $1$ on
  $X$ are exactly the positive Radon measures on the locally compact space
  $X$;
\item\label{itm:ch15-thu-identification} under the injection
  $\operatorname{SH}(X)\hookrightarrow\mathrm D^0(X)$, $u\mapsto u|_{I(X)}$, of
  \cite[Lemme~3.4.1]{Thu05}, the real cone $\operatorname{SH}(X)$ is exactly the set of
  currents $T\in\mathrm D^0(X)$ such that $\dprimeddprime T$ is positive.
\end{enumerate}
In particular a \emph{smooth} $\varphi\in\mathcal A^0(X)$ is subharmonic if
and only if $\dprimeddprime\varphi\ge0$ \cite[Proposition~3.4.4]{Thu05}.
The current identification is imported in full from
\cite[Th\'eor\`eme~3.4.12]{Thu05}, with the injection supplied by
\cite[Lemme~3.4.1]{Thu05}; no proof is claimed here.
\end{quote}

\refstepcounter{theorem}\label{thm:ch15-thuillier-approximation}
\begin{quote}
\textbf{Foundational input \thetheorem\ (Thuillier's approximation
theorem; source boundary).}
Let $X$ be a smooth strictly $k$-analytic curve, let $\Omega\subseteq X$ be
open, and let $u\in\operatorname{SH}(\Omega)$.  Then for every relatively compact open
$\Omega'\subseteq\Omega$ the restriction $u|_{\Omega'}$ is the lower envelope
of a decreasing family of \emph{smooth subharmonic} functions.
This is imported in full from \cite[Th\'eor\`eme~3.4.2]{Thu05}; no proof
is claimed here.
\end{quote}

It is worth recording what
produces the approximants in \Cref{thm:ch15-thuillier-approximation}, because
the shape of the family matters below.  One reduces to
$\Omega'=Y\setminus\partial Y$ for a compact $k$-analytic domain $Y$; for
each finite subset $S\subseteq I(Y)$ containing $\partial Y$ one solves the
Dirichlet problem for the function $u_S$ that agrees with $u$ on $S$ and is
harmonic on $Y\setminus S$; the family $(u_S)_S$, indexed by the filtered set
of such finite subsets, is decreasing, consists of smooth subharmonic
functions, and has lower envelope $u$.  The maximum principle
\cite[Proposition~3.1.11]{Thu05} is used twice, and the identification of
$u$ with its restriction to $I(Y)$ is \cite[Lemme~3.4.1]{Thu05}.

Three further facts from the same source are used elsewhere in this book and
are recorded here for reference.  The maximum principle: a subharmonic
function on an open $\Omega\subseteq X$ attaining a local maximum at a point
of $\Omega$ is locally constant near that point
\cite[Proposition~3.1.11]{Thu05}; this is the one-dimensional input to
\Cref{cor:ch15-gr-maximum}, and it is unconditional, with no continuity
hypothesis.  Harnack's principle: for an increasing sequence of harmonic
functions on a connected open set, either the sequence tends to $+\infty$
locally uniformly or it converges locally uniformly to a harmonic function;
this is the special case, for a constant family of domains, of
\cite[Proposition~3.1.2]{Thu05}, which allows the $n$-th function to be
defined on a domain $\Omega_n$ subject to a local eventual-containment
condition.  And a rigidity statement: a subharmonic
function on a smooth strictly $k$-analytic curve is entirely determined by
its restriction to the set
\begin{equation}
\label{eq:ch15-thuillier-X0}
 X_0=\{x\in X\;:\;[\kappa(x):k]<\infty\}
\end{equation}
of points whose residue field is finite over $k$
\cite[Th\'eor\`eme~3.6.26]{Thu05}, a consequence of the fine-topology
density of $X_0$ \cite[Proposition~3.6.27]{Thu05} and of the capacity theory
of \cite[\S3.6]{Thu05}.

\paragraph{What this does and does not say about stage one.}

\begin{warning}{Thuillier's approximation is not the collapse of the CLD
hierarchy}
\index{CLD hierarchy!early collapse}
\Cref{thm:ch15-thuillier-approximation} is, in form, exactly a stage-one
statement: a member of a psh class exhibited as a decreasing limit of smooth
members.  It is nevertheless not a proof that
$\operatorname{Psh}_{S,1}=\operatorname{Psh}_S$, even on a curve, and it is
not evidence about dimension $\ge2$.  Three separate gaps must be crossed,
and none is crossed by Thuillier's theorem alone.  In the continuous
smooth-proper algebraic curve scope, all three are crossed by the separate
construction of \Cref{thm:p063-current-psh-ordinary-regularization}.

\emph{A family, not a sequence.}  The approximants $(u_S)_S$ are indexed by
the filtered set of finite subsets of $I(Y)$ containing $\partial Y$.  This
is a decreasing \emph{net}, and $I(Y)$ is in general uncountable.  The
hierarchy of \Cref{def:ch15-hierarchy} is defined by decreasing limits of
\emph{sequences}, deliberately so, for the reason stated in
\Cref{sec:ch15-psh}: a Berkovich space need not be first countable and a
decreasing net of psh functions is not automatically the limit of a
decreasing sequence.  Extracting a cofinal sequence from $(u_S)_S$ is an
additional argument, not a restatement.

\emph{A different class of approximants.}  Thuillier's smooth functions are
piecewise affine after retraction to a skeleton
\cite[Proposition~3.2.4]{Thu05}, whereas
$\operatorname{Psh}_{S,0}$ consists of the CLD-smooth functions, locally
pullbacks of $C^\infty$ functions along a moment map, with positive Hessian.
Neither class contains the other: a piecewise affine function with a corner
is not CLD-smooth, and a CLD-smooth function is not in general locally in
$\mathbf R\log|\mathcal O_X^\times|$.  A decreasing approximation by
functions of one class is not a decreasing approximation by functions of the
other.

\emph{A different class of functions.}  $\operatorname{SH}(\Omega)$ is Thuillier's class,
defined by comparison \eqref{eq:ch15-thuillier-sh}.  For continuous functions
on an algebraic curve, Wanner identifies this class with the condition
$\dprimeddprime[u]\ge0$; see \Cref{thm:wanner-curve-comparison}.  Wanner's
terminology calls the latter condition ``psh'', but it does not thereby
identify it with the approximation-generated CLD hierarchy.  The missing
identification for continuous functions on smooth proper curves is proved in
\Cref{thm:p063-current-psh-ordinary-regularization}.  For functions taking
the value $-\infty$ this book asserts no such
comparison, and \Cref{thm:cld-to-classical-psh} supplies only one implication.

Finally, and most importantly, \Cref{thm:ch15-thuillier-approximation} is a
statement \emph{about curves}.  Transferring it to the CLD class in
dimension $\ge2$ is not licensed by anything proved here, and the open
problem \Cref{prob:frontier-cld-hierarchy} is stated in that generality.
\end{warning}

\begin{remark}[Why tree geometry is decisive]
\label{rem:ch15-thuillier-dimension-one}
The following observation explains the geometry behind
\Cref{thm:p063-current-psh-ordinary-regularization}; the complete proof is in
Chapter~19.  It is not attributed to \cite{Thu05} or to another source.

The construction sketched after
\Cref{thm:ch15-thuillier-approximation} works because, in dimension one,
every stratum at which the approximant $u_S$ can fail to be harmonic is a
\emph{vertex}, that is, a single point of the skeleton.  Prescribing the
approximant there is prescribing finitely many real numbers, the Dirichlet
problem on the complementary edges has a unique solution, and refining $S$
only adds further point conditions, which by the maximum principle can only
decrease the solution.  Both the interpolation and the monotonicity are
therefore finite-dimensional and combinatorial.

In dimension $\ge2$ a skeleton has strata of positive dimension, and the
corresponding local model problem is a genuine boundary-value problem on a
polyhedron rather than a finite interpolation.  Monotonicity under
refinement is no longer supplied by a maximum principle on edges, and the
patching of solutions across faces is exactly where a regularized-maximum
argument would be needed and where, as noted after
\Cref{cor:ch15-uniform-approachable}, no analogue is available.  This is
consistent with the fact that the collapse results available here and in
Chapter~19---\Cref{thm:ch15-convexity-model} on
$\mathbf R^n$, \Cref{prop:ch15-log-approachable} for $\log|f|$, the
one-dimensional interval and star models of
\Cref{thm:p071-cld-constant-collar-endpoint,%
thm:p071-cld-infinite-collar-endpoint,%
thm:ch16-collapse-finite-tropical-curves,%
thm:p071-infinite-star-collapse-candidate}, and the continuous
smooth-proper-curve result---have a finite-dimensional, order-theoretic, or
globally convex model behind them.
\end{remark}

\section{Convexity: the model case on \texorpdfstring{$\mathbf R^n$}{Rn}}
\label{sec:ch15-convexity}

On a real affine space every question raised in \Cref{sec:ch15-hierarchy} has
a positive answer, and the whole theory reduces to classical convex analysis.
This is the basic model case and the one whose collapse proof uses
translation-invariant radial mollification.  It is worth stating in full
because it supplies the intuition behind the general definition and because
the different collapse mechanisms on tropical toric partial compactifications,
finite ambient-induced real-affine stars, constant-tail intervals, and the
exact CLD stars $X_B$ are treated separately in
\Cref{sec:ch16-collapse-positive-cases}.

\begin{theorem}
\label{thm:ch15-convexity-model}
Let $X$ be the tropical space of level $n$ associated with $\mathbf R^n$, let
$U\subseteq\mathbf R^n$ be open and convex, and let
$u:U\to\mathbf R\cup\{-\infty\}$ be a function that is not identically
$-\infty$.  The following are equivalent:
\begin{enumerate}
\item\label{itm:ch15-conv-psh} $u$ is psh, that is, a section of
  $\operatorname{Psh}_X$ over $U$;
\item\label{itm:ch15-conv-convex} $u$ is real-valued and convex;
\item\label{itm:ch15-conv-current} $u$ is locally integrable and
  $\dprimeddprime[u]\ge0$;
\item\label{itm:ch15-conv-stage1} $u$ is a section of
  $\operatorname{CPsh}_{X,1}$, that is, $u$ is continuous and locally a
  uniform limit of smooth psh functions.
\end{enumerate}
Consequently, on $\mathbf R^n$ the sheaves of psh functions, of continuous
psh functions, and of locally convex functions coincide, and the hierarchy
collapses: $\operatorname{Psh}_{X,1}=\operatorname{Psh}_X$.
\end{theorem}

\begin{proof}
We first record that a smooth function $w$ on a convex open subset with
$\dprimeddprime w\ge0$ is convex: by \eqref{eq:appJ-real-hessian} the
coefficient matrix of $\dprimeddprime w$ is the Hessian of $w$, and
positivity in bidegree $(1,1)$ is positive semidefiniteness of that matrix
(\Cref{prop:ch09-positivity-dictionary}), so
\Cref{thm:convex-ddc-positive} applies.  Thus
$\operatorname{Psh}_{X,0}$ consists of the smooth convex functions.

\cref{itm:ch15-conv-psh}$\Rightarrow$\cref{itm:ch15-conv-convex}.  Let
$\mathscr C$ be the subsheaf of functions that are, locally, either finite
and convex or identically $-\infty$.  By the previous paragraph it contains
$\operatorname{Psh}_{X,0}$.  It is stable under decreasing sequential limits:
let $(u_k)$ be decreasing with limit $u$ on a convex open $V$, all $u_k$
finite convex.  If $u(a)=-\infty$ for some $a\in V$, then for every $b\in V$
and $t\in[0,1[$ convexity of $u_k$ gives
$u_k((1-t)a+tb)\le(1-t)u_k(a)+tu_k(b)$, and letting $k\to\infty$ gives
$u((1-t)a+tb)=-\infty$; hence $u\equiv-\infty$ on the half-open segment
$[a,b[$ for every $b$, and therefore on $V$.  Otherwise $u$ is finite
everywhere and the convexity inequalities pass to the limit.  By
\Cref{lem:ch15-minimality}, $\operatorname{Psh}_X\subseteq\mathscr C$; since
$u\not\equiv-\infty$ and $U$ is connected, $u$ is finite and convex.

\cref{itm:ch15-conv-convex} or
\cref{itm:ch15-conv-current}$\Rightarrow$\cref{itm:ch15-conv-stage1}.  A finite
convex function on an open convex set is continuous, indeed locally Lipschitz
\cite[\S10]{Roc70}, hence locally integrable; so under either hypothesis $u$
is locally integrable.  The assertion is local, so fix a ball $B(a,r)$ with
$\overline{B(a,s)}\subset U$ for some $s>r$.  Let $\rho$ be a radial
mollifier as in \Cref{prop:ch09-radial-mollification} and put
$u_\varepsilon=u*\rho_\varepsilon$.  If $u$ is convex, $u_\varepsilon$ is
smooth and convex by \Cref{prop:superform-convex-mollification}, and
$u_\varepsilon\to u$ uniformly on $B(a,r)$.  If instead
$\dprimeddprime[u]\ge0$, then
$\dprimeddprime u_\varepsilon=\dprimeddprime[u]*\rho_\varepsilon\ge0$, so
$u_\varepsilon$ is smooth with positive semidefinite Hessian, hence convex by
\Cref{thm:convex-ddc-positive}, and $u_\varepsilon\to u$ in $L^1_{\loc}$ and
--- by \Cref{thm:ch09-distributional-convexity} --- pointwise to the convex
representative, uniformly on $B(a,r)$.  Either way $u$ is a uniform limit on
$B(a,r)$ of smooth psh functions, so
$u\in\operatorname{CPsh}_{X,1}(B(a,r))$.

\cref{itm:ch15-conv-stage1}$\Rightarrow$\cref{itm:ch15-conv-current} is
\Cref{prop:cld-psh-positive-current} together with continuity, and
\cref{itm:ch15-conv-stage1}$\Rightarrow$\cref{itm:ch15-conv-psh} is
\Cref{prop:cld-psh-calculus}.  This closes the cycle.  The theorem is
\cite[Proposition~8.3.2]{CLD25}.

For the final assertion, let $u$ be psh on an open $V\subseteq\mathbf R^n$.
The question is local, so shrink $V$ to a ball, on which either
$u\equiv-\infty$ --- and then $u$ is the decreasing limit of the smooth psh
constants $-k$, hence in $\operatorname{Psh}_{X,1}$ --- or $u$ is finite and
convex by the equivalence just proved, hence in
$\operatorname{CPsh}_{X,1}\subseteq\operatorname{Psh}_{X,1}$.
\end{proof}

\begin{remark}[Where the monotonicity comes from]
\label{rem:ch15-radial}
The proof above produces a \emph{uniform} approximation and then invokes
\Cref{prop:ch15-uniform-limit} to convert it into a decreasing one.  There is
a more direct route, and it is worth seeing because it explains the role of
radiality.  If $\rho$ is radial and $u$ is convex, then by
\Cref{prop:ch09-radial-mollification} the family $u*\rho_\varepsilon$ is
smooth, convex, dominates $u$, and is \emph{nondecreasing in
$\varepsilon$}; hence $u*\rho_{1/k}\downarrow u$ is already a decreasing
sequence of smooth psh functions, and no shift device is needed.  Radiality
enters twice: once through the evenness of the kernel, which makes its
barycentre the origin and gives $u*\rho_\varepsilon\ge u$ by Jensen, and once
in the convex-even argument that gives monotonicity in the parameter.  For a
non-radial kernel both conclusions fail; see \Cref{ex:ch09-nonradial-kernel}.

The same convex-even argument reappears in
\Cref{lem:ch15-regmax}\,\cref{itm:ch15-regmax-mono}.  This is not a
coincidence: the regularized maximum \emph{is} a mollification, of the convex
function $\max$, with a radial --- indeed product-even --- kernel.
\end{remark}

\begin{example}[A convex function with a nontrivial polar locus, and one
without]
\label{ex:ch15-convex-polar}
On $\mathbf R^n$ a psh function that is not identically $-\infty$ on a
connected open set is finite there, by
\Cref{thm:ch15-convexity-model}.  \emph{There are no nontrivial polar loci in
the model case.}  This is a genuine difference from complex geometry, where
$\log|z|$ is psh on $\mathbf C$ with polar locus a point, and it is a genuine
difference from the Berkovich case, where $\log|f|$ has polar locus
$\{f=0\}$ by \Cref{prop:ch15-log-approachable}.  The reason is dimensional:
on $\mathbf R^n$ the psh condition is convexity along \emph{real} lines, and
a convex function on an interval that is $-\infty$ at one point is
$-\infty$ on a half-line, hence on everything by connectedness.  In the
complex and Berkovich settings the polar locus has real codimension two,
respectively is a Zariski-closed set not meeting any skeleton in top
dimension, and no such propagation occurs.

This is a warning about transporting intuition.  The statement
``psh functions are convex'' is exactly true on $\mathbf R^n$; the statement
``psh functions have no polar locus'' is true on $\mathbf R^n$ and false on
a Berkovich space, although both are consequences of the same definition
applied on different spaces.
\end{example}

\section{Polar loci}
\label{sec:ch15-polar}

\begin{definition}[Polar locus]
\label{def:ch15-polar}
\index{polar locus}
The \emph{polar locus} of a psh function $u$ on $S$ is the set
$P(u)=\{x\in S:u(x)=-\infty\}$.  It is a $G_\delta$ subset, being the
intersection of the open sets $\{u<-k\}$, and it is closed in $S$ when $u$ is
continuous as a map to $\mathbf R\cup\{-\infty\}$.
\end{definition}

The basic example is $\log|f|$ for a holomorphic $f$.  It is also the basic
example of a locally approachable function that is not continuous as a
real-valued function, and the construction of its approximants is the place
where the truncation property of the regularized maximum is indispensable.

\begin{proposition}
\label{prop:ch15-log-approachable}
Let $f$ be a holomorphic function on an analytic space $S$, regarded as an
affine $G$-tropical space.  Then $\log|f|$, extended by $-\infty$ on
$\{f=0\}$, is a locally approachable psh function, and its polar locus is
$\{f=0\}$.  If $f$ is moreover a non-zero-divisor, then $\log|f|$ is locally
integrable and $\{f=0\}$ is a Zariski-closed subset of codimension one.
\end{proposition}

\begin{proof}
Write $u=\log|f|$ and, for $t>0$,
\[
 u_t=\log\max\{|f|,t\}=\max\{u,\log t\},
\]
a continuous real-valued function on $S$, equal to the constant $\log t$ on
the open set $\{|f|<t\}$ and to $u$ on $\{|f|>t\}$.

\emph{Step 1: $u_t$ is smooth psh after a regularized maximum.}  Choose
$\varepsilon>0$ and let $M_\varepsilon$ be the regularized maximum of
\Cref{lem:ch15-regmax} with $m=2$.  Put
\[
 w_{t,\varepsilon}=M_\varepsilon\bigl(\log|f|,\log t\bigr).
\]
We claim that $w_{t,\varepsilon}$ is a globally defined smooth psh function
on $S$.  Two open sets cover $S$.  On $\{f\neq0\}$ the function $\log|f|$ is
smooth with $\dprimeddprime\log|f|=0$ by \Cref{lem:ch15-log-smooth}, and
$\log t$ is a smooth psh constant; \Cref{prop:ch15-composition} with
$\alpha=0$ shows that $w_{t,\varepsilon}$ is smooth psh there.  On the open
set $\{|f|<t\rme^{-2\varepsilon}\}$ we have
$\log|f|\le\log t-2\varepsilon$, so by
\Cref{lem:ch15-regmax}\,\cref{itm:ch15-regmax-trunc} the value
$M_\varepsilon(\log|f|,\log t)$ does not depend on the first argument;
it therefore equals the constant $M_\varepsilon(\log t-2\varepsilon,\log t)$
on that set, which is smooth psh.  Since the two open sets cover $S$ and the
two descriptions agree on the overlap, $w_{t,\varepsilon}$ is smooth psh on
$S$.  It is here that the truncation property is used: the log-sum-exp
regularization would give a function that is not smooth across $\{f=0\}$,
because every argument contributes to it.

\emph{Step 2: a decreasing diagonal.}  Choose $t_k\downarrow0$ and
$\varepsilon_k\downarrow0$, for instance $t_k=\varepsilon_k=1/k$, and set
$w_k=w_{t_k,\varepsilon_k}$.  Then $w_k$ is smooth psh by Step 1, and
\[
 \begin{aligned}
  w_{k+1}
   &=M_{\varepsilon_{k+1}}(\log|f|,\log t_{k+1})\\
   &\le M_{\varepsilon_k}(\log|f|,\log t_{k+1})\\
   &\le M_{\varepsilon_k}(\log|f|,\log t_k)=w_k .
 \end{aligned}
\]
The first inequality is the monotonicity in \Cref{lem:ch15-regmax},
specifically \cref{itm:ch15-regmax-mono}; the second is monotonicity in the
second argument.  Finally
\Cref{lem:ch15-regmax}\,\cref{itm:ch15-regmax-sandwich} gives
\[
 \max\{u,\log t_k\}\le w_k\le\max\{u,\log t_k\}+\varepsilon_k ,
\]
and $\max\{u,\log t_k\}\downarrow u$ pointwise, including at points of
$\{f=0\}$, where both sides tend to $-\infty$.  Hence $w_k\downarrow u$ and
$u\in\operatorname{Psh}_{S,1}$.

\emph{Step 3: the polar locus and integrability.}  By construction
$u(x)=-\infty$ exactly when $|f(x)|=0$, that is, on $\{f=0\}$.  If $f$ is a
non-zero-divisor, local integrability is
\Cref{prop:ch15-log-current}\,\cref{itm:ch15-logcur-int} and the codimension
statement is \Cref{lem:ch15-divisor-basic}\,\cref{itm:ch15-div-supp}.
\end{proof}

\begin{corollary}
\label{cor:ch15-meromorphic-difference}
Let $f$ be a regular meromorphic function on $S$.  Locally, $\log|f|$ is a
difference of two locally approachable psh functions, and
$\dprimeddprime[\log|f|]=\delta_{\divis(f)}$ is a difference of two positive
currents.  In particular $\log|f|$ is psh if and only if the polar part of
$\divis(f)$ is empty, that is, if and only if $f$ is holomorphic on the
locus where it is defined.
\end{corollary}

\begin{proof}
Locally $f=g/h$ with $g,h$ non-zero-divisors, and
$\log|f|=\log|g|-\log|h|$; apply \Cref{prop:ch15-log-approachable} to each.
The current statement is \Cref{thm:cld-poincare-lelong} together with the
decomposition $\divis(f)=\divis(g)-\divis(h)$ of
\Cref{lem:ch15-divisor-basic}\,\cref{itm:ch15-div-add}.  For the last
assertion, if $f$ is holomorphic then $\log|f|$ is psh by
\Cref{prop:ch15-log-approachable}; conversely if $\log|f|$ is psh then
$\dprimeddprime[\log|f|]\ge0$ by \Cref{prop:cld-psh-positive-current},
which by \Cref{thm:cld-poincare-lelong} forces $\divis(f)$ to be effective.
\end{proof}

\begin{example}[A psh function with a nontrivial polar locus]
\label{ex:ch15-polar-example}
Let $K$ be algebraically closed and nontrivially valued, let
$X=\mathbf A^{2,\An}_K$ with coordinates $T_1,T_2$, and put
\[
 u=\log\max\{|T_1|,|T_2|\}.
\]
By \Cref{prop:ch15-log-approachable} and \Cref{prop:ch15-max-stability}, $u$
is locally approachable, and its polar locus is
$P(u)=\{T_1=T_2=0\}$, the single rigid point $o=(0,0)$.  This is a polar
locus of codimension two.

Three features are worth noting.  First, $u$ is continuous as a map to
$\mathbf R\cup\{-\infty\}$ but not as a real-valued function, so it is psh
without being cpsh; the class $\operatorname{CPsh}$ is strictly smaller.
Second, $u$ is locally integrable: on a tropical chart adapted to
$(T_1,T_2)$ the function $u$ becomes $-\min\{x_1,x_2\}$ on
$\mathbf R^2$, which is continuous and hence locally integrable, and the
polar locus corresponds to the point at infinity of the chart, of measure
zero.  Third, $\dprimeddprime u$ is the positive current computed in
\Cref{ex:ch15-bm-dn} for the truncated potentials; the Bochner--Martinelli
formula \Cref{thm:cld-local-bochner-martinelli} says that its second power,
suitably interpreted, is $\delta_o$.  Thus the polar locus of $u$ carries
all of its Monge--Amp\`ere mass, which is the exact opposite of what happens
for the potentials appearing in \Cref{thm:cld-polar-negligibility}.  The
distinction is that the negligibility theorem concerns the measure attached
to \emph{other} potentials, not to $u$ itself.
\end{example}

\begin{example}[A polar locus that is not analytic]
\label{ex:ch15-polar-nonanalytic}
Let $\Gamma\subset\mathbf R$ be an open interval, regarded as a
one-dimensional tropical space, and let $u$ be the convex function
$u\equiv-\infty$.  Its polar locus is all of $\Gamma$.  By
\Cref{thm:ch15-convexity-model} this is the only way a psh function on a
connected open subset of $\mathbf R^n$ can have a nonempty polar locus.  The
example is trivial but it makes the point that \Cref{def:ch15-polar} does
\emph{not} assert analyticity of the polar locus: on a Berkovich space the
polar locus of a general psh function is only a $G_\delta$ set, and it is
Zariski-closed in the examples above because they are built from
$\log|f|$.  A Siu-type analyticity theorem for polar loci or Lelong sublevel
sets is not available in this theory.
\end{example}

\subsection{Negligibility and compact-sublevel barriers}

Polar loci are negligible for the Monge--Amp\`ere measures of \emph{other}
continuous psh potentials.  The precise statement requires the
Chern--Levine--Nirenberg inequalities and the closure argument of
\Cref{chap:cld-bedford-taylor}; we recall its analytic specialization here.
The curvature factors may belong to the full continuous psh sheaf.  Local
approachability is retained for the auxiliary, possibly unbounded polar
function, while goodness, strictness, and the compact-sublevel barrier are
the geometric inputs that make the weighted estimate available.

% Production verification: Rethlas v3.3.3-r1,
% run 20260901T075502Z_0210bc4a60f4_8aed4fe9.
\begin{theorem}
\label{thm:ch15-negligibility-recalled}
Let $K$ be any complete non-Archimedean field, possibly trivially valued,
and let $X$ be a good, separated, strict, boundaryless, pure
$n$-dimensional $K$-analytic space.  Fix $0\leq m\leq n$.  Let
$u_1,\dots,u_m$ be arbitrary continuous psh functions on $X$, and let
$\alpha$ be a smooth form of type
$(n-m,n-m)$, and let $V$ be a locally approachable psh function on $X$ that
is locally integrable for smooth top-degree measures.  Then $V$ is locally
integrable for the total variation of the signed measure
\[
 \dprimeddprime u_1\wedge\cdots\wedge
 \dprimeddprime u_m\wedge\alpha,
\]
and this measure has zero total variation on $P(V)$.  Moreover it has zero
total variation on every closed analytic subset of local dimension $<n$.
\end{theorem}

\begin{lemma}
\label{lem:ch15-negligibility-zero-dimensional-edge}
In the setting of \Cref{thm:ch15-negligibility-recalled}, suppose that
$n=0$.  Then necessarily $m=0$, and the conclusions of that theorem hold:
$V$ is locally integrable for the total variation of the smooth signed
top-degree measure $\alpha$, that measure has zero total variation on
$P(V)$, and it has zero total variation on every closed analytic subset of
local dimension $<0$.
\end{lemma}

\begin{proof}
The inequalities $0\leq m\leq n=0$ force $m=0$.  Thus the displayed wedge
product in \Cref{thm:ch15-negligibility-recalled} has no curvature factor
and is the smooth $(0,0)$-form $\alpha$, viewed as the signed measure
$\alpha\,\delta_X$ through the integration current of $X$.

This interpretation is genuinely measure-theoretic in dimension zero.
Every point has a zero-dimensional affinoid neighbourhood; its affinoid
algebra is Artinian, so its Berkovich spectrum is finite.  On such a
neighbourhood the integration current $\delta_X$ is therefore represented
by a finite positive atomic Radon measure $\mu_X$, with the usual positive
analytic multiplicities.  In particular, every analytic point $y$ in that
neighbourhood has strictly positive finite weight
$0<\mu_X(\{y\})<\infty$.  The constant smooth $(0,0)$-form $1$ gives this
smooth top-degree measure.  Hence the assumed local integrability of $V$
for smooth top-degree measures gives
\[
                       \int |V|\,\mathrm d\mu_X<\infty.
\]
The smooth scalar $\alpha$ is bounded on the finite neighbourhood, and
the total variation of $\alpha\,\mu_X$ is $|\alpha|\,\mu_X$.  Consequently
\[
 \int |V|\,\mathrm d|\alpha\mu_X|
 \leq \|\alpha\|_{\infty}\int |V|\,\mathrm d\mu_X<\infty.
\]
Since $|V|=\infty$ on $P(V)=\{V=-\infty\}$, the positive finite weight of
each analytic point shows directly that local integrability excludes
$V(y)=-\infty$ at every $y$.  Equivalently $\mu_X(P(V))=0$, and therefore
$|\alpha\mu_X|(P(V))=0$.  Finally, a nonempty analytic space has
nonnegative local dimension.  Thus a closed analytic subset of local
dimension $<0$ is empty, and the remaining assertion is vacuous.
\end{proof}

\begin{proof}[of \Cref{thm:ch15-negligibility-recalled}]
When $n=0$, the result is
\Cref{lem:ch15-negligibility-zero-dimensional-edge}.  Assume henceforth
that $n\geq1$.  The compact-sublevel hypothesis of
\Cref{thm:cld-polar-negligibility} is supplied inside every prescribed
ordinary neighbourhood by \Cref{thm:ch16-barrier}; goodness and separation
also give the required local differential separation.  The polar-locus
assertion is therefore exactly \Cref{thm:cld-polar-negligibility}.  The last
assertion is \Cref{cor:ch16-cpsh-analytic-nullity}, whose proof constructs the
locally approachable logarithmic pole by prime avoidance and smooth convex
truncation.  For first-stage curvature factors these conclusions are
\cite[Theorem~16.9.2 and Corollaries~16.9.3--16.9.4]{CLD25}; the passage to
the full continuous psh sheaf is proved in
\Cref{chap:cld-bedford-taylor}.
\end{proof}

The next statement is the local compact-sublevel barrier used in the proof of
\Cref{thm:ch15-negligibility-recalled}.  Its two levels must be kept
distinct: the value at the chosen point is negative, whereas compactness is
required for a strictly positive sublevel.  Requiring a function to be
negative on its entire domain and its nonpositive sublevel to be compact
would instead require that entire open domain to be compact.

% Production verification: Rethlas v3.3.2-r1,
% run 20260901T065725Z_06ee37b2e57b.
\begin{proposition}
\label{prop:ch15-exhaustion}
Let $K$ be any complete non-Archimedean field, possibly trivially valued,
and let $X$ be a good, separated, strict, boundaryless, pure
$n$-dimensional $K$-analytic space.  Given $x\in X$ and an ordinary open
neighbourhood $\Omega$ of $x$, there are an ordinary open
$Y\subseteq\Omega$, a smooth psh function $\rho$ on $Y$, and $\delta>0$
such that
\[
             \rho(x)<0,
 \qquad K_\delta:=\{\rho\leq\delta\}\Subset Y.
\]
In particular, $U=\{\rho<0\}$ is a nonempty relatively compact ordinary
open subset of $Y$.
\end{proposition}

\begin{proof}
Choose an affinoid neighbourhood $A$ of $x$ contained in $\Omega$ and a
rational affinoid neighbourhood $W$ of $x$ contained in
$\operatorname{Int}(A)$.  Thus
\[
 W=\bigl\{|f_i|\leq\lambda_i|g|\ (1\leq i\leq s)\bigr\},
\]
where $f_1,\ldots,f_s,g$ are analytic on $A$, have no common zero, and
$\lambda_i>0$.  For $\mu>1$ put
$W_\mu=\{|f_i|\leq\mu\lambda_i|g|\}$.  For some $\mu>1$ one has
$W_\mu\subseteq\operatorname{Int}(A)$: otherwise the nested compact sets
$W_\mu\cap(A\setminus\operatorname{Int}(A))$, as $\mu\downarrow1$, would
have a point in $W\setminus\operatorname{Int}(A)$.  Since the $f_i$ and $g$
have no common zero, $g$ has no zero on $W_\mu$.

On $Y=\operatorname{Int}(A)\cap\{g\ne0\}$ set
\[
 a_i=\log(|f_i/g|/\lambda_i),\qquad a_0=-1,
 \qquad t=\log\mu>0.
\]
Choose $0<\varepsilon<t/8$ and a smooth convex coordinatewise
nondecreasing regularized maximum $M_\varepsilon$ satisfying
\[
 \max_i s_i\leq M_\varepsilon(s_0,\ldots,s_s)
       \leq\max_i s_i+\varepsilon
\]
and ignoring a coordinate more than $2\varepsilon$ below the maximum.
Define
\[
             \rho=M_\varepsilon(a_0,a_1,\ldots,a_s)-t/2.
\]
The constant branch keeps the maximum finite when all numerators vanish;
near a zero of one $f_i$, that coordinate is ignored, so $\rho$ extends as
a smooth psh function.  At $x$ all $a_i\leq0$, and therefore
$\rho(x)\leq\varepsilon-t/2<0$.

Take $\delta=t/4$.  The inequality $\rho\leq\delta$ implies
$a_i\leq3t/4$ for all $i$, whence
\[
 K_\delta\subseteq W_{\exp(3t/4)}\subseteq W_\mu\subseteq Y.
\]
The set $K_\delta$ is closed in the compact rational domain on the right,
and hence compact.  The construction uses arbitrary positive real radii,
so it also applies when the valuation of $K$ is trivial.  This is the
compact-positive-sublevel form of
\cite[Proposition~16.9.1]{CLD25}, proved again in
\Cref{thm:ch16-barrier}.
\end{proof}

\section{Maximum principles}
\label{sec:ch15-maximum}

A maximum principle asserts that a psh function attaining an interior maximum
is locally constant.  In this theory the status of such a statement depends
sharply on the class of functions and on the space, and it is worth being
explicit about what is proved and what is not.

\begin{itemize}
\item On $\mathbf R^n$ the maximum principle holds and is elementary; see
  \Cref{prop:ch15-max-convex}.
\item On a one-dimensional boundaryless tropical space it holds for
  continuous psh functions; see \Cref{thm:ch15-max-onedim}.  The proof uses
  balancing in an essential way, and the statement fails without it
  (\Cref{ex:ch15-unbalanced-max}).
\item On a smooth strictly $k$-analytic curve it holds for Thuillier's
  subharmonic functions, with no continuity and no boundedness hypothesis;
  this is \cite[Proposition~3.1.11]{Thu05}, recorded in
  \Cref{sec:ch15-thuillier-curves}.
\item For CLD-psh functions on a general analytic space of dimension $\ge2$
  it is \emph{not known}.  CLD say so explicitly in the introduction to their
  memoir, contrasting their approximation-defined class with the
  curve-test class of Gubler and Rabinoff, for which a maximum principle is
  proved \cite[\S0.3.5]{CLD25}.
\item For classically psh functions in the sense of
  \Cref{sec:ch15-classical}, a maximum principle is available and is stated
  there.
\end{itemize}

\begin{proposition}
\label{prop:ch15-max-convex}
Let $\Omega\subseteq\mathbf R^n$ be open and connected and let
$u:\Omega\to\mathbf R\cup\{-\infty\}$ be psh.  If $u$ attains its supremum at
a point of $\Omega$, then $u$ is constant.
\end{proposition}

\begin{proof}
If $u\equiv-\infty$ there is nothing to prove, so by
\Cref{thm:ch15-convexity-model} we may assume $u$ finite and locally convex.
Let $M=\sup_\Omega u$ be attained at $x_0$ and let
$A=\{x\in\Omega:u(x)=M\}$, a nonempty set, closed in $\Omega$ because $u$ is
continuous.  We show $A$ is open.  Let $a\in A$ and take $r>0$ with
$B(a,r)\subseteq\Omega$ and $u$ convex on $B(a,r)$.  Suppose $u(b)<M$ for
some $b\in B(a,r)$, and set $c=a+t(a-b)$ for $t>0$ small enough that
$c\in B(a,r)$.  Then $a=\frac{t}{1+t}b+\frac{1}{1+t}c$ is a convex
combination, so
\[
 M=u(a)\le\frac{t}{1+t}u(b)+\frac{1}{1+t}u(c)
 <\frac{t}{1+t}M+\frac{1}{1+t}M=M,
\]
a contradiction.  Hence $u\equiv M$ on $B(a,r)$ and $A$ is open.  Since
$\Omega$ is connected, $A=\Omega$.
\end{proof}

\subsection{Dimension one: the atom at a vertex}

Let $(\mathscr C,m)$ be a one-dimensional weighted rational polyhedral
complex in $N_{\mathbf R}$, with support $X=|\mathscr C|$, positive weights
$m_\sigma$ on the edges, and lattice arclength on each edge as in
\Cref{def:ch11-lattice-measure}.  For a vertex $p$ we write $T_pX$ for the
set of edges of $\mathscr C$ containing $p$, $u_{\sigma/p}\in N$ for the
primitive generator of $\sigma$ pointing away from $p$, and $d_\sigma u(p)$
for the outgoing slope of a function $u$ along $\sigma$ at $p$, computed in
the lattice arclength of $\sigma$.  By \Cref{def:ch11-defect} the boundary
defect at $p$ is $\mathbf N(p)=-\sum_{\sigma\in T_pX}m_\sigma u_{\sigma/p}$,
and $X$ is balanced --- equivalently boundaryless --- exactly when
$\mathbf N(p)=0$ for every $p$.

\begin{lemma}
\label{lem:ch15-vertex-atom}
Assume $X$ balanced.  Let $u$ be a continuous function on $X$, of class
$C^2$ on each edge, and let $p$ be a vertex.  Then, on a neighbourhood of
$p$,
\[
 \dprimeddprime[u]
 =\Bigl(\sum_{\sigma\in T_pX}m_\sigma\,d_\sigma u(p)\Bigr)\delta_p
  +\sum_{\sigma\in T_pX}m_\sigma\,u''\,\rd s|_\sigma .
\]
Without the balancing hypothesis the same computation produces the extra term
$u(p)\,\bigl\langle \rd\Psi(p),\mathbf N(p)\bigr\rangle\,\delta_p$, where
$\Psi$ is any smooth ambient representative of the test function at $p$; the
sign is fixed by
$\sum_\sigma m_\sigma d_\sigma\Psi(p)=\langle\rd\Psi(p),\sum_\sigma
m_\sigma u_{\sigma/p}\rangle=-\langle\rd\Psi(p),\mathbf N(p)\rangle$.
\end{lemma}

\begin{proof}
Let $\varphi$ be a smooth function supported in a small neighbourhood of $p$
meeting no other vertex, and let $\Psi$ be a smooth function on
$N_{\mathbf R}$ representing $\varphi$ near $p$.  Identify each edge
$\sigma\in T_pX$ near $p$ with $[0,\ell)$ by lattice arclength $s$, so that
$s=0$ is $p$ and $\partial_s$ corresponds to $u_{\sigma/p}$.  Then, by the
definition of the pairing and \Cref{prop:current-differential-calculus},
\[
 \langle\dprimeddprime[u],\varphi\rangle
 =\langle[u],\dprimeddprime\varphi\rangle
 =\sum_{\sigma\in T_pX}m_\sigma\int_0^\ell u(s)\,\varphi_\sigma''(s)\,\rd s ,
\]
where $\varphi_\sigma(s)=\varphi$ restricted to $\sigma$, and the weights and
lattice measures are those of \Cref{def:ch11-lattice-measure}.  Two
integrations by parts on each edge, all boundary terms at $s=\ell$ vanishing
by the support hypothesis, give
\[
 \int_0^\ell u\varphi_\sigma''\,\rd s
 =-u(p)\,\varphi_\sigma'(0)+u'(0^+)\varphi(p)
   +\int_0^\ell u''\varphi_\sigma\,\rd s .
\]
Now $\varphi_\sigma'(0)=\langle \rd\Psi(p),u_{\sigma/p}\rangle$ and
$u'(0^+)=d_\sigma u(p)$.  Summing over $\sigma$ with the weights,
\[
 \begin{aligned}
 \langle\dprimeddprime[u],\varphi\rangle
 &=-u(p)\Bigl\langle \rd\Psi(p),
      \sum_\sigma m_\sigma u_{\sigma/p}\Bigr\rangle
   +\varphi(p)\sum_\sigma m_\sigma d_\sigma u(p)\\
 &\quad+\sum_\sigma m_\sigma\int u''\varphi_\sigma\,\rd s .
 \end{aligned}
\]
The first term is $u(p)\langle \rd\Psi(p),\mathbf N(p)\rangle$, which
vanishes when $X$ is balanced.  This gives both assertions.
\end{proof}

\Cref{lem:ch15-vertex-atom} is the reason \Cref{prop:cld-psh-positive-current}
carries a boundarylessness hypothesis: positivity of $\dprimeddprime[u]$
compares two things, the edgewise second derivative and the weighted sum of
outgoing slopes, and only balancing removes the third, sign-indefinite term.
For weights equal to one it is \Cref{thm:ch11-graph-green}\,\cref{itm:ch11-graph:lap}
combined with \eqref{eq:graph-laplacian}.

\begin{theorem}
\label{thm:ch15-max-onedim}
Let $X$ be a connected one-dimensional balanced weighted rational polyhedral
complex, or more generally a connected one-dimensional boundaryless tropical
space, and let $u$ be a continuous psh function on $X$.  If $u$ attains its
supremum at a point of $X$, then $u$ is constant.
\end{theorem}

\begin{proof}
Let $M=\sup_Xu$, attained at some point, and let $A=\{u=M\}$, nonempty and
closed by continuity.  We show that $A$ is open; connectedness then gives
$A=X$.

Let $a\in A$.  If $a$ lies in the interior of an edge, then a neighbourhood
of $a$ in $X$ is an open interval of $\mathbf R$ with its standard tropical
structure, $u$ is convex there by \Cref{thm:ch15-convexity-model}, and
\Cref{prop:ch15-max-convex} gives that $u$ is constant near $a$.

Let $a=p$ be a vertex.  On each edge $\sigma\in T_pX$ the restriction of $u$
to the interior is convex, again by \Cref{thm:ch15-convexity-model}, and $u$
is continuous up to $p$; hence the outgoing slope
\[
 d_\sigma u(p)=\lim_{s\downarrow0}\frac{u(s)-u(p)}{s}
\]
exists in $[-\infty,+\infty)$, the difference quotient of a convex function
being nondecreasing in $s$.  Since $u\le M=u(p)$ on $X$, every difference
quotient is $\le0$, so
\begin{equation}
\label{eq:ch15-max-slopes}
 d_\sigma u(p)\le0\qquad\text{for every }\sigma\in T_pX .
\end{equation}

On the other hand $u$ is psh and continuous, hence locally integrable, so
$\mu:=\dprimeddprime[u]\ge0$ by \Cref{prop:cld-psh-positive-current}.  We
claim that
\begin{equation}
\label{eq:ch15-max-atom}
 \sum_{\sigma\in T_pX}m_\sigma\,d_\sigma u(p)\ \ge\ 0 ,
\end{equation}
the sum being taken in $[-\infty,+\infty)$.  For edgewise $C^2$ functions
this is \Cref{lem:ch15-vertex-atom} together with positivity of $\mu$; for a
merely edgewise convex $u$ we argue with a family of cutoffs, since the
outgoing slope may a priori be $-\infty$.

Fix a Euclidean norm on $N_{\mathbf R}$ and a smooth
$\chi:\mathbf R_{\ge0}\to[0,1]$ with $\chi\equiv1$ on $[0,1]$,
$\chi\equiv0$ on $[2,\infty)$.  For $\eta>0$ let $\Psi_\eta(x)=\chi(|x|/\eta)$
and $\varphi_\eta=\Psi_\eta|_X$, a smooth compactly supported function on $X$
which is identically $1$ near $p$; for $\eta$ small its support meets no
vertex other than $p$.  On the edge $\sigma$, identified with $[0,\ell)$ by
lattice arclength, one has $\varphi_\eta=\chi\bigl(s|u_{\sigma/p}|/\eta\bigr)$,
so with $\eta_\sigma=\eta/|u_{\sigma/p}|$ the function
$\varphi_{\eta,\sigma}$ equals $1$ on $[0,\eta_\sigma]$ and $0$ beyond
$2\eta_\sigma$.  Write $u'_\sigma$ for the right derivative of $u$ along
$\sigma$, a nondecreasing real function on $(0,\ell)$ with
$u'_\sigma(0^+)=d_\sigma u(p)$, and $\rd u'_\sigma\ge0$ for the associated
Stieltjes measure.  Since $\varphi_{\eta,\sigma}''$ vanishes on
$[0,\eta_\sigma]$, two integrations by parts on $[\eta_\sigma,\ell)$ ---
the second in the Stieltjes sense, legitimate because $u$ is convex hence
locally Lipschitz on $(0,\ell)$ --- give
\[
 \begin{aligned}
 \int_0^{\ell}u\,\varphi_{\eta,\sigma}''\,\rd s
 &=u'_\sigma(\eta_\sigma)
   +\int_{\eta_\sigma}^{2\eta_\sigma}
        \varphi_{\eta,\sigma}\,\rd u'_\sigma
 \\
 &\le u'_\sigma(\eta_\sigma)
   +\bigl(u'_\sigma(2\eta_\sigma)-u'_\sigma(\eta_\sigma)\bigr)
  =u'_\sigma(2\eta_\sigma),
 \end{aligned}
\]
using $0\le\varphi_{\eta,\sigma}\le1$ and $\rd u'_\sigma\ge0$.  No term
involving $u(p)$ arises, because $\varphi_\eta$ is constant near $p$; this
is the same cancellation as in \Cref{lem:ch15-vertex-atom}, obtained here by
the choice of cutoff instead of by balancing.  Summing over the edges with
the weights,
\[
 0\ \le\ \langle\mu,\varphi_\eta\rangle
 =\sum_{\sigma}m_\sigma\int_0^\ell u\,\varphi_{\eta,\sigma}''\,\rd s
 \ \le\ \sum_\sigma m_\sigma\,u'_\sigma(2\eta_\sigma).
\]
As $\eta\downarrow0$ the right-hand side decreases to
$\sum_\sigma m_\sigma\,d_\sigma u(p)$ in $[-\infty,+\infty)$, because each
$u'_\sigma$ is nondecreasing.  A nonincreasing sequence of nonnegative
numbers has nonnegative limit, which is \eqref{eq:ch15-max-atom}.

Since all $m_\sigma>0$, \eqref{eq:ch15-max-slopes} and
\eqref{eq:ch15-max-atom} together force $d_\sigma u(p)=0$ for every
$\sigma$; in particular no slope is $-\infty$.

Fix $\sigma$ and identify it near $p$ with $[0,\ell)$.  The function
$s\mapsto u(s)$ is convex with right derivative $0$ at $s=0$, hence
nondecreasing on $[0,\ell)$, so $u\ge u(p)=M$ there; combined with $u\le M$
this gives $u\equiv M$ on $\sigma$ near $p$.  As $\sigma$ was arbitrary,
$u\equiv M$ on a neighbourhood of $p$, and $A$ is open.
\end{proof}

Balancing is used exactly once in this proof, and not where one might expect:
it enters through the hypothesis of
\Cref{prop:cld-psh-positive-current}, which converts psh-ness into positivity
of the current $\dprimeddprime[u]$.  The cutoff argument then needs nothing
further.  \Cref{ex:ch15-unbalanced-max} shows that this single step is the
whole content: on an unbalanced star there is a smooth function with positive
$\dprimeddprime$ as a superform whose current is negative, and the maximum
principle fails for it.

\begin{example}[Failure without balancing]
\label{ex:ch15-unbalanced-max}
Let $N=\mathbf Z^2$ and let
$X=\mathbf R_{\ge0}e_1\cup\mathbf R_{\ge0}e_2$ with weight one on each ray,
the unbalanced fan of \Cref{ex:ch10-unbalanced}: its defect at the origin is
$\mathbf N(0)=-(e_1+e_2)\neq0$.  Let
\[
 u=-(x_1+x_2)\big|_X ,
\]
the restriction of an affine function on $\mathbf R^2$.  Then:
\begin{enumerate}
\item\label{itm:ch15-unbal-smooth} $u$ is smooth on $X$ and
  $\dprimeddprime u=0$, so $u$ is smooth psh in the sense of
  \Cref{def:ch15-hierarchy};
\item\label{itm:ch15-unbal-max} $u(0)=0$ and $u=-t<0$ at the point $te_i$
  for $t>0$, so $u$ attains a \emph{strict} global maximum at the origin and
  is not locally constant there;
\item\label{itm:ch15-unbal-current} the current $\dprimeddprime[u]$ equals
  $-2\delta_0$, hence is \emph{not} positive.
\end{enumerate}
For \cref{itm:ch15-unbal-current}, take a test function $\varphi$ supported
near the origin.  Each ray is parametrized by lattice arclength $t$, and
$u(te_i)=-t$, so
\[
 \langle\dprimeddprime[u],\varphi\rangle
 =\sum_{i=1,2}\int_0^\infty(-t)\,\varphi_i''(t)\,\rd t
 =\sum_{i=1,2}\bigl(-\varphi(0)\bigr)
 =-2\varphi(0),
\]
by the same two integrations by parts as in
\Cref{lem:ch15-pl-constant}.  Equivalently, by
\Cref{lem:ch15-vertex-atom} the atom is $\sum_im_i\,d_{e_i}u(0)=-1-1=-2$,
the correction term $u(0)\langle \rd\Psi(0),\mathbf N(0)\rangle$ vanishing
here only because $u(0)=0$.

Three lessons.  First, the maximum principle genuinely needs balancing.
Second, on an unbalanced space ``smooth with positive $\dprimeddprime$'' and
``positive $\dprimeddprime$ as a current'' are different conditions: item
\cref{itm:ch15-unbal-smooth} and item \cref{itm:ch15-unbal-current} contradict
each other under the identification $\dprimeddprime[u]=[\dprimeddprime u]$,
which is exactly the identification that
\Cref{prop:current-differential-calculus} makes only for boundaryless
spaces.  Third, the failure is quantitative: the missing mass $-2$ is the
pairing of the defect with the differential of $u$, which is the
one-dimensional case of the defect Stokes formula
\Cref{thm:ch11-defect-stokes}.
\end{example}

\begin{example}[Failure at an analytic or polyhedral boundary]
\label{ex:ch15-boundary-max}
Let $\Gamma=[0,\ell]$ be a segment, regarded as a one-dimensional weighted
complex with weight one; its defect is nonzero at each endpoint, so $\Gamma$
has boundary.  Let $u(s)=s$, a smooth convex function.  Then $u$ attains its
maximum at the endpoint $s=\ell$ and is not constant, while by
\Cref{ex:ch11-graph-segment-check} the \emph{graph Laplacian} is
\[
 \Delta_\Gamma u=\delta_0-\delta_\ell ,
\]
which is not a positive measure; the negative atom sits exactly at the point
where the maximum is attained.  The current $\dprimeddprime[u]$ is not even a
measure.  By \Cref{lem:ch15-vertex-atom}, applied at both endpoints with
$\mathbf N(0)=-1$ and $\mathbf N(\ell)=+1$,
\[
 \langle\dprimeddprime[u],\varphi\rangle
 =\varphi(0)-\varphi(\ell)+\ell\,\varphi'(\ell) ,
\]
the last term being the boundary correction
$u(\ell)\langle\rd\Psi(\ell),\mathbf N(\ell)\rangle$, which vanished in the
previous example only because $u$ vanished at the defect point.  It is not
positive either: for $\varphi\geq0$ equal to $1$ near $\ell$ the value is
$-1$.  This is the polyhedral shadow of the
following analytic situation: on the closed unit disc $\cM(K\{T\})$, whose
analytic boundary is the single Gauss point $\zeta_{0,1}$, the continuous
function $\log\max\{|T|,r\}$ is psh, attains its maximum $0$ precisely on
$\{|T|=1\}$, and is not constant.  It was the same endpoint term that had to
be excluded in the proof of \Cref{lem:ch15-pl-constant}, which is why that
computation was carried out on the open disc.
\end{example}

\begin{warning}{Do not extrapolate to higher dimension}
\Cref{thm:ch15-max-onedim} is a one-dimensional theorem, and its proof is
one-dimensional in an essential way: it converts positivity of a current into
an inequality between finitely many numbers, the outgoing slopes at a vertex.
In dimension $\ge2$ there is no such finite list, and no proof of a maximum
principle for CLD-psh functions is known.  What is available is either the
curve-test class of \Cref{sec:ch15-classical}, where a maximum principle is
proved by a genuinely different argument, or the local compact-sublevel barrier
\Cref{prop:ch15-exhaustion}, which substitutes for a maximum principle in
the estimates of \Cref{chap:cld-bedford-taylor} without proving one.
\end{warning}

\section{Pullbacks, scalar extension, and functoriality}
\label{sec:ch15-pullback}

\begin{proposition}
\label{prop:ch15-pullback-stagewise}
Let $g:S'\to S$ be a morphism of affine $G$-tropical spaces --- in
particular, a morphism of analytic spaces --- and let $u$ be psh on $S$.
Then $g^*u=u\circ g$ is psh on $S'$; and if
$u\in\operatorname{Psh}_{S,\alpha}$ then
$g^*u\in\operatorname{Psh}_{S',\alpha}$.  In particular the pullback of a
locally approachable function is locally approachable.
\end{proposition}

\begin{proof}
By \Cref{prop:ch15-hierarchy-stabilizes} it suffices to prove the stagewise
assertion, and we argue by transfinite induction on $\alpha$.

For $\alpha=0$: if $u$ is smooth with $\dprimeddprime u\ge0$, then $g^*u$ is
smooth by \Cref{prop:cld-form-pullback}, and
$\dprimeddprime(g^*u)=g^*(\dprimeddprime u)$ because pullback commutes with
both differentials; pullback preserves the positivity cones, so
$\dprimeddprime(g^*u)\ge0$.

For a limit ordinal, membership is local and the inductive hypothesis applies
on the preimage of a suitable neighbourhood.

For a successor $\alpha=\beta+1$: near a point of $S'$, choose a
neighbourhood $U$ of its image on which $u$ is the pointwise limit of a
decreasing sequence $(u_k)$ in $\operatorname{Psh}_{S,\beta}(U)$.  Then
$(g^*u_k)$ is a decreasing sequence on $g^{-1}(U)$, lies in
$\operatorname{Psh}_{S',\beta}$ by induction, and converges pointwise to
$g^*u$; hence $g^*u\in\operatorname{Psh}_{S',\beta+1}$.  This is
\cite[Proposition~8.1.5]{CLD25}.
\end{proof}

\begin{corollary}
\label{cor:ch15-pullback-cases}
Let $u$ be psh on an analytic space $X$, at stage $\alpha$.  Then:
\begin{enumerate}
\item\label{itm:ch15-pb-open} the restriction of $u$ to an open subset or to
  an analytic domain is psh at stage $\alpha$;
\item\label{itm:ch15-pb-closed} the restriction of $u$ to a closed analytic
  subspace $Z\subseteq X$ is psh at stage $\alpha$ on $Z$;
\item\label{itm:ch15-pb-ext} for a complete extension $K'/K$, the pullback of
  $u$ to $X_{K'}$ is psh at stage $\alpha$.
\end{enumerate}
\end{corollary}

\begin{proof}
Each case is \Cref{prop:ch15-pullback-stagewise} applied to the corresponding
morphism: an open immersion, a closed immersion, and the base-change morphism
$X_{K'}\to X$.
\end{proof}

The corresponding functoriality of the Poincar\'e--Lelong formula is worth
recording separately, since the naive statement is false.

\begin{proposition}
\label{prop:ch15-pl-pullback}
Let $g:X'\to X$ be a morphism of boundaryless pure-dimensional analytic
spaces and let $f$ be a regular meromorphic function on $X$ such that $g^*f$
is again a regular meromorphic function on $X'$.  Then
\[
 \dprimeddprime[\log|g^*f|]=\delta_{g^*\divis(f)} .
\]
\end{proposition}

\begin{proof}
Apply \Cref{thm:cld-poincare-lelong} on $X'$ to the regular meromorphic
function $g^*f$, and identify $\divis(g^*f)=g^*\divis(f)$ by
\Cref{lem:ch15-divisor-basic}\,\cref{itm:ch15-div-pullback}.
\end{proof}

\begin{warning}{Currents do not pull back}
\Cref{prop:ch15-pl-pullback} is \emph{not} obtained by pulling back the
identity on $X$.  There is no pullback operation on currents in general
(\Cref{chap:cld-currents}); what is available is the pullback of the
\emph{function} $\log|f|$, which is again of the same type, plus the pullback
of the \emph{Cartier divisor}, which requires that $g^*f$ be a
non-zero-divisor.  If that hypothesis fails --- for instance if $g$ maps
$X'$ into $\{f=0\}$ --- then $\log|g^*f|\equiv-\infty$, the left-hand side is
not defined as a current, and the right-hand side has no meaning either.
The hypothesis is thus not removable, and the proposition is a statement
about two functorialities agreeing, not about one operation.
\end{warning}

\section{Curve-test, or classically psh, functions}
\label{sec:ch15-classical}

Gubler--Rabinoff introduce a notion patterned on the restriction-to-curves
definition in complex geometry.

The contrast with \Cref{sec:ch15-psh} is worth stating before the definition.
The CLD class is defined \emph{from below}, by generating from the smooth
psh functions under an operation; membership is therefore proved by
exhibiting approximants, and every property of the class is proved by the
minimality device \Cref{lem:ch15-minimality}.  The curve-test class is
defined \emph{from above}, by a family of necessary conditions; membership is
proved by verifying a test on every curve, and properties of the class are
inherited from the corresponding properties of subharmonic functions on
curves --- for which the primary source is Thuillier's thesis \cite{Thu05},
summarized in \Cref{sec:ch15-thuillier-curves}.  Neither definition is
obviously stronger in general, and the two closure
mechanisms are genuinely different: the first uses smooth approximation and
no field extension, the second uses arbitrary analytic extensions of the
ground field and no approximation at all.  One implication is a theorem
(\Cref{thm:cld-to-classical-psh}); the general converse is
\Cref{prob:frontier-classical-regularization}.  For continuous functions on
smooth proper curves, the current-positive bridge of
\Cref{thm:p063-current-psh-ordinary-regularization} supplies the relevant
local ordinary-smooth converse.

\begin{definition}[Classically psh]
An upper semicontinuous function
$u:X\to\mathbf R\cup\{-\infty\}$ is \emph{classically psh} if for every
analytic extension field $K'/K$, every $K'$-analytic curve $C$, and every
morphism $g:C\to X_{K'}$, the function
\[
                    u\circ\pi_{K'/K}\circ g
\]
is subharmonic on $C$, in the sense of \eqref{eq:ch15-thuillier-sh}, that is
of \cite[D\'efinition~3.1.5]{Thu05}.  It is enough to test connected smooth
strictly $K'$-analytic curves.
\end{definition}

\begin{theorem}\label{thm:cld-to-classical-psh}
Every CLD-psh function is classically psh.
\end{theorem}
\begin{proof}
CLD-psh functions are upper semicontinuous and stable under analytic pullback
and scalar extension.  It remains to work on a smooth analytic curve.  A
smooth CLD-psh function is subharmonic by
\cite[Theorem~4.7]{Wanner19}.  Subharmonic functions are stable under
decreasing limits \cite[Proposition~3.1.8]{Thu05}, so transfinite induction
through the CLD definition gives the result.  This argument is recorded in the introduction of
\cite{GR25Classical}.
\end{proof}

No general converse is asserted.  The definitions have different closure
mechanisms: CLD uses approximation from smooth Hessian-positive functions;
classical psh uses all scalar extensions and all analytic curve tests.

\begin{theorem}\label{thm:gr-pl-comparison}
Assume that $K$ is nontrivially valued.  Let $X$ be a good strictly
$K$-analytic space and let $h:X\to\mathbf R$ be an $\mathbf R$-piecewise
linear function.  Then $h$ is classically psh if and only if it is
semipositive, meaning that its residue line bundle on every reduction of a
germ is nef.
\end{theorem}
\begin{proof}
Assume first that $h$ is semipositive.  Semipositivity of an
$\mathbf R$-PL function is preserved by analytic scalar extension and
pullback.  On a smooth strictly analytic curve, the residue degree of the PL
metric at a vertex is the sum of the outgoing slopes; nefness of every
residue line bundle is therefore equivalent to subharmonicity.  Hence the
pullback of $h$ to every curve in the definition is subharmonic, so $h$ is
classically psh.

Conversely, assume that $h$ is classically psh and fix a germ $(X,x)$.  If
its residue line bundle $L_h(x)$ were not nef, there would be a proper curve
in the reduction of the germ on which $L_h(x)$ has negative degree.  The
local curve-lifting theorem, after an algebraically closed analytic scalar
extension, realizes this residual curve by a smooth strictly analytic curve
$C$ mapping to a sufficiently small neighborhood of $x$.  The slope formula
identifies the degree of $L_h(x)$ on the residual curve with the mass of the
Laplacian of $h|_C$ over the corresponding compact subcurve.  Classical
plurisubharmonicity makes $h|_C$ subharmonic, so that mass is nonnegative, a
contradiction.  Thus $L_h(x)$ is nef for every $x$, which is semipositivity.

The curve criterion and its lifting step are
\cite[Theorem~5.3 and Corollary~5.4]{GR25PL}; the assembled equivalence is
\cite[Theorem~4.3]{GR25Classical}.  The pointwise-closedness result
\cite[Theorem~5.6]{GR25PL} is not needed for this equivalence.  Goodness,
strictness, nontriviality of the valuation, and the PL hypothesis enter the
curve-lifting and residual-line-bundle steps and are therefore retained.
\end{proof}

\begin{proposition}
\label{prop:classical-psh-calculus}
Classically psh functions form a sheaf and a convex cone, are stable under
finite maxima and decreasing nets, scalar extension, and analytic pullback.
If $u$ attains a local maximum at an interior point, then $u$ is locally
constant there.  A real-valued $h$ is \emph{pluriharmonic} in this theory
when both $h$ and $-h$ are classically psh.
\end{proposition}
\begin{proof}
Subharmonic functions on analytic curves form a sheaf and a convex cone and
are stable under finite maxima and decreasing nets; these are
\cite[Corollaire~3.1.13 and Proposition~3.1.8]{Thu05}, and the last of them
is the stability invoked in \Cref{thm:ch15-thuillier-approximation}.
Pulling each asserted
operation back after every scalar extension therefore proves the first four
claims.  Scalar extension is built into the definition.  For analytic
pullback, compose a test curve with the given analytic morphism; the result
is again a test curve for the original function.  These arguments are
formalized in \cite[Proposition~4.2]{GR25Classical}.

Suppose that $u$ attains a local maximum at an interior point $x$.  Shrink to
a connected boundaryless neighborhood on which this is a global maximum.
Every nearby point can, after analytic scalar extension, be joined to $x$ by
a connected overconvergent analytic curve.  The pullback of $u$ to that curve
is subharmonic and attains an interior maximum; the one-dimensional maximum
principle makes it constant on the curve.  Hence $u$ is constant near $x$.
The curve-connectivity theorem and this argument are
\cite[Theorems~5.4 and 6.1]{GR25Classical}.

The final sentence in the statement is a definition.  Since the preceding
properties apply to both $h$ and $-h$, pluriharmonic functions are
intrinsically stable under pullback and scalar extension.  Vanishing of the
current $\dprimeddprime[h]$ alone does not supply the required curve tests.
\end{proof}

\begin{proposition}[Pointwise recovery of classical pluriharmonicity]
\label{prop:ch15-pointwise-pluriharmonic-closure}
Let $X$ be a $K$-analytic space and let $h:X\to\mathbf R$ be continuous.
If there is a sequence $(h_j)_{j\geq1}$ of classically pluriharmonic
functions on $X$ such that $h_j(x)\to h(x)$ for every $x\in X$, then $h$ is
classically pluriharmonic.  Moreover, the convergence is locally uniform.
\end{proposition}
\begin{proof}
This is \cite[Theorem~7.6]{GR25Classical}.  Its local proof first reduces to
a quasicompact neighbourhood.  There the space of classically
pluriharmonic functions is finite dimensional, and restriction to the finite
union of the Shilov boundaries of an affinoid cover is injective.  Pointwise
convergence on that finite set therefore converges inside this finite
dimensional space; the maximum principle identifies the resulting norm with
the sup norm.  Thus $(h_j)$ converges uniformly to a pluriharmonic function,
which agrees pointwise with $h$.
\end{proof}

\begin{remark}[A recovery criterion, not a current-regularity theorem]
\label{rem:ch15-pointwise-ph-boundary}
The proposition requires a sequence of functions already known to pass all
curve tests.  In particular, it does not turn a continuous equality
$\dprimeddprime[h]=0$ of ordinary CLD currents into such a sequence, nor
does it furnish ordinary CLD-smooth harmonic approximants.  It is therefore
an exact final step for a prospective recovery argument, rather than a
solution of the current-zero regularity problem.
\end{remark}

\begin{corollary}
\label{cor:ch15-gr-maximum}
Let $K$ be nontrivially valued and let $X$ be a good, separated, strictly
$K$-analytic space without boundary.  If $u$ is classically psh on $X$ and
attains a local maximum at a point of $X$, then $u$ is constant near that
point.  Consequently, if $X$ is connected and compact and $u$ is classically
psh on $X$, then $u$ is constant.
\end{corollary}

\begin{proof}
The first assertion is the second paragraph of the proof of the preceding
proposition, that is, \cite[Theorems~5.4 and 6.1]{GR25Classical}.  For the
second, an upper semicontinuous function on a compact space attains its
supremum; the first assertion makes the locus where the supremum is attained
open, and it is closed by upper semicontinuity, so by connectedness it is all
of $X$.
\end{proof}

\begin{corollary}
\label{cor:ch15-cld-curve-max}
Under the hypotheses of \Cref{cor:ch15-gr-maximum}, a CLD-psh function on $X$
attaining a local maximum is locally constant there.
\end{corollary}

\begin{proof}
Combine \Cref{thm:cld-to-classical-psh} with \Cref{cor:ch15-gr-maximum}.
\end{proof}

\begin{warning}{What this does and does not settle}
\Cref{cor:ch15-cld-curve-max} is a maximum principle for CLD-psh functions,
but only under the hypotheses that the curve-test theory requires: $K$
nontrivially valued, $X$ good, separated, strict and boundaryless.  These
are precisely the hypotheses of the local curve-lifting theorem
\cite[Theorem~5.3]{GR25PL}, which is the analytic engine of
\Cref{cor:ch15-gr-maximum}.  Outside them --- over a trivially valued field,
on a non-strict space, or on an abstract tropical space that is not the
tropicalization of anything --- no maximum principle for the CLD class is
known, which is the statement recorded in \cite[\S0.3.5]{CLD25}.
\Cref{thm:ch15-max-onedim} is unconditional but one-dimensional; the two
results overlap and neither contains the other.
\end{warning}

\subsection{Failure of regularization across a reducible node}

The next example shows that upper-semicontinuous regularization of a local
semipositive envelope can create a branch jump.  The obstacle is continuous,
the domain is a strict affinoid curve, and the failure concerns all three
targets: classical plurisubharmonicity, CLD-v2 plurisubharmonicity, and
continuity.

\begin{proposition}
\label{prop:ch15-nodal-local-envelope}
Let \(K\) be a complete nontrivially valued non-Archimedean field.  Choose
\(\varpi\in K\) with \(0<|\varpi|<1\), and put
\[
                  r:=|\varpi|,
          \qquad  a:=-\log r>0.
\]
Consider the reduced strict \(K\)-affinoid curve
\[
 U:=\mathcal M\bigl(K\langle x,y\rangle/(xy)\bigr),
 \qquad
 A:=V(y),\quad B:=V(x),\quad o:=V(x,y).
\]
Thus \(A\) and \(B\) are closed unit discs meeting only at \(o\).  Define
\(f:U\to\mathbf R\) by
\begin{equation}
\label{eq:ch15-nodal-obstacle}
 f(z):=
 \begin{cases}
  0,
       & z\in A,\\[2mm]
  -a\min\{r^{-1}|y(z)|,1\},
       & z\in B.
 \end{cases}
\end{equation}
Then \(f\) is continuous.  If
\[
 P_U^{\mathrm{raw}}f
 :=\sup\{g:g\text{ is locally Fubini--Study or semipositive
                    \(\mathbf R\)-PL on \(U\), and }g\leq f\},
\]
then
\begin{equation}
\label{eq:ch15-nodal-raw-envelope}
 P_U^{\mathrm{raw}}f(z)=
 \begin{cases}
  0,  &z\in A\setminus\{o\},\\
  -a, &z\in B,
 \end{cases}
\end{equation}
whereas
\begin{equation}
\label{eq:ch15-nodal-regularized-envelope}
 \bigl(P_U^{\mathrm{raw}}f\bigr)^*(z)=
 \begin{cases}
  0,  &z\in A,\\
  -a, &z\in B\setminus\{o\}.
 \end{cases}
\end{equation}
In particular, the upper-semicontinuous regularization is neither
classically psh nor CLD-v2 psh, and it is not continuous.
\end{proposition}

\begin{proof}
Since \(xy=0\), every point of \(U\) belongs to \(A\) or to \(B\), and
\(A\cap B=\{o\}\).  The two expressions in
\eqref{eq:ch15-nodal-obstacle} are continuous on their respective closed
branches and both take the value \(0\) at \(o\).  The finite closed-pasting
lemma therefore shows that \(f\) is continuous.

We first construct admissible minorants which detect the whole branch
\(A\setminus\{o\}\).  For \(m\geq1\), set, with the convention
\(\log0=-\infty\),
\begin{equation}
\label{eq:ch15-nodal-hm}
                  h_m:=\max\{m^{-1}\log|x|,-a\}.
\end{equation}
Although \(x\) vanishes identically on \(B\), the function \(h_m\) is a
real-valued \(\mathbf R\)-PL function.  Indeed, the two strict affinoid
domains
\[
 U_m^-:=\{|x|\leq r^m\},
 \qquad
 U_m^+:=\{|x|\geq r^m\}
\]
form a Laurent \(G\)-cover of \(U\), since \(r^m=|\varpi^m|\).  On \(U_m^-\)
one has
\[
                         h_m=-a=\log|\varpi|.
\]
On \(U_m^+\), the element \(x\) is an analytic unit: in the rational-domain
algebra its inverse is
\[
                  x^{-1}=\varpi^{-m}(\varpi^m/x).
\]
Consequently
\[
                         h_m=m^{-1}\log|x|
                         \qquad\text{on }U_m^+.
\]
These two \(\mathbf R\)-PL expressions agree on the overlap.  Notice in
particular that the zero locus of \(x\) is contained entirely in \(U_m^-\),
where the expression involving \(\log|x|\) has been replaced by the finite
constant \(-a\).

We next verify semipositivity, including at the node.  The normalization is
the finite surjective morphism
\[
 \nu:D_x\amalg D_y\longrightarrow U,
 \qquad
 D_x:=\mathcal M(K\langle x\rangle),
 \quad
 D_y:=\mathcal M(K\langle y\rangle).
\]
Its pullback is
\[
 h_m\circ\nu=
 \begin{cases}
  \max\{m^{-1}\log|x|,-a\},&\text{on }D_x,\\
  -a,&\text{on }D_y.
 \end{cases}
\]
On \(D_x\), this is subharmonic by the logarithm, positive-scaling, and
finite-maximum properties
\cite[Proposition~3.5(2),(4)]{GR25Classical}; on \(D_y\) it is constant and
hence subharmonic.  Subharmonicity descends under a finite surjective
morphism by \cite[Proposition~3.5(8)]{GR25Classical}.  Thus \(h_m\) is
subharmonic on the possibly singular curve \(U\).  On analytic curves,
subharmonicity agrees with classical plurisubharmonicity
\cite[\S4.1.2]{GR25Classical}.  Since \(h_m\) is real-valued
\(\mathbf R\)-PL, \cite[Theorem~4.3]{GR25Classical} now shows that \(h_m\)
is semipositive.

For \(z\in A\), one has \(|x(z)|\leq1\), so \(h_m(z)\leq0=f(z)\).
On \(B\), one has \(x=0\), hence \(h_m=-a\), while
\(-a\leq f\leq0\).  Therefore every \(h_m\) is admissible.  The constant
function \(-a=\log|\varpi|\) is also an admissible semipositive
\(\mathbf R\)-PL function.

It remains to show that no admissible function can exceed \(-a\) anywhere
on \(B\).  Put
\[
 D_r:=\{z\in B:|y(z)|\leq r\}
       \simeq\mathcal M\bigl(K\langle r^{-1}y\rangle\bigr).
\]
This closed disc has the unique Shilov point \(\eta_r\), namely the
\(r\)-Gauss point
\begin{equation}
\label{eq:ch15-nodal-r-gauss}
 \left|\sum_{n\geq0}c_ny^n\right|_{\eta_r}
       =\max_{n\geq0}|c_n|r^n.
\end{equation}
For completeness, after writing \(y=\varpi T\), the algebra of \(D_r\)
becomes \(K\langle T\rangle\).  The Gauss norm of every element of
\(K\langle T\rangle\) is its supremum norm on the closed unit disc and is
its value at the Gauss point.  Hence the singleton consisting of that point
is already a boundary for every analytic function.  Since the Shilov
boundary is the smallest nonempty closed boundary, it is precisely that
singleton.  Transporting it by \(y=\varpi T\) gives \(\eta_r\).

Every function occurring in the defining supremum is classically psh.  For
semipositive \(\mathbf R\)-PL functions this is
\cite[Theorem~4.3]{GR25Classical}.  Locally Fubini--Study functions are,
locally, positive scalar multiples of finite maxima of logarithms of
analytic functions, up to additive constants; hence their classical
plurisubharmonicity follows from
\cite[Proposition~4.2(1)--(3)]{GR25Classical}.  Classical psh functions are
functorial under analytic pullback, so an admissible \(g\) restricts to a
classically psh function on \(D_r\).  The affinoid maximum principle
\cite[Theorem~6.4]{GR25Classical} and the uniqueness of the Shilov point give
\begin{equation}
\label{eq:ch15-nodal-disc-maximum}
             g(z)\leq g(\eta_r)
                    \leq f(\eta_r)
                    =-a
             \qquad (z\in D_r).
\end{equation}
On \(B\setminus D_r=\{|y|>r\}\), the obstacle itself equals \(-a\).
Thus \(g\leq-a\) on all of \(B\).  Since the constant \(-a\) is admissible,
\begin{equation}
\label{eq:ch15-nodal-raw-on-b}
                         P_U^{\mathrm{raw}}f=-a
                         \quad\text{on }B.
\end{equation}

Now fix \(z\in A\setminus\{o\}\).  Then \(|x(z)|>0\), and hence
\[
 \lim_{m\to\infty}h_m(z)
 =\lim_{m\to\infty}
       \max\{m^{-1}\log|x(z)|,-a\}
 =0.
\]
Since every candidate is bounded above by \(f=0\) on \(A\), the raw
envelope equals \(0\) at \(z\).  Together with
\eqref{eq:ch15-nodal-raw-on-b}, this proves
\eqref{eq:ch15-nodal-raw-envelope}.

The sets
\[
 A\setminus\{o\}=\{|x|>0\},
 \qquad
 B\setminus\{o\}=\{|y|>0\}
\]
are open in \(U\).  Hence the raw envelope is locally constant away from
\(o\), and its upper-semicontinuous regularization has the same values
there.  Globally the raw envelope is at most \(0\).  On the other hand, the
rigid points
\[
                  z_n=(x=\varpi^n,y=0)\in A\setminus\{o\}
\]
converge to \(o\), and the raw envelope takes the value \(0\) at each
\(z_n\).  Therefore
\[
                 \bigl(P_U^{\mathrm{raw}}f\bigr)^*(o)=0,
\]
which proves \eqref{eq:ch15-nodal-regularized-envelope}.

Let \(u=(P_U^{\mathrm{raw}}f)^*\).  If \(u\) were classically psh, then its
restriction to \(D_r\) would be classically psh.  But
\[
                 \max_{D_r}u=u(o)=0,
 \qquad
                 u(\eta_r)=-a<0,
\]
contradicting the affinoid maximum principle
\cite[Theorem~6.4]{GR25Classical}.  Thus \(u\) is not classically psh.
Every CLD-v2 psh function is classically psh by
\Cref{thm:cld-to-classical-psh}; this is the implication recorded in
\cite[\S1.2]{GR25Classical}.  Hence \(u\) is not CLD-v2 psh either.
Finally, the points
\[
                  w_n=(x=0,y=\varpi^n)\in B\setminus\{o\}
\]
converge to \(o\), while \(u(w_n)=-a\) and \(u(o)=0\).  Therefore \(u\) is
not continuous.
\end{proof}

This obstruction is not confined to curves.  Taking its product with a
closed strict polydisc propagates the same maximum-principle contradiction
to every dimension at least two.  The point is to use the zero section in
both directions: pullback of one-dimensional competitors gives one
inequality for the raw envelope, while restriction of ambient competitors
gives the reverse inequality.  No interchange of pullback with
upper-semicontinuous regularization is needed.

% Quarantined P03.1 product-obstruction candidate, 30 August 2026.
% Not production verified and not included by the live manuscript.

\begin{proposition}
\label{prop:p031-nodal-product-envelope-obstruction-candidate}
Let \(K\) be a complete nontrivially valued non-Archimedean field.  Choose
\(\varpi\in K\) with \(0<|\varpi|<1\), and put
\[
                  r:=|\varpi|,
          \qquad  a:=-\log r>0.
\]
Consider the reduced strict \(K\)-affinoid curve
\[
 U:=\mathcal M\bigl(K\langle x,y\rangle/(xy)\bigr),
 \qquad
 A:=V(y),\quad B:=V(x),\quad o:=V(x,y),
\]
and the continuous function
\[
 f(z):=
 \begin{cases}
  0,
       & z\in A,\\[2mm]
  -a\min\{r^{-1}|y(z)|,1\},
       & z\in B.
 \end{cases}
\]
For an integer \(d\geq1\), let
\[
 \mathbb D^d:=\mathcal M\bigl(K\langle t_1,\ldots,t_d\rangle\bigr),
 \qquad
 X_d:=U\times_K\mathbb D^d,
\]
write \(\pi:X_d\to U\) for the first projection, let
\[
 e:=V(t_1,\ldots,t_d)\in\mathbb D^d(K),
 \qquad
 j:U\longrightarrow X_d,\quad z\longmapsto(z,e),
\]
and set \(F:=f\circ\pi\).  Let \(\mathcal C_{X_d}(F)\) be the set of all
locally Fubini--Study or semipositive \(\mathbf R\)-PL functions \(g\) on
\(X_d\) such that \(g\leq F\), and define
\[
            G:=P_{X_d}^{\mathrm{raw}}F
              :=\sup_{g\in\mathcal C_{X_d}(F)}g,
            \qquad E:=G^*.
\]
Let
\[
 D_r:=\{z\in B:|y(z)|\leq r\},
\]
and denote by \(\eta_r\) its \(r\)-Gauss point.  Then
\[
                         E(j(o))=0,
              \qquad E(j(\eta_r))\leq-a.
\]
Consequently, \(E\) is not classically psh and is not CLD-v2 psh on
\(X_d\).  Moreover, \(X_d\) is reduced and strict of pure dimension
\(d+1\); hence such failures occur in every analytic dimension at least
two.
\end{proposition}

% Quarantined P03.1 product-obstruction candidate, 30 August 2026.
% Not production verified and not included by the live manuscript.

\begin{proof}[Proof of
\cref{prop:p031-nodal-product-envelope-obstruction-candidate}]
The projection and section satisfy
\[
                         \pi\circ j=\operatorname{id}_U.
\]
Analytic pullback preserves locally Fubini--Study functions, real PL
functions, and semipositivity.  Thus, if \(h\) is admissible for
\(P_U^{\mathrm{raw}}f\), then \(h\circ\pi\) is admissible for \(G\), and
\[
                         G\circ j\geq P_U^{\mathrm{raw}}f.
\]
Conversely, if \(g\) is admissible for \(G\), then \(g\circ j\) is admissible
for \(P_U^{\mathrm{raw}}f\), because \(F\circ j=f\).  Taking suprema gives
\begin{equation}
\label{eq:p031-product-section-identity-candidate}
                 G\circ j=P_U^{\mathrm{raw}}f.
\end{equation}

By \cref{prop:ch15-nodal-local-envelope},
\begin{equation}
\label{eq:p031-product-base-values-candidate}
 P_U^{\mathrm{raw}}f=0\quad\text{on }A\setminus\{o\},
 \qquad
 P_U^{\mathrm{raw}}f=-a\quad\text{on }B.
\end{equation}
Every competitor is bounded above by \(F\), so \(G\leq F\).  Since \(F\) is
continuous, upper-semicontinuous regularization preserves this inequality:
for every \(q\in X_d\),
\[
 E(q)=\limsup_{q'\to q}G(q')
     \leq\limsup_{q'\to q}F(q')
     =F(q).
\]
In particular,
\begin{equation}
\label{eq:p031-product-obstacle-upper-candidate}
                              E\leq F.
\end{equation}

For \(m\geq1\), let \(z_m\in A\setminus\{o\}\) be the rigid point
\[
                       x(z_m)=\varpi^m,\qquad y(z_m)=0.
\]
Then \(z_m\to o\), and
\eqref{eq:p031-product-section-identity-candidate}--%
\eqref{eq:p031-product-base-values-candidate} give
\[
                          G(j(z_m))=0.
\]
It follows from the definition of \(E=G^*\) that \(E(j(o))\geq0\).
On the other hand, \eqref{eq:p031-product-obstacle-upper-candidate} and
\(F(j(o))=f(o)=0\) give \(E(j(o))\leq0\).  Hence
\begin{equation}
\label{eq:p031-product-node-value-candidate}
                              E(j(o))=0.
\end{equation}

The obstacle has \(f(\eta_r)=-a\).  Therefore
\eqref{eq:p031-product-obstacle-upper-candidate} yields
\begin{equation}
\label{eq:p031-product-gauss-upper-candidate}
                         E(j(\eta_r))\leq-a.
\end{equation}

Suppose that \(E\) were classically psh on \(X_d\).  This property is
preserved by analytic pullback and restriction to analytic domains.
Consequently, \(E\circ j\) would be classically psh on the strict closed
disc \(D_r\).  This disc has the singleton Shilov boundary
\(\{\eta_r\}\), so the affinoid maximum principle would imply
\[
                         E(j(o))\leq E(j(\eta_r)).
\]
This contradicts \eqref{eq:p031-product-node-value-candidate} and
\eqref{eq:p031-product-gauss-upper-candidate}, since \(a>0\).  Thus \(E\)
is not classically psh.  Every CLD-v2 psh function is classically psh, so
\(E\) is not CLD-v2 psh either.

Finally,
\[
 X_d=\mathcal M\bigl(
 K\langle x,y,t_1,\ldots,t_d\rangle/(xy)\bigr).
\]
The ideals \((x)\) and \((y)\) are prime, their intersection is \((xy)\),
and their quotients are Tate algebras in \(d+1\) variables.  Thus \((xy)\)
is radical, its two irreducible components both have dimension \(d+1\),
and \(X_d\) is reduced and pure-dimensional.  It is strict by construction.
\end{proof}


The exact statement and proof passed an independent exact-byte hostile audit
and then a dependency-closed production verification using
\texttt{gpt-5.6-sol} at \texttt{xhigh}; the final verdict was
\texttt{correct}, with zero critical errors and zero gaps.  The conclusion
deliberately does not assert a formula for the raw envelope away from the
zero section and does not assert continuity or discontinuity of the
higher-dimensional regularized envelope.

The mechanism in both propositions is reducibility at the rigid node: the
obstacle pins every admissible minorant to the negative boundary value on one
branch, whereas minorants on the other branch approach zero away from the
node.  Thus
\Cref{prop:ch15-nodal-local-envelope,prop:p031-nodal-product-envelope-obstruction-candidate}
disprove an unrestricted local envelope theorem in every positive analytic
dimension.  The next result shows that this branching obstruction is sharp
in dimension one after the purely inseparable and nilpotent defects are
removed: bijectivity of the perfected normalization restores the classical
psh envelope property, including at relative-boundary points.  Neither the
CLD-v2 conclusion nor continuity of the envelope is inferred.

\begin{lemma}
\label{lem:p031-rigsmooth-curve-envelope}
Let \(K\) be a complete nontrivially valued non-Archimedean field, and let
\(C\) be a separated, rig-smooth, strict \(K\)-analytic curve.  Let
\((u_i)_{i\in I}\) be a nonempty family of classically plurisubharmonic
functions on \(C\) which is locally bounded above.  Then
\[
                 u:=\left(\sup_{i\in I}u_i\right)^*
\]
is classically plurisubharmonic on \(C\).
\end{lemma}

\begin{theorem}
\label{thm:p031-perfected-unibranch-curve-envelope}
Let \(k\) be a complete nontrivially valued non-Archimedean field, and let
\(U\) be a reduced, separated, strict \(k\)-analytic curve.  If
\(\operatorname{char}k=0\), put \(k^\dagger=k\).  If
\(\operatorname{char}k=p>0\), let \(k^\dagger\) be the completion of the
perfect closure of \(k\).  Set
\[
 U^\dagger:=U_{k^\dagger},
 \qquad Z:=(U^\dagger)_{\mathrm{red}},
\]
and let
\[
                   \nu:\widetilde Z\longrightarrow Z
\]
be the normalization.  Assume that \(\nu\) is bijective on underlying
Berkovich spaces.

If \((\varphi_i)_{i\in I}\) is a nonempty family of classically
plurisubharmonic functions on \(U\) which is locally bounded above, then
\begin{equation}
\label{eq:p031-perfected-unibranch-envelope}
             \left(\sup_{i\in I}\varphi_i\right)^*
\end{equation}
is classically plurisubharmonic on \(U\).

In particular, when \(k\) is perfect, the conclusion holds if \(U\) is
normal, and more generally if the normalization of \(U\) is bijective.
Thus it holds for reduced analytically unibranch strict analytic curves in
the explicit normalization-fibre sense.
\end{theorem}

\begin{corollary}
\label{cor:p031-local-fs-envelope-on-unibranch-curves}
Under the hypotheses of
\Cref{thm:p031-perfected-unibranch-curve-envelope}, let
\(f:U\to\mathbf R\cup\{-\infty\}\) be upper semicontinuous and set
\[
 \mathcal A_f:=
 \left\{
 h:\begin{array}{l}
 h\text{ is locally Fubini--Study or semipositive real PL on }U,\\
 h\leq f
 \end{array}
 \right\}.
\]
With the convention that the supremum of the empty family is \(-\infty\),
the function
\begin{equation}
\label{eq:p031-local-fs-envelope-on-unibranch-curves}
              \left(\sup_{h\in\mathcal A_f}h\right)^*
\end{equation}
is classically plurisubharmonic on \(U\).
The assertion includes points of the relative boundary of \(U\).
\end{corollary}


% The frozen proof headings already contain the words ``Proof of''.  This
% local style adjustment suppresses ntheorem's redundant outer
% ``Proof (...)'' wrapper without changing any verified packet byte.
\begingroup
\makeatletter
\let\pzeroonethree@original@nonumberplain\th@nonumberplain
\def\th@nonumberplain{%
  \pzeroonethree@original@nonumberplain
  \def\@opargbegintheorem##1##2##3{%
    \item[\theorem@headerfont\hskip\labelsep ##3\theorem@separator]}}
\makeatother
\begin{proof}[Proof of \Cref{lem:p031-rigsmooth-curve-envelope}]
Put \(s:=\sup_i u_i\) and \(u:=s^*\).  Local boundedness from above implies
that \(u\) takes values in \(\mathbf R\cup\{-\infty\}\).  Upper
semicontinuous regularization is local on open subsets: if \(W\subset C\)
is open, then
\begin{equation}
\label{eq:p031-usc-open-restriction}
                         u|_W=(s|_W)^*.
\end{equation}
Indeed, the neighborhoods of a point \(x\in W\) that are contained in
\(W\) are cofinal in the neighborhoods of \(x\) in \(C\), and both sides
of \eqref{eq:p031-usc-open-restriction} equal the corresponding
\(\limsup_{y\to x}s(y)\).

We first work on a connected component \(W\) of
\(\operatorname{Int}(C)\).  It is a smooth, boundaryless, strict analytic
curve.  By \cite[\S4.1.2 and \S3.4.1]{GR25Classical}, every \(u_i|_W\)
is either subharmonic in Thuillier's sense or identically \(-\infty\) on
\(W\).  If \(u|_W\equiv-\infty\), it is subharmonic in the extended sense
of \cite[Definition~3.4]{GR25Classical}.  We may therefore suppose that it
is not identically \(-\infty\).

Let \(D\subset W\) be a strict affinoid domain, and let
\(q:D\to\mathbf R\) be harmonic with
\begin{equation}
\label{eq:p031-harmonic-boundary-majorant}
                              u\leq q
                              \quad\text{on }\partial D.
\end{equation}
Since \(u_i\leq s\leq u\), every \(u_i\) satisfies the same boundary
inequality.  Thuillier's defining comparison property gives
\[
                              u_i\leq q
                              \quad\text{on }D
\]
for every \(i\); the assertion is tautological when \(u_i\equiv-\infty\).
Hence \(s\leq q\) on \(D\).  The complement of the relative boundary of
\(D\) is its topological interior in \(W\): this is the combination of
the relative-interior identities in \cite[\S2.11(3),(4)]{GR25Classical},
using that \(W\) is boundaryless and separated.  Therefore, for
\(x\in D\setminus\partial D\), upper semicontinuous regularization and
continuity of \(q\) give
\[
                         u(x)=\limsup_{y\to x}s(y)
                              \leq q(x).
\]
On \(\partial D\) the same inequality is
\eqref{eq:p031-harmonic-boundary-majorant}.  Thus \(u\leq q\) on \(D\).
Thuillier's comparison criterion proves that \(u|_W\) is subharmonic.

We have proved that \(u\) is subharmonic on
\(\operatorname{Int}(C)\).  The boundary-extension property
\cite[Proposition~3.5(3)]{GR25Classical} makes \(u\) subharmonic on all of
\(C\), including its relative boundary.  Finally,
\cite[\S4.1.2]{GR25Classical} identifies subharmonicity and classical
plurisubharmonicity on analytic curves.  This proves the lemma.
\end{proof}

\begin{proof}[Proof of
\Cref{thm:p031-perfected-unibranch-curve-envelope}]
Write
\[
                  v:=\sup_{i\in I}\varphi_i,
                  \qquad u:=v^*.
\]
We first record a topological identity.  If \(a:X\to
\mathbf R\cup\{-\infty\}\) is locally bounded above and
\(q:Y\to X\) is a homeomorphism, then
\begin{equation}
\label{eq:p031-homeomorphism-usc}
                          a^*\circ q=(a\circ q)^*.
\end{equation}
This follows directly because \(q\) identifies the two neighborhood
filters in the \(\limsup\) definition of upper semicontinuous
regularization.  The identity also commutes with a pointwise supremum:
\((\sup_i\varphi_i)\circ q=\sup_i(\varphi_i\circ q)\).

Let
\[
 \pi:U^\dagger\longrightarrow U,
 \qquad
 \iota:Z\lhook\joinrel\longrightarrow U^\dagger
\]
be the scalar-extension projection and the reduction inclusion.  In
positive characteristic, \(\pi\) is a homeomorphism by
\cite[Lemma~3.7]{GR25Classical}; in characteristic zero it is the identity.
The map \(\iota\) is a homeomorphism by
\cite[Lemma~3.8]{GR25Classical}.  The normalization map \(\nu\) is a
topological quotient by \cite[Lemma~3.9(4)]{GR25Classical}; it is therefore
a homeomorphism under the bijectivity hypothesis of the theorem.
Repeated use of \eqref{eq:p031-homeomorphism-usc} yields
\begin{equation}
\label{eq:p031-upstairs-envelope-identity}
 u\circ\pi\circ\iota\circ\nu
 =
 \left(
   \sup_{i\in I}\varphi_i\circ\pi\circ\iota\circ\nu
 \right)^*.
\end{equation}

The field \(k^\dagger\) is perfect.  Since \(Z\) is a reduced strict
\(k^\dagger\)-analytic curve, its normalization \(\widetilde Z\) is
rig-smooth by \cite[Lemma~3.9(1)]{GR25Classical}.  Each function
\(\varphi_i\circ\pi\circ\iota\circ\nu\) is classically
plurisubharmonic: scalar extension is
\cite[Proposition~4.2(5)]{GR25Classical}, and analytic pullback is the
forward implication of \cite[Proposition~4.2(6)]{GR25Classical}.  The
family remains locally bounded above because all three underlying maps are
homeomorphisms.  Applying
\Cref{lem:p031-rigsmooth-curve-envelope} to
\eqref{eq:p031-upstairs-envelope-identity} shows that
\(u\circ\pi\circ\iota\circ\nu\) is classically plurisubharmonic on
\(\widetilde Z\).

The map \(\nu\) is finite and surjective.  The map \(\iota\) is a finite
closed immersion and is surjective on underlying points.  The finite
surjective descent implication in
\cite[Proposition~4.2(6)]{GR25Classical}, applied first to \(\nu\) and then
to \(\iota\), gives
\[
                         u\circ\pi
                         \quad\text{classically psh on }U^\dagger.
\]
The converse scalar-extension implication in
\cite[Proposition~4.2(5)]{GR25Classical} now shows that \(u\) is
classically psh on \(U\).  This proves
\eqref{eq:p031-perfected-unibranch-envelope}.

If \(k\) is perfect, then \(k^\dagger=k\) and \(Z=U\).  For a normal
curve the normalization is the identity.  More generally, if the
normalization has one point over every point, then it is bijective, which
is precisely the explicit analytic-unibranch condition used in the final
assertion.
\end{proof}

\begin{proof}[Proof of
\Cref{cor:p031-local-fs-envelope-on-unibranch-curves}]
A locally Fubini--Study function is classically plurisubharmonic by the
logarithm, positive-scaling, finite-maximum, and sheaf properties in
\cite[Proposition~4.2(1)--(3)]{GR25Classical}.  A semipositive real-PL
function is classically plurisubharmonic by
\cite[Theorem~4.3]{GR25Classical}.  Thus every member of
\(\mathcal A_f\) is classically psh.

The upper semicontinuous function \(f\), which never takes the value
\(+\infty\), is locally bounded above.  Hence the family
\(\mathcal A_f\) is locally bounded above.  If it is nonempty, apply
\Cref{thm:p031-perfected-unibranch-curve-envelope}.  If it is empty, the
regularized supremum is the constant function \(-\infty\), which is
classically psh by \cite[Definitions~3.4 and~4.1]{GR25Classical}.
Relative-boundary points were retained in
\Cref{lem:p031-rigsmooth-curve-envelope} through
\cite[Proposition~3.5(3)]{GR25Classical}.
\end{proof}

\endgroup

Pointwise bijectivity of that normalization is a one-dimensional condition,
not a product-stable substitute for geometric unibranchness.  The following
explicit surface counterexample has a curve factor in precisely the stated
perfected-unibranch range, but a separable residue-field extension at its
hidden node splits after taking a product with an open disc.

% Book-local P03.1 nonsplit-node product obstruction, 10 September 2026.
% It records an explicit counterexample; no priority claim is made.
\begin{proposition}[A pointwise-unibranch curve whose disc product has a bad envelope]
\label{prop:p031-nonsplit-node-open-disc-envelope-obstruction}
Let \(K=\mathbf Q_3\), choose \(s\) with \(s^2=2\), and put
\[
 L:=K(s),\qquad
 C:=\mathcal M\bigl(K\langle x,y\rangle/(x^2-2y^2)\bigr),
 \qquad B:=\{t:|t|<1\}.
\]
Then \(C\) is a reduced irreducible separated strict \(K\)-analytic curve
whose normalization is bijective on underlying Berkovich spaces.  Nevertheless,
on the good strict analytic surface \(X:=C\times_K B\) there is a finite
continuous function \(f:X\to\mathbf R\).  Let \(\mathcal A_f\) be the
collection of locally Fubini--Study or semipositive real-PL functions \(h\)
on \(X\) satisfying \(h\leq f\), and put
\[
 E:=\left(\sup_{h\in\mathcal A_f}h\right)^*,
\]
the family \(\mathcal A_f\) is nonempty but \(E\) is neither classically
plurisubharmonic nor CLD-v2 psh.
\end{proposition}

\begin{proof}
Write \(r:=|3|\) and \(\alpha:=-\log r>0\).  The element \(2\) is a
nonsquare in \(K\), and \(L/K\) is its unramified quadratic extension.  The
normalization of \(C\) is
\[
 \nu:\widetilde C:=\mathcal M(L\langle u\rangle)\longrightarrow C,
 \qquad y=u,\quad x=su.
 \tag{\(\heartsuit\)}
\]
It is finite, and it is an isomorphism over \(\{|y|>0\}\), since there
\(s=x/y\).  The remaining fibre is \(\mathcal M(L)\) over the unique
point \(o\) cut out by \(x=y=0\).  It is a singleton.  Thus \(\nu\) is
bijective on underlying Berkovich spaces; as \(K\) has characteristic zero,
this is exactly the normalization hypothesis in
\Cref{thm:p031-perfected-unibranch-curve-envelope}.

Put \(N:=\widetilde C\times_K B\), let \(z:=t/3\), and define
\[
 \widetilde f(u,t):=
 \max\left\{\log|s-z|,\ -\alpha\min\{|u|/r,1\}\right\}.
 \tag{\(\heartsuit_1\)}
\]
This is finite and continuous.  At \(u=0\), its second entry is zero and
\[
 \widetilde f(0,t)=\max\{\log|s-z|,0\}
                    =\log\max\{1,|z|\}.
 \tag{\(\heartsuit_2\)}
\]
The final expression is independent of a point of a normalization fibre:
if \(|z|\ne1\), the ultrametric inequality gives the equality immediately,
and if \(|z|=1\), both sides of \((\heartsuit_2)\) are zero even in the
cancellation case.  Away from \(u=0\), \(\nu\times\mathrm{id}_B\) is an
isomorphism.  Hence \(\widetilde f\) is constant on every fibre of the
finite surjection
\[
 \rho:=\nu\times\mathrm{id}_B:N\longrightarrow X.
\]
Finite surjective analytic maps are topological quotient maps, so
\(\widetilde f=f\circ\rho\) for one finite continuous function \(f\) on
\(X\).

For every \(m\ge1\), write uniquely
\[
 (s-z)^m=P_m(z)+sQ_m(z),\qquad P_m,Q_m\in K[z],
\]
and set
\[
 A_m:=yP_m(t/3)+xQ_m(t/3),
 \qquad
 h_m:=\frac1m\log\max\{|A_m|,|3|^m\}.
 \tag{\(\heartsuit_3\)}
\]
The \(A_m\) are global \(K\)-analytic functions on \(X\); the nonzero
constant \(3^m\) makes \(h_m\) finite.  Thus \(h_m\) is a global locally
Fubini--Study function, hence an allowed competitor type.  Pullback along
\(\rho\) gives
\[
 h_m\circ\rho=
 \max\left\{m^{-1}\log|u|+\log|s-z|,\ -\alpha\right\}.
 \tag{\(\heartsuit_4\)}
\]
Because \(|u|\le1\) on \(\widetilde C\), this is at most
\[
 \max\{\log|s-z|,-\alpha\}
 \leq\widetilde f.
\]
Hence \(h_m\leq f\).  The constant \(-\alpha\) also lies below \(f\), so
\(\mathcal A_f\) is nonempty.

Fix a root \(s_0\in L\) of \(s_0^2=2\).  The \(L\)-rigid point
\(b\in B\) with \(t=3s_0\) has completed residue field \(L\).  On the
closed unit \(L\)-disc, consider the two analytic maps to \(X\)
\[
 \gamma_\pm(u)=(\pm s_0u,u,3s_0).
\]
Their origins have the same image \(q=(o,b)\).  On the plus branch,
\[
 f\circ\gamma_+(u)=-\alpha\min\{|u|/r,1\}.
 \tag{\(\heartsuit_5\)}
\]
Let \(D_r=\{|u|\le r\}\) and let \(\eta_r\) be its Gauss point.  Every
\(g\in\mathcal A_f\) is classically psh, and so is its pullback to \(D_r\).
Since \(g(\gamma_+(\eta_r))\leq-\alpha\), the affinoid maximum principle
\cite[Theorem~6.4]{GR25Classical} gives
\[
              g\circ\gamma_+\leq-\alpha\quad\text{on }D_r.
\]
The constant competitor proves that the raw envelope equals \(-\alpha\) on
this disc, in particular at \(q\).

On the minus branch, \(|-s_0-s_0|=1\), and therefore
\[
 f\circ\gamma_-=0,
 \qquad
 h_m\circ\gamma_-=
 \max\{m^{-1}\log|u|,-\alpha\}.
 \tag{\(\heartsuit_6\)}
\]
For every \(u\ne0\), the latter tends to zero.  Thus the raw envelope is
zero at every point \(\gamma_-(u)\) with \(u\ne0\).  Its upper
semicontinuous regularization satisfies \(E\leq f\), because \(f\) is
upper semicontinuous.  Letting \(u=3^n\to0\) in the preceding equality now
gives
\[
 E(q)=0,
 \qquad E(\gamma_+(\eta_r))=-\alpha.
 \tag{\(\heartsuit_7\)}
\]
Indeed, the first equality follows from the approaching raw values and
\(f(q)=0\), and the second from the constant competitor and
\((\heartsuit_5)\).

If \(E\) were classically psh, then \(E\circ\gamma_+\) would be
subharmonic on \(D_r\).  Its value at the centre would be zero, whereas its
only Shilov-boundary value would be \(-\alpha\), contradicting the same
affinoid maximum principle.  Thus \(E\) is not classically psh.  Since every
CLD-v2 psh function is classically psh by \Cref{thm:cld-to-classical-psh}, it
is not CLD-v2 psh either.
\end{proof}


Thus the curve theorem cannot be promoted by merely taking a standard open
disc product.  It leaves untouched the separate normal-curve product range,
or any possible higher-dimensional theorem under a genuinely geometric
unibranch hypothesis; the present obstruction only rules out the weaker
pointwise normalization condition.

There is also a higher-dimensional criterion that isolates precisely one way
to avoid the raw-envelope regularization gap: require a common local modulus
for the actual competitor family.  It is a condition on that family, not a
consequence of normality, unibranchness, or continuity of the obstacle.

% Candidate only.  It does not retire any P03 atom.
\begin{proposition}[A modulus-controlled local envelope criterion]
\label{prop:p031-modulus-controlled-envelope-candidate}
Let $k$ be a complete non-Archimedean field, let $X$ be a $k$-analytic
space, and let $\mathcal F$ be a nonempty family of finite-valued
classically plurisubharmonic functions on $X$.  Put
\[
 s(x):=\sup_{h\in\mathcal F}h(x).
\]
Assume that $s(x)<+\infty$ for every $x\in X$ and that $\mathcal F$ has a
common local modulus: for every $x\in X$ and every $\varepsilon>0$, there is
an open neighbourhood $W$ of $x$ such that
\[
 |h(y)-h(x)|<\varepsilon
 \qquad(y\in W,\ h\in\mathcal F).
\]
Then $s$ is finite, continuous, and classically plurisubharmonic.  In
particular $s^*=s$.

Consequently, suppose in addition that $k$ is complete and nontrivially
valued and that $X$ is good and strictly $k$-analytic.  If
$f:X\to\mathbf R\cup\{-\infty\}$ is upper semicontinuous and its nonempty
local-Fubini--Study/semipositive-$\mathbf R$-PL minorant family
\[
 \mathcal A_f:=\{h:h\text{ is locally Fubini--Study or semipositive
 $\mathbf R$-PL on }X,\ h\leq f\}
\]
consists of finite-valued functions and has this common local modulus, then
the raw envelope $P_X^{\mathrm{raw}}f:=\sup_{h\in\mathcal A_f}h$ is itself
continuous and classically plurisubharmonic; hence
$(P_X^{\mathrm{raw}}f)^*=P_X^{\mathrm{raw}}f$.
\end{proposition}

% Candidate only.  The source interface is recorded in SOURCE_LEDGER.md.
\begin{proof}
Fix $x\in X$ and $\varepsilon>0$, and choose $W$ as in the common-modulus
hypothesis.  For $y\in W$ and $h\in\mathcal F$, we have
\[
 h(y)\leq h(x)+\varepsilon\leq s(x)+\varepsilon,
 \qquad
 h(y)\geq h(x)-\varepsilon.
\]
Taking suprema over $h$ in the two inequalities gives
\[
 s(x)-\varepsilon\leq s(y)\leq s(x)+\varepsilon.
\]
Thus $s$ is continuous at $x$, and hence it is finite and continuous on
$X$ by the hypotheses.

Let $\mathcal J$ be the directed set of nonempty finite subsets of
$\mathcal F$, ordered by inclusion, and put
\[
 s_J:=\max_{h\in J}h\qquad(J\in\mathcal J).
\]
By the cone property for classically psh functions,
\cite[Proposition~4.2(3)]{GR25Classical}, every $s_J$ is classically psh.
For every $z\in X$, the net $(s_J(z))_{J\in\mathcal J}$ converges to
$s(z)$.  Since $s$ is continuous, it is upper semicontinuous.  The
pointwise-net limit theorem \cite[Theorem~6.5]{GR25Classical} therefore
shows that $s$ is classically psh.  Its upper-semicontinuous regularization
is consequently $s$ itself.

For the final assertion, every member of $\mathcal A_f$ is classically psh:
this is \cite[Theorem~4.3]{GR25Classical} for a semipositive
$\mathbf R$-PL function, while locally Fubini--Study functions are locally
positive multiples of finite maxima of logarithms of analytic functions, up
to constants, and hence use \cite[Proposition~4.2(1)--(3)]{GR25Classical}.
The first part, applied to $\mathcal F=\mathcal A_f$, gives the stated
conclusion.
\end{proof}


The criterion gives neither CLD-v2 pshness nor a classification of domains
for which every obstacle family has a common modulus.  In particular it does
not replace the unresolved higher-dimensional part of
\Cref{prob:frontier-unibranch-local-envelope}.

There is a second higher-dimensional positive mechanism which needs no
common modulus.  On an ordinary open subset of analytic affine space, a
small closed polydisc tube around the graph of a test curve has a finite
relative-Gauss boundary.  The affinoid maximum principle controls every
competitor on that tube, while openness of the tube projection near the
zero section converts the common tube bound into an ambient-neighbourhood
bound for the raw supremum.  This proves all curve tests defined over the
ground field.  The remaining scalar-extension issue is stated explicitly:
upper-semicontinuous regularization need not commute formally with scalar
projection.

% Book-local P03.1 affine ground-field curve-test theorem,
% 12 September 2026.  No novelty or priority claim is made.

\begin{definition}
\label{def:p031-ground-field-curvewise-psh}
Let $K$ be a complete nontrivially valued non-Archimedean field and let
$X$ be a $K$-analytic space.  An upper-semicontinuous function
$v:X\to\mathbf R\cup\{-\infty\}$ is called
\emph{$K$-curvewise plurisubharmonic} if, for every smooth connected
strict $K$-analytic curve $C$ and every $K$-analytic morphism
$g:C\to X$, the function $v\circ g$ is subharmonic in the sense of
\cite[Definition~3.4]{GR25Classical}.  Equivalently, by
\cite[\S3.4.1]{GR25Classical}, its restriction to the relative interior of
$C$ satisfies Thuillier's harmonic comparison, with the identically
$-\infty$ case allowed.

This is deliberately weaker than classical plurisubharmonicity: the latter
also tests curves after every analytic extension of the ground field.
\end{definition}

\begin{lemma}
\label{lem:p031-open-pullback-usc}
Let $q:Y\to X$ be a continuous open surjection of topological spaces and
let $a:X\to\mathbf R\cup\{-\infty\}$ be locally bounded above.  Then
\[
                         a^*\circ q=(a\circ q)^*.
\]
\end{lemma}

\begin{theorem}
\label{thm:p031-affine-ground-field-perron}
Let $K$ be a complete nontrivially valued non-Archimedean field, let
$n\geq 1$, and let $V\subseteq\mathbf A_K^{n,\mathrm{an}}$ be an
ordinary open subset.  Let $(u_i)_{i\in I}$ be a nonempty locally
upper-bounded family of classically plurisubharmonic functions on $V$.
Then
\[
                   u:=\left(\sup_{i\in I}u_i\right)^*
\]
is $K$-curvewise plurisubharmonic on $V$.
\end{theorem}

\begin{corollary}
\label{cor:p031-scalar-compatibility-upgrade}
Keep the hypotheses and notation of
\Cref{thm:p031-affine-ground-field-perron}, and put
$s:=\sup_i u_i$.  Assume that for every analytic extension field
$K'/K$, with scalar projection
$
 \pi_{K'/K}:V_{K'}\to V,
$
one has
\begin{equation}
\label{eq:p031-scalar-compatibility}
 u\circ\pi_{K'/K}
       =\left(\sup_i u_i\circ\pi_{K'/K}\right)^*
       =(s\circ\pi_{K'/K})^*.
\end{equation}
Then $u$ is classically plurisubharmonic on $V$.

For a fixed extension $K'/K$, the identity
\eqref{eq:p031-scalar-compatibility} holds whenever
$
 \pi_{K'/K}:V_{K'}\to V
$
is open.  Thus openness of every scalar projection is a sufficient
topological hypothesis for the classical conclusion.
\end{corollary}

\begin{corollary}
\label{cor:p031-affine-obstacle-ground-field}
Let $f:V\to\mathbf R\cup\{-\infty\}$ be upper semicontinuous and let
\[
 \mathcal A_f=
 \{h:h\text{ is locally Fubini--Study or semipositive real PL on }V,
          \ h\leq f\}.
\]
With the convention $\sup\varnothing=-\infty$, the regularized envelope
\[
                       P_f:=\left(\sup_{h\in\mathcal A_f}h\right)^*
\]
is $K$-curvewise plurisubharmonic.  If the scalar-compatibility identity
\eqref{eq:p031-scalar-compatibility} holds for this family after every
analytic extension, then $P_f$ is classically plurisubharmonic.
\end{corollary}

\begin{proposition}
\label{prop:p031-scalar-openness-failure}
There are complete non-Archimedean extensions $K'/K$ for which the
scalar projection
\[
        \mathbf A_{K'}^{1,\mathrm{an}}
              \longrightarrow\mathbf A_K^{1,\mathrm{an}}
\]
is not open.  Consequently
\Cref{thm:p031-affine-ground-field-perron} does not, by topology alone,
upgrade to the unrestricted classical-psh assertion in P03.1.

The proposition is only a failure of this automatic upgrade.  It neither
produces a normal higher-dimensional counterexample to P03.1 nor proves
that scalar compatibility fails for a particular psh envelope.
\end{proposition}

% Book-local P03.1 affine ground-field curve-test proof,
% 12 September 2026.  No novelty or priority claim is made.

\begin{proof}[\Cref{lem:p031-open-pullback-usc}]
For $x\in X$, write the upper-semicontinuous regularization as the
neighborhood-filter limsup
\[
                  a^*(x)=\inf_{U\ni x}\sup_{z\in U}a(z).
\]
Fix $y\in Y$ and put $x=q(y)$.  If $N$ is a neighborhood of $y$,
then $q(N)$ is a neighborhood of $x$ because $q$ is open.
Conversely, if $U$ is a neighborhood of $x$, then $q^{-1}(U)$ is a
neighborhood of $y$ and, by surjectivity, $q(q^{-1}(U))=U$.  Hence the
sets $q(N)$, as $N$ runs over the neighborhoods of $y$, are cofinal
among the neighborhoods of $x$.  Therefore
\[
 \begin{split}
 (a\circ q)^*(y)
   &=\inf_{N\ni y}\sup_{z\in N}a(q(z))\\
   &=\inf_{N\ni y}\sup_{x'\in q(N)}a(x')
    =a^*(x),
 \end{split}
\]
which is the claimed identity.
\end{proof}

\begin{proof}[
\Cref{thm:p031-affine-ground-field-perron}]
Put $s:=\sup_i u_i$ and $u:=s^*$.  Local upper boundedness makes $u$
an upper-semicontinuous function with values in
$\mathbf R\cup\{-\infty\}$.

Let $C$ be a smooth connected strict $K$-analytic curve and let
$g:C\to V$ be a morphism.  We prove that $u\circ g$ is subharmonic.  It
is enough, by \cite[Proposition~3.5(3)]{GR25Classical}, to work on each
connected component of $\operatorname{Int}(C)$.  If $u\circ g$ is
identically $-\infty$ there, there is nothing to prove.  Otherwise let
$D$ be a strict affinoid domain in that component and let
$h:D\to\mathbf R$ be harmonic with
\begin{equation}
\label{eq:p031-ground-field-boundary-hypothesis}
                         u\circ g\leq h
                         \qquad\text{on }\partial D.
\end{equation}
We must prove the same inequality on $D$.

We first construct a small analytic tube around the graph of $g|_D$.
Write $z_1,\ldots,z_n$ for the affine coordinates and
$g_j:=z_j\circ g\in\mathcal O(D)$.  On
$D\times\mathbf A_K^{n,\mathrm{an}}$, the translation
\[
  \Psi(t,w_1,\ldots,w_n)
     =(t,g_1(t)+w_1,\ldots,g_n(t)+w_n)
\]
is an analytic automorphism.  The inverse image
$\Psi^{-1}(D\times V)$ is an ordinary open neighborhood of the zero
section.  Choose a sequence $\rho_m\in|K^\times|$ decreasing to zero.
Compactness shows that, for all sufficiently large $m$,
\begin{equation}
\label{eq:p031-ground-field-tube-inclusion}
 D\times\mathbb B^n_{\rho_m}
       \subseteq\Psi^{-1}(D\times V).
\end{equation}
Indeed, otherwise the nonempty nested compact sets
\[
 \bigl(D\times\mathbb B^n_{\rho_m}\bigr)
        \setminus\Psi^{-1}(D\times V)
\]
would have nonempty intersection, whereas their intersection is contained
in the zero section and the zero section lies in
$\Psi^{-1}(D\times V)$.

Fix a radius $\rho$ satisfying
\eqref{eq:p031-ground-field-tube-inclusion} and set
\[
 T_\rho:=D\times\mathbb B^n_\rho,
 \qquad
 p_\rho(t,w):=g(t)+w\in V,
 \qquad
 H_\rho(t,w):=h(t).
\]
Thus $T_\rho$ is a $K$-affinoid space.  Let
$\sigma_\rho:D\to T_\rho$ be the relative Gauss section.  We need two
elementary properties of this tube.

First,
\begin{equation}
\label{eq:p031-relative-gauss-boundary}
                  \sigma_\rho(\Gamma(D))
                  \quad\text{is a boundary of }T_\rho,
\end{equation}
where $\Gamma(D)=\partial D$ is the finite Shilov boundary of the strict
affinoid curve $D$, as recalled in
\cite[\S3.1(3)]{GR25Classical}.  To see this, write an analytic function on
the tube as $F=\sum_{\alpha}a_\alpha w^\alpha$, with
$a_\alpha\in\mathcal O(D)$.  The defining property of the relative
Gauss point gives
\[
 |F(\sigma_\rho(x))|
        =\max_\alpha |a_\alpha(x)|\rho^{|\alpha|}.
\]
For every point of the fiber over $x$, the value of $|F|$ is bounded
above by the right side.  Since every coefficient $a_\alpha$ attains its
supremum on $\Gamma(D)$, and the weighted coefficient norms tend to
zero, we get
\[
 \sup_{T_\rho}|F|
   =\max_{b\in\Gamma(D)}|F(\sigma_\rho(b))|.
\]
This identifies the displayed finite set as a boundary of $T_\rho$.
Consequently the Shilov boundary of $T_\rho$ is contained in that set.

Second, if $t_0\in D\setminus\partial D$, then $p_\rho$ is open on an
ordinary neighborhood of $(t_0,0)$.  In the $w$-coordinates, put
\[
 U_\rho^\circ
    :=\operatorname{Int}(D)\times\mathbb B_{\rho}^{n,\circ}
       \subseteq T_\rho,
\]
where the second factor is the open polydisc.  Before imposing the closed
tube, $p_\rho$ is, through the translation $\Psi$, the projection
\[
             D\times\mathbf A_K^{n,\mathrm{an}}
                     \longrightarrow\mathbf A_K^{n,\mathrm{an}}.
\]
The structural map of the smooth curve $\operatorname{Int}(D)$ is
quasi-smooth and boundaryless.  Base change, the translation $\Psi$, and
restriction to the ordinary open set $U_\rho^\circ$ preserve these two
properties.  Thus
\[
                   p_\rho|_{U_\rho^\circ}:U_\rho^\circ\longrightarrow V
\]
is quasi-smooth and boundaryless, and is open by
\cite[Corollary~6.2.5]{Ducros18Families}.  Since
$(t_0,0)\in U_\rho^\circ$, this proves the local assertion.

Return to the boundary comparison.  As $\rho\to0$, for every
$b\in\Gamma(D)$ one has
\begin{equation}
\label{eq:p031-gauss-convergence}
                 p_\rho(\sigma_\rho(b))\longrightarrow g(b).
\end{equation}
This follows directly from the Gauss formula: on every fixed large tube,
the value of a convergent power series at the relative Gauss point tends
to the value of its constant coefficient at the zero section.  The set
$\Gamma(D)$ is finite.  Upper semicontinuity of $u$, together with
\eqref{eq:p031-ground-field-boundary-hypothesis}, therefore shows that for every
$\varepsilon>0$ one may choose $\rho\in|K^\times|$ small enough that
\begin{equation}
\label{eq:p031-gauss-bound}
 u(p_\rho(\sigma_\rho(b)))
       \leq h(b)+\varepsilon
       \qquad(b\in\Gamma(D)).
\end{equation}

The harmonic function $h$ is pluriharmonic on $D$: this is
\cite[Definition~3.13 and \S4.1.2]{GR25Classical}.  Hence $H_\rho$ is
pluriharmonic by analytic pullback, using
\cite[Proposition~7.2(7)]{GR25Classical}.  For every $i$, the function
\[
                         u_i\circ p_\rho-H_\rho
\]
is consequently classically psh on $T_\rho$.  Since
$u_i\leq s\leq u$, equation \eqref{eq:p031-gauss-bound} bounds this
function by $\varepsilon$ on $\sigma_\rho(\Gamma(D))$.  The global
affinoid maximum principle \cite[Theorem~6.4]{GR25Classical}, together
with \eqref{eq:p031-relative-gauss-boundary}, gives
\[
                    u_i\circ p_\rho-H_\rho\leq\varepsilon
                    \qquad\text{on }T_\rho
\]
for every $i$, and hence
\begin{equation}
\label{eq:p031-raw-tube-bound}
                    s\circ p_\rho\leq H_\rho+\varepsilon
                    \qquad\text{on }T_\rho.
\end{equation}

Fix $t_0\in D\setminus\partial D$.  By the local openness proved above,
we may choose an ordinary neighborhood $N$ of $(t_0,0)$ in $T_\rho$
such that $p_\rho(N)$ is an ordinary open neighborhood of $g(t_0)$.
Shrinking $N$, continuity of $h$ also gives
$H_\rho<h(t_0)+\varepsilon$ on $N$.  It follows from
\eqref{eq:p031-raw-tube-bound} that
\[
                    s<h(t_0)+2\varepsilon
                    \qquad\text{on }p_\rho(N).
\]
The definition of usc regularization now yields
\[
                         u(g(t_0))\leq h(t_0)+2\varepsilon.
\]
Letting $\varepsilon\downarrow0$ proves $u\circ g\leq h$ at every
point of $D\setminus\partial D$; the boundary inequality was assumed in
\eqref{eq:p031-ground-field-boundary-hypothesis}.  Thuillier's comparison criterion
shows that $u\circ g$ is subharmonic on each component of
$\operatorname{Int}(C)$, and
\cite[Proposition~3.5(3)]{GR25Classical} extends it across the relative
boundary of $C$.  This is exactly the property in
\Cref{def:p031-ground-field-curvewise-psh}.
\end{proof}

\begin{proof}[
\Cref{cor:p031-scalar-compatibility-upgrade}]
Fix an analytic extension $K'/K$.  Pointwise supremum commutes with
pullback, so
\[
             s\circ\pi_{K'/K}=\sup_i(u_i\circ\pi_{K'/K}).
\]
The pulled-back functions are classically psh by
\cite[Proposition~4.2(5)]{GR25Classical} and remain locally upper bounded.
Applying \Cref{thm:p031-affine-ground-field-perron} over $K'$ shows
that $(s\circ\pi_{K'/K})^*$ is $K'$-curvewise psh.  Under
\eqref{eq:p031-scalar-compatibility}, this function equals
$u\circ\pi_{K'/K}$.  Since this holds for every $K'/K$, the smooth
strict curve test in \cite[Definition~4.1 and \S4.1.3]{GR25Classical}
shows that $u$ is classically psh.

If $\pi_{K'/K}$ is open, it is also continuous and surjective; the latter
two properties of scalar projection are recalled in
\cite[\S2.9]{GR25Classical}.  Apply
\Cref{lem:p031-open-pullback-usc} to $a=s$ and use $u=s^*$ to get
\[
 u\circ\pi_{K'/K}
   =s^*\circ\pi_{K'/K}
   =(s\circ\pi_{K'/K})^*,
\]
which is the desired identity for that extension.
\end{proof}

\begin{proof}[
\Cref{cor:p031-affine-obstacle-ground-field}]
A locally Fubini--Study function is classically psh by the logarithm and
sheaf properties, together with closure under positive linear combinations
and finite maxima; see
\cite[Proposition~4.2(1)--(3)]{GR25Classical}.  A semipositive real-PL
function is classically psh by \cite[Theorem~4.3]{GR25Classical}.  The
upper-semicontinuous obstacle $f$, which has no $+\infty$ value, is
locally bounded above, and therefore bounds the whole family
$\mathcal A_f$ locally.  If this family is nonempty, apply
\Cref{thm:p031-affine-ground-field-perron}; the scalar-compatible
conclusion follows from
\Cref{cor:p031-scalar-compatibility-upgrade}.  If the family is empty,
then $P_f=-\infty$, which is already classically psh by
\cite[Definitions~3.4 and~4.1]{GR25Classical}.
\end{proof}

\begin{proof}[
\Cref{prop:p031-scalar-openness-failure}]
Poineau--Turchetti consider an algebraically closed complete field $K$
which is not spherically complete, a maximally complete algebraically
closed extension $K'/K$, and a type--$4$ point of
$\mathbf A_K^{1,\mathrm{an}}$.  They identify its inverse image under
scalar projection as a positive-radius closed disc
\cite[Lemma~5.6.2]{PT21SchottkyI}; their
\cite[Remark~5.6.4]{PT21SchottkyI} records explicitly that the projection
is not open in general.  Thus the open-map hypothesis in
\Cref{lem:p031-open-pullback-usc} is not automatic, even for the
analytic affine line.  The final scope assertions follow because this
topological example does not supply a psh family for which the envelope
identity fails.
\end{proof}


The remaining equality has an exact fiberwise form.  Scalar projection is
a closed surjection with compact fibers, so the downstairs regularization
is the maximum, over the complete scalar fiber, of the regularization of
the pulled-back raw function.  The following results prove this formula,
characterize when topology alone gives pointwise equality, and keep the
psh and purely topological issues separate.

% Book-local P03.1 scalar-fiber maximum statement,
% 19 September 2026.  No novelty or priority claim is made.

\begin{lemma}
\label{lem:p031-perfect-pullback-usc}
Let $q:Y\to X$ be a continuous closed surjection of topological spaces,
and assume that every fiber of $q$ is compact.  Let
$a:X\to\mathbf R\cup\{-\infty\}$ be locally bounded above and put
\[
             b:=(a\circ q)^*.
\]
Then, for every $x\in X$,
\begin{equation}
\label{eq:p031-perfect-pullback-usc}
             a^*(x)=\max_{y\in q^{-1}(x)} b(y).
\end{equation}
Equivalently, for every $y\in Y$,
\begin{equation}
\label{eq:p031-fiber-maximum-pullback}
 a^*(q(y))
   =\max_{z\in q^{-1}(q(y))}(a\circ q)^*(z).
\end{equation}
Thus the possible failure of usc regularization to commute with pullback is
exactly the failure of the pulled regularization to attain the same value at
every point of a fiber.
\end{lemma}

\begin{proposition}
\label{prop:p031-open-map-characterization}
Let $q:Y\to X$ be a continuous surjection.  The following are equivalent.
\begin{enumerate}
\item The map $q$ is open.
\item For every bounded function $a:X\to\mathbf R$, one has
      $a^*\circ q=(a\circ q)^*$.
\item The same identity holds for every function $a:X\to\{0,1\}$.
\end{enumerate}
In particular, non-openness always produces an actual strict
regularization defect for bounded raw data; it is not merely a failure of
one proof of the pullback identity.
\end{proposition}

\begin{corollary}
\label{cor:p031-scalar-fiber-maximum}
Let $K'/K$ be an analytic extension of complete non-Archimedean fields,
let $X$ be a $K$-analytic space, and let
$
 \pi:X_{K'}\to X
$
be scalar projection.  If $a:X\to\mathbf R\cup\{-\infty\}$ is locally
bounded above, then
\begin{equation}
\label{eq:p031-scalar-fiber-maximum}
 a^*(x)
   =\max_{y\in\pi^{-1}(x)}(a\circ\pi)^*(y)
       \qquad(x\in X).
\end{equation}
The same formula holds after restricting $\pi$ to the inverse image of
an ordinary open subset of $X$.
\end{corollary}

\begin{corollary}
\label{cor:p031-fiber-constancy-upgrade}
Keep the notation and hypotheses of
\Cref{thm:p031-affine-ground-field-perron}.  Thus
\[
       s:=\sup_i u_i,
       \qquad
       u:=s^*.
\]
For an analytic extension $K'/K$, put
\[
       v_{K'}:=(s\circ\pi_{K'/K})^*
       \quad\text{on }V_{K'}.
\]
Then the scalar-compatibility identity
\begin{equation}
\label{eq:p031-fiber-constancy}
       u\circ\pi_{K'/K}=v_{K'}
\end{equation}
holds if and only if $v_{K'}$ is constant on every fiber of
$\pi_{K'/K}$.  More precisely,
\begin{equation}
\label{eq:p031-perron-fiber-maximum}
 u\circ\pi_{K'/K}(y)
   =\max_{z\in\pi_{K'/K}^{-1}(\pi_{K'/K}(y))}v_{K'}(z).
\end{equation}
Consequently, if $v_{K'}$ is fiberwise constant for every analytic
extension $K'/K$, then $u$ is classically plurisubharmonic.  The same
conclusion applies to the local-Fubini--Study/semipositive-real-PL obstacle
envelope in \Cref{cor:p031-affine-obstacle-ground-field}.
\end{corollary}

\begin{proposition}
\label{prop:p031-type-four-topological-defect}
There are algebraically closed complete non-Archimedean fields $K\subset
K'$, a type--$4$ point $x\in\mathbf A_K^{1,\mathrm{an}}$, a point
$y\in\pi_{K'/K}^{-1}(x)$, and a bounded function
\[
             a:\mathbf A_K^{1,\mathrm{an}}\longrightarrow\{0,1\}
\]
such that
\begin{equation}
\label{eq:p031-type-four-strict-defect}
 (a\circ\pi_{K'/K})^*(y)
       <a^*(x)
       =\max_{z\in\pi_{K'/K}^{-1}(x)}
                    (a\circ\pi_{K'/K})^*(z).
\end{equation}
This is a sharp topological obstruction to universal scalar compatibility.
The function $a$ is not asserted to be psh or to arise as a supremum of
psh competitors.  Hence the proposition is not a counterexample to P03.1.
\end{proposition}

% Book-local P03.1 scalar-fiber maximum proof,
% 19 September 2026.  No novelty or priority claim is made.

\begin{proof}[\Cref{lem:p031-perfect-pullback-usc}]
Write
\[
 a^*(x)=\inf_{U\ni x}\sup_U a,
 \qquad
 b(y)=\inf_{N\ni y}\sup_N(a\circ q),
\]
where the infima may be taken over open neighborhoods.  Fix $x\in X$.
For $y\in q^{-1}(x)$ and every neighborhood $U$ of $x$, the set
$q^{-1}(U)$ is a neighborhood of $y$, and surjectivity gives
\[
 \sup_{q^{-1}(U)}(a\circ q)=\sup_Ua.
\]
Taking first the infimum over the neighborhoods of $y$, and then the
infimum over $U$, yields
\begin{equation}
\label{eq:p031-fiber-easy-inequality}
                         b(y)\leq a^*(x).
\end{equation}

The function $b$ is upper semicontinuous.  Its restriction to the compact
fiber $F_x:=q^{-1}(x)$ therefore attains a maximum; denote that maximum by
$M$.  We prove $a^*(x)\leq M$.  Let $r\in\mathbf R$ satisfy $r>M$.
For each $y\in F_x$, the defining neighborhood formula for $b(y)<r$
gives an open neighborhood $N_y$ such that
\[
                         a\circ q<r\quad\text{on }N_y.
\]
Compactness supplies $y_1,\ldots,y_m\in F_x$ for which
$
 N:=N_{y_1}\cup\cdots\cup N_{y_m}
$
contains $F_x$.  Since $q$ is closed,
\[
                         U:=X\setminus q(Y\setminus N)
\]
is an open neighborhood of $x$, and $q^{-1}(U)\subseteq N$.  If
$x'\in U$, surjectivity lets us choose $y'\in q^{-1}(x')$; then
$y'\in N$, so $a(x')=(a\circ q)(y')<r$.  Thus
$
 \sup_Ua\leq r
$
and $a^*(x)\leq r$.  Letting $r\downarrow M$ proves the desired
inequality when $M\in\mathbf R$.  If $M=-\infty$, the same argument
works for every real $r$ and gives $a^*(x)=-\infty$.

Together with \eqref{eq:p031-fiber-easy-inequality}, this proves
\eqref{eq:p031-perfect-pullback-usc}; substituting $x=q(y)$ gives
\eqref{eq:p031-fiber-maximum-pullback}.
\end{proof}

\begin{proof}[\Cref{prop:p031-open-map-characterization}]
If $q$ is open, the neighborhood images $q(N)$, with $N$ ranging over
the neighborhoods of $y$, are cofinal among the neighborhoods of
$q(y)$.  The neighborhood-filter formula therefore gives
$
 (a\circ q)^*(y)=a^*(q(y))
$
for every locally bounded-above $a$; this is
\Cref{lem:p031-open-pullback-usc}.  Hence (1) implies (2), and (2)
implies (3).

Conversely, suppose that $q$ is not open.  There are an open set
$N\subseteq Y$ and a point $y\in N$ such that $q(N)$ is not a
neighborhood of $x:=q(y)$.  Define
\[
                  a:=\mathbf 1_{X\setminus q(N)}.
\]
Every neighborhood of $x$ meets $X\setminus q(N)$, so $a^*(x)=1$.
On the other hand, $a\circ q=0$ throughout the neighborhood $N$ of
$y$, and hence
\[
                  (a\circ q)^*(y)=0<a^*(q(y)).
\]
Thus (3) fails.  This proves that (3) implies (1) and completes the
equivalence.
\end{proof}

\begin{proof}[\Cref{cor:p031-scalar-fiber-maximum}]
Scalar projection is continuous, closed, and surjective by
\cite[\S2.9]{GR25Classical}.  Its fiber over $x$ is canonically
\[
       \mathcal M\bigl(\mathcal H(x)\widehat\otimes_KK'\bigr),
\]
as recalled in \Cref{sec:berk-functoriality}; it is compact by
\Cref{prop:berk-compact-spectrum}.  The hypotheses of
\Cref{lem:p031-perfect-pullback-usc} therefore hold and give
\eqref{eq:p031-scalar-fiber-maximum}.

If $W\subseteq X$ is ordinary open, the restricted map
$
 \pi^{-1}(W)\to W
$
is still continuous and surjective and has the same fibers.  It is closed:
if $F=G\cap\pi^{-1}(W)$ is relatively closed, with $G$ closed in
$X_{K'}$, then
$
 \pi(F)=\pi(G)\cap W
$
is relatively closed in $W$.  The same lemma applies.
\end{proof}

\begin{proof}[\Cref{cor:p031-fiber-constancy-upgrade}]
Apply \Cref{cor:p031-scalar-fiber-maximum} to $a=s$.  Since
$u=s^*$, the result is exactly
\eqref{eq:p031-perron-fiber-maximum}.  If $v_{K'}$ is constant on a
fiber, its common value is its maximum there, so
\eqref{eq:p031-fiber-constancy} holds.  Conversely, if that identity
holds, its left side is constant on each scalar fiber; hence $v_{K'}$ is
fiberwise constant.  This proves the equivalence.

If the condition holds after every analytic extension, the exact
compatibility hypothesis of
\Cref{cor:p031-scalar-compatibility-upgrade} holds.  That corollary gives
classical plurisubharmonicity of $u$, and its obstacle specialization gives
the final assertion.
\end{proof}

\begin{proof}[
\Cref{prop:p031-type-four-topological-defect}]
Poineau--Turchetti take an algebraically closed complete field $K$ which
is not spherically complete, a maximally complete algebraically closed
extension $K'/K$, and a type--$4$ point $x$ of the Berkovich affine
line.  By \cite[Lemma~5.6.2]{PT21SchottkyI}, the fiber over $x$ is a
positive-radius closed disc $D^+(\gamma,r(x))$.  The open disc
$D^-(\gamma,r(x))$ lies in that fiber and therefore has image $\{x\}$,
which is not a neighborhood of $x$; this is the non-openness recorded in
\cite[Remark~5.6.4]{PT21SchottkyI}.

Apply the contrapositive construction in the proof of
\Cref{prop:p031-open-map-characterization} to this open disc and a point
$y$ in it.  It gives a bounded $\{0,1\}$-valued function $a$ such that
\[
                 (a\circ\pi_{K'/K})^*(y)=0<1=a^*(x).
\]
Finally, the affine-line scalar projection is proper and surjective by
\cite[Proposition~5.6.1]{PT21SchottkyI}, so
\Cref{lem:p031-perfect-pullback-usc} identifies $a^*(x)$ with the
fiber maximum and proves \eqref{eq:p031-type-four-strict-defect}.  The
construction imposes no psh property on $a$, which proves the stated
scope boundary.
\end{proof}


One scalar extension is nevertheless harmless.  After passing to the
completion of an algebraic closure, scalar projection is open on every
strict affinoid chart.  Relative interiors globalize this fact to the
ordinary affine open set, so the pulled-back regularized envelope is
exactly the regularization of the pulled-back raw supremum.

% Book-local P03.1 completed-algebraic-closure reduction,
% 19 September 2026.  No novelty or priority claim is made.

\begin{theorem}
\label{thm:p031-completed-algebraic-closure-reduction}
Keep the hypotheses and notation of
\Cref{thm:p031-affine-ground-field-perron}; in particular,
\[
 V\subseteq\mathbf A_K^{n,\mathrm{an}},
 \qquad
 s:=\sup_i u_i,
 \qquad
 u:=s^*.
\]
Let $L=\widehat{K^{\mathrm{alg}}}$ be the completion of an algebraic
closure of $K$, and let
$
 \pi_{L/K}:V_L\to V
$
be scalar projection.  Then:
\begin{enumerate}
\item the map $\pi_{L/K}$ is open;
\item with $s_L:=s\circ\pi_{L/K}$, one has
\begin{equation}
\label{eq:p031-no-completed-algebraic-closure-defect}
                    u\circ\pi_{L/K}=s_L^*;
\end{equation}
this common function, denoted by $u_L$, is
$L$-curvewise plurisubharmonic;
\item $u$ is classically plurisubharmonic on $V$ if and only if
$u_L$ is classically plurisubharmonic on $V_L$.
\end{enumerate}
Thus passage to the completed algebraic closure creates no
upper-regularization defect.  For the stronger affine Perron problem
treated by \Cref{thm:p031-affine-ground-field-perron}, it is enough to
resolve the remaining scalar-extension obstruction over complete
algebraically closed fields.
\end{theorem}

% Book-local P03.1 completed-algebraic-closure reduction proof,
% 19 September 2026.  No novelty or priority claim is made.

\begin{proof}[\Cref{thm:p031-completed-algebraic-closure-reduction}]
We first prove openness on the ordinary open set $V$.  Fix $y\in V_L$,
put $x=\pi_{L/K}(y)$, and let $N\subseteq V_L$ be an ordinary open
neighborhood of $y$.  Since $K$ is nontrivially valued and $V$ is open in
analytic affine space, there is a strict affinoid domain $X\subseteq V$
such that $x$ lies in the relative interior $\operatorname{Int}(X/V)$.
Compatibility of relative interiors with scalar extension gives
\[
 y\in\pi_{L/K}^{-1}\bigl(\operatorname{Int}(X/V)\bigr)
       \subseteq\operatorname{Int}(X_L/V_L);
\]
see \cite[\S2.11(5)]{GR25Classical}.

The completed-algebraic-closure openness theorem is recorded and applied
in the proof of \cite[Proposition~3.5(9)]{GR25Classical}, where it is
traced to Berkovich's Corollary~1.3.6.  It makes the scalar map
$X_L\to X$ open.  Hence
\[
 \pi_{L/K}\bigl(N\cap\operatorname{Int}(X_L/V_L)\bigr)
       \cap\operatorname{Int}(X/V)
\]
is an ordinary open neighborhood of $x$ in $V$, contained in
$\pi_{L/K}(N)$.  Therefore $\pi_{L/K}:V_L\to V$ is open.

Scalar projection is continuous and surjective by
\cite[\S2.9]{GR25Classical}.  Applying
\Cref{lem:p031-open-pullback-usc} to $a=s$ gives
\[
              u\circ\pi_{L/K}
                 =s^*\circ\pi_{L/K}
                 =(s\circ\pi_{L/K})^*=s_L^*,
\]
which proves \eqref{eq:p031-no-completed-algebraic-closure-defect}.

By \cite[Proposition~4.2(5)]{GR25Classical}, each
$u_i\circ\pi_{L/K}$ is classically psh on $V_L$.  The pulled-back family
remains locally upper bounded and has pointwise supremum $s_L$.
Applying \Cref{thm:p031-affine-ground-field-perron} over $L$ proves that
$u_L=s_L^*$ is $L$-curvewise psh.  Finally, another application of
\cite[Proposition~4.2(5)]{GR25Classical} gives
\[
 u\text{ is classically psh on }V
 \quad\Longleftrightarrow\quad
 u\circ\pi_{L/K}=u_L\text{ is classically psh on }V_L.
\]
This proves all assertions.
\end{proof}


The raw-supremum gap can also be localized without any scalar-extension
hypothesis.  Away from the closure of the strict gap between a raw
supremum and its upper-semicontinuous regularization, the finite maxima of
the competitors converge to an upper-semicontinuous limit.  The general
pointwise-limit theorem then makes that limit classically psh.  For the
completed algebraic closure, openness shows that the entire closed defect
support pulls back exactly.

% Book-local P03.1 upper-regularization defect localization,
% 19 September 2026.  No novelty or priority claim is made.

\begin{theorem}
\label{thm:p031-usc-defect-localization}
Let $X$ be a $K$-analytic space, and let $(\varphi_i)_{i\in I}$ be a
nonempty locally upper-bounded family of classically plurisubharmonic
functions on $X$.  Put
\[
                  s:=\sup_{i\in I}\varphi_i,\qquad u:=s^*,
\]
and define the upper-regularization defect set and its complement by
\[
 D(s):=\{x\in X:u(x)>s(x)\},
 \qquad
 G(s):=X\setminus\overline{D(s)}.
\]
Then $G(s)$ is the largest ordinary open subset of $X$ on which
$u=s$, and
\[
                         u|_{G(s)}=s|_{G(s)}
\]
is classically plurisubharmonic.

In particular, if the raw supremum $s$ is upper semicontinuous, then
$D(s)=\varnothing$ and $u=s$ is classically plurisubharmonic on all of
$X$.
\end{theorem}

\begin{corollary}
\label{cor:p031-completed-algebraic-closure-defect-support}
Keep the notation of
\Cref{thm:p031-completed-algebraic-closure-reduction}, and put
\[
 D_K:=\{x\in V:u(x)>s(x)\},
 \qquad
 D_L:=\{y\in V_L:u_L(y)>s_L(y)\}.
\]
Then
\[
 D_L=\pi_{L/K}^{-1}(D_K),
 \qquad
\overline{D_L}=\pi_{L/K}^{-1}(\overline{D_K}).
\]
Consequently,
\[
 V_L\setminus\overline{D_L}
   =\pi_{L/K}^{-1}\bigl(V\setminus\overline{D_K}\bigr).
\]
Thus passage to the completed algebraic closure preserves the full
upper-regularization defect support.
\end{corollary}

% Book-local P03.1 upper-regularization defect localization proof,
% 19 September 2026.  No novelty or priority claim is made.

\begin{proof}[\Cref{thm:p031-usc-defect-localization}]
By construction $u=s^*\geq s$.  An ordinary open set $U\subseteq X$
satisfies $u=s$ on $U$ exactly when $U\cap D(s)=\varnothing$.
It follows that the union of all such open sets is
\[
                         X\setminus\overline{D(s)}=G(s),
\]
which proves the maximality assertion.  Since $u$ is upper
semicontinuous, the equality $s=u$ on $G(s)$ also shows that
$s|_{G(s)}$ is upper semicontinuous.

Let $\mathcal F(I)$ be the directed set of nonempty finite subsets of
$I$, ordered by inclusion.  For $F\in\mathcal F(I)$, put
\[
                            m_F:=\max_{i\in F}\varphi_i.
\]
The cone property
\cite[Proposition~4.2(3)]{GR25Classical} makes every $m_F$ classically
psh.  On $G(s)$, the net $(m_F)_F$ converges pointwise to
\[
                         s|_{G(s)}=u|_{G(s)}.
\]
Its limit is upper semicontinuous, so the pointwise-limit theorem
\cite[Theorem~6.5]{GR25Classical} proves that $u|_{G(s)}$ is classically
psh.

Finally, $s$ is upper semicontinuous if and only if $s=s^*=u$, which is
equivalent to $D(s)=\varnothing$.  The preceding conclusion then applies
with $G(s)=X$.
\end{proof}

\begin{proof}[\Cref{cor:p031-completed-algebraic-closure-defect-support}]
Equation \eqref{eq:p031-no-completed-algebraic-closure-defect} gives
\[
                     u_L=u\circ\pi_{L/K},
 \qquad
                     s_L=s\circ\pi_{L/K}.
\]
The strict inequality defining the defect set therefore gives
$D_L=\pi_{L/K}^{-1}(D_K)$.

Continuity of $\pi_{L/K}$ gives
\[
 \overline{\pi_{L/K}^{-1}(D_K)}
       \subseteq\pi_{L/K}^{-1}(\overline{D_K}).
\]
For the reverse inclusion, fix $y$ over a point of
$\overline{D_K}$.  Every ordinary open neighborhood $N$ of $y$ has
open image by \Cref{thm:p031-completed-algebraic-closure-reduction}; this
image meets $D_K$, so $N$ meets $\pi_{L/K}^{-1}(D_K)$.  Hence
$y\in\overline{D_L}$, which proves the second equality.  Taking
complements gives the final identity.
\end{proof}


There is an exact descent statement in the genuine covering-space range.
If a finite analytic map is already a disjoint union of analytic
isomorphisms over an ordinary neighborhood of every base point, then one
may descend an arbitrary locally Fubini--Study function by taking its
fiberwise maximum.  This uses the full real-coefficient LFS calculus, not a
rational or invariant subfamily.

% P03.1 ordinary-locally-split LFS descent, live integration version.

\begin{definition}
\label{def:p031-ordinary-locally-split}
Let $\pi:E\to X$ be a finite surjective morphism of $K$-analytic spaces.
We call $\pi$ \emph{ordinary-locally split} if every $x\in X$ has an
ordinary open neighborhood $U$ for which
\begin{equation}
\label{eq:p031-ordinary-local-splitting}
 \pi^{-1}(U)=\coprod_{i=1}^{d}U_i,
 \qquad
 \pi_i:=\pi|_{U_i}:U_i\xrightarrow{\ \sim\ }U
\end{equation}
as $K$-analytic spaces.  In particular, the fiber cardinality is locally
constant on $X$.
\end{definition}

\begin{theorem}
\label{thm:p031-lfs-fiber-maximum}
Let $\pi:E\to X$ be an ordinary-locally-split finite surjective morphism
of good strict $K$-analytic spaces.  If $u$ is locally Fubini--Study on
$E$, then
\begin{equation}
\label{eq:p031-lfs-fiber-maximum}
             \pi_{\max}u(x):=\max_{e\in\pi^{-1}(x)}u(e)
\end{equation}
is locally Fubini--Study on $X$.  No rationality assumption on the
coefficients of $u$ is required.
\end{theorem}

\begin{theorem}
\label{thm:p031-lfs-envelope-pullback}
Keep the setting of \Cref{thm:p031-lfs-fiber-maximum}.  For an arbitrary
function $f:X\to\mathbf R\cup\{-\infty\}$, put
\begin{align*}
 s_{\mathrm{LFS},X}
   &:=\sup\{b:b\text{ is locally Fubini--Study on }X,\ b\leq f\},\\
 s_{\mathrm{LFS},E}
   &:=\sup\{u:u\text{ is locally Fubini--Study on }E,
                    \ u\leq f\circ\pi\},
\end{align*}
with the empty-family supremum equal to $-\infty$.  Then
\begin{equation}
\label{eq:p031-lfs-envelope-pullback}
               s_{\mathrm{LFS},E}=s_{\mathrm{LFS},X}\circ\pi.
\end{equation}
The competitor family on $E$ is nonempty if and only if the competitor
family on $X$ is nonempty.
\end{theorem}

% P03.1 ordinary-locally-split LFS descent proof, live integration version.

\begin{proof}[\Cref{thm:p031-lfs-fiber-maximum}]
Fix $x\in X$ and a splitting neighborhood $U$ as in
\eqref{eq:p031-ordinary-local-splitting}.  Write
$\sigma_i:=\pi_i^{-1}:U\to U_i$.  The local description of
locally Fubini--Study functions in \cite[\S7.9]{GJR26} is invariant under
analytic isomorphisms: composing the analytic units in a local presentation
with $\sigma_i$ again gives analytic units.  Hence
\[
                         u_i:=u|_{U_i}\circ\sigma_i
\]
is locally Fubini--Study on $U$ for every $i$.

For $y\in U$, the points
$\sigma_1(y),\ldots,\sigma_d(y)$ are exactly the fiber $\pi^{-1}(y)$, so
\begin{equation}
\label{eq:p031-local-fiber-maximum}
                    \pi_{\max}u|_U=\max_{1\leq i\leq d}u_i.
\end{equation}
Locally Fubini--Study functions are closed under finite maxima, again by
the full real-coefficient local calculus of \cite[\S7.9]{GJR26}.  Thus
$\pi_{\max}u$ is locally Fubini--Study on $U$.  Such splitting
neighborhoods cover $X$, proving the theorem.
\end{proof}

\begin{proof}[\Cref{thm:p031-lfs-envelope-pullback}]
Pullback preserves the local presentation of a locally Fubini--Study
function.  Consequently every base competitor $b$ gives the cover
competitor $b\circ\pi$, and hence
\begin{equation}
\label{eq:p031-lfs-envelope-easy-inequality}
              s_{\mathrm{LFS},E}
                    \geq s_{\mathrm{LFS},X}\circ\pi.
\end{equation}

Conversely, let $u$ be a cover competitor and set
$b:=\pi_{\max}u$.  By \Cref{thm:p031-lfs-fiber-maximum}, the function $b$
is locally Fubini--Study on $X$.  For $x\in X$, every
$e\in\pi^{-1}(x)$ satisfies $u(e)\leq f(x)$; therefore $b\leq f$ and $b$
is a base competitor.  For every $e\in E$ one then has
\[
              u(e)\leq b(\pi(e))
                    \leq s_{\mathrm{LFS},X}(\pi(e)).
\]
Taking the supremum over all cover competitors gives the reverse of
\eqref{eq:p031-lfs-envelope-easy-inequality} whenever the cover family is
nonempty.

Finally, a base competitor pulls back to a cover competitor, while a cover
competitor produces a base competitor by fiberwise maximization.  Thus the
two families are empty simultaneously, and the asserted equality also holds
under the stated empty-family convention.
\end{proof}


Thus the raw LFS envelope commutes exactly with pullback along every
ordinary-locally-split finite cover, including the empty-family case.  The
hypothesis is an explicit local splitting, not a consequence asserted here
for a general finite flat or finite etale morphism.  Nor does this result
compare the LFS subfamily with every semipositive real-PL competitor, produce
such a cover for scalar extension, or prove regularity of the common raw
envelope.

Ordinary local splitting is not the only mechanism that can control a
fiberwise maximum.  A finite group may organize each fiber as a single
orbit even when no local disjoint-union presentation has been chosen.  Under
the explicit finite-flat orbit-presentation hypothesis below, symmetrization
descends the rational semipositive branch exactly and places the full LFS
branch below the base real-PL envelope.

% P03.1 finite analytic orbit descent, live integration version.

\begin{definition}
\label{def:p031-finite-analytic-orbit-presentation}
Let $E$ and $X$ be $K$-analytic spaces.  A morphism
$\pi:E\to X$ has a \emph{finite analytic orbit presentation} if it is
finite, flat, and surjective, and if a finite group $G$ acts on $E$ by
$K$-analytic automorphisms such that
\begin{equation}
\label{eq:p031-action-over-base}
              \pi\circ g=\pi\quad(g\in G),
              \qquad
              \pi^{-1}(x)=G\cdot e
\end{equation}
for every $x\in X$ and every $e\in\pi^{-1}(x)$.  This is only a
set-theoretic orbit condition on the fibers; no categorical quotient is
included in the definition.
\end{definition}

\begin{theorem}
\label{thm:p031-orbit-fiber-maximum}
Let $K$ be complete and nontrivially valued, and let
$\pi:E\to X$ have a finite analytic orbit presentation between good strict
$K$-analytic spaces.
\begin{enumerate}
\item If $u$ is a semipositive $\mathbf Q$-PL function on $E$, then
\begin{equation}
\label{eq:p031-rational-fiber-maximum}
       \pi_{\max}u(x):=\max_{e\in\pi^{-1}(x)}u(e)
\end{equation}
is a semipositive $\mathbf Q$-PL function on $X$.
\item If $u$ is locally Fubini--Study on $E$, then $\pi_{\max}u$ is a
semipositive real-PL function on $X$.
\end{enumerate}
The second conclusion does not assert that the descended function is
locally Fubini--Study.
\end{theorem}

\begin{theorem}
\label{thm:p031-orbit-envelope-comparison}
Keep the hypotheses of \Cref{thm:p031-orbit-fiber-maximum}.  For an
arbitrary function $f:X\to\mathbf R\cup\{-\infty\}$, let
$s_{\mathbf Q+,X}$ and $s_{\mathbf Q+,E}$ be the raw envelopes of
semipositive $\mathbf Q$-PL minorants of $f$ and $f\circ\pi$,
respectively.  Define $s_{\mathrm{PL},X}$ and $s_{\mathrm{PL},E}$ using
semipositive real-PL minorants, and define $s_{\mathrm{LFS},X}$ and
$s_{\mathrm{LFS},E}$ using locally Fubini--Study minorants.  Every
empty-family supremum is $-\infty$.  Then
\begin{equation}
\label{eq:p031-rational-semipositive-envelope-pullback}
                  s_{\mathbf Q+,E}
                    =s_{\mathbf Q+,X}\circ\pi,
\end{equation}
and
\begin{equation}
\label{eq:p031-lfs-pl-envelope-chain}
 s_{\mathrm{LFS},X}\circ\pi
       \leq s_{\mathrm{LFS},E}
       \leq s_{\mathrm{PL},X}\circ\pi
       \leq s_{\mathrm{PL},E}.
\end{equation}
The two rational competitor families are empty simultaneously.
\end{theorem}

\begin{proposition}
\label{prop:p031-real-pl-fiber-maximum-failure}
The rational conclusion of \Cref{thm:p031-orbit-fiber-maximum} cannot be
extended to every semipositive real-PL function.  More precisely, there are
a finite analytic orbit presentation $\pi:E\to X$ and a semipositive
real-PL function $u$ on $E$ such that $\pi_{\max}u$ is not real PL.
Consequently, the preceding results give no identity
\[
                    s_{\mathrm{PL},E}
                       =s_{\mathrm{PL},X}\circ\pi
\]
for the full real-PL branch.
\end{proposition}

% P03.1 finite analytic orbit descent proof, live integration version.

\begin{proof}[\Cref{thm:p031-orbit-fiber-maximum}]
Let $G$ be the group in
\Cref{def:p031-finite-analytic-orbit-presentation}, and put
\begin{equation}
\label{eq:p031-orbit-symmetrization}
                         m:=\max_{g\in G}u\circ g.
\end{equation}
Pullback by an analytic automorphism preserves piecewise linearity and
semipositivity, the latter by
\cite[Proposition~7.6\textup{(1')}]{GJR26}.

Suppose first that $u$ is semipositive $\mathbf Q$-PL.  Every translate
$u\circ g$ is semipositive $\mathbf Q$-PL, and $m$ is $\mathbf Q$-PL by
the finite-maximum property in \cite[\S5.2\textup{(7)}]{GJR26}.  The
translates are classically psh by
\cite[Theorem~4.3]{GR25Classical}; their finite maximum is classically psh
by \cite[Proposition~4.2\textup{(3)}]{GR25Classical}.  The converse
direction of that same Theorem~4.3 makes $m$ semipositive.

Suppose instead that $u$ is locally Fubini--Study.  Each translate has the
same property, and the full real-coefficient local calculus of
\cite[\S7.9]{GJR26} makes their finite maximum locally Fubini--Study.
Hence $m$ is, in particular, semipositive real PL.

Equation \eqref{eq:p031-action-over-base} makes $m$ invariant under $G$,
and each fiber is one orbit.  There is therefore a unique function
$b:X\to\mathbf R$ satisfying
\begin{equation}
\label{eq:p031-orbit-symmetrization-descends}
                         m=b\circ\pi.
\end{equation}
In the rational case apply \cite[Proposition~5.11]{GJR26} with
$\Lambda=\mathbf Q$; in the locally Fubini--Study case apply it with
$\Lambda=\mathbf R$.  Since these coefficient groups are divisible and
$\pi$ is finite, flat, and surjective, the proposition shows that $b$ is
$\mathbf Q$-PL or real PL, respectively.  Semipositivity reflects through
the proper surjection $\pi$ by
\cite[Proposition~7.6\textup{(1')}]{GJR26}.

For $e\in\pi^{-1}(x)$, the orbit condition finally gives
\[
 b(x)=m(e)=\max_{g\in G}u(g(e))
              =\max_{e'\in\pi^{-1}(x)}u(e').
\]
Thus $b=\pi_{\max}u$.  In the second case the argument descends real
piecewise linearity and nef semipositivity, but constructs no semiample
residue model and hence proves no locally Fubini--Study conclusion.
\end{proof}

\begin{proof}[\Cref{thm:p031-orbit-envelope-comparison}]
Pullback preserves semipositive $\mathbf Q$-PL functions.  Hence every
base competitor gives a cover competitor and
\[
                  s_{\mathbf Q+,X}\circ\pi
                         \leq s_{\mathbf Q+,E}.
\]
Conversely, if $u\leq f\circ\pi$ is a semipositive $\mathbf Q$-PL cover
competitor, then \Cref{thm:p031-orbit-fiber-maximum} makes
$b:=\pi_{\max}u$ a semipositive $\mathbf Q$-PL base competitor.  Thus
\[
          u(e)\leq b(\pi(e))
              \leq s_{\mathbf Q+,X}(\pi(e)).
\]
Taking the supremum over the cover competitors proves the reverse
inequality in
\eqref{eq:p031-rational-semipositive-envelope-pullback}.  Pullback and
fiberwise maximization also show that the two rational families are empty
simultaneously.

Pullback of locally Fubini--Study base competitors gives
\[
                 s_{\mathrm{LFS},X}\circ\pi
                        \leq s_{\mathrm{LFS},E}.
\]
If $u\leq f\circ\pi$ is a locally Fubini--Study cover competitor, the
second part of \Cref{thm:p031-orbit-fiber-maximum} makes
$b:=\pi_{\max}u$ a semipositive real-PL base competitor below $f$.  Hence
\[
          u(e)\leq b(\pi(e))
              \leq s_{\mathrm{PL},X}(\pi(e)),
\]
and taking the supremum proves the middle inequality in
\eqref{eq:p031-lfs-pl-envelope-chain}.  Finally, pullback of semipositive
real-PL base competitors gives the last inequality in that chain.
\end{proof}

\begin{proof}[\Cref{prop:p031-real-pl-fiber-maximum-failure}]
Work over a complete nontrivially valued field whose additive value group
is $\mathbf Q$.  On $X=\mathbf G_{m,K}^{\mathrm{an}}$, Example~5.3 of
\cite{GJR26} gives the real-PL function
\[
                         h=\pi\log|t|+1
\]
for which $\max\{h,0\}$ is not real PL: under tropicalization it has its
unique bend at the irrational point $1/\pi$.  The function $h$ is
classically psh by \cite[Proposition~4.2\textup{(2),(3)}]{GR25Classical},
and therefore semipositive by \cite[Theorem~4.3]{GR25Classical}.

Let $E=X\amalg X$, let $\pi_E:E\to X$ be the fold map, and let the group
of order two exchange the components.  This is a finite analytic orbit
presentation.  Define $u$ to equal $h$ on the first component and $0$ on
the second.  Then $u$ is semipositive real PL, whereas
\[
                         (\pi_E)_{\max}u=\max\{h,0\}
\]
is not real PL.  This proves the asserted failure.
\end{proof}


The distinction between rational and real slopes is essential: the fold-map
example in \Cref{prop:p031-real-pl-fiber-maximum-failure} already prevents a
full real-PL fiber-maximum theorem.  The orbit result does not construct such
a presentation for scalar extension, descend the LFS class itself, compare
the full local-FS/semipositive-PL union, or prove regularity of an envelope.

Although maxima leave the full real-PL class, finite orbit averaging stays
inside it.  Averaging need not dominate the original competitor, so it does
not by itself prove descent of the full envelope.  It does, however, identify
the invariant subenvelope exactly and isolate the entire remaining gain as a
nonnegative symmetry-breaking defect.

% P03.1 finite-orbit symmetry defect, live integration version.

\begin{definition}
\label{def:p031-orbit-average-competitors}
Keep the finite analytic orbit presentation $\pi:E\to X$ and its finite
group $G$ from \Cref{def:p031-finite-analytic-orbit-presentation}.  For a
real-valued function $u$ on $E$, its \emph{orbit average} is
\begin{equation}
\label{eq:p031-orbit-average}
                 \mathsf A_Gu
                    :=\frac{1}{|G|}\sum_{g\in G}u\circ g.
\end{equation}
For an obstacle $f:X\to\mathbf R\cup\{-\infty\}$, set
\begin{align*}
 \mathcal P_X(f)
  &:=\{b:b\text{ is semipositive real PL on }X,\ b\leq f\},\\
 \mathcal P_E(f\circ\pi)
  &:=\{u:u\text{ is semipositive real PL on }E,\ u\leq f\circ\pi\},
\end{align*}
and write $\mathcal P_E(f\circ\pi)^G$ for the invariant subfamily.
\end{definition}

\begin{theorem}
\label{thm:p031-invariant-real-pl-descent}
Let $K$ be complete and nontrivially valued, and let $\pi:E\to X$ have a
finite analytic orbit presentation between good strict $K$-analytic spaces.
For every obstacle $f:X\to\mathbf R\cup\{-\infty\}$, pullback is an
order-preserving bijection
\begin{equation}
\label{eq:p031-invariant-real-pl-bijection}
 \mathcal P_X(f)
   \xrightarrow{\ \sim\ }
 \mathcal P_E(f\circ\pi)^G,
 \qquad b\longmapsto b\circ\pi.
\end{equation}
Moreover, $\mathsf A_Gu\in\mathcal P_E(f\circ\pi)^G$ for every
$u\in\mathcal P_E(f\circ\pi)$.  Consequently the base and cover competitor
families are empty simultaneously, and
\begin{equation}
\label{eq:p031-invariant-real-pl-subenvelope}
 s_{\mathrm{PL},X}\circ\pi
   =\sup_{v\in\mathcal P_E(f\circ\pi)^G}v.
\end{equation}
\end{theorem}

\begin{theorem}
\label{thm:p031-centered-symmetry-defect}
Keep the hypotheses of \Cref{thm:p031-invariant-real-pl-descent}, and suppose
that $\mathcal P_X(f)$ is nonempty.  The \emph{PL symmetry-breaking defect}
\begin{equation}
\label{eq:p031-pl-symmetry-defect}
 \Sigma_{\pi,f}(e)
   :=s_{\mathrm{PL},E}(e)-s_{\mathrm{PL},X}(\pi(e))
\end{equation}
is finite, nonnegative, invariant under $G$, and constant on every fiber of
$\pi$.  For $u\in\mathcal P_E(f\circ\pi)$ define the nonnegative average
headroom
\begin{equation}
\label{eq:p031-average-headroom}
 H_u(e):=s_{\mathrm{PL},X}(\pi(e))-\mathsf A_Gu(e).
\end{equation}
Then
\begin{equation}
\label{eq:p031-centered-symmetry-formula}
 \Sigma_{\pi,f}(e)
   =\sup_{u\in\mathcal P_E(f\circ\pi)}
       \bigl(u(e)-\mathsf A_Gu(e)-H_u(e)\bigr).
\end{equation}
If
\begin{equation}
\label{eq:p031-orbit-asymmetry-budget}
 \Omega_{\pi,f}(e)
   :=\sup_{u\in\mathcal P_E(f\circ\pi)}
        \bigl(u(e)-\mathsf A_Gu(e)\bigr),
\end{equation}
then, with values in $[0,+\infty]$,
\begin{equation}
\label{eq:p031-symmetry-oscillation-bounds}
 0\leq\Sigma_{\pi,f}(e)
   \leq\Omega_{\pi,f}(e)
   \leq
   \sup_{u\in\mathcal P_E(f\circ\pi)}
       \left(
        \max_{g\in G}u(g(e))-\min_{g\in G}u(g(e))
       \right).
\end{equation}
In particular, $\Sigma_{\pi,f}(e)=0$ if and only if every cover competitor
satisfies
\begin{equation}
\label{eq:p031-pointwise-symmetry-criterion}
                 u(e)-\mathsf A_Gu(e)\leq H_u(e).
\end{equation}

For the raw envelope $s_{\mathrm{P03}}$ over the union of locally
Fubini--Study and semipositive real-PL minorants, one has
\begin{equation}
\label{eq:p031-full-union-symmetry-defect}
 s_{\mathrm{P03},E}-s_{\mathrm{P03},X}\circ\pi
     =\Sigma_{\pi,f}.
\end{equation}
Thus exact full-union compatibility is equivalent to
$\Sigma_{\pi,f}\equiv0$, and it holds in particular if
$\Omega_{\pi,f}\equiv0$.  No regularity of the raw envelopes, finiteness of
$\Omega_{\pi,f}$, or vanishing of $\Sigma_{\pi,f}$ is asserted.
\end{theorem}

% P03.1 finite-orbit symmetry defect proof, live integration version.

\begin{proof}[\Cref{thm:p031-invariant-real-pl-descent}]
Pullback sends every member of $\mathcal P_X(f)$ to a $G$-invariant member
of $\mathcal P_E(f\circ\pi)$, and it is injective because $\pi$ is
surjective.

Conversely, let $v\in\mathcal P_E(f\circ\pi)^G$.  Every fiber of $\pi$ is a
single $G$-orbit, so $v$ is constant on fibers and there is a unique
set-theoretic function $b:X\to\mathbf R$ such that
\begin{equation}
\label{eq:p031-invariant-real-pl-descends}
                              v=b\circ\pi.
\end{equation}
Apply \cite[Proposition~5.11]{GJR26} with the divisible coefficient group
$\mathbf R$.  Since $\pi$ is finite, flat, and surjective, the proposition
shows directly that $b$ is real PL.  Semipositivity reflects through the
proper surjection $\pi$ by
\cite[Proposition~7.6\textup{(1')}]{GJR26}.  Finally,
$v\leq f\circ\pi$ and surjectivity give $b\leq f$.  Thus
$b\in\mathcal P_X(f)$, and the construction is inverse to pullback and
preserves pointwise order.

Now let $u\in\mathcal P_E(f\circ\pi)$.  Every translate $u\circ g$ is
semipositive real PL.  Such functions form a convex cone by
\cite[\S7.3\textup{(5)}]{GJR26}, so $\mathsf A_Gu$ is semipositive real PL.
It is $G$-invariant and remains below $f\circ\pi$.  This proves the averaging
claim.  A nonempty cover family therefore has a nonempty invariant subfamily
and hence a nonempty base family; the converse follows by pullback.  Taking
pointwise suprema in \eqref{eq:p031-invariant-real-pl-bijection} proves
\eqref{eq:p031-invariant-real-pl-subenvelope}, including the simultaneous
empty-family case.
\end{proof}

\begin{proof}[\Cref{thm:p031-centered-symmetry-defect}]
A base competitor and its pullback bound both raw envelopes from below.
They also force $f$ to be real-valued, so the obstacle bounds both envelopes
from above.  Hence the two envelope values and their difference are finite.
Pullback of base competitors gives
\[
                 s_{\mathrm{PL},X}\circ\pi
                     \leq s_{\mathrm{PL},E},
\]
and therefore $\Sigma_{\pi,f}\geq0$.

For $g_0\in G$, composition $u\mapsto u\circ g_0$ permutes the cover
competitor family.  Consequently
\[
 s_{\mathrm{PL},E}(g_0(e))
  =\sup_{u\in\mathcal P_E(f\circ\pi)}u(g_0(e))
  =s_{\mathrm{PL},E}(e).
\]
The pullback term in \eqref{eq:p031-pl-symmetry-defect} is invariant as
well.  Thus the defect is invariant and, since each fiber is one orbit,
constant on fibers.

By \Cref{thm:p031-invariant-real-pl-descent}, the average
$\mathsf A_Gu$ is an invariant cover competitor.  Equation
\eqref{eq:p031-invariant-real-pl-subenvelope} therefore gives
$\mathsf A_Gu\leq s_{\mathrm{PL},X}\circ\pi$, so $H_u\geq0$.  For every
cover competitor, there is the exact pointwise cancellation
\[
 u(e)-\mathsf A_Gu(e)-H_u(e)
   =u(e)-s_{\mathrm{PL},X}(\pi(e)).
\]
Taking the supremum proves
\eqref{eq:p031-centered-symmetry-formula}.

Discarding the nonnegative headroom gives
$\Sigma_{\pi,f}\leq\Omega_{\pi,f}$.  The pullback of any base competitor is
invariant and has $u-\mathsf A_Gu=0$, so $\Omega_{\pi,f}\geq0$.  Moreover,
for each individual competitor,
\[
 u(e)-\mathsf A_Gu(e)
   \leq
      \max_{g\in G}u(g(e))-\min_{g\in G}u(g(e)).
\]
Taking suprema proves \eqref{eq:p031-symmetry-oscillation-bounds}.  The
centered formula is a supremum of its displayed summands, while the defect
is nonnegative.  It vanishes exactly when every summand is nonpositive,
which is \eqref{eq:p031-pointwise-symmetry-criterion}.

Finally, every locally Fubini--Study function is semipositive real PL by
\cite[Definition~7.8 and \S7.9]{GJR26}.  Adjoining LFS minorants to the
semipositive real-PL family therefore does not enlarge it.  The P03 union
envelopes are precisely the two PL envelopes, and
\eqref{eq:p031-full-union-symmetry-defect} follows from the definition of
$\Sigma_{\pi,f}$.  Its vanishing criterion is immediate, and the sufficient
condition $\Omega_{\pi,f}\equiv0$ follows from
\eqref{eq:p031-symmetry-oscillation-bounds}.  None of these pointwise
identities supplies any of the regularity or vanishing statements excluded
in the theorem.
\end{proof}


Thus the finite-orbit part of the full P03 union reduces to a precise
pointwise cofinality question: can a non-invariant cover competitor improve
the invariant subenvelope after its average headroom is charged?  The theorem
does not show that this defect vanishes or is regular, and it does not produce
a finite orbit presentation for scalar extension.

The defect does vanish when finitely many global analytic sections cover
every fiber.  One selects a section only after fixing a cover competitor and
a fiber point, and then restricts that competitor along the selected section.
This transfers the required value to a global base competitor without taking
a maximum in the real-PL class.

% P03.1 global-section-palette descent, live integration version.

\begin{definition}
\label{def:p031-global-section-palette}
Let $\pi:E\to X$ be a finite surjective morphism of $K$-analytic spaces.
A \emph{global section palette} for $\pi$ is a nonempty finite family of
analytic sections
\[
                 \sigma_i:X\longrightarrow E
                    \qquad(1\leq i\leq N)
\]
such that, for every $x\in X$,
\begin{equation}
\label{eq:p031-global-section-palette}
             \pi^{-1}(x)
                =\{\sigma_1(x),\ldots,\sigma_N(x)\}.
\end{equation}
Repetitions among the displayed points are allowed.
\end{definition}

\begin{theorem}
\label{thm:p031-section-palette-envelope-descent}
Let $K$ be complete and nontrivially valued, and let $\pi:E\to X$ be a
finite surjective morphism of good strict $K$-analytic spaces admitting a
global section palette.  Fix an obstacle
$f:X\to\mathbf R\cup\{-\infty\}$ and a semipositive real-PL cover
competitor $u\leq f\circ\pi$.  For every $x\in X$, there is a semipositive
real-PL base competitor $b_{x,u}\leq f$ such that
\begin{equation}
\label{eq:p031-section-palette-exact-repair}
 (b_{x,u}\circ\pi)(e)
   =\max_{e'\in\pi^{-1}(x)}u(e')
       \qquad(e\in\pi^{-1}(x)).
\end{equation}
The competitor $b_{x,u}$ may depend on $x$ and $u$.

Consequently, including when both competitor families are empty,
\begin{equation}
\label{eq:p031-section-palette-envelope-pullback}
                 s_{\mathrm{PL},E}
                    =s_{\mathrm{PL},X}\circ\pi.
\end{equation}
The same identity holds for the raw P03 envelope over the union of locally
Fubini--Study and semipositive real-PL minorants.
\end{theorem}

\begin{corollary}
\label{cor:p031-sectioned-orbit-defect-vanishing}
Keep the finite analytic orbit presentation of
\Cref{def:p031-finite-analytic-orbit-presentation}.  If $\pi$ admits one
global analytic section $\sigma:X\to E$, then
$\{g\circ\sigma:g\in G\}$ is a global section palette.  Hence
\eqref{eq:p031-section-palette-envelope-pullback} and the full P03 envelope
identity hold for every obstacle.  If $\mathcal P_X(f)$ is nonempty, then
the symmetry-breaking defect from
\Cref{thm:p031-centered-symmetry-defect} vanishes identically.
\end{corollary}

\begin{corollary}
\label{cor:p031-fold-envelope-without-maximum-closure}
In the fold-map example from
\Cref{prop:p031-real-pl-fiber-maximum-failure}, write
$q:X\amalg X\to X$ for the fold map, let $h$ be the semipositive real-PL
function on $X$, and set $f:=\max\{h,0\}$.  Let $u$ equal $h$ on the first
component and $0$ on the second.  Then
\begin{equation}
\label{eq:p031-fold-envelope-values}
 s_{\mathrm{PL},X}=f,
 \qquad
 s_{\mathrm{PL},X\amalg X}=f\circ q,
 \qquad
 \Sigma_{q,f}\equiv0,
\end{equation}
although
\[
                         q_{\max}u=f
\]
is not real PL.  Thus failure of real-PL fiber-maximum closure does not by
itself produce a nonzero envelope defect.
\end{corollary}

% P03.1 global-section-palette descent proof, live integration version.

\begin{proof}[\Cref{thm:p031-section-palette-envelope-descent}]
Fix $u$ and $x$.  The finite nonempty fiber has a point $e_x$ at which
$u$ attains its maximum.  Choose $i=i(x,u)$ in
\eqref{eq:p031-global-section-palette} such that $\sigma_i(x)=e_x$, and set
\[
                            b_{x,u}:=u\circ\sigma_i.
\]
Real piecewise linearity is preserved by analytic pullback by
\cite[Lemma~5.4\textup{(1)}, p.~171]{GJR26}, and semipositivity is
preserved by
\cite[Proposition~7.6\textup{(1)}, p.~189]{GJR26}.  Thus $b_{x,u}$ is
semipositive real PL on $X$.  Since $\pi\circ\sigma_i$ is the identity and
$u\leq f\circ\pi$, one has $b_{x,u}\leq f$.  Its pullback is constant on
$\pi^{-1}(x)$, with value
\[
 b_{x,u}(x)=u(\sigma_i(x))=u(e_x)
          =\max_{e'\in\pi^{-1}(x)}u(e').
\]
This proves \eqref{eq:p031-section-palette-exact-repair}.  The index was
selected after $x$ and $u$; no single section is asserted to maximize on
other fibers.

The same two pullback results show that every base competitor pulls back to
a cover competitor.  Hence
\[
                 s_{\mathrm{PL},X}\circ\pi
                    \leq s_{\mathrm{PL},E}.
\]
Conversely, fix $e\in E$, put $x:=\pi(e)$, and let $u$ be a cover
competitor.  Choose $j$ with $\sigma_j(x)=e$.  The preceding argument,
without choosing a maximizing point, makes $u\circ\sigma_j$ a base
competitor.  Therefore
\[
                   u(e)=(u\circ\sigma_j)(x)
                         \leq s_{\mathrm{PL},X}(x).
\]
Taking the supremum over $u$ proves the reverse inequality.  Moreover, a
base competitor pulls back to a cover competitor, while composition with
any palette section sends a cover competitor to a base competitor.  The
two families are therefore empty simultaneously, so the equality includes
the value $-\infty$.

Every locally Fubini--Study function is semipositive real PL by
\cite[Definition~7.8 and \S7.9]{GJR26}.  Adjoining LFS minorants therefore
does not enlarge either real-PL competitor family, and the P03 union
identity follows from
\eqref{eq:p031-section-palette-envelope-pullback}.
\end{proof}

\begin{proof}[\Cref{cor:p031-sectioned-orbit-defect-vanishing}]
For every $g\in G$,
\[
                  \pi\circ g\circ\sigma
                    =\pi\circ\sigma
                    =\operatorname{id}_X,
\]
so $g\circ\sigma$ is a section.  If $e\in\pi^{-1}(x)$, fiber
transitivity gives $e=g(\sigma(x))$ for some $g\in G$.  The translates of
$\sigma$ therefore form a global section palette, with repetitions allowed
when a stabilizer is nontrivial.  The two envelope identities follow from
\Cref{thm:p031-section-palette-envelope-descent}.  Under the stated
nonemptiness hypothesis, subtracting the PL envelope identity in
\eqref{eq:p031-pl-symmetry-defect} gives $\Sigma_{\pi,f}=0$.
\end{proof}

\begin{proof}[\Cref{cor:p031-fold-envelope-without-maximum-closure}]
The two canonical inclusions $X\to X\amalg X$ form a global section
palette for $q$.  Hence
\Cref{thm:p031-section-palette-envelope-descent} gives exact semipositive
real-PL envelope pullback.

Both $h$ and $0$ are base competitors below $f$.  Their pointwise supremum
is $f$, while the envelope is bounded above by the obstacle.  Hence
$s_{\mathrm{PL},X}=f$, and the pullback identity gives
$s_{\mathrm{PL},X\amalg X}=f\circ q$.  The constant zero function is an
LFS minorant of $f$, so the defect is defined and vanishes.

The two values of $u$ on the fiber over $x$ are $h(x)$ and $0$.  Thus
$q_{\max}u=\max\{h,0\}=f$, which is not real PL by
\Cref{prop:p031-real-pl-fiber-maximum-failure}.  The same example therefore
has a failed real-PL maximum operation and exact envelope descent.
\end{proof}


The fold example shows that failure of real-PL maximum closure is therefore
not itself an envelope obstruction.  The section-palette theorem does not
construct global sections for scalar projection or for a general finite flat
cover, and it proves no regularity of the common raw envelope.

A global palette is stronger than what is needed to compare local envelopes.
One analytic section through a chosen cover point transfers every local
competitor value at that point.  Separating this local identity from global
cofinality yields an exact balance for the remaining pullback gain.

% P03.1 local-section PL defect balance, live integration version.

\begin{definition}
\label{def:p031-analytically-sectioned-point}
Let $\pi:E\to X$ be a morphism of $K$-analytic spaces and let $e\in E$,
with $x:=\pi(e)$.  The point $e$ is \emph{analytically sectioned over
$X$} if there are an ordinary open neighborhood $U$ of $x$ and an analytic
morphism $\sigma:U\to E$ such that
\begin{equation}
\label{eq:p031-local-analytic-section}
        \pi\circ\sigma=\operatorname{id}_U,
        \qquad
        \sigma(x)=e.
\end{equation}
\end{definition}

\begin{definition}
\label{def:p031-local-pl-globalization-defect}
Let $Y$ be a good strict $K$-analytic space and let
$g:Y\to\mathbf R\cup\{-\infty\}$.  The local raw semipositive real-PL
envelope is
\begin{equation}
\label{eq:p031-local-pl-envelope}
 \begin{aligned}
 P_{\mathrm{PL}}^{\mathrm{loc}}g(y)
  :=\sup\{v(y):\;&V\text{ is open},\ y\in V,\\
                 &v\text{ is semipositive real PL on }V,
                   \ v\leq g|_V\}.
 \end{aligned}
\end{equation}
The supremum of an empty family is $-\infty$.

Suppose that $g$ has a global semipositive real-PL minorant, and write
$s_{\mathrm{PL},Y}$ for its global raw semipositive real-PL envelope.  The
\emph{PL globalization defect} is
\begin{equation}
\label{eq:p031-pl-globalization-defect}
 \Delta^{\mathrm{PL}}_{g,Y}(y)
   :=P_{\mathrm{PL}}^{\mathrm{loc}}g(y)-s_{\mathrm{PL},Y}(y).
\end{equation}
It is a finite nonnegative real number.
\end{definition}

\begin{theorem}
\label{thm:p031-sectioned-local-pl-pullback}
Let $K$ be complete and nontrivially valued, let $\pi:E\to X$ be a
morphism of good strict $K$-analytic spaces, and let $e\in E$ be
analytically sectioned over $X$.  For every function
$f:X\to\mathbf R\cup\{-\infty\}$,
\begin{equation}
\label{eq:p031-sectioned-local-pl-pullback}
 P_{\mathrm{PL}}^{\mathrm{loc}}(f\circ\pi)(e)
    =P_{\mathrm{PL}}^{\mathrm{loc}}f(\pi(e)).
\end{equation}
This identity includes the simultaneous empty-family case.
\end{theorem}

\begin{theorem}
\label{thm:p031-sectioned-globalization-balance}
Keep the morphism and sectioned point from
\Cref{thm:p031-sectioned-local-pl-pullback}, put $x:=\pi(e)$, and suppose
that $f$ has a global semipositive real-PL minorant.  Define the global PL
pullback gain by
\begin{equation}
\label{eq:p031-global-pl-pullback-gain}
 \Gamma^{\mathrm{PL}}_{\pi,f}(e)
   :=s_{\mathrm{PL},E}(e)-s_{\mathrm{PL},X}(x).
\end{equation}
Then
\begin{align}
 \Gamma^{\mathrm{PL}}_{\pi,f}(e)
   &=\Delta^{\mathrm{PL}}_{f,X}(x)
      -\Delta^{\mathrm{PL}}_{f\circ\pi,E}(e),
\label{eq:p031-sectioned-globalization-balance}\\
 0\leq\Delta^{\mathrm{PL}}_{f\circ\pi,E}(e)
   &\leq\Delta^{\mathrm{PL}}_{f,X}(x),
\label{eq:p031-sectioned-defect-contraction}\\
 0\leq\Gamma^{\mathrm{PL}}_{\pi,f}(e)
   &\leq\Delta^{\mathrm{PL}}_{f,X}(x).
\label{eq:p031-sectioned-gain-bound}
\end{align}
The same gain measures the pullback failure of the full P03 union envelope.
If $\pi$ has a finite analytic orbit presentation, then
$\Gamma^{\mathrm{PL}}_{\pi,f}=\Sigma_{\pi,f}$.
\end{theorem}

\begin{corollary}
\label{cor:p031-local-section-orbit-fiber-closure}
Let $\pi:E\to X$ have a finite analytic orbit presentation, let $f$ have a
global semipositive real-PL minorant, and suppose that $\pi$ admits an
analytic section on a neighborhood of $x\in X$.  Then every point of
$\pi^{-1}(x)$ is
analytically sectioned.  If
\begin{equation}
\label{eq:p031-zero-base-globalization-defect}
                  \Delta^{\mathrm{PL}}_{f,X}(x)=0,
\end{equation}
then, for every $e\in\pi^{-1}(x)$,
\begin{equation}
\label{eq:p031-zero-fiber-globalization-defects}
 \Delta^{\mathrm{PL}}_{f\circ\pi,E}(e)=0,
 \qquad
 \Sigma_{\pi,f}(e)=0.
\end{equation}
Hence the full P03 union envelope has exact pullback compatibility on the
entire fiber.
\end{corollary}

\begin{corollary}
\label{cor:p031-local-section-defect-preservation}
Let $\pi\colon E\to X$ have a finite analytic orbit presentation, let $f$
admit a global semipositive real PL minorant, and suppose that every base
point has an ordinary open neighborhood on which $\pi$ admits an analytic
section.  Then
\[
                 s_{\mathrm{PL},E}
                    =s_{\mathrm{PL},X}\circ\pi
\]
if and only if
\begin{equation}
\label{eq:p031-globalization-defect-preservation}
 \Delta^{\mathrm{PL}}_{f\circ\pi,E}
    =\Delta^{\mathrm{PL}}_{f,X}\circ\pi.
\end{equation}
In particular, base local-to-global cofinality
\[
           P_{\mathrm{PL}}^{\mathrm{loc}}f=s_{\mathrm{PL},X}
\]
forces both cover local-to-global cofinality and exact full-union
compatibility.
\end{corollary}

% P03.1 local-section PL defect balance proof, live integration version.

The defect in \eqref{eq:p031-pl-globalization-defect} is finite whenever it
is defined.  Indeed, let $b_0$ be a global semipositive real-PL minorant of
$g$.  Then
\[
 b_0(y)\leq s_{\mathrm{PL},Y}(y)
   \leq P_{\mathrm{PL}}^{\mathrm{loc}}g(y)
   \leq g(y)<+\infty.
\]
Here $b_0\leq g$ forces the $\mathbf R\cup\{-\infty\}$-valued obstacle
$g$ to be real-valued.  The displayed inequalities also prove
nonnegativity.

\begin{proof}[\Cref{thm:p031-sectioned-local-pl-pullback}]
Put $x:=\pi(e)$.  Pullback of a local semipositive real-PL minorant of $f$
gives such a minorant of $f\circ\pi$ on the inverse image of its domain.
Therefore
\begin{equation}
\label{eq:p031-easy-local-pl-inequality}
 P_{\mathrm{PL}}^{\mathrm{loc}}f(x)
    \leq P_{\mathrm{PL}}^{\mathrm{loc}}(f\circ\pi)(e).
\end{equation}

For the reverse inequality, let $u$ be a semipositive real-PL minorant of
$f\circ\pi$ on an ordinary open neighborhood $V$ of $e$.  Choose
$\sigma:U\to E$ as in \eqref{eq:p031-local-analytic-section}.  The set
\[
                         W:=U\cap\sigma^{-1}(V)
\]
is an ordinary open neighborhood of $x$, because analytic morphisms are
continuous and $\sigma(x)=e\in V$.  On $W$, put $b:=u\circ\sigma$.
Real piecewise linearity is preserved by analytic pullback by
\cite[Lemma~5.4\textup{(1)}, p.~171]{GJR26}, and semipositivity is
preserved by
\cite[Proposition~7.6\textup{(1)}, p.~189]{GJR26}.  Moreover,
\[
 b\leq f\circ\pi\circ\sigma=f|_W,
 \qquad
 b(x)=u(e).
\]
Thus $b$ is a local base competitor with the same central value as $u$.
Taking the supremum over $u$ gives the reverse of
\eqref{eq:p031-easy-local-pl-inequality}.  The same two constructions show
that one local family is empty if and only if the other is empty.
\end{proof}

\begin{proof}[\Cref{thm:p031-sectioned-globalization-balance}]
The global semipositive real-PL minorant of $f$ and its pullback make all
four local and global envelope values finite.  By
\Cref{thm:p031-sectioned-local-pl-pullback},
\[
 P_{\mathrm{PL}}^{\mathrm{loc}}(f\circ\pi)(e)
     =P_{\mathrm{PL}}^{\mathrm{loc}}f(x).
\]
Subtracting the cover and base globalization defects therefore gives
\begin{align*}
 \Delta^{\mathrm{PL}}_{f,X}(x)
  -\Delta^{\mathrm{PL}}_{f\circ\pi,E}(e)
   &=s_{\mathrm{PL},E}(e)-s_{\mathrm{PL},X}(x)\\
   &=\Gamma^{\mathrm{PL}}_{\pi,f}(e),
\end{align*}
which proves \eqref{eq:p031-sectioned-globalization-balance}.

Both globalization defects are nonnegative.  Pullback of global base
competitors gives
\[
                 s_{\mathrm{PL},X}\circ\pi
                       \leq s_{\mathrm{PL},E},
\]
so $\Gamma^{\mathrm{PL}}_{\pi,f}\geq0$.  The balance identity now gives
\eqref{eq:p031-sectioned-defect-contraction} and
\eqref{eq:p031-sectioned-gain-bound}.

The LFS class is contained in the semipositive real-PL class by
\cite[Definition~7.8 and \S7.9]{GJR26}, so the P03 union envelopes are the
PL envelopes.  In the orbit setting, the definition
\eqref{eq:p031-pl-symmetry-defect} identifies the PL pullback gain with
$\Sigma_{\pi,f}$.
\end{proof}

\begin{proof}[\Cref{cor:p031-local-section-orbit-fiber-closure}]
Let $\sigma:U\to E$ be a local section near $x$.  For every $g\in G$, the
map $g\circ\sigma$ is another section on $U$, since
\[
                         \pi\circ g\circ\sigma
                            =\pi\circ\sigma
                            =\operatorname{id}_U.
\]
Fiber transitivity shows that the points $g(\sigma(x))$, for $g\in G$,
exhaust $\pi^{-1}(x)$.  Thus every point of that fiber is analytically
sectioned.  Apply \eqref{eq:p031-sectioned-defect-contraction} and
\eqref{eq:p031-sectioned-gain-bound} at each such point.  If the base defect
vanishes, the nonnegative cover defect and symmetry defect both vanish.
The full-union assertion follows from
\eqref{eq:p031-full-union-symmetry-defect}.
\end{proof}

\begin{proof}[\Cref{cor:p031-local-section-defect-preservation}]
The group translates of a local section through one point over $x$ give
local sections through every point over $x$, as in the preceding proof.
Hence \eqref{eq:p031-sectioned-globalization-balance} holds on all of $E$.
Its left-hand side vanishes identically exactly when
\eqref{eq:p031-globalization-defect-preservation} holds, proving the
equivalence.

If the base local and global PL envelopes agree, then
$\Delta^{\mathrm{PL}}_{f,X}=0$.  The two nonnegative quantities in the
balance identity must then vanish at every cover point.  This gives cover
local-to-global cofinality and exact full-union compatibility.
\end{proof}


At an interior point, the PL and LFS local families are in fact cofinal for
every obstacle.  The one-sided approximation is local and therefore does not
solve the globalization problem; it does, however, split that problem into
the LFS globalization defect and the global PL--LFS branch gap.

% P03.1 local PL--LFS defect split, live integration version.

\begin{definition}
\label{def:p031-local-lfs-envelope}
Let $Y$ be a good strict $K$-analytic space and let
$g:Y\to\mathbf R\cup\{-\infty\}$.  The local raw LFS envelope is
\begin{equation}
\label{eq:p031-local-lfs-envelope}
 P_{\mathrm{LFS}}^{\mathrm{loc}}g(y)
  :=\sup\{v(y):v\text{ is LFS on an open }V\ni y,
                    \ v\leq g|_V\}.
\end{equation}
The supremum of an empty family is $-\infty$.
\end{definition}

\begin{theorem}
\label{thm:p031-local-pl-lfs-cofinality}
Let $K$ be complete and nontrivially valued, let $Y$ be a good strict
$K$-analytic space, and let $y\in\operatorname{Int}(Y)$.  For every
function $g:Y\to\mathbf R\cup\{-\infty\}$,
\begin{equation}
\label{eq:p031-local-pl-lfs-cofinality}
             P_{\mathrm{PL}}^{\mathrm{loc}}g(y)
                =P_{\mathrm{LFS}}^{\mathrm{loc}}g(y).
\end{equation}
The equality includes the simultaneous empty-family case.  The
neighborhood and the approximating LFS functions may depend on the chosen
PL competitor.
\end{theorem}

\begin{definition}
\label{def:p031-lfs-globalization-defect}
Suppose that $g$ has a global LFS minorant.  Write $s_{\mathrm{LFS},Y}$
and $s_{\mathrm{PL},Y}$ for the global raw LFS and semipositive real-PL
envelopes.  At $y\in\operatorname{Int}(Y)$, define
\begin{align}
 \Delta^{\mathrm{LFS}}_{g,Y}(y)
   &:=P_{\mathrm{LFS}}^{\mathrm{loc}}g(y)
        -s_{\mathrm{LFS},Y}(y),
\label{eq:p031-lfs-globalization-defect}\\
 B_{g,Y}(y)
   &:=s_{\mathrm{PL},Y}(y)-s_{\mathrm{LFS},Y}(y).
\label{eq:p031-pl-lfs-branch-gap}
\end{align}
These are finite nonnegative real numbers.  The second quantity is the
\emph{global PL--LFS branch gap}.
\end{definition}

\begin{theorem}
\label{thm:p031-pl-lfs-defect-split}
Under the hypotheses of
\Cref{def:p031-lfs-globalization-defect}, one has at every
$y\in\operatorname{Int}(Y)$
\begin{equation}
\label{eq:p031-pl-lfs-defect-split}
 \Delta^{\mathrm{PL}}_{g,Y}(y)
   =\Delta^{\mathrm{LFS}}_{g,Y}(y)-B_{g,Y}(y).
\end{equation}
Consequently,
\begin{equation}
\label{eq:p031-branch-gap-globalization-bound}
             0\leq B_{g,Y}(y)
                \leq\Delta^{\mathrm{LFS}}_{g,Y}(y).
\end{equation}
Moreover, $\Delta^{\mathrm{LFS}}_{g,Y}(y)=0$ if and only if
\begin{equation}
\label{eq:p031-lfs-zero-four-envelope-collapse}
 s_{\mathrm{LFS},Y}(y)
  =s_{\mathrm{PL},Y}(y)
  =P_{\mathrm{LFS}}^{\mathrm{loc}}g(y)
  =P_{\mathrm{PL}}^{\mathrm{loc}}g(y).
\end{equation}
Thus local LFS cofinality simultaneously removes the PL globalization
defect and the global PL--LFS branch gap.
\end{theorem}

\begin{corollary}
\label{cor:p031-sectioned-four-defect-balance}
Let $\pi:E\to X$ have a finite analytic orbit presentation between good
boundaryless strict $K$-analytic spaces, let $f$ have a global LFS
minorant, and let $e\in E$ be analytically sectioned over $X$.  Put
$x:=\pi(e)$.  Then
\begin{equation}
\label{eq:p031-sectioned-four-defect-balance}
 \Sigma_{\pi,f}(e)
  =\bigl(\Delta^{\mathrm{LFS}}_{f,X}(x)-B_{f,X}(x)\bigr)
   -\bigl(
      \Delta^{\mathrm{LFS}}_{f\circ\pi,E}(e)
        -B_{f\circ\pi,E}(e)
    \bigr).
\end{equation}
If $\pi$ admits an analytic section on a neighborhood of $x$ and
$\Delta^{\mathrm{LFS}}_{f,X}(x)=0$, then, for every
$e'\in\pi^{-1}(x)$,
\[
 \Delta^{\mathrm{LFS}}_{f\circ\pi,E}(e')
  =\Delta^{\mathrm{PL}}_{f\circ\pi,E}(e')
  =B_{f\circ\pi,E}(e')
  =\Sigma_{\pi,f}(e')=0.
\]
In particular, the full P03 union envelope has exact pullback compatibility
on $\pi^{-1}(x)$.
\end{corollary}

\begin{corollary}
\label{cor:p031-global-lfs-cofinality-closes-orbit-cover}
Keep the finite orbit presentation and obstacle from
\Cref{cor:p031-sectioned-four-defect-balance}.  Suppose that every base
point has an ordinary open neighborhood on which $\pi$ admits an analytic
section.  If
\[
            P_{\mathrm{LFS}}^{\mathrm{loc}}f=s_{\mathrm{LFS},X},
\]
then
\begin{equation}
\label{eq:p031-global-lfs-cofinality-collapse}
 P_{\mathrm{LFS}}^{\mathrm{loc}}(f\circ\pi)
  =P_{\mathrm{PL}}^{\mathrm{loc}}(f\circ\pi)
  =s_{\mathrm{LFS},E}
  =s_{\mathrm{PL},E}
  =s_{\mathrm{LFS},X}\circ\pi
  =s_{\mathrm{PL},X}\circ\pi.
\end{equation}
Hence the full P03 union envelope commutes with pullback.
\end{corollary}

% P03.1 local PL--LFS defect split proof, live integration version.

\begin{proof}[\Cref{thm:p031-local-pl-lfs-cofinality}]
Every LFS function is semipositive real PL, so
\[
 P_{\mathrm{LFS}}^{\mathrm{loc}}g(y)
    \leq P_{\mathrm{PL}}^{\mathrm{loc}}g(y).
\]

Conversely, let $h$ be a semipositive real-PL minorant of $g$ on an
ordinary open neighborhood $V$ of $y$.  Interiority is local on ordinary
open subspaces, so $y\in\operatorname{Int}(V)$.  By
\cite[Proposition~7.10, pp.~190--191]{GJR26}, after shrinking inside $V$
there are a strictly analytic neighborhood $U$ of $y$ and LFS functions
$v_j$ on $U$ converging uniformly to $h|_U$.  Choose an ordinary open
neighborhood $W$ of $y$ contained in $U$ and put
\[
 \varepsilon_j:=\lVert v_j-h\rVert_{U,\infty},
 \qquad
 u_j:=v_j-\varepsilon_j.
\]
The finite-unit description in \cite[\S7.9, p.~190]{GJR26} and the
nontriviality of the valuation show that subtraction of a nonnegative real
constant preserves the LFS class.  On $W$ one has
\[
                 u_j\leq h\leq g,
 \qquad
                 \lVert u_j-h\rVert_{W,\infty}
                     \leq2\varepsilon_j.
\]
Thus $u_j(y)\to h(y)$ through local LFS minorants of $g$, which gives the
reverse inequality in \eqref{eq:p031-local-pl-lfs-cofinality} after taking
the supremum over $h$.

If the local PL family is empty, its LFS subfamily is empty.  If it is
nonempty, the preceding construction produces a local LFS minorant.  Hence
the two empty-family cases are simultaneous.
\end{proof}

The quantities in \Cref{def:p031-lfs-globalization-defect} are finite.
Indeed, a global LFS minorant is also a global semipositive real-PL
minorant.  It forces the obstacle to be real-valued and gives a finite lower
bound for all four envelopes, while the obstacle gives their finite upper
bound.  Inclusion of the global and local competitor families gives
nonnegativity.

\begin{proof}[\Cref{thm:p031-pl-lfs-defect-split}]
By \Cref{thm:p031-local-pl-lfs-cofinality},
\[
       P_{\mathrm{PL}}^{\mathrm{loc}}g(y)
          =P_{\mathrm{LFS}}^{\mathrm{loc}}g(y).
\]
Insert the global LFS envelope between this common local value and the
global PL envelope.  The definitions give
\begin{align*}
 \Delta^{\mathrm{PL}}_{g,Y}(y)
  &=P_{\mathrm{LFS}}^{\mathrm{loc}}g(y)-s_{\mathrm{PL},Y}(y)\\
  &=\bigl(
      P_{\mathrm{LFS}}^{\mathrm{loc}}g(y)-s_{\mathrm{LFS},Y}(y)
    \bigr)
    -\bigl(s_{\mathrm{PL},Y}(y)-s_{\mathrm{LFS},Y}(y)\bigr)\\
  &=\Delta^{\mathrm{LFS}}_{g,Y}(y)-B_{g,Y}(y),
\end{align*}
which is \eqref{eq:p031-pl-lfs-defect-split}.  Since the PL globalization
defect and the branch gap are nonnegative, the same identity gives
\eqref{eq:p031-branch-gap-globalization-bound}.

If the LFS globalization defect vanishes, then its decomposition
\[
       \Delta^{\mathrm{LFS}}_{g,Y}(y)
         =\Delta^{\mathrm{PL}}_{g,Y}(y)+B_{g,Y}(y)
\]
forces both terms on the right to vanish.  The two definitions and local
cofinality give \eqref{eq:p031-lfs-zero-four-envelope-collapse}.  The
converse follows from the first and third terms in that equality.
\end{proof}

\begin{proof}[\Cref{cor:p031-sectioned-four-defect-balance}]
The spaces are boundaryless, so the defect split applies at $x$ and $e$.
Substitute it on both spaces into
\eqref{eq:p031-sectioned-globalization-balance}.  In the orbit setting the
left-hand side is $\Sigma_{\pi,f}(e)$, which proves
\eqref{eq:p031-sectioned-four-defect-balance}.

Now suppose that a local section exists near $x$ and that the base LFS
globalization defect vanishes there.  By
\Cref{thm:p031-pl-lfs-defect-split},
\[
          \Delta^{\mathrm{PL}}_{f,X}(x)=B_{f,X}(x)=0.
\]
The group translates of the local section analytically section every point
of the fiber.  Hence
\Cref{cor:p031-local-section-orbit-fiber-closure} makes the cover PL defect
and $\Sigma_{\pi,f}$ vanish at every $e'\in\pi^{-1}(x)$.

The base branch gap is zero, so
$s_{\mathrm{LFS},X}(x)=s_{\mathrm{PL},X}(x)$.  The envelope chain
\eqref{eq:p031-lfs-pl-envelope-chain}, together with
$\Sigma_{\pi,f}(e')=0$, gives
\[
 s_{\mathrm{LFS},E}(e')
   =s_{\mathrm{PL},E}(e')
   =s_{\mathrm{PL},X}(x)
   =s_{\mathrm{LFS},X}(x).
\]
Thus the cover branch gap vanishes.  The defect split then makes the cover
LFS globalization defect vanish as well.  Exact full-union compatibility
follows from \eqref{eq:p031-full-union-symmetry-defect}.
\end{proof}

\begin{proof}[\Cref{cor:p031-global-lfs-cofinality-closes-orbit-cover}]
Base LFS local-to-global cofinality says that the base LFS globalization
defect vanishes everywhere.  Apply
\Cref{cor:p031-sectioned-four-defect-balance} on each fiber.  It gives
cover LFS and PL local-to-global cofinality, equality of the cover LFS and
PL envelopes, and exact pullback from the equal base envelopes.  Local
PL--LFS cofinality on $E$ supplies the first equality in
\eqref{eq:p031-global-lfs-cofinality-collapse}; the remaining equalities
are the pointwise conclusions just obtained.
\end{proof}


The local-section hypothesis can be removed for finite flat covers.  First,
rational LFS formulas survive fiberwise maximum: the characteristic
polynomial of each analytic unit records its spectral maximum, while common
denominators compress finite sums and maxima.  Second, every real-coefficient
LFS formula is locally cofinal from below by rational formulas.  Together
these facts descend a competitor defined near one cover point after an LFS
baseline fills the rest of the complete fiber.

% P03.1 finite-flat local PL pullback, live integration version.

\begin{definition}
\label{def:p031-rational-lfs}
Let $Y$ be a good strict $K$-analytic space.  A function $u:Y\to\mathbf R$
is \emph{$\mathbf Q$-LFS} if every point has an ordinary open neighborhood
on which
\begin{equation}
\label{eq:p031-rational-lfs-normal-form}
 u=\sum_{\nu=1}^{r}a_\nu
       \max_{1\leq j\leq N_\nu}\{-\log|f_{\nu j}|\},
 \qquad
 a_\nu\in\mathbf Q_{\geq0},\quad
 f_{\nu j}\in\mathcal O_Y^\times.
\end{equation}
The empty sum is allowed.  By \cite[\S7.9, p.~190]{GJR26}, every
$\mathbf Q$-LFS function is LFS.

For $g:Y\to\mathbf R\cup\{-\infty\}$, define the local raw
$\mathbf Q$-LFS envelope by
\begin{equation}
\label{eq:p031-local-rational-lfs-envelope}
 P_{\mathbf Q}^{\mathrm{loc}}g(y)
   :=\sup\{v(y):v\text{ is $\mathbf Q$-LFS on an open }V\ni y,
                    \ v\leq g|_V\}.
\end{equation}
The supremum of an empty family is $-\infty$.
\end{definition}

\begin{lemma}
\label{lem:p031-characteristic-spectral-maximum}
Let $\pi:V\to U$ be a finite locally free morphism of constant positive
degree $d$ between good strict $K$-analytic spaces, and let
$h\in\mathcal O(V)^\times$.  Write the characteristic polynomial of
multiplication by $h$ as
\begin{equation}
\label{eq:p031-characteristic-polynomial}
 \chi_h(T)=T^d+c_1T^{d-1}+\cdots+c_d,
 \qquad c_i\in\mathcal O(U).
\end{equation}
Then, for every $y\in U$,
\begin{equation}
\label{eq:p031-characteristic-spectral-maximum}
 \max_{v\in\pi^{-1}(y)}|h(v)|
       =\max_{1\leq i\leq d}|c_i(y)|^{1/i}.
\end{equation}
Moreover, the function
\begin{equation}
\label{eq:p031-log-spectral-maximum}
 y\longmapsto\log\max_{v\in\pi^{-1}(y)}|h(v)|
\end{equation}
is $\mathbf Q$-LFS on $U$.
\end{lemma}

\begin{theorem}
\label{thm:p031-rational-lfs-fiber-maximum}
Let $\pi:E\to X$ be a finite, flat, surjective morphism of separated good
strict $K$-analytic spaces.  If $u$ is $\mathbf Q$-LFS on $E$, then
\begin{equation}
\label{eq:p031-rational-lfs-fiber-maximum}
       \pi_{\max}u(x):=\max_{e\in\pi^{-1}(x)}u(e)
\end{equation}
is $\mathbf Q$-LFS on $X$.

For any $f:X\to\mathbf R\cup\{-\infty\}$, let $s_{\mathbf Q,X}$ and
$s_{\mathbf Q,E}$ be the global raw $\mathbf Q$-LFS envelopes of $f$ and
$f\circ\pi$, respectively.  Then
\begin{equation}
\label{eq:p031-rational-lfs-envelope-pullback}
                     s_{\mathbf Q,E}
                       =s_{\mathbf Q,X}\circ\pi.
\end{equation}
This equality includes simultaneous empty competitor families.
\end{theorem}

\begin{proposition}
\label{prop:p031-local-rational-lfs-cofinality}
Assume that $K$ is nontrivially valued.  Let $u$ be LFS on a good strict
$K$-analytic space $Y$, and fix $y\in Y$.  There are a compact strictly
analytic neighborhood $U$ of $y$, chosen arbitrarily small, and
$\mathbf Q$-LFS functions $v_n$ on $U$ such that
\begin{equation}
\label{eq:p031-one-sided-local-rationalization}
                    u-\frac1n<v_n\leq u
                    \qquad\text{on }U.
\end{equation}
Consequently, for every $g:Y\to\mathbf R\cup\{-\infty\}$,
\begin{equation}
\label{eq:p031-local-real-rational-lfs-cofinality}
             P_{\mathrm{LFS}}^{\mathrm{loc}}g
                  =P_{\mathbf Q}^{\mathrm{loc}}g.
\end{equation}
The equality includes simultaneous empty competitor families.
\end{proposition}

\begin{theorem}
\label{thm:p031-finite-flat-local-envelope-pullback}
Let $K$ be complete and nontrivially valued, and let
$\pi:E\to X$ be a finite, flat, surjective morphism of separated good
strict $K$-analytic spaces.  Let
$f:X\to\mathbf R\cup\{-\infty\}$, fix $x\in X$, and choose
$e\in\pi^{-1}(x)$.  Suppose that $f\circ\pi$ has an LFS minorant on an
ordinary open neighborhood of the complete fiber $\pi^{-1}(x)$.  Then
\begin{equation}
\label{eq:p031-finite-flat-local-lfs-pullback}
 P_{\mathrm{LFS}}^{\mathrm{loc}}(f\circ\pi)(e)
       =P_{\mathrm{LFS}}^{\mathrm{loc}}f(x).
\end{equation}
If $E$ and $X$ are boundaryless, then also
\begin{equation}
\label{eq:p031-finite-flat-local-pl-pullback}
 P_{\mathrm{PL}}^{\mathrm{loc}}(f\circ\pi)(e)
       =P_{\mathrm{PL}}^{\mathrm{loc}}f(x).
\end{equation}
No analytic section through $e$ is assumed.
\end{theorem}

\begin{theorem}
\label{thm:p031-finite-flat-pl-defect-balance}
Keep the finite flat morphism above, assume that $E$ and $X$ are
boundaryless, and suppose that $f$ has a global LFS minorant.  For
$e\in E$, put $x:=\pi(e)$ and
\[
 \Gamma^{\mathrm{PL}}_{\pi,f}(e)
   :=s_{\mathrm{PL},E}(e)-s_{\mathrm{PL},X}(x).
\]
Then $\Gamma^{\mathrm{PL}}_{\pi,f}(e)$ is finite and nonnegative, and
\begin{align}
 \Gamma^{\mathrm{PL}}_{\pi,f}(e)
   &=\Delta^{\mathrm{PL}}_{f,X}(x)
       -\Delta^{\mathrm{PL}}_{f\circ\pi,E}(e),
\label{eq:p031-finite-flat-pl-defect-balance}\\
 0\leq\Delta^{\mathrm{PL}}_{f\circ\pi,E}(e)
   &\leq\Delta^{\mathrm{PL}}_{f,X}(x),
\label{eq:p031-finite-flat-pl-defect-contraction}\\
 0\leq\Gamma^{\mathrm{PL}}_{\pi,f}(e)
   &\leq\Delta^{\mathrm{PL}}_{f,X}(x).
\label{eq:p031-finite-flat-pl-gain-bound}
\end{align}
The full P03 raw-envelope pullback defect is
$\Gamma^{\mathrm{PL}}_{\pi,f}$.  Hence exact pullback compatibility holds
if and only if
\begin{equation}
\label{eq:p031-finite-flat-pl-defect-preservation}
 \Delta^{\mathrm{PL}}_{f\circ\pi,E}
     =\Delta^{\mathrm{PL}}_{f,X}\circ\pi.
\end{equation}
If $\Delta^{\mathrm{PL}}_{f,X}(x)=0$, then both the cover PL defect and
$\Gamma^{\mathrm{PL}}_{\pi,f}$ vanish at every point of $\pi^{-1}(x)$.
\end{theorem}

\begin{corollary}
\label{cor:p031-finite-orbit-four-defect-balance}
Suppose that $\pi:E\to X$ has a finite analytic orbit presentation between
separated good boundaryless strict $K$-analytic spaces, and let $f$ have a
global LFS minorant.  For every $e\in E$, with $x:=\pi(e)$,
\begin{equation}
\label{eq:p031-finite-orbit-four-defect-balance}
 \Sigma_{\pi,f}(e)
  =\bigl(\Delta^{\mathrm{LFS}}_{f,X}(x)-B_{f,X}(x)\bigr)
   -\bigl(
      \Delta^{\mathrm{LFS}}_{f\circ\pi,E}(e)
        -B_{f\circ\pi,E}(e)
    \bigr).
\end{equation}
If $\Delta^{\mathrm{LFS}}_{f,X}(x)=0$, then, for every
$e'\in\pi^{-1}(x)$,
\[
 \Delta^{\mathrm{LFS}}_{f\circ\pi,E}(e')
  =\Delta^{\mathrm{PL}}_{f\circ\pi,E}(e')
  =B_{f\circ\pi,E}(e')
  =\Sigma_{\pi,f}(e')=0.
\]
Thus the full P03 raw envelope has exact pullback compatibility on the
complete fiber, without a local-section hypothesis.
\end{corollary}

% P03.1 finite-flat local PL pullback proof, live integration version.

\begin{proof}[\Cref{lem:p031-characteristic-spectral-maximum}]
Fix $y\in U$, put $H:=\mathcal H(y)$, and base change
\eqref{eq:p031-characteristic-polynomial} to the finite $H$-algebra of the
analytic fiber.  After extension to an algebraic closure of $H$, its roots,
counted with multiplicity, are the residue-field images of $h$ at the points
of $\pi^{-1}(y)$.  Indeed, the fiber algebra is the product of its local
Artinian factors.  On each factor multiplication by $h$ is triangular after
scalar extension, with diagonal entries given by the residue-field
conjugates of $h$.  Nilpotents change multiplicities but introduce no new
eigenvalues.  Uniqueness of the extended absolute value on every finite
extension of the complete field $H$ shows that the largest root modulus is
\[
                       R_y:=\max_{v\in\pi^{-1}(y)}|h(v)|.
\]

Write the roots as $\alpha_1,\ldots,\alpha_d$.  Since $c_i(y)$ is, up to
sign, their $i$-th elementary symmetric function, the ultrametric inequality
gives
\[
                           |c_i(y)|\leq R_y^i.
\]
Conversely, choose a root $\alpha$ with $|\alpha|=R_y$.  If
$|c_i(y)|^{1/i}<R_y$ for every $i$, then in
\[
       \alpha^d+c_1(y)\alpha^{d-1}+\cdots+c_d(y)=0
\]
the first term would have strictly larger absolute value than every other
term, contradicting the ultrametric triangle law.  This proves
\eqref{eq:p031-characteristic-spectral-maximum}.

Multiplication by the unit $h$ is an automorphism of the finite locally free
module $\pi_*\mathcal O_V$, so $c_d$ is a unit.  Fix $y\in U$ and let
\[
                         I_y:=\{i:c_i(y)\ne0\}.
\]
After shrinking around $y$, each $c_i$ with $i\in I_y$ is a unit.  For
$i\notin I_y$, continuity and the nonvanishing of $c_d$ give
\[
                 |c_i(z)|^{1/i}<|c_d(z)|^{1/d}
\]
on a smaller neighborhood.  These indices do not contribute to the
spectral maximum there.  With $N:=\operatorname{lcm}(1,\ldots,d)$,
\begin{equation}
\label{eq:p031-local-characteristic-lfs-formula}
 \log R_z
  =\max_{i\in I_y}\frac1i\log|c_i(z)|
  =\frac1N\max_{i\in I_y}
       \{-\log|c_i(z)^{-N/i}|\}.
\end{equation}
This is a $\mathbf Q$-LFS presentation, proving the final assertion.
\end{proof}

\begin{proof}[\Cref{thm:p031-rational-lfs-fiber-maximum}]
Fix $x\in X$ and write $\pi^{-1}(x)=\{e_1,\ldots,e_q\}$.  Choose pairwise
disjoint ordinary open neighborhoods $V_k$ of the $e_k$ on which $u$ has
a presentation \eqref{eq:p031-rational-lfs-normal-form}.  Since a finite
morphism is proper,
\[
 U_0:=X\setminus\pi\bigl(E\setminus(V_1\cup\cdots\cup V_q)\bigr)
\]
is an open neighborhood of $x$, and $\pi^{-1}(U_0)$ is the disjoint union
of the $V_k\cap\pi^{-1}(U_0)$.

A finite flat analytic morphism is open by
\cite[proof of Proposition~5.11, p.~175]{GJR26}.  Shrink $U_0$ to the
intersection of the images of these finitely many branches and then along
their locally constant rank strata.  We obtain an open neighborhood $U$ of
$x$ and a decomposition
\begin{equation}
\label{eq:p031-local-finite-flat-branch-decomposition}
                 \pi^{-1}(U)=E_1\amalg\cdots\amalg E_q
\end{equation}
such that each $E_k\to U$ is finite locally free, surjective, and of a
fixed positive degree $d_k$.

On $E_k$, take a common denominator in the selected rational presentation
and distribute the finite maxima.  This gives
\begin{equation}
\label{eq:p031-rational-lfs-block-compression}
 u=\frac1{D_k}\max_{F\in\mathcal F_k}\{-\log|F|\},
 \qquad D_k\in\mathbf Z_{>0},
\end{equation}
where $\mathcal F_k$ is a finite family of analytic units on $E_k$.  For
the empty sum, take $\mathcal F_k=\{1\}$.  Applying
\Cref{lem:p031-characteristic-spectral-maximum} to $F^{-1}$ on every
branch gives
\begin{equation}
\label{eq:p031-branchwise-spectral-compression}
 \pi_{\max}u(z)
  =\max_{1\leq k\leq q}\max_{F\in\mathcal F_k}
      \frac1{D_k}
      \log\max_{e\in(E_k)_z}|F(e)^{-1}|.
\end{equation}
Each function inside the finite maxima is $\mathbf Q$-LFS.  A finite
maximum of rational presentations is again $\mathbf Q$-LFS: compress each
presentation as in \eqref{eq:p031-rational-lfs-block-compression}, take a
common denominator, and raise the finitely many analytic units to the
corresponding positive integral powers.  Hence
\eqref{eq:p031-branchwise-spectral-compression} is $\mathbf Q$-LFS on
$U$.  Since $x$ was arbitrary, the fiber-maximum assertion follows.

Pullback sends a base $\mathbf Q$-LFS minorant of $f$ to a cover minorant,
so $s_{\mathbf Q,X}\circ\pi\leq s_{\mathbf Q,E}$.  Conversely, if
$u\leq f\circ\pi$ is a cover competitor, then $\pi_{\max}u$ is a base
$\mathbf Q$-LFS competitor and
\[
 u(e)\leq\pi_{\max}u(\pi(e))
       \leq s_{\mathbf Q,X}(\pi(e)).
\]
Taking the supremum over $u$ proves the reverse inequality in
\eqref{eq:p031-rational-lfs-envelope-pullback}.  It also proves that the
two competitor families are empty simultaneously.
\end{proof}

\begin{proof}[\Cref{prop:p031-local-rational-lfs-cofinality}]
By \cite[\S7.9, p.~190]{GJR26}, after shrinking around $y$ one may write
\[
 u=\sum_{\nu=1}^{r}a_\nu\phi_\nu,
 \qquad
 \phi_\nu=\max_{1\leq j\leq N_\nu}\{-\log|f_{\nu j}|\},
 \qquad a_\nu\in\mathbf R_{\geq0},
\]
with analytic units $f_{\nu j}$.  Choose an arbitrarily small compact
strictly analytic neighborhood $U$ of $y$ inside this presentation domain.
The finitely many functions $\phi_\nu$ are bounded on $U$.

Put $\eta_n:=1/(4n)$ and choose $a_{\nu,n}\in\mathbf Q_{\geq0}$ so that
\[
 \left\|\sum_{\nu=1}^{r}(a_{\nu,n}-a_\nu)\phi_\nu
 \right\|_{U,\infty}<\eta_n.
\]
Set $w_n:=\sum_\nu a_{\nu,n}\phi_\nu$.  Nontriviality of the valuation
gives $\gamma\in K^\times$ with $\tau:=-\log|\gamma|>0$.  Choose
$q_n\in\mathbf Q_{>0}$ with
\[
                         \eta_n<q_n\tau<2\eta_n.
\]
The constant $-q_n\tau$ is $\mathbf Q$-LFS because
\[
              -q_n\tau=q_n\max\{-\log|\gamma^{-1}|\}.
\]
Thus $v_n:=w_n-q_n\tau$ is $\mathbf Q$-LFS and satisfies
$u-3\eta_n<v_n<u$ on $U$, which proves
\eqref{eq:p031-one-sided-local-rationalization}.

The rational local competitor family is contained in the LFS family.  For
the reverse inequality, apply the ordered approximation to an arbitrary
local LFS minorant of $g$.  Restrict $v_n$ to an ordinary open neighborhood
of $y$ contained in the interior of $U$.  It remains below that minorant and
hence below $g$, while $v_n(y)\to u(y)$.  Taking the supremum proves
\eqref{eq:p031-local-real-rational-lfs-cofinality}.  If the LFS family is
empty, so is its rational subfamily, proving the final clause.
\end{proof}

\begin{proof}[\Cref{thm:p031-finite-flat-local-envelope-pullback}]
Pullback of local LFS minorants gives
\begin{equation}
\label{eq:p031-finite-flat-easy-local-inequality}
 P_{\mathrm{LFS}}^{\mathrm{loc}}f(x)
   \leq P_{\mathrm{LFS}}^{\mathrm{loc}}(f\circ\pi)(e).
\end{equation}

For the reverse inequality, let $u$ be an LFS minorant of $f\circ\pi$ on
an ordinary open neighborhood of $e$.  Write
$\pi^{-1}(x)=\{e_1,\ldots,e_q\}$ with $e_1=e$, and let $b_0$ be the
assumed LFS minorant on a neighborhood of the complete fiber.  Choose
pairwise disjoint ordinary open neighborhoods $W_k$ of the $e_k$ such that
$u$ is defined on $W_1$ and $b_0$ is defined on every $W_k$ for $k>1$.
On their disjoint union, set
\[
 \widetilde u|_{W_1}:=u,
 \qquad
 \widetilde u|_{W_k}:=b_0\quad(k>1).
\]
This is an LFS minorant of $f\circ\pi$ near the complete fiber.

Fix $\varepsilon>0$.  Apply
\Cref{prop:p031-local-rational-lfs-cofinality} at each $e_k$, and restrict
the resulting compact neighborhoods to pairwise disjoint ordinary open
neighborhoods.  Their union is a neighborhood $W'$ of the complete fiber
carrying a $\mathbf Q$-LFS function $v$ such that
\begin{equation}
\label{eq:p031-complete-fiber-rational-minorant}
             v\leq\widetilde u\leq f\circ\pi,
             \qquad v(e)>u(e)-\varepsilon.
\end{equation}
Since $\pi$ is finite and hence proper, after shrinking to an ordinary open
neighborhood $U$ of $x$ one has $\pi^{-1}(U)\subset W'$.  By
\Cref{thm:p031-rational-lfs-fiber-maximum},
$m:=\pi_{\max}(v|_{\pi^{-1}(U)})$ is $\mathbf Q$-LFS on $U$, and
\[
                         m\leq f|_U,
             \qquad m(x)\geq v(e)>u(e)-\varepsilon.
\]
Therefore $u(e)-\varepsilon<P_{\mathrm{LFS}}^{\mathrm{loc}}f(x)$.
Letting $\varepsilon\downarrow0$ and taking the supremum over $u$ proves
the reverse of \eqref{eq:p031-finite-flat-easy-local-inequality}.

If both spaces are boundaryless, every point is interior.  Apply
\Cref{thm:p031-local-pl-lfs-cofinality} on $E$ and $X$ to the two sides of
\eqref{eq:p031-finite-flat-local-lfs-pullback}; this gives
\eqref{eq:p031-finite-flat-local-pl-pullback}.
\end{proof}

\begin{proof}[\Cref{thm:p031-finite-flat-pl-defect-balance}]
A global LFS minorant pulls back to a global cover LFS minorant, so
\Cref{thm:p031-finite-flat-local-envelope-pullback} applies at every cover
point.  It also makes all local and global envelope values finite.  Pullback
of global PL competitors gives
\[
                   s_{\mathrm{PL},X}\circ\pi
                      \leq s_{\mathrm{PL},E},
\]
so $\Gamma^{\mathrm{PL}}_{\pi,f}\geq0$.

Fix $e\in E$, put $x:=\pi(e)$, and denote the common value in
\eqref{eq:p031-finite-flat-local-pl-pullback} by $P_x$.  Then
\begin{align*}
 \Delta^{\mathrm{PL}}_{f,X}(x)
   -\Delta^{\mathrm{PL}}_{f\circ\pi,E}(e)
  &=(P_x-s_{\mathrm{PL},X}(x))
      -(P_x-s_{\mathrm{PL},E}(e))\\
  &=s_{\mathrm{PL},E}(e)-s_{\mathrm{PL},X}(x)\\
  &=\Gamma^{\mathrm{PL}}_{\pi,f}(e),
\end{align*}
which proves \eqref{eq:p031-finite-flat-pl-defect-balance}.  Since the
cover defect and the gain are nonnegative, the same identity gives
\eqref{eq:p031-finite-flat-pl-defect-contraction} and
\eqref{eq:p031-finite-flat-pl-gain-bound}.

Every LFS function is semipositive real PL, so the full P03 union envelope
is the PL envelope.  Its pullback defect is therefore the displayed gain,
which vanishes exactly when
\eqref{eq:p031-finite-flat-pl-defect-preservation} holds.  If the base
defect vanishes at $x$, the two nonnegative summands
\[
 \Delta^{\mathrm{PL}}_{f\circ\pi,E}(e)
       +\Gamma^{\mathrm{PL}}_{\pi,f}(e)
       =\Delta^{\mathrm{PL}}_{f,X}(x)
\]
vanish for every $e\in\pi^{-1}(x)$.
\end{proof}

\begin{proof}[\Cref{cor:p031-finite-orbit-four-defect-balance}]
For a finite analytic orbit presentation, the global PL pullback gain is
$\Sigma_{\pi,f}$ by \eqref{eq:p031-full-union-symmetry-defect}.  Substitute
the split identity \eqref{eq:p031-pl-lfs-defect-split} on $X$ and $E$ into
\eqref{eq:p031-finite-flat-pl-defect-balance}.  This gives
\eqref{eq:p031-finite-orbit-four-defect-balance} at every cover point.

If the base LFS globalization defect vanishes at $x$, then
\Cref{thm:p031-pl-lfs-defect-split} gives
\[
             \Delta^{\mathrm{PL}}_{f,X}(x)=B_{f,X}(x)=0.
\]
The finite-flat defect balance makes the cover PL defect and
$\Sigma_{\pi,f}$ vanish throughout the fiber.  The envelope chain
\eqref{eq:p031-lfs-pl-envelope-chain}, together with the zero base branch
gap and zero symmetry defect, then forces the cover branch gap to vanish.
The split identity makes the cover LFS defect vanish as well.  Exact
full-envelope pullback compatibility follows from
\eqref{eq:p031-full-union-symmetry-defect}.
\end{proof}


The rational subclass also carries a conserved globalization ledger.  Its
local and global envelopes both commute with finite flat pullback, so its
local-to-global defect is unchanged.  The full LFS and PL pullback gains are
therefore exactly the changes in the two successive global branch gaps.

% P03.1 rational globalization-defect conservation, live integration version.

\begin{definition}
\label{def:p031-rational-globalization-ledger}
Let $Y$ be a good strict $K$-analytic space, let
$g:Y\to\mathbf R\cup\{-\infty\}$, and suppose that $g$ has a global
$\mathbf Q$-LFS minorant.  Write $s_{\mathbf Q,Y}$ for the global raw
$\mathbf Q$-LFS envelope.  Define the \emph{global rational-coefficient
gap} by
\begin{equation}
\label{eq:p031-rational-lfs-global-gap}
 A_{g,Y}:=s_{\mathrm{LFS},Y}-s_{\mathbf Q,Y}
\end{equation}
and the \emph{rational globalization defect} by
\begin{equation}
\label{eq:p031-rational-globalization-defect}
 \Delta^{\mathbf Q}_{g,Y}
   :=P_{\mathbf Q}^{\mathrm{loc}}g-s_{\mathbf Q,Y}.
\end{equation}
Both are finite nonnegative real-valued functions.
\end{definition}

\begin{theorem}
\label{thm:p031-rational-defect-conservation}
Let $K$ be complete and nontrivially valued, and let
$\pi:E\to X$ be a finite, flat, surjective morphism of separated good
boundaryless strict $K$-analytic spaces.  Suppose that
$f:X\to\mathbf R\cup\{-\infty\}$ has a global $\mathbf Q$-LFS
minorant.  On either $(Y,g)=(X,f)$ or
$(Y,g)=(E,f\circ\pi)$, one has
\begin{equation}
\label{eq:p031-three-level-defect-ledger}
 \Delta^{\mathbf Q}_{g,Y}
   =\Delta^{\mathrm{LFS}}_{g,Y}+A_{g,Y}
   =\Delta^{\mathrm{PL}}_{g,Y}+A_{g,Y}+B_{g,Y}.
\end{equation}
For every $e\in E$, with $x:=\pi(e)$,
\begin{equation}
\label{eq:p031-rational-defect-conservation}
 \Delta^{\mathbf Q}_{f\circ\pi,E}(e)
       =\Delta^{\mathbf Q}_{f,X}(x).
\end{equation}

Define the global LFS pullback gain by
\[
 \Gamma^{\mathrm{LFS}}_{\pi,f}(e)
   :=s_{\mathrm{LFS},E}(e)-s_{\mathrm{LFS},X}(x).
\]
Then
\begin{align}
 0\leq\Gamma^{\mathrm{LFS}}_{\pi,f}(e)
   &=A_{f\circ\pi,E}(e)-A_{f,X}(x) \notag\\
   &=\Delta^{\mathrm{LFS}}_{f,X}(x)
      -\Delta^{\mathrm{LFS}}_{f\circ\pi,E}(e),
\label{eq:p031-finite-flat-lfs-gap-balance}\\
 0\leq\Gamma^{\mathrm{PL}}_{\pi,f}(e)
   &=\bigl(A_{f\circ\pi,E}(e)+B_{f\circ\pi,E}(e)\bigr)
      -\bigl(A_{f,X}(x)+B_{f,X}(x)\bigr) \notag\\
   &=\Delta^{\mathrm{PL}}_{f,X}(x)
      -\Delta^{\mathrm{PL}}_{f\circ\pi,E}(e).
\label{eq:p031-finite-flat-pl-two-gap-balance}
\end{align}
Consequently,
\begin{align}
 s_{\mathrm{LFS},E}=s_{\mathrm{LFS},X}\circ\pi
 &\quad\Longleftrightarrow\quad
 A_{f\circ\pi,E}=A_{f,X}\circ\pi,
\label{eq:p031-lfs-compatibility-gap-criterion}\\
 s_{\mathrm{PL},E}=s_{\mathrm{PL},X}\circ\pi
 &\quad\Longleftrightarrow\quad
 A_{f\circ\pi,E}+B_{f\circ\pi,E}
   =(A_{f,X}+B_{f,X})\circ\pi.
\label{eq:p031-p03-compatibility-total-gap-criterion}
\end{align}
The second equivalence is exact pullback compatibility for the full P03
raw envelope.
\end{theorem}

\begin{corollary}
\label{cor:p031-rational-cofinality-fiber-collapse}
Keep the hypotheses above.  If
$\Delta^{\mathbf Q}_{f,X}(x)=0$, then, at $x$ and at every
$e\in\pi^{-1}(x)$, the rational, LFS, and PL globalization defects and
the two global branch gaps $A$ and $B$ all vanish.  In particular,
\begin{equation}
\label{eq:p031-rational-cofinality-fiber-collapse}
 \begin{split}
 s_{\mathbf Q,E}(e)
   &=s_{\mathrm{LFS},E}(e)
    =s_{\mathrm{PL},E}(e)\\
   &=s_{\mathbf Q,X}(x)
    =s_{\mathrm{LFS},X}(x)
    =s_{\mathrm{PL},X}(x),
 \end{split}
\end{equation}
and their common value is also the value of each local rational, LFS, and
PL envelope at the corresponding point.  Thus pointwise rational LFS local-to-global cofinality
on the base forces exact full P03 compatibility on the complete fiber.
\end{corollary}

\begin{corollary}
\label{cor:p031-orbit-rational-gap-decomposition}
If, in addition, $\pi$ has a finite analytic orbit presentation, then
for every $e\in E$, with $x:=\pi(e)$,
\begin{equation}
\label{eq:p031-orbit-rational-gap-decomposition}
 \Sigma_{\pi,f}(e)
   =\bigl(A_{f\circ\pi,E}(e)-A_{f,X}(x)\bigr)
      +\bigl(B_{f\circ\pi,E}(e)-B_{f,X}(x)\bigr).
\end{equation}
Hence the symmetry-breaking defect vanishes exactly when the total global
gain from rational LFS to semipositive real PL is preserved.
\end{corollary}

% P03.1 rational globalization-defect conservation proof, live integration.

The quantities in \Cref{def:p031-rational-globalization-ledger} are finite.
Indeed, a global $\mathbf Q$-LFS minorant is also a global LFS and
semipositive real-PL minorant.  It forces the obstacle to be real-valued and
gives a finite lower bound for every envelope under consideration, while
the obstacle gives a finite upper bound.  Inclusion of the rational
competitor family in the LFS family gives $A_{g,Y}\geq0$, and restriction
of global rational competitors gives
$\Delta^{\mathbf Q}_{g,Y}\geq0$.

\begin{proof}[\Cref{thm:p031-rational-defect-conservation}]
By \eqref{eq:p031-local-real-rational-lfs-cofinality}, the local rational
and LFS envelopes agree.  Boundarylessness and
\Cref{thm:p031-local-pl-lfs-cofinality} identify this common value with the
local PL envelope.  Insert the three global envelopes
\[
             s_{\mathbf Q,Y}\leq s_{\mathrm{LFS},Y}
                    \leq s_{\mathrm{PL},Y}
\]
below the common local value.  The definitions of the three defects and
the two gaps then give
\eqref{eq:p031-three-level-defect-ledger}.

Let $q_0$ be a global $\mathbf Q$-LFS minorant of $f$.  Its pullback is a
global $\mathbf Q$-LFS minorant of $f\circ\pi$ and, in particular, an
LFS baseline near every complete fiber.  Therefore
\Cref{thm:p031-finite-flat-local-envelope-pullback} and
\eqref{eq:p031-local-real-rational-lfs-cofinality} give
\[
 P_{\mathbf Q}^{\mathrm{loc}}(f\circ\pi)(e)
   =P_{\mathrm{LFS}}^{\mathrm{loc}}(f\circ\pi)(e)
   =P_{\mathrm{LFS}}^{\mathrm{loc}}f(x)
   =P_{\mathbf Q}^{\mathrm{loc}}f(x).
\]
The global rational envelope identity
\eqref{eq:p031-rational-lfs-envelope-pullback} gives
\[
                         s_{\mathbf Q,E}(e)
                              =s_{\mathbf Q,X}(x).
\]
Subtracting the global identity from the local identity proves
\eqref{eq:p031-rational-defect-conservation}.

Since $s_{\mathrm{LFS}}=s_{\mathbf Q}+A$, exact rational pullback gives
\[
 \Gamma^{\mathrm{LFS}}_{\pi,f}(e)
   =A_{f\circ\pi,E}(e)-A_{f,X}(x).
\]
Pullback of every global base LFS competitor is a cover competitor, so this
gain is nonnegative.  Subtracting the two LFS decompositions in
\eqref{eq:p031-three-level-defect-ledger} from the conserved rational
defect gives the last equality in
\eqref{eq:p031-finite-flat-lfs-gap-balance}.

Similarly, $s_{\mathrm{PL}}=s_{\mathbf Q}+A+B$, so exact rational
pullback gives the first equality in
\eqref{eq:p031-finite-flat-pl-two-gap-balance}.  Pullback of global PL
competitors makes this gain nonnegative.  Subtracting the two PL
decompositions in \eqref{eq:p031-three-level-defect-ledger} gives the
last equality, equivalently recovering
\eqref{eq:p031-finite-flat-pl-defect-balance}.

The LFS envelopes commute exactly when their gain vanishes, which is the
first gap-preservation criterion.  The full P03 raw envelope is the PL
envelope, and its gain vanishes exactly when the total gap $A+B$ is
preserved.  This proves
\eqref{eq:p031-lfs-compatibility-gap-criterion} and
\eqref{eq:p031-p03-compatibility-total-gap-criterion}.
\end{proof}

\begin{proof}[\Cref{cor:p031-rational-cofinality-fiber-collapse}]
By \eqref{eq:p031-rational-defect-conservation}, the cover rational defect
vanishes at every point over $x$.  Each instance of
\eqref{eq:p031-three-level-defect-ledger} writes
\[
       0=\Delta^{\mathbf Q}
          =\Delta^{\mathrm{LFS}}+A
          =\Delta^{\mathrm{PL}}+A+B,
\]
where all displayed terms are nonnegative.  Hence every term vanishes on
the base and the complete fiber.  The definitions, local cofinality, and
exact rational pullback now give
\eqref{eq:p031-rational-cofinality-fiber-collapse} and the asserted local
equalities.
\end{proof}

\begin{proof}[\Cref{cor:p031-orbit-rational-gap-decomposition}]
For a finite analytic orbit presentation,
$\Gamma^{\mathrm{PL}}_{\pi,f}=\Sigma_{\pi,f}$ by
\eqref{eq:p031-full-union-symmetry-defect}.  Substitute this identity in
\eqref{eq:p031-finite-flat-pl-two-gap-balance}.  The vanishing criterion is
the zero case of the resulting equality.
\end{proof}


The conservation law is additive in finite flat towers.  Since every
one-step gain is nonnegative, exact compatibility for the composite forces
exact compatibility at each intermediate stage; later covers cannot cancel
an earlier defect.

\medskip
% P03.1 finite-flat tower additivity and no-cancellation, live integration.

\begin{theorem}
\label{thm:p031-finite-flat-tower-defect-additivity}
Let $K$ be complete and nontrivially valued, and let
\[
 X_r\xrightarrow{\pi_r}X_{r-1}\longrightarrow\cdots
   \longrightarrow X_1\xrightarrow{\pi_1}X_0
\]
be a finite tower, with $r\geq1$, of finite, flat, surjective morphisms of
separated good boundaryless strict $K$-analytic spaces.  Put
$\Pi_i:=\pi_1\circ\cdots\circ\pi_i$, with $\Pi_0:=\mathrm{id}_{X_0}$,
and let $f_i:=f\circ\Pi_i$, where
$f:X_0\to\mathbf R\cup\{-\infty\}$ has a global $\mathbf Q$-LFS
minorant.

For $\mathcal C\in\{\mathbf Q,\mathrm{LFS},\mathrm{PL}\}$, write
$s_{\mathcal C,i}$ for the global raw $\mathcal C$-envelope of $f_i$.
Set
\[
 A_i:=A_{f_i,X_i},\qquad B_i:=B_{f_i,X_i},\qquad
 \Delta_i^{\mathcal C}:=\Delta_{f_i,X_i}^{\mathcal C}.
\]
Choose a compatible chain of points
$x_i\in X_i$ with $\pi_i(x_i)=x_{i-1}$.  Then
\begin{equation}
\label{eq:p031-tower-rational-defect-conservation}
 \Delta_i^{\mathbf Q}(x_i)=\Delta_0^{\mathbf Q}(x_0)
 \qquad(0\leq i\leq r).
\end{equation}
If
\begin{align*}
 \Gamma_i^{\mathrm{LFS}}(x_i)
   &:=s_{\mathrm{LFS},i}(x_i)-s_{\mathrm{LFS},i-1}(x_{i-1}),\\
 \Gamma_i^{\mathrm{PL}}(x_i)
   &:=s_{\mathrm{PL},i}(x_i)-s_{\mathrm{PL},i-1}(x_{i-1}),
\end{align*}
then every step gain is nonnegative and
\begin{align}
 0\leq\sum_{i=1}^{r}\Gamma_i^{\mathrm{LFS}}(x_i)
   &=s_{\mathrm{LFS},r}(x_r)-s_{\mathrm{LFS},0}(x_0) \notag\\
   &=A_r(x_r)-A_0(x_0) \notag\\
   &=\Delta_0^{\mathrm{LFS}}(x_0)
       -\Delta_r^{\mathrm{LFS}}(x_r),
\label{eq:p031-tower-lfs-gain-additivity}\\
 0\leq\sum_{i=1}^{r}\Gamma_i^{\mathrm{PL}}(x_i)
   &=s_{\mathrm{PL},r}(x_r)-s_{\mathrm{PL},0}(x_0) \notag\\
   &=\bigl(A_r(x_r)+B_r(x_r)\bigr)
       -\bigl(A_0(x_0)+B_0(x_0)\bigr) \notag\\
   &=\Delta_0^{\mathrm{PL}}(x_0)
       -\Delta_r^{\mathrm{PL}}(x_r).
\label{eq:p031-tower-pl-gain-additivity}
\end{align}
In either line, the total gain vanishes if and only if every step gain
vanishes.
\end{theorem}

\begin{corollary}
\label{cor:p031-finite-flat-tower-no-cancellation}
Keep the hypotheses above.  Exact compatibility at the top of the tower is
equivalent to exact compatibility at every stage:
\begin{align}
 s_{\mathrm{LFS},r}=s_{\mathrm{LFS},0}\circ\Pi_r
 &\quad\Longleftrightarrow\quad
 s_{\mathrm{LFS},i}=s_{\mathrm{LFS},i-1}\circ\pi_i
 \quad(1\leq i\leq r),
\label{eq:p031-tower-lfs-no-cancellation}\\
 s_{\mathrm{PL},r}=s_{\mathrm{PL},0}\circ\Pi_r
 &\quad\Longleftrightarrow\quad
 s_{\mathrm{PL},i}=s_{\mathrm{PL},i-1}\circ\pi_i
 \quad(1\leq i\leq r).
\label{eq:p031-tower-pl-no-cancellation}
\end{align}
Thus neither an LFS defect nor a full P03 defect at an intermediate finite
flat stage can be cancelled by later stages.
\end{corollary}

\begin{corollary}
\label{cor:p031-rational-cofinality-tower-collapse}
If $\Delta_0^{\mathbf Q}(x_0)=0$, then for every compatible lift chain and
every $0\leq i\leq r$, the three globalization defects
$\Delta_i^{\mathbf Q}$, $\Delta_i^{\mathrm{LFS}}$, and
$\Delta_i^{\mathrm{PL}}$, as well as $A_i$ and $B_i$, vanish at $x_i$.
Moreover,
\begin{equation}
\label{eq:p031-rational-cofinality-tower-collapse}
 \begin{split}
 s_{\mathbf Q,i}(x_i)
  &=s_{\mathrm{LFS},i}(x_i)=s_{\mathrm{PL},i}(x_i)\\
  &=P_{\mathbf Q}^{\mathrm{loc}}f_i(x_i)
    =P_{\mathrm{LFS}}^{\mathrm{loc}}f_i(x_i)
    =P_{\mathrm{PL}}^{\mathrm{loc}}f_i(x_i)\\
  &=s_{\mathbf Q,0}(x_0).
 \end{split}
\end{equation}
\end{corollary}

% P03.1 finite-flat tower additivity and no-cancellation proof.

\begin{proof}[\Cref{thm:p031-finite-flat-tower-defect-additivity}]
The chosen global rational LFS minorant of $f$ pulls back to a global
rational LFS minorant of every $f_i$.  Apply
\Cref{thm:p031-rational-defect-conservation} to
$\pi_i:X_i\to X_{i-1}$ with obstacle $f_{i-1}$.  Its rational-defect
identity gives
\[
 \Delta_i^{\mathbf Q}(x_i)=\Delta_{i-1}^{\mathbf Q}(x_{i-1}).
\]
Iteration proves \eqref{eq:p031-tower-rational-defect-conservation}.

The same theorem gives, at every stage,
\begin{align*}
 0\leq\Gamma_i^{\mathrm{LFS}}(x_i)
   &=A_i(x_i)-A_{i-1}(x_{i-1})\\
   &=\Delta_{i-1}^{\mathrm{LFS}}(x_{i-1})
       -\Delta_i^{\mathrm{LFS}}(x_i),\\
 0\leq\Gamma_i^{\mathrm{PL}}(x_i)
   &=\bigl(A_i(x_i)+B_i(x_i)\bigr)
       -\bigl(A_{i-1}(x_{i-1})+B_{i-1}(x_{i-1})\bigr)\\
   &=\Delta_{i-1}^{\mathrm{PL}}(x_{i-1})
       -\Delta_i^{\mathrm{PL}}(x_i).
\end{align*}
Summing telescopes the envelope values, the gaps, and the defects, which
proves
\eqref{eq:p031-tower-lfs-gain-additivity} and
\eqref{eq:p031-tower-pl-gain-additivity}.  A finite sum of nonnegative real
numbers is zero exactly when every summand is zero.
\end{proof}

\begin{proof}[\Cref{cor:p031-finite-flat-tower-no-cancellation}]
If every stage is compatible, composition gives compatibility at the top.
Conversely, suppose first that
$s_{\mathrm{LFS},r}=s_{\mathrm{LFS},0}\circ\Pi_r$.  Fix
$1\leq i\leq r$ and $x_i\in X_i$.  If $i=r$, take $x_r=x_i$; otherwise,
surjectivity of the later maps supplies a compatible lift $x_r$ of $x_i$.
The preceding points are forced by the tower.  The composite LFS gain along
this chain is zero.  Every step gain is
nonnegative by
\eqref{eq:p031-tower-lfs-gain-additivity}, so the $i$-th gain vanishes at
$x_i$.  Since $x_i$ was arbitrary, the $i$-th LFS envelopes commute.  This
holds for every $i$ and proves
\eqref{eq:p031-tower-lfs-no-cancellation}.

The same argument, using
\eqref{eq:p031-tower-pl-gain-additivity}, proves
\eqref{eq:p031-tower-pl-no-cancellation}.  The PL envelope is the full P03
raw envelope, so the second assertion is precisely the claimed P03
no-cancellation statement.
\end{proof}

\begin{proof}[\Cref{cor:p031-rational-cofinality-tower-collapse}]
Equation \eqref{eq:p031-tower-rational-defect-conservation} makes the rational
defect zero at every $x_i$.  Apply
\Cref{cor:p031-rational-cofinality-fiber-collapse} successively to the edges
of the compatible chain.  It makes the LFS and PL defects and both gaps
vanish at every $x_i$, identifies the three local and three global envelope
values there, and identifies the global values across each edge.  Iteration
gives \eqref{eq:p031-rational-cofinality-tower-collapse}.
\end{proof}


The conserved rational defect itself admits an exact boundary description.
A local rational competitor with excess $\delta$ over the global rational
envelope must exceed every global rational competitor by at least $\delta$
on the boundary of every sufficiently small domain.  Conversely, continuity
recovers the pointwise excess as an infimum of these boundary loads.  This
turns rational local-to-global cofinality into a collar-absorption criterion
and gives a concrete boundary-saturation hypothesis that forces the whole
defect ledger to vanish.

\medskip
% P03.1 rational collar minimax and boundary saturation, live integration.

\begin{lemma}
\label{lem:p031-rational-protected-collar-gluing}
Let $X$ be a good strict $K$-analytic space, let
$f:X\to\mathbf R\cup\{-\infty\}$, and let $b$ be a global
$\mathbf Q$-LFS minorant of $f$.  Suppose that $D,V\subset X$ are
ordinary open subsets with $\overline V\subset D$, and that $u$ is a
$\mathbf Q$-LFS minorant of $f|_D$ satisfying
\begin{equation}
\label{eq:p031-rational-protected-collar-order}
                         u\leq b
                 \qquad\text{on }D\setminus\overline V.
\end{equation}
Then
\begin{equation}
\label{eq:p031-rational-protected-collar-gluing}
 H:=
 \begin{cases}
    \max\{b,u\},&\text{on }D,\\
    b,&\text{on }X\setminus\overline V
 \end{cases}
\end{equation}
is a global $\mathbf Q$-LFS minorant of $f$.
\end{lemma}

\begin{lemma}
\label{lem:p031-small-nonempty-boundary-domain}
Let $X$ be a connected separated good $K$-analytic space with at least two
points.  If $x\in W\subset X$ and $W$ is ordinary open, then there is an
ordinary open subset $D$ such that
\begin{equation}
\label{eq:p031-small-nonempty-boundary-domain}
 x\in D,\qquad \overline D\text{ is compact},\qquad
 \overline D\subset W,\qquad \partial D\ne\varnothing.
\end{equation}
Moreover, $D$ may be chosen inside any prescribed ordinary neighborhood
of $x$ contained in $W$.
\end{lemma}

\begin{theorem}
\label{thm:p031-rational-collar-minimax}
Let $K$ be complete and nontrivially valued, let $X$ be a connected
separated good strict $K$-analytic space with at least two points, and let
$f:X\to\mathbf R\cup\{-\infty\}$ be upper semicontinuous with a global
$\mathbf Q$-LFS minorant.  For an ordinary open set $U\subset X$, put
\begin{equation}
\label{eq:p031-rational-minorant-family}
 \mathcal Q_f(U):=
 \{v:v\text{ is $\mathbf Q$-LFS on }U,\ v\leq f|_U\}.
\end{equation}
Fix $x\in X$, an ordinary open neighborhood $W$ of $x$, and
$v\in\mathcal Q_f(W)$ such that
\begin{equation}
\label{eq:p031-positive-rational-globalization-gap}
             \delta_x(v):=v(x)-s_{\mathbf Q,X}(x)>0.
\end{equation}
Define its rational collar load by
\begin{equation}
\label{eq:p031-rational-local-collar-load}
 \Lambda_x(v;W):=
 \inf_{\substack{x\in D,\ \overline D\subset W,\ \overline D\ \mathrm{compact}\\
                  D\ \mathrm{ordinary\ open},\ \partial D\ne\varnothing}}
 \ \inf_{b\in\mathcal Q_f(X)}
                  \max_{\partial D}(v-b).
\end{equation}
Then
\begin{equation}
\label{eq:p031-rational-collar-minimax}
                  \Lambda_x(v;W)=\delta_x(v).
\end{equation}
Consequently the complete rational globalization defect has the exact
collar representation
\begin{equation}
\label{eq:p031-rational-defect-collar-formula}
 \Delta^{\mathbf Q}_{f,X}(x)
  =\sup_{\substack{x\in W\ \mathrm{ordinary\ open},\ v\in\mathcal Q_f(W)\\
                    v(x)>s_{\mathbf Q,X}(x)}}
                         \Lambda_x(v;W),
\end{equation}
where the supremum of the empty family is declared to be $0$.
\end{theorem}

\begin{corollary}
\label{cor:p031-rational-collar-absorption}
Under the hypotheses of \Cref{thm:p031-rational-collar-minimax}, the
following are equivalent:
\begin{enumerate}[label=\textup{(\roman*)}]
\item
$\Delta^{\mathbf Q}_{f,X}(x)=0$;
\item
for every ordinary open neighborhood $W$ of $x$, every
$v\in\mathcal Q_f(W)$, and every $\varepsilon>0$, there are an ordinary
open $D$ and a global $b\in\mathcal Q_f(X)$ such that
\begin{equation}
\label{eq:p031-rational-collar-absorption}
 x\in D,\quad \overline D\subset W,\quad \overline D\text{ is compact},
 \quad\partial D\ne\varnothing,
 \quad\max_{\partial D}(v-b)<\varepsilon.
\end{equation}
\end{enumerate}
\end{corollary}

\begin{definition}
\label{def:p031-rational-boundary-saturation}
In the setting of \Cref{thm:p031-rational-collar-minimax}, the obstacle
$f$ is \emph{rationally boundary-saturated from below at $x$} if, for every
ordinary open neighborhood $W$ of $x$ and every $\varepsilon>0$, there are
an ordinary open $D$ and a global $b\in\mathcal Q_f(X)$ such that
\begin{equation}
\label{eq:p031-rational-boundary-saturation}
 x\in D,\quad \overline D\subset W,\quad \overline D\text{ is compact},
 \quad\partial D\ne\varnothing,
 \quad\max_{\partial D}(f-b)<\varepsilon.
\end{equation}
The maximum exists because $f-b$ is real-valued upper semicontinuous on
the compact set $\partial D$.
\end{definition}

\begin{corollary}
\label{cor:p031-rational-boundary-saturation-collapse}
Keep the hypotheses of \Cref{thm:p031-rational-collar-minimax}, and assume
in addition that $X$ is boundaryless.  If $f$ is rationally
boundary-saturated from below at $x$, then
\begin{equation}
\label{eq:p031-rational-boundary-saturation-zero-ledger}
 \Delta^{\mathbf Q}_{f,X}(x)
 =\Delta^{\mathrm{LFS}}_{f,X}(x)
 =\Delta^{\mathrm{PL}}_{f,X}(x)
 =A_{f,X}(x)=B_{f,X}(x)=0,
\end{equation}
and all six envelope values agree:
\begin{equation}
\label{eq:p031-rational-boundary-saturation-collapse}
 \begin{split}
 s_{\mathbf Q,X}(x)&=s_{\mathrm{LFS},X}(x)=s_{\mathrm{PL},X}(x)\\
 &=P_{\mathbf Q}^{\mathrm{loc}}f(x)
   =P_{\mathrm{LFS}}^{\mathrm{loc}}f(x)
   =P_{\mathrm{PL}}^{\mathrm{loc}}f(x).
 \end{split}
\end{equation}
\end{corollary}

\begin{corollary}
\label{cor:p031-rational-boundary-saturation-tower-collapse}
Let $K$ be complete and nontrivially valued, and let
\[
 X_r\xrightarrow{\pi_r}X_{r-1}\longrightarrow\cdots
   \longrightarrow X_1\xrightarrow{\pi_1}X_0
\]
be a finite tower, with $r\geq1$, of finite, flat, surjective morphisms of
separated good boundaryless strict $K$-analytic spaces.  Suppose that $X_0$
is connected
with at least two points, that $f:X_0\to\mathbf R\cup\{-\infty\}$ is upper
semicontinuous with a global $\mathbf Q$-LFS minorant, and that $f$ is
rationally boundary-saturated from below at $x_0\in X_0$.  With the notation
of \Cref{thm:p031-finite-flat-tower-defect-additivity}, for every
$0\leq i\leq r$ and every $y\in\Pi_i^{-1}(x_0)$ one has
\begin{equation}
\label{eq:p031-rational-boundary-saturation-tower-zero}
 \Delta_i^{\mathbf Q}(y)=\Delta_i^{\mathrm{LFS}}(y)
 =\Delta_i^{\mathrm{PL}}(y)=A_i(y)=B_i(y)=0,
\end{equation}
and
\begin{equation}
\label{eq:p031-rational-boundary-saturation-tower-collapse}
 \begin{split}
 s_{\mathbf Q,i}(y)&=s_{\mathrm{LFS},i}(y)=s_{\mathrm{PL},i}(y)\\
 &=P_{\mathbf Q}^{\mathrm{loc}}f_i(y)
   =P_{\mathrm{LFS}}^{\mathrm{loc}}f_i(y)
   =P_{\mathrm{PL}}^{\mathrm{loc}}f_i(y)
   =s_{\mathbf Q,0}(x_0).
 \end{split}
\end{equation}
Thus the rational, LFS, and full P03 envelopes have exact pullback
compatibility throughout the complete inverse image of $x_0$ at every
stage of the tower.
\end{corollary}

% P03.1 rational collar minimax and boundary saturation proof, live integration.

\begin{proof}[\Cref{lem:p031-rational-protected-collar-gluing}]
The two displayed formulas agree on the overlap
$D\setminus\overline V$ by
\eqref{eq:p031-rational-protected-collar-order}, so they define a function
$H$ on the ordinary open cover
\[
                         X=D\cup(X\setminus\overline V).
\]
The $\mathbf Q$-LFS class is local by
\Cref{def:p031-rational-lfs}.  It is closed under finite maxima: locally
compress the two rational presentations as in
\eqref{eq:p031-rational-lfs-block-compression}, take a common denominator,
and combine the resulting finite families of analytic units.  Hence $H$ is
$\mathbf Q$-LFS.  Both $b$ and $u$ lie below $f$ on their domains, so the
same is true of $H$.
\end{proof}

\begin{proof}[\Cref{lem:p031-small-nonempty-boundary-domain}]
A separated good analytic space is locally compact Hausdorff and therefore
regular.  Since a connected Hausdorff space with at least two points has no
isolated point, the prescribed neighborhood contains a point $y\ne x$.
Regularity and local compactness give an ordinary open neighborhood $D$ of
$x$ whose compact closure is contained in the prescribed neighborhood
minus $\{y\}$.  If $\partial D$ were empty, then $D$ would be a nonempty
clopen proper subset of $X$, contradicting connectedness.  Thus
$\partial D\ne\varnothing$.
\end{proof}

\begin{proof}[\Cref{thm:p031-rational-collar-minimax}]
All functions in $\mathcal Q_f(U)$ are continuous.  A global
$\mathbf Q$-LFS minorant forces $f$ to be real-valued, and it also makes
$s_{\mathbf Q,X}$ and $P_{\mathbf Q}^{\mathrm{loc}}f$ finite.  In
particular, every maximum and difference used below is finite.

Fix an admissible $D$ and $b\in\mathcal Q_f(X)$.  We first prove
\begin{equation}
\label{eq:p031-rational-collar-lower-bound}
                       \max_{\partial D}(v-b)\geq\delta_x(v).
\end{equation}
Suppose instead that this maximum is smaller than $\delta_x(v)$.  Choose
$\gamma\in K^\times$ with $\tau:=-\log|\gamma|>0$.  Since
$\mathbf Q\tau$ is dense in $\mathbf R$, there is a number
\begin{equation}
\label{eq:p031-rational-collar-shift-choice}
 c=q\tau,\qquad q\in\mathbf Q_{>0},\qquad
 \max\bigl\{0,\max_{\partial D}(v-b)\bigr\}<c<\delta_x(v).
\end{equation}
The constant $-c$ is $\mathbf Q$-LFS because
\[
                   -c=q\max\{-\log|\gamma^{-1}|\}.
\]

Put $\phi:=v-b$ on $\overline D$.  Compactness of $\partial D$ and
continuity give a relative open neighborhood $N$ of $\partial D$ in
$\overline D$ on which $\phi<c$.  On the other hand,
\[
 \phi(x)=v(x)-b(x)\geq v(x)-s_{\mathbf Q,X}(x)
                         =\delta_x(v)>c,
\]
so $x\notin N$.  The compact set $\overline D\setminus N$ lies in $D$.
Choose an ordinary open $V$ with
\[
 x\in\overline D\setminus N\subset V,
                    \qquad\overline V\subset D.
\]
Then $v-c<b$ on $D\setminus\overline V$.  Apply
\Cref{lem:p031-rational-protected-collar-gluing} to $u:=v-c$.  It produces
a global $\mathbf Q$-LFS minorant $H$ of $f$ satisfying
\[
                    H(x)\geq v(x)-c>s_{\mathbf Q,X}(x),
\]
contrary to the definition of the global rational envelope.  This proves
\eqref{eq:p031-rational-collar-lower-bound} and hence
$\Lambda_x(v;W)\geq\delta_x(v)$.

For the reverse inequality, fix $\varepsilon>0$.  Choose
$b\in\mathcal Q_f(X)$ with
\[
                 b(x)>s_{\mathbf Q,X}(x)-\frac{\varepsilon}{3}.
\]
Continuity of $v-b$ gives an ordinary neighborhood $N$ of $x$, with
$N\subset W$, such that
\[
 (v-b)(z)<(v-b)(x)+\frac{\varepsilon}{3}
       <\delta_x(v)+\frac{2\varepsilon}{3}
                  \qquad(z\in N).
\]
By \Cref{lem:p031-small-nonempty-boundary-domain}, choose an admissible
$D$ with $\overline D\subset N$.  Then
\[
                  \max_{\partial D}(v-b)
                       <\delta_x(v)+\varepsilon.
\]
Taking the two infima and then letting $\varepsilon\downarrow0$ proves
\eqref{eq:p031-rational-collar-minimax}.

Finally, the definition of the local rational envelope gives
\[
 \Delta^{\mathbf Q}_{f,X}(x)
   =\sup_{\substack{x\in W\ \mathrm{ordinary\ open}\\
                     v\in\mathcal Q_f(W)}}
          \bigl(v(x)-s_{\mathbf Q,X}(x)\bigr).
\]
The left side is nonnegative.  If it is positive, the subfamily with
positive gap has the same supremum; on that subfamily
\eqref{eq:p031-rational-collar-minimax} replaces the gap by the collar
load.  If it is zero, the positive-gap subfamily is empty.  This proves
\eqref{eq:p031-rational-defect-collar-formula} with the stated convention.
\end{proof}

\begin{proof}[\Cref{cor:p031-rational-collar-absorption}]
Assume first that $\Delta^{\mathbf Q}_{f,X}(x)=0$.  Then every
$v\in\mathcal Q_f(W)$ satisfies $v(x)\leq s_{\mathbf Q,X}(x)$.  Given
$\varepsilon>0$, choose $b\in\mathcal Q_f(X)$ with
$b(x)>s_{\mathbf Q,X}(x)-\varepsilon/3$.  Continuity gives a neighborhood
$N\subset W$ of $x$ on which $v-b<2\varepsilon/3$.  Choose $D$ with
$\overline D\subset N$ and nonempty boundary by
\Cref{lem:p031-small-nonempty-boundary-domain}.  This gives
\eqref{eq:p031-rational-collar-absorption}.

Conversely, if the defect were positive, there would be a pair $(W,v)$
with $\delta_x(v)>0$.  By
\eqref{eq:p031-rational-collar-minimax}, every admissible $D$ and $b$ would
satisfy
\[
                         \max_{\partial D}(v-b)\geq\delta_x(v),
\]
contradicting \eqref{eq:p031-rational-collar-absorption} with
$\varepsilon:=\delta_x(v)/2$.
\end{proof}

\begin{proof}[\Cref{cor:p031-rational-boundary-saturation-collapse}]
For $v\in\mathcal Q_f(W)$ one has $v\leq f$, so rational boundary
saturation implies the collar-absorption condition
\eqref{eq:p031-rational-collar-absorption}.  Hence
\Cref{cor:p031-rational-collar-absorption} gives
$\Delta^{\mathbf Q}_{f,X}(x)=0$.  Apply the three-level ledger
\eqref{eq:p031-three-level-defect-ledger}.  It writes zero as sums of the
five nonnegative defects and gaps in
\eqref{eq:p031-rational-boundary-saturation-zero-ledger}; all of them
therefore vanish.  Their definitions and the local cofinality identities
\eqref{eq:p031-local-real-rational-lfs-cofinality} and
\Cref{thm:p031-local-pl-lfs-cofinality} give
\eqref{eq:p031-rational-boundary-saturation-collapse}.
\end{proof}

\begin{proof}[\Cref{cor:p031-rational-boundary-saturation-tower-collapse}]
By \Cref{cor:p031-rational-boundary-saturation-collapse},
$\Delta_0^{\mathbf Q}(x_0)=0$.  Given
$y\in\Pi_i^{-1}(x_0)$, its successive images form a compatible lift chain.
Apply \Cref{cor:p031-rational-cofinality-tower-collapse} to that chain.
This proves \eqref{eq:p031-rational-boundary-saturation-tower-zero} and
\eqref{eq:p031-rational-boundary-saturation-tower-collapse}.  Since $y$
was arbitrary, the conclusion holds on the complete inverse image at every
stage.
\end{proof}


The collar formula also forces a topological propagation law.  A positive
rational globalization gap cannot be carried by an isolated or trapped
superlevel component: below every positive gap level and inside every
smaller relatively compact domain, its connected high-defect carrier must
reach the surrounding boundary.  Conversely, one such scale-by-scale escape
system, for one competitor and one positive level across every smaller
domain, characterizes positivity of the complete rational defect.

\medskip
% P03.1 rational-defect escape criterion, live integration.

\begin{lemma}
\label{lem:p031-compact-component-collar}
Let $Y$ be a locally compact Hausdorff space, let $F\subset Y$ be compact,
and let $C$ be a connected component of $F$.  If $O\subset Y$ is open,
$C\subset O$, and $\overline O$ is compact, then there is an open subset
$D\subset Y$ such that
\begin{equation}
\label{eq:p031-compact-component-collar}
 C\subset D,\qquad \overline D\subset O,\qquad
 \overline D\text{ is compact},\qquad \partial D\cap F=\varnothing.
\end{equation}
\end{lemma}

\begin{theorem}
\label{thm:p031-rational-defect-escape}
Keep the hypotheses and notation of
\Cref{thm:p031-rational-collar-minimax}.  Fix an ordinary open
neighborhood $W$ of $x$ and $v\in\mathcal Q_f(W)$ with
\begin{equation}
\label{eq:p031-rational-defect-escape-positive-gap}
             \delta_x(v)=v(x)-s_{\mathbf Q,X}(x)>0.
\end{equation}
Put
\begin{equation}
\label{eq:p031-rational-defect-profile}
              g_v:=v-s_{\mathbf Q,X}\quad\text{on }W.
\end{equation}
Let $0<\alpha<\delta_x(v)$, and let $U$ be any ordinary open subset
satisfying
\begin{equation}
\label{eq:p031-rational-defect-test-domain}
 x\in U,\qquad \overline U\subset W,\qquad
 \overline U\text{ is compact},\qquad \partial U\ne\varnothing.
\end{equation}
Then the connected component $C_{\alpha,U}(x)$ of $x$ in the compact set
\begin{equation}
\label{eq:p031-rational-defect-superlevel-carrier}
 E_{\alpha,U}(v):=
 \{z\in\overline U:g_v(z)\geq\alpha\}
\end{equation}
meets $\partial U$.  In particular, every positive rational globalization
gap carries, at every smaller scale and at every level below the gap, a
compact connected subset joining $x$ to the surrounding boundary on which
$v-s_{\mathbf Q,X}\geq\alpha$.
\end{theorem}

\begin{corollary}
\label{cor:p031-rational-defect-escape-characterization}
Under the hypotheses of \Cref{thm:p031-rational-collar-minimax}, the
following are equivalent:
\begin{enumerate}[label=\textup{(\roman*)}]
\item
$\Delta^{\mathbf Q}_{f,X}(x)>0$;
\item
there are an ordinary open neighborhood $W$ of $x$, a function
$v\in\mathcal Q_f(W)$, and a number $\alpha>0$ such that
$g_v:=v-s_{\mathbf Q,X}$ satisfies $g_v(x)\geq\alpha$ and, for every $U$
satisfying
\eqref{eq:p031-rational-defect-test-domain}, the component
$C_{\alpha,U}(x)$ in \eqref{eq:p031-rational-defect-superlevel-carrier}
meets $\partial U$.
\end{enumerate}
Thus vanishing of the rational globalization defect is equivalent to the
absence of such a scale-by-scale escaping high-defect carrier system.  When
$X$ is also
boundaryless, this topological condition is equivalent, through
\eqref{eq:p031-three-level-defect-ledger}, to simultaneous vanishing of the
rational, LFS, and PL defects and both global branch gaps at $x$.
\end{corollary}

% P03.1 rational-defect escape criterion proof, live integration.

\begin{proof}[\Cref{lem:p031-compact-component-collar}]
We first recall the needed compact-Hausdorff separation fact.  Fix $c\in C$
and let $Q$ be the intersection of all clopen subsets of $F$ containing
$c$.  Every one of those subsets contains $C$, so $C\subset Q$.  The set
$Q$ is connected.  Indeed, if $Q=A\amalg B$ were a separation with
$c\in A$, compact Hausdorff normality would give disjoint open subsets
$G,H\subset F$ containing $A,B$, respectively.  The compact set
$F\setminus(G\cup H)$ is disjoint from $Q$.  For each of its points, the
definition of $Q$ supplies a clopen neighborhood of $c$ omitting that
point.  A finite intersection gives a clopen neighborhood $N$ of $c$ with
$N\subset G\cup H$.  Then $N\cap G$ and $N\cap H$ are complementary open
subsets of the clopen set $N$.  Thus $N\cap G$ is clopen in $F$, contains
$c$, and omits $B$, contradicting the definition of $Q$.  Hence $Q$ is connected,
so maximality of the component gives $Q=C$.

If $F\setminus O$ is empty, put $A:=F$.  Otherwise, for each point of the
compact set $F\setminus O$, choose a clopen subset of $F$ that contains $C$
and omits that point.  A finite intersection gives a clopen set
$A\subset F$ with
\[
                         C\subset A\subset O.
\]
Both $A$ and $F\setminus A$ are compact.  The open set
$O\setminus(F\setminus A)$ contains $A$.  Local compactness and regularity,
followed by a finite subcover of $A$, give an open set $D$ such that
\[
 A\subset D\subset\overline D
             \subset O\setminus(F\setminus A)
\]
and $\overline D$ is compact.  Its boundary misses $A$ because $A\subset D$
and misses $F\setminus A$ because its closure does.  This proves
\eqref{eq:p031-compact-component-collar}.
\end{proof}

\begin{proof}[\Cref{thm:p031-rational-defect-escape}]
The global family $\mathcal Q_f(X)$ is nonempty, and its members are
continuous.  Hence its pointwise supremum $s_{\mathbf Q,X}$ is lower
semicontinuous.  The function $g_v=v-s_{\mathbf Q,X}$ is therefore upper
semicontinuous on $W$, so $E_{\alpha,U}(v)$ is a compact subset of
$\overline U$.  It contains $x$ by
\eqref{eq:p031-rational-defect-escape-positive-gap}.

Suppose that $C_{\alpha,U}(x)$ does not meet $\partial U$.  It is then
contained in $U$.  Apply \Cref{lem:p031-compact-component-collar} with
$F=E_{\alpha,U}(v)$ and $O=U$.  We obtain an ordinary open set $D$ such
that
\begin{equation}
\label{eq:p031-rational-defect-trapping-domain}
 x\in D,\qquad \overline D\subset U,\qquad
 \overline D\text{ is compact},\qquad
 \partial D\cap E_{\alpha,U}(v)=\varnothing.
\end{equation}
The boundary of $D$ is nonempty.  Otherwise $D$ would be a nonempty clopen
proper subset of the connected space $X$; it is proper because
$\overline D\subset U$ and $\partial U\ne\varnothing$.

For each $y\in\partial D$, the last condition in
\eqref{eq:p031-rational-defect-trapping-domain} gives
\[
                    s_{\mathbf Q,X}(y)>v(y)-\alpha.
\]
Choose $b_y\in\mathcal Q_f(X)$ with $b_y(y)>v(y)-\alpha$.  Since
$v-b_y$ is continuous on $W$, it is smaller than $\alpha$ on a neighborhood
of $y$.  Finitely many of these
neighborhoods cover the compact set $\partial D$.  The maximum $b$ of the
corresponding finitely many functions $b_y$ is again a global
$\mathbf Q$-LFS minorant of $f$, and it satisfies
\begin{equation}
\label{eq:p031-rational-defect-subcritical-boundary-load}
                         \max_{\partial D}(v-b)<\alpha.
\end{equation}
On the other hand, the lower-bound half of
\Cref{thm:p031-rational-collar-minimax} applies to $D$ and $b$ and gives
\[
             \max_{\partial D}(v-b)\geq\delta_x(v)>\alpha,
\]
contradicting
\eqref{eq:p031-rational-defect-subcritical-boundary-load}.  Thus
$C_{\alpha,U}(x)$ meets $\partial U$.  It is a connected component of a
compact Hausdorff space and is therefore compact, which gives the final
assertion.
\end{proof}

\begin{proof}[\Cref{cor:p031-rational-defect-escape-characterization}]
Suppose first that $\Delta^{\mathbf Q}_{f,X}(x)>0$.  By the definition of
the local rational envelope, there are $W$ and $v\in\mathcal Q_f(W)$ with
$\delta_x(v)>0$.  Choose $0<\alpha<\delta_x(v)$ and apply
\Cref{thm:p031-rational-defect-escape} to every admissible $U$.  This proves
\textup{(ii)}.

Conversely, \textup{(ii)} includes $g_v(x)\geq\alpha>0$.  Since the local
rational envelope dominates $v$ at $x$, one has
\[
 \Delta^{\mathbf Q}_{f,X}(x)
   =P_{\mathbf Q}^{\mathrm{loc}}f(x)-s_{\mathbf Q,X}(x)
   \geq v(x)-s_{\mathbf Q,X}(x)
   =g_v(x)>0.
\]
This proves \textup{(i)}.  The last assertion follows from the equivalence
just proved and the nonnegative three-level defect ledger
\eqref{eq:p031-three-level-defect-ledger}.
\end{proof}


The escape property has a comparison form that no longer fixes the central
point.  The positive part of the difference between any local rational
minorant and the global rational envelope satisfies the boundary maximum
principle on every relatively compact domain.  Moreover, the truncated
components from the preceding theorem all lie in one intrinsic component;
that single carrier crosses every smaller boundary and is not relatively
compact in its local domain.

\medskip
% P03.1 rational-defect maximum principle, live integration.

\begin{theorem}
\label{thm:p031-rational-defect-maximum-principle}
Keep the hypotheses and notation of
\Cref{thm:p031-rational-collar-minimax}.  Fix an ordinary open subset
$W\subset X$ and $v\in\mathcal Q_f(W)$, and put
\[
                       g_v:=v-s_{\mathbf Q,X}.
\]
For $t\in\mathbf R$, write $t_+:=\max\{t,0\}$.  If $U\subset X$ is
ordinary open and satisfies
\begin{equation}
\label{eq:p031-rational-defect-maximum-principle-domain}
 \overline U\subset W,\qquad \overline U\text{ is compact},\qquad
 \partial U\ne\varnothing,
\end{equation}
then
\begin{equation}
\label{eq:p031-rational-defect-maximum-principle}
        \max_{\overline U}(g_v)_+
                  =\max_{\partial U}(g_v)_+.
\end{equation}
Consequently, for every $c\geq0$,
\begin{equation}
\label{eq:p031-rational-defect-boundary-comparison}
 v\leq s_{\mathbf Q,X}+c\text{ on }\partial U
 \quad\Longrightarrow\quad
 v\leq s_{\mathbf Q,X}+c\text{ on }\overline U.
\end{equation}
In particular, a local rational minorant that is dominated by the global
rational envelope on the boundary is dominated by it throughout the
domain.
\end{theorem}

\begin{corollary}
\label{cor:p031-rational-defect-coherent-carrier}
In the setting of \Cref{thm:p031-rational-defect-maximum-principle}, fix
$x\in W$ and assume that $g_v(x)>0$.  For $0<\alpha<g_v(x)$, let
$\Gamma_{\alpha,W}(x)$ be the connected component of $x$ in the closed
subset
\begin{equation}
\label{eq:p031-rational-defect-intrinsic-superlevel}
             \{z\in W:g_v(z)\geq\alpha\}\subset W.
\end{equation}
Then
\begin{equation}
\label{eq:p031-rational-defect-coherent-boundary-crossing}
       \Gamma_{\alpha,W}(x)\cap\partial U\ne\varnothing
\end{equation}
for every ordinary open $U$ such that
\[
 x\in U,\qquad \overline U\subset W,\qquad
 \overline U\text{ is compact},\qquad \partial U\ne\varnothing.
\]
Consequently $\Gamma_{\alpha,W}(x)$ is not relatively compact in $W$.
If $\overline W$ is compact, then
\begin{equation}
\label{eq:p031-rational-defect-coherent-boundary-limit}
       \overline{\Gamma_{\alpha,W}(x)}^{\,X}\cap\partial W
                         \ne\varnothing.
\end{equation}
Thus one intrinsic connected high-defect carrier, rather than a
scale-dependent collection of components, reaches the boundary of every
smaller test domain.
\end{corollary}

\begin{corollary}
\label{cor:p031-rational-defect-end-characterization}
Under the hypotheses of \Cref{thm:p031-rational-collar-minimax}, the
following are equivalent:
\begin{enumerate}[label=\textup{(\roman*)}]
\item
$\Delta^{\mathbf Q}_{f,X}(x)>0$;
\item
there are an ordinary open neighborhood $W$ of $x$, a function
$v\in\mathcal Q_f(W)$, and a number $\alpha>0$ such that, with
$g_v:=v-s_{\mathbf Q,X}$, one has $\alpha<g_v(x)$ and the component
$\Gamma_{\alpha,W}(x)$ in
\eqref{eq:p031-rational-defect-intrinsic-superlevel} is not relatively
compact in $W$.
\end{enumerate}
Hence vanishing of the complete rational globalization defect is equivalent
to the absence of a connected positive-defect carrier that is not relatively
compact.  If $X$ is boundaryless, this absence is equivalent, through
\eqref{eq:p031-three-level-defect-ledger}, to simultaneous vanishing of the
rational, LFS, and PL defects and both global branch gaps at $x$.
\end{corollary}

% P03.1 rational-defect maximum principle proof, live integration.

\begin{proof}[\Cref{thm:p031-rational-defect-maximum-principle}]
The function $g_v$ is upper semicontinuous: $v$ is continuous, while
$s_{\mathbf Q,X}$ is the pointwise supremum of the nonempty family of
continuous global rational minorants and is therefore lower
semicontinuous.  Thus both maxima in
\eqref{eq:p031-rational-defect-maximum-principle} exist, and the left one
dominates the right one.

Put $M:=\max_{\overline U}(g_v)_+$.  If $M=0$, equality follows.  Suppose
that $M>0$, and choose $z\in\overline U$ with $(g_v(z))_+=M$.
If $z\in\partial U$, the right-hand maximum is at least $M$, and equality
follows.  Otherwise $z\in U$ and
$g_v(z)=M>0$.  For every $0<\alpha<M$, apply
\Cref{thm:p031-rational-defect-escape} with base point $z$ and test domain
$U$.  The component supplied there meets $\partial U$, so
\[
                       \max_{\partial U}(g_v)_+\geq\alpha.
\]
Letting $\alpha\uparrow M$ proves the reverse inequality in
\eqref{eq:p031-rational-defect-maximum-principle}.

If the left condition in
\eqref{eq:p031-rational-defect-boundary-comparison} holds, then
$\max_{\partial U}(g_v)_+\leq c$.  The maximum principle gives
$(g_v)_+\leq c$ on $\overline U$, and hence $g_v\leq c$ there.  This is
the asserted comparison.
\end{proof}

\begin{proof}[\Cref{cor:p031-rational-defect-coherent-carrier}]
Fix a test domain $U$ satisfying the displayed conditions in the statement.
The connected component of $x$ in
\[
       \{g_v\geq\alpha\}\cap\overline U
\]
is a connected subset of
\eqref{eq:p031-rational-defect-intrinsic-superlevel} containing $x$; it is
therefore contained in $\Gamma_{\alpha,W}(x)$.  By
\Cref{thm:p031-rational-defect-escape}, that smaller component meets
$\partial U$.  This proves
\eqref{eq:p031-rational-defect-coherent-boundary-crossing}.

Suppose that $\Gamma_{\alpha,W}(x)$ is relatively compact in $W$.  The
superlevel set in
\eqref{eq:p031-rational-defect-intrinsic-superlevel} is closed in $W$, and
each of its connected components is closed.  Hence
$\Gamma_{\alpha,W}(x)$ is compact.  It cannot equal $W$: otherwise $W$
would be compact and hence closed in the Hausdorff space $X$.  Since it is
also open and $X$ is connected, one would have $W=X$; but then $v$ would be
a global rational minorant and
$g_v(x)\leq0$.

Choose $y\in W\setminus\Gamma_{\alpha,W}(x)$.  Local compactness and
regularity, followed by a finite subcover of the compact carrier, give an
ordinary open set $U$ such that
\[
 \Gamma_{\alpha,W}(x)\subset U,\quad
 \overline U\text{ is compact},\quad\text{and}\quad
 \overline U\subset W\setminus\{y\}.
\]
Its boundary is nonempty, since otherwise $U$ would be a nonempty clopen
proper subset of the connected space $X$.  The boundary-crossing property
now gives
$\Gamma_{\alpha,W}(x)\cap\partial U\ne\varnothing$, contradicting
$\Gamma_{\alpha,W}(x)\subset U$.  Thus the carrier is not relatively
compact in $W$.

Finally assume that $\overline W$ is compact.  If the closure of the
carrier in $X$ missed $\partial W$, it would be a compact subset of $W$.
This would make the carrier relatively compact in $W$, which has just been
excluded.  This proves
\eqref{eq:p031-rational-defect-coherent-boundary-limit}.
\end{proof}

\begin{proof}[\Cref{cor:p031-rational-defect-end-characterization}]
If $\Delta^{\mathbf Q}_{f,X}(x)>0$, the definition of the local rational
envelope supplies $W$ and $v\in\mathcal Q_f(W)$ with $g_v(x)>0$.  Choose
$0<\alpha<g_v(x)$ and apply
\Cref{cor:p031-rational-defect-coherent-carrier}.  This proves
\textup{(ii)}.

Conversely, \textup{(ii)} includes $g_v(x)>0$, and therefore
\[
 \Delta^{\mathbf Q}_{f,X}(x)
   =P_{\mathbf Q}^{\mathrm{loc}}f(x)-s_{\mathbf Q,X}(x)
   \geq v(x)-s_{\mathbf Q,X}(x)
   =g_v(x)>0.
\]
The last assertion follows from this equivalence and the nonnegative
three-level defect ledger
\eqref{eq:p031-three-level-defect-ledger}.
\end{proof}


On a relatively compact local domain, the same carrier argument reaches the
actual boundary and records its full height.  The positive defect has an
exact upper boundary trace: its supremum in the domain is the maximum of its
interior upper limits at boundary points.  This gives a removable-boundary
criterion and an open-set characterization of rational cofinality.  Under a
finite flat pullback, openness preserves every boundary trace pointwise, while
properness and compact-Hausdorff component separation bound the pullback of
each compact carrier by the size of one finite fiber.

\medskip
% P03.1 rational-defect boundary trace, live integration.

\begin{definition}
\label{def:p031-rational-defect-boundary-trace}
Keep the hypotheses and notation of
\Cref{thm:p031-rational-collar-minimax}.  Let $W\subset X$ be ordinary
open with compact closure and nonempty boundary, and let
$v\in\mathcal Q_f(W)$.  Put $g_v:=v-s_{\mathbf Q,X}$.  For
$\xi\in\partial W$, define the \emph{positive boundary trace} by
\begin{equation}
\label{eq:p031-rational-defect-boundary-trace}
 \tau_{W,v}(\xi)
  :=\inf_{\substack{\xi\in V\\ V\ \mathrm{ordinary\ open}}}
       \ \sup_{z\in V\cap W}(g_v(z))_+
   =\limsup_{\substack{z\to\xi\\ z\in W}}(g_v(z))_+.
\end{equation}
The neighborhood formula specifies the last expression even when $X$ is
not first countable.
\end{definition}

\begin{theorem}
\label{thm:p031-rational-defect-boundary-trace}
In the setting of
\Cref{def:p031-rational-defect-boundary-trace}, the function
$\tau_{W,v}:\partial W\to\mathbf R_{\geq0}$ is finite and upper
semicontinuous, and
\begin{equation}
\label{eq:p031-rational-defect-boundary-trace-maximum}
       \sup_{z\in W}(g_v(z))_+
             =\max_{\xi\in\partial W}\tau_{W,v}(\xi).
\end{equation}
Consequently, for every $c\geq0$,
\begin{equation}
\label{eq:p031-rational-defect-asymptotic-comparison}
 \begin{split}
 v&\leq s_{\mathbf Q,X}+c\quad\text{on }W\\
 &\Longleftrightarrow\quad
 \limsup_{\substack{z\to\xi\\ z\in W}}
       \bigl(v(z)-s_{\mathbf Q,X}(z)\bigr)\leq c
       \quad\text{for every }\xi\in\partial W.
 \end{split}
\end{equation}
Thus domination by the global rational envelope is removable across an
ordinary boundary whenever it holds in the interior-limit sense.
\end{theorem}

\begin{corollary}
\label{cor:p031-rational-defect-end-height}
Keep the setting above and fix $x\in W$.  For $\alpha>0$ with
$g_v(x)\geq\alpha$, let $\Gamma_{\alpha,W}(x)$ be the component of $x$ in
\[
                       \{g_v\geq\alpha\}\subset W.
\]
Then, with the supremum of the empty set declared to be $0$,
\begin{equation}
\label{eq:p031-rational-defect-end-height}
 (g_v(x))_+
  =\sup\left\{\alpha>0:
      g_v(x)\geq\alpha,
      \overline{\Gamma_{\alpha,W}(x)}^{\,X}\cap\partial W
          \ne\varnothing\right\}.
\end{equation}
In particular,
\begin{equation}
\label{eq:p031-rational-defect-end-height-formula}
 \Delta^{\mathbf Q}_{f,X}(x)
  =
  \sup_{\substack{x\in W,\ \overline W\ \mathrm{compact}\\
                  W\ \mathrm{ordinary\ open},\ \partial W\ne\varnothing\\
                  v\in\mathcal Q_f(W)}}
  \ \sup\left\{\alpha>0:
      g_v(x)\geq\alpha,\
      \overline{\Gamma_{\alpha,W}(x)}^{\,X}\cap\partial W
          \ne\varnothing\right\}.
\end{equation}
Hence the complete rational globalization defect is exactly the supremal
height of a connected positive-defect carrier that reaches a genuine
boundary.
\end{corollary}

\begin{corollary}
\label{cor:p031-rational-defect-open-set-trace-characterization}
Let $\Omega\subset X$ be ordinary open.  The following are equivalent:
\begin{enumerate}[label=\textup{(\roman*)}]
\item
$\Delta^{\mathbf Q}_{f,X}=0$ on $\Omega$;
\item
for every ordinary open $W$ with
\[
 \overline W\subset\Omega,\qquad
 \overline W\text{ compact},\qquad
 \partial W\ne\varnothing,
\]
and every $v\in\mathcal Q_f(W)$, one has
$\tau_{W,v}=0$ on $\partial W$.
\end{enumerate}
If $X$ is boundaryless, these conditions are also equivalent to
simultaneous vanishing on $\Omega$ of the rational, LFS, and PL
globalization defects and both global branch gaps.
\end{corollary}

\begin{lemma}
\label{lem:p031-proper-open-compact-component-anchor}
Let $p:Y\to Z$ be a continuous, proper, open, surjective map of locally
compact Hausdorff spaces.  Let $K\subset Z$ be nonempty, compact, and
connected.  If $z\in K$ and
\[
                         p^{-1}(z)=\{y_1,\ldots,y_q\}
\]
is finite, then every connected component of $p^{-1}(K)$ contains at
least one $y_j$.  Consequently,
\begin{equation}
\label{eq:p031-proper-open-compact-component-anchor}
                  \#\pi_0\bigl(p^{-1}(K)\bigr)\leq q.
\end{equation}
No metrizability, local connectedness, path lifting, or covering-space
structure is required.
\end{lemma}

\begin{corollary}
\label{cor:p031-finite-flat-boundary-trace-and-carrier-anchor}
Keep the hypotheses of
\Cref{thm:p031-rational-defect-conservation}.  Let $W\subset X$ be
ordinary open with compact closure and nonempty boundary, let
$v\in\mathcal Q_f(W)$, and put
\[
 \widetilde W:=\pi^{-1}(W),\qquad
 \widetilde v:=v\circ\pi,\qquad
 \widetilde g:=\widetilde v-s_{\mathbf Q,E}.
\]
Then
\begin{equation}
\label{eq:p031-finite-flat-boundary-pullback}
 \overline{\widetilde W}=\pi^{-1}(\overline W),
 \qquad
 \partial\widetilde W=\pi^{-1}(\partial W),
 \qquad
 \widetilde g=g_v\circ\pi.
\end{equation}
Define $\tau_{\widetilde W,\widetilde v}$ by the same neighborhood
formula as in
\eqref{eq:p031-rational-defect-boundary-trace}, applied on the possibly
disconnected space $E$.  For every $\eta\in\partial\widetilde W$,
\begin{equation}
\label{eq:p031-finite-flat-boundary-trace-pullback}
       \tau_{\widetilde W,\widetilde v}(\eta)
          =\tau_{W,v}\bigl(\pi(\eta)\bigr).
\end{equation}
In particular, finite flat pullback preserves the exact maximum boundary
trace in \eqref{eq:p031-rational-defect-boundary-trace-maximum}.

More precisely, fix $x\in W$, $0<\alpha<g_v(x)$, and an ordinary open
$U$ satisfying
\[
 x\in U,\qquad \overline U\subset W,\qquad
 \overline U\text{ compact},\qquad \partial U\ne\varnothing.
\]
Let $C_{\alpha,U}(x)$ be the compact connected carrier in
\eqref{eq:p031-rational-defect-superlevel-carrier}.  Every component of
$\pi^{-1}(C_{\alpha,U}(x))$ meets $\pi^{-1}(x)$, and
\begin{equation}
\label{eq:p031-finite-flat-carrier-branch-bound}
 \#\pi_0\bigl(\pi^{-1}(C_{\alpha,U}(x))\bigr)
       \leq\#\pi^{-1}(x).
\end{equation}
Moreover, some $e\in\pi^{-1}(x)$ has a connected component in
\[
 \{\widetilde g\geq\alpha\}
       \cap\overline{\pi^{-1}(U)}
\]
that meets $\partial\pi^{-1}(U)$.  Thus a finite flat cover neither lowers
the boundary trace nor replaces one downstairs carrier by infinitely many
unanchored compact branches.
\end{corollary}

% P03.1 rational-defect boundary trace proof, live integration.

\begin{proof}[\Cref{thm:p031-rational-defect-boundary-trace}]
Choose a global $b\in\mathcal Q_f(X)$.  Since $v\leq f$ and
$s_{\mathbf Q,X}\geq b$ on $W$,
\[
                    (g_v)_+\leq(f-b)_+\quad\text{on }W.
\]
The global minorant makes $f$ real-valued.  The upper semicontinuous
function $f-b$ is bounded above on the compact set $\overline W$.
Consequently the trace in
\eqref{eq:p031-rational-defect-boundary-trace} is finite.

It is upper semicontinuous on $\partial W$.  Indeed, if
$\tau_{W,v}(\xi)<a$, the definition supplies an ordinary open
neighborhood $V$ of $\xi$ such that
\[
                     \sup_{V\cap W}(g_v)_+<a.
\]
The same neighborhood shows that
$\tau_{W,v}(\xi')<a$ for every
$\xi'\in V\cap\partial W$.  Compactness of $\partial W$ therefore makes
the right-hand side of
\eqref{eq:p031-rational-defect-boundary-trace-maximum} a maximum.  Its
value is at most the left-hand side by definition.

For the reverse inequality, fix $z\in W$.  If $(g_v(z))_+=0$, there is
nothing to prove.  Otherwise choose $0<\alpha<g_v(z)$.  By
\Cref{cor:p031-rational-defect-coherent-carrier},
\[
      \overline{\Gamma_{\alpha,W}(z)}^{\,X}
                      \cap\partial W\ne\varnothing.
\]
Choose $\xi$ in this intersection.  Every ordinary neighborhood of $\xi$
meets $\Gamma_{\alpha,W}(z)$, and $(g_v)_+\geq\alpha$ on that carrier.
Hence $\tau_{W,v}(\xi)\geq\alpha$.  Letting
$\alpha\uparrow g_v(z)$ gives
\[
             (g_v(z))_+
                  \leq\max_{\partial W}\tau_{W,v}.
\]
Taking the supremum over $z\in W$ proves
\eqref{eq:p031-rational-defect-boundary-trace-maximum}.

If $v\leq s_{\mathbf Q,X}+c$ on $W$, the boundary upper limits in
\eqref{eq:p031-rational-defect-asymptotic-comparison} are at most $c$.
Conversely, assume those upper limits are at most $c$.  Since $c\geq0$,
the positive traces satisfy $\tau_{W,v}\leq c$ on $\partial W$.
Equation \eqref{eq:p031-rational-defect-boundary-trace-maximum} then gives
$(g_v)_+\leq c$ on $W$, and hence $v\leq s_{\mathbf Q,X}+c$ there.
\end{proof}

\begin{proof}[\Cref{cor:p031-rational-defect-end-height}]
If $g_v(x)\leq0$, the set on the right-hand side of
\eqref{eq:p031-rational-defect-end-height} is empty.  Suppose that
$g_v(x)>0$.  For every $0<\alpha<g_v(x)$,
\Cref{cor:p031-rational-defect-coherent-carrier} gives
\[
      \overline{\Gamma_{\alpha,W}(x)}^{\,X}
                         \cap\partial W\ne\varnothing.
\]
No $\alpha>g_v(x)$ belongs to the displayed set.  Its supremum is
therefore $g_v(x)$, proving
\eqref{eq:p031-rational-defect-end-height}.

Every local rational competitor near $x$ may be restricted to a domain
$W$ with compact closure and nonempty boundary by
\Cref{lem:p031-small-nonempty-boundary-domain}.  The definition of the
local rational envelope thus gives
\[
 \Delta^{\mathbf Q}_{f,X}(x)
  =\sup_{\substack{x\in W,\ \overline W\ \mathrm{compact}\\
                   W\ \mathrm{ordinary\ open},\ \partial W\ne\varnothing\\
                   v\in\mathcal Q_f(W)}}(g_v(x))_+.
\]
Substitution of
\eqref{eq:p031-rational-defect-end-height} proves
\eqref{eq:p031-rational-defect-end-height-formula}.
\end{proof}

\begin{proof}[\Cref{cor:p031-rational-defect-open-set-trace-characterization}]
Assume first that $\Delta^{\mathbf Q}_{f,X}=0$ on $\Omega$.  For
$W$ and $v$ as in \textup{(ii)}, every $z\in W$ satisfies
\[
 v(z)\leq P_{\mathbf Q}^{\mathrm{loc}}f(z)
        =s_{\mathbf Q,X}(z).
\]
Thus $(g_v)_+=0$ on $W$, and
$\tau_{W,v}=0$ on $\partial W$.

Conversely, fix $x\in\Omega$ and a rational competitor $v$ on an ordinary
open neighborhood $V$ of $x$ with $V\subset\Omega$.  Choose $W$ with
\[
 x\in W,\qquad \overline W\subset V,\qquad
 \overline W\text{ compact},\qquad \partial W\ne\varnothing
\]
by \Cref{lem:p031-small-nonempty-boundary-domain}.  Apply \textup{(ii)}
to $v|_W$.  The trace maximum formula gives
$v(x)\leq s_{\mathbf Q,X}(x)$.  Taking the supremum over all local
competitors proves
$\Delta^{\mathbf Q}_{f,X}(x)=0$.  Since $x$ was arbitrary, this proves
\textup{(i)}.

When $X$ is boundaryless, apply the nonnegative three-level ledger
\eqref{eq:p031-three-level-defect-ledger} pointwise on $\Omega$.  The
rational defect vanishes exactly when every summand in that ledger
vanishes.
\end{proof}

\begin{proof}[\Cref{lem:p031-proper-open-compact-component-anchor}]
Restrict $p$ to $p^{-1}(K)\to K$.  Its source is compact by properness,
and the restricted map remains continuous, open, and surjective.

Let $C$ be a connected component of $p^{-1}(K)$.  In a compact Hausdorff
space a component is the intersection of its clopen neighborhoods; the
compact-Hausdorff proof of this fact is given at the start of the proof of
\Cref{lem:p031-compact-component-collar}.  Suppose that
$C\cap p^{-1}(z)=\varnothing$.  For every $1\leq j\leq q$, choose a
clopen subset $A_j$ of $p^{-1}(K)$ that contains $C$ and omits $y_j$.
Then
\[
                         A:=A_1\cap\cdots\cap A_q
\]
is nonempty and clopen, and $z\notin p(A)$.  Since the restricted map is
open, $p(A)$ is open in $K$; since $A$ is compact, $p(A)$ is also closed
in the Hausdorff space $K$.  This gives a nonempty proper clopen subset
of connected $K$, a contradiction.  Hence every component contains an
anchor point $y_j$.  Distinct components are disjoint, so
\eqref{eq:p031-proper-open-compact-component-anchor} follows.
\end{proof}

\begin{proof}[\Cref{cor:p031-finite-flat-boundary-trace-and-carrier-anchor}]
A finite morphism is proper, and a finite flat analytic morphism is open
by \cite[proof of Proposition~5.11, p.~175]{GJR26}.  Continuity gives
\[
       \overline{\widetilde W}\subset\pi^{-1}(\overline W).
\]
For the reverse inclusion, let $\eta\in\pi^{-1}(\overline W)$ and let
$N$ be an ordinary open neighborhood of $\eta$.  The open neighborhood
$\pi(N)$ of $\pi(\eta)$ meets $W$, so $N$ meets $\widetilde W$.
Therefore
\[
 \overline{\widetilde W}=\pi^{-1}(\overline W),
 \qquad
 \partial\widetilde W=\pi^{-1}(\partial W).
\]
The rational envelope pullback identity
\eqref{eq:p031-rational-lfs-envelope-pullback} gives
$\widetilde g=g_v\circ\pi$.  This proves
\eqref{eq:p031-finite-flat-boundary-pullback}.

Fix $\eta\in\partial\widetilde W$ and put $\xi:=\pi(\eta)$.  If $N$ is
an ordinary open neighborhood of $\eta$, openness and surjectivity onto
its image give
\[
 \sup_{\zeta\in N\cap\widetilde W}(g_v\circ\pi)(\zeta)_+
   =\sup_{z\in\pi(N)\cap W}(g_v(z))_+.
\]
The right-hand side is at least $\tau_{W,v}(\xi)$.  Taking the infimum
over $N$ proves
\[
                  \tau_{\widetilde W,\widetilde v}(\eta)
                         \geq\tau_{W,v}(\xi).
\]
Conversely, if $V$ is an ordinary open neighborhood of $\xi$, then
$\pi^{-1}(V)$ is a neighborhood of $\eta$ and
\[
 \sup_{\zeta\in\pi^{-1}(V)\cap\widetilde W}
        (g_v\circ\pi)(\zeta)_+
   =\sup_{z\in V\cap W}(g_v(z))_+.
\]
Taking the infimum over $V$ gives the reverse inequality and proves
\eqref{eq:p031-finite-flat-boundary-trace-pullback}.

Put $C:=C_{\alpha,U}(x)$.  By
\Cref{thm:p031-rational-defect-escape}, it is compact, connected, and
meets $\partial U$.  Apply
\Cref{lem:p031-proper-open-compact-component-anchor} to
$\pi^{-1}(C)\to C$.  This proves the anchoring assertion and
\eqref{eq:p031-finite-flat-carrier-branch-bound}.

Choose $z\in C\cap\partial U$ and
$\eta\in\pi^{-1}(z)$.  The component of $\eta$ in $\pi^{-1}(C)$ contains
some $e\in\pi^{-1}(x)$.  It is a connected subset of
\[
 \{\widetilde g\geq\alpha\}
       \cap\pi^{-1}(\overline U)
  =\{\widetilde g\geq\alpha\}
       \cap\overline{\pi^{-1}(U)}
\]
that contains $e$ and meets
$\pi^{-1}(\partial U)=\partial\pi^{-1}(U)$.  The component of $e$ in the
displayed compact superlevel therefore meets that boundary, as asserted.
\end{proof}


The boundary contact can be made independent of the truncation level.
Nested compactness produces one terminal boundary core reached by every
strict superlevel component through the root point.  Its rooted trace has
maximum exactly the pointwise defect.  Under finite flat pullback, every
lift of a terminal endpoint admits a common central-fiber anchor that works
simultaneously at all strict levels; no covering-space or path-lifting
structure is used.

\medskip
% P03.1 rational-defect terminal boundary core, live integration.

\begin{definition}
\label{def:p031-rational-defect-rooted-terminal-trace}
Keep the setting of
\Cref{def:p031-rational-defect-boundary-trace}, and fix $x\in W$.  For
$0<\alpha<g_v(x)$, put
\begin{equation}
\label{eq:p031-rational-defect-rooted-boundary-contact}
 B_{\alpha,W,v}(x)
   :=\overline{\Gamma_{\alpha,W}(x)}^{\,X}\cap\partial W.
\end{equation}
For $\xi\in\partial W$, define the \emph{rooted terminal trace}
\begin{equation}
\label{eq:p031-rational-defect-rooted-terminal-trace}
 \rho_{W,v,x}(\xi)
  :=\sup\Bigl(
       \{\alpha\in(0,g_v(x)):\xi\in B_{\alpha,W,v}(x)\}
       \cup\{0\}\Bigr).
\end{equation}
If $g_v(x)>0$, define the \emph{terminal boundary core}
\begin{equation}
\label{eq:p031-rational-defect-terminal-boundary-core}
 \mathfrak T_{W,v}(x)
   :=\bigcap_{0<\alpha<g_v(x)}B_{\alpha,W,v}(x).
\end{equation}
Thus the root $x$, rather than an arbitrary high-defect point near the
boundary, selects the trace.
\end{definition}

\begin{theorem}
\label{thm:p031-rational-defect-terminal-boundary-core}
In the setting of
\Cref{def:p031-rational-defect-rooted-terminal-trace}, the function
$\rho_{W,v,x}:\partial W\to\mathbf R_{\geq0}$ is finite and upper
semicontinuous, and
\begin{equation}
\label{eq:p031-rational-defect-rooted-terminal-maximum}
             (g_v(x))_+
                =\max_{\xi\in\partial W}\rho_{W,v,x}(\xi).
\end{equation}
Moreover,
\begin{equation}
\label{eq:p031-rational-defect-rooted-trace-bound}
                \rho_{W,v,x}\leq\tau_{W,v}
                \quad\text{on }\partial W.
\end{equation}
If $g_v(x)>0$, then $\mathfrak T_{W,v}(x)$ is a nonempty compact set and
\begin{equation}
\label{eq:p031-rational-defect-terminal-argmax}
 \mathfrak T_{W,v}(x)
   =\operatorname*{arg\,max}_{\xi\in\partial W}
        \rho_{W,v,x}(\xi),
 \qquad
 \tau_{W,v}(\xi)\geq g_v(x)
       \quad(\xi\in\mathfrak T_{W,v}(x)).
\end{equation}
Hence all strict superlevels through $x$ reach one common boundary point;
the endpoint need not change as the level increases to $g_v(x)$.
If, in addition, $(g_v(x))_+=\sup_W(g_v)_+$, then
\begin{equation}
\label{eq:p031-rational-defect-terminal-max-trace}
 \mathfrak T_{W,v}(x)
   \subset\operatorname*{arg\,max}_{\partial W}\tau_{W,v}.
\end{equation}
\end{theorem}

\begin{corollary}
\label{cor:p031-rational-defect-rooted-terminal-formula}
For every $x\in X$,
\begin{equation}
\label{eq:p031-rational-defect-rooted-terminal-formula}
 \Delta^{\mathbf Q}_{f,X}(x)
  =
  \sup_{\substack{x\in W,\ \overline W\ \mathrm{compact}\\
                  W\ \mathrm{ordinary\ open},\ \partial W\ne\varnothing\\
                  v\in\mathcal Q_f(W)}}
       \ \max_{\xi\in\partial W}\rho_{W,v,x}(\xi).
\end{equation}
Thus every positive witnessing competitor for the complete rational
globalization defect admits a common terminal endpoint for all its
truncation levels, not merely one boundary contact point at each level.
\end{corollary}

\begin{corollary}
\label{cor:p031-finite-flat-persistent-terminal-anchor}
Keep the hypotheses and notation of
\Cref{cor:p031-finite-flat-boundary-trace-and-carrier-anchor}.  Fix $x\in W$
with $g_v(x)>0$, choose
$\xi\in\mathfrak T_{W,v}(x)$, and let
$\eta\in\pi^{-1}(\xi)$.  For $0<\alpha<g_v(x)$, put
\[
 K_\alpha:=\overline{\Gamma_{\alpha,W}(x)}^{\,X},
\]
and let $\widetilde K_\alpha(\eta)$ be the connected component of $\eta$ in
$\pi^{-1}(K_\alpha)$.  Then the compact connected sets
$\widetilde K_\alpha(\eta)$ decrease with $\alpha$, each meets
$\pi^{-1}(x)$, and
\begin{equation}
\label{eq:p031-finite-flat-persistent-terminal-anchor}
 \varnothing\ne
 \bigcap_{0<\alpha<g_v(x)}
       \bigl(\widetilde K_\alpha(\eta)\cap\pi^{-1}(x)\bigr).
\end{equation}
Consequently one may choose a single $e\in\pi^{-1}(x)$ such that $e$ and
$\eta$ lie in the same component of $\pi^{-1}(K_\alpha)$ for every
$0<\alpha<g_v(x)$.  In addition,
\begin{equation}
\label{eq:p031-finite-flat-persistent-terminal-trace}
 \tau_{\widetilde W,\widetilde v}(\eta)
   =\tau_{W,v}(\xi)\geq g_v(x).
\end{equation}
Thus every lift of a terminal point admits at least one central-fiber anchor
that remains valid as the level approaches the rooted defect height.
\end{corollary}

% P03.1 rational-defect terminal boundary core proof, live integration.

\begin{proof}[\Cref{thm:p031-rational-defect-terminal-boundary-core}]
Write $h:=g_v(x)$.  If $h\leq0$, the set in
\eqref{eq:p031-rational-defect-rooted-terminal-trace} is empty for every
$\xi$, so $\rho_{W,v,x}=0$ and
\eqref{eq:p031-rational-defect-rooted-terminal-maximum} is immediate.
Suppose henceforth that $h>0$.

For $0<\alpha<h$, the set
$B_{\alpha,W,v}(x)$ is nonempty by
\Cref{cor:p031-rational-defect-coherent-carrier}.  It is compact because it
is closed in the compact set $\partial W$.  If
$0<\alpha<\beta<h$, then
\[
 \Gamma_{\beta,W}(x)\subset\Gamma_{\alpha,W}(x).
\]
Indeed, the set on the left is a connected subset of
$\{g_v\geq\alpha\}$ that contains $x$, so it lies in the component of $x$
there.  Therefore
\begin{equation}
\label{eq:p031-rational-defect-rooted-contact-nesting}
             B_{\beta,W,v}(x)\subset B_{\alpha,W,v}(x).
\end{equation}
The nested nonempty compact subsets $B_{\alpha,W,v}(x)$ of the compact
boundary have nonempty intersection.  Hence
$\mathfrak T_{W,v}(x)$ is nonempty and compact.

The rooted trace is bounded above by $h$.  For $0<r\leq h$, nesting gives
\begin{equation}
\label{eq:p031-rational-defect-rooted-trace-superlevel}
 \{\xi\in\partial W:\rho_{W,v,x}(\xi)\geq r\}
   =\bigcap_{0<\alpha<r}B_{\alpha,W,v}(x),
\end{equation}
while the left-hand side is empty for $r>h$.  Thus every positive closed
superlevel of $\rho_{W,v,x}$ is closed, and the function is upper
semicontinuous.  Every point of $\mathfrak T_{W,v}(x)$ has rooted trace
$h$, whereas no rooted trace exceeds $h$.  This proves
\eqref{eq:p031-rational-defect-rooted-terminal-maximum} and identifies its
maximizer set as in
\eqref{eq:p031-rational-defect-terminal-argmax}.

If $\xi\in B_{\alpha,W,v}(x)$, every ordinary neighborhood of $\xi$ meets
$\Gamma_{\alpha,W}(x)$, on which $(g_v)_+\geq\alpha$.  Hence
$\tau_{W,v}(\xi)\geq\alpha$.  Taking the supremum over all levels selected
by $\xi$ proves
\eqref{eq:p031-rational-defect-rooted-trace-bound}.  In particular, every
$\xi\in\mathfrak T_{W,v}(x)$ satisfies
$\tau_{W,v}(\xi)\geq h$, completing
\eqref{eq:p031-rational-defect-terminal-argmax}.

Finally, if $h=\sup_W(g_v)_+$, then
\Cref{thm:p031-rational-defect-boundary-trace} gives
$\tau_{W,v}\leq h$ on $\partial W$.  Equality holds at every point of the
terminal core, proving
\eqref{eq:p031-rational-defect-terminal-max-trace}.
\end{proof}

\begin{proof}[\Cref{cor:p031-rational-defect-rooted-terminal-formula}]
For each admissible pair $(W,v)$ through $x$,
\Cref{thm:p031-rational-defect-terminal-boundary-core} gives
\[
       (v(x)-s_{\mathbf Q,X}(x))_+
          =\max_{\xi\in\partial W}\rho_{W,v,x}(\xi).
\]
Take the supremum over the same relatively compact test domains and local
rational competitors as in
\eqref{eq:p031-rational-defect-end-height-formula}.
\end{proof}

\begin{proof}[\Cref{cor:p031-finite-flat-persistent-terminal-anchor}]
Each $K_\alpha$ is compact and connected: it is the closure of the connected
set $\Gamma_{\alpha,W}(x)$ inside the compact set $\overline W$.  It contains
$x$ and $\xi$.  The nesting argument in
\eqref{eq:p031-rational-defect-rooted-contact-nesting} also gives
\begin{equation}
\label{eq:p031-rational-defect-rooted-carrier-nesting}
                    K_\beta\subset K_\alpha
                    \qquad(0<\alpha<\beta<h).
\end{equation}

For each $\alpha$, apply
\Cref{lem:p031-proper-open-compact-component-anchor} to
$\pi^{-1}(K_\alpha)\to K_\alpha$, with anchor point $x$.  The component
$\widetilde K_\alpha(\eta)$ is compact and connected and meets the finite
fiber $\pi^{-1}(x)$.  By
\eqref{eq:p031-rational-defect-rooted-carrier-nesting},
$\widetilde K_\beta(\eta)$ is a connected subset of
$\pi^{-1}(K_\alpha)$ containing $\eta$, and therefore
\[
             \widetilde K_\beta(\eta)
                  \subset\widetilde K_\alpha(\eta)
             \qquad(0<\alpha<\beta<h).
\]
Consequently the sets
\[
 A_\alpha:=\widetilde K_\alpha(\eta)\cap\pi^{-1}(x)
\]
form a decreasing family of nonempty subsets of the finite set
$\pi^{-1}(x)$.  Their intersection is nonempty, which is exactly
\eqref{eq:p031-finite-flat-persistent-terminal-anchor}.  Any point $e$ in
that intersection is the asserted common anchor.

The trace identity
\eqref{eq:p031-finite-flat-boundary-trace-pullback} and
\eqref{eq:p031-rational-defect-terminal-argmax} give
\eqref{eq:p031-finite-flat-persistent-terminal-trace}.
\end{proof}


The finite-fiber hypothesis is unnecessary for the component anchor itself.
For a proper open surjection of locally compact Hausdorff spaces, every
component over a compact connected set maps onto that set.  Compactness of
one fiber then gives a common anchor for any decreasing family of components
through a fixed lifted endpoint.  Applied to the terminal carriers above,
this works even for infinite or nondiscrete fibers.  Finiteness is needed
only for a finite numerical component bound, while the finite-flat trace
identity remains an analytic input.

\medskip
% P03.1 proper-open persistent component anchors, live integration.

\begin{theorem}
\label{thm:p031-proper-open-component-surjectivity}
Let $p:Y\to Z$ be a continuous, proper, open, surjective map of locally
compact Hausdorff spaces, and let $K\subset Z$ be nonempty, compact, and
connected.  Every connected component $C$ of $p^{-1}(K)$ satisfies
\begin{equation}
\label{eq:p031-proper-open-component-surjectivity}
                         p(C)=K.
\end{equation}
Consequently every component meets $p^{-1}(z)$ for every $z\in K$.  If one
such fiber is finite, then
\begin{equation}
\label{eq:p031-proper-open-component-cardinality}
        \#\pi_0\bigl(p^{-1}(K)\bigr)\leq\#p^{-1}(z).
\end{equation}
No metrizability, discreteness of the other fibers, local connectedness, or
path-lifting hypothesis is required.
\end{theorem}

\begin{theorem}
\label{thm:p031-proper-open-persistent-component-anchor}
Keep the hypotheses on $p$ in
\Cref{thm:p031-proper-open-component-surjectivity}.  Let $I$ be a nonempty
linearly ordered set, and let $(K_i)_{i\in I}$ be nonempty compact connected
subsets of $Z$ such that
\[
                   K_j\subset K_i\qquad(i\leq j).
\]
Suppose that $x,\xi\in\bigcap_{i\in I}K_i$, fix
$\eta\in p^{-1}(\xi)$, and let $C_i(\eta)$ be the connected component of
$\eta$ in $p^{-1}(K_i)$.  Then
\begin{equation}
\label{eq:p031-proper-open-persistent-component-nesting}
 C_j(\eta)\subset C_i(\eta)\quad(i\leq j),
 \qquad
 p\bigl(C_i(\eta)\bigr)=K_i\quad(i\in I),
\end{equation}
and
\begin{equation}
\label{eq:p031-proper-open-persistent-component-anchor}
 \varnothing\ne
 p^{-1}(x)\cap\bigcap_{i\in I}C_i(\eta).
\end{equation}
Thus one point of the compact fiber $p^{-1}(x)$ anchors the entire decreasing
family, even when that fiber is infinite or nondiscrete.
\end{theorem}

\begin{corollary}
\label{cor:p031-proper-open-persistent-terminal-anchor}
Keep the setting of
\Cref{def:p031-rational-defect-rooted-terminal-trace}.  Suppose that
$g_v(x)>0$, choose $\xi\in\mathfrak T_{W,v}(x)$, and let
$p:Y\to X$ be a continuous, proper, open, surjective map of locally compact
Hausdorff spaces.  Fix $\eta\in p^{-1}(\xi)$ and, for
$0<\alpha<g_v(x)$, let $C_\alpha(\eta)$ be the connected component of
$\eta$ in
\[
       p^{-1}\!\left(
          \overline{\Gamma_{\alpha,W}(x)}^{\,X}
       \right).
\]
Then every $C_\alpha(\eta)$ maps onto
$\overline{\Gamma_{\alpha,W}(x)}^{\,X}$, the family decreases with
$\alpha$, and
\begin{equation}
\label{eq:p031-proper-open-persistent-terminal-anchor}
 \varnothing\ne
 p^{-1}(x)\cap
 \bigcap_{0<\alpha<g_v(x)}C_\alpha(\eta).
\end{equation}
In particular, the persistent component-anchor conclusion does not require
finite fibers.  The finite-flat trace identity in
\eqref{eq:p031-finite-flat-persistent-terminal-trace} remains a separate
analytic statement.
\end{corollary}

% Keep the frozen proof unchanged; allow its prose an emergency stretch.
\begingroup
\emergencystretch=1em
% P03.1 proper-open persistent component anchors proof, live integration.

\begin{proof}[\Cref{thm:p031-proper-open-component-surjectivity}]
Put $L:=p^{-1}(K)$ and let $q:L\to K$ be the restricted map.  Properness
makes $L$ compact.  The map $q$ is continuous and surjective, and it is open:
if $U=V\cap L$ with $V\subset Y$ open, then
\[
                         q(U)=p(V)\cap K.
\]

Let $C$ be a connected component of $L$ and fix $z\in K$.  Suppose that
$C\cap q^{-1}(z)=\varnothing$.  In a compact Hausdorff space a connected
component is the intersection of its clopen neighborhoods, as recalled in
the proof of \Cref{lem:p031-compact-component-collar}.  Hence, for every
$y\in q^{-1}(z)$, there is a clopen set $A_y\subset L$ which contains $C$
and omits $y$.

The fiber $q^{-1}(z)$ is compact.  The open sets $L\setminus A_y$ therefore
admit a finite subcover of that fiber, say the sets indexed by
$y_1,\ldots,y_m$.  The set
\[
                         A:=\bigcap_{r=1}^m A_{y_r}
\]
is nonempty and clopen in $L$, contains $C$, and is disjoint from
$q^{-1}(z)$.  Its image $q(A)$ is open in $K$ because $q$ is open.  It is
also compact, hence closed in the Hausdorff space $K$.  Thus $q(A)$ is a
nonempty proper clopen subset of the connected space $K$, a contradiction.
Since $z$ was arbitrary, this proves
\eqref{eq:p031-proper-open-component-surjectivity}.

If $p^{-1}(z)$ is finite, every component of $L$ contains a point of that
fiber, and distinct components contain disjoint subsets of the fiber.  This
gives \eqref{eq:p031-proper-open-component-cardinality}.
\end{proof}

\begin{proof}[\Cref{thm:p031-proper-open-persistent-component-anchor}]
Each $C_i(\eta)$ is compact and connected.  By
\Cref{thm:p031-proper-open-component-surjectivity}, it maps onto $K_i$ and
therefore meets the compact fiber $F:=p^{-1}(x)$.  If $i\leq j$, then
$C_j(\eta)$ is a connected subset of $p^{-1}(K_i)$ containing $\eta$, so
it lies in the component $C_i(\eta)$.  This proves
\eqref{eq:p031-proper-open-persistent-component-nesting}.

The sets
\[
                         A_i:=C_i(\eta)\cap F
\]
are nonempty compact subsets of $F$ and decrease with $i$.  Every finite
subfamily has nonempty intersection: choose its largest index and use the
corresponding $A_i$.  Compactness of $F$ now gives
\[
                   \varnothing\ne\bigcap_{i\in I}A_i,
\]
which is \eqref{eq:p031-proper-open-persistent-component-anchor}.
\end{proof}

\begin{proof}[\Cref{cor:p031-proper-open-persistent-terminal-anchor}]
For $0<\alpha<g_v(x)$, put
\[
             K_\alpha
               :=\overline{\Gamma_{\alpha,W}(x)}^{\,X}.
\]
By the proof of
\Cref{thm:p031-rational-defect-terminal-boundary-core}, each $K_\alpha$ is
compact and connected, contains $x$ and $\xi$, and satisfies
$K_\beta\subset K_\alpha$ when $0<\alpha<\beta<g_v(x)$.  Apply
\Cref{thm:p031-proper-open-persistent-component-anchor} to this linearly
ordered family.  Its conclusions are exactly the asserted surjectivity,
nesting, and \eqref{eq:p031-proper-open-persistent-terminal-anchor}.
\end{proof}

\par\endgroup

Global openness is not needed when the endpoint lift is selected by a
continuous section.  The section image of every connected base carrier lies
in the component through the lifted endpoint, so that component maps onto the
carrier and contains the section image of every base point.  Over an
algebraically closed base, Poineau's universal-point theorem supplies a
canonical continuous section of every scalar projection.  It gives both
pointwise openness along the section image and an explicit persistent anchor
for every canonical terminal lift, even when scalar projection is not
globally open.

\medskip
% P03.1 canonical-section persistent anchors, live integration.

\begin{lemma}
\label{lem:p031-continuous-section-pointwise-openness}
Let $q:Y\to Z$ be a continuous map with a continuous section
$\sigma:Z\to Y$.  Then $q$ is open at every point of $\sigma(Z)$: if $N$
is a neighborhood of $\sigma(z)$, then $q(N)$ is a neighborhood of $z$.
Consequently, for every locally bounded-above function
$a:Z\to\mathbf R\cup\{-\infty\}$,
\begin{equation}
\label{eq:p031-continuous-section-usc-regularization}
                  (a\circ q)^*(\sigma(z))=a^*(z).
\end{equation}
No openness of $q$ away from the section image is asserted.
\end{lemma}

\begin{theorem}
\label{thm:p031-continuous-section-persistent-component-anchor}
Let $q:Y\to Z$ be a continuous map with a continuous section
$\sigma:Z\to Y$.  Let $I$ be a nonempty linearly ordered set and let
$(K_i)_{i\in I}$ be nonempty connected subsets of $Z$ such that
\[
                    K_j\subset K_i\qquad(i\leq j).
\]
Suppose that $x,\xi\in\bigcap_{i\in I}K_i$, and let $C_i$ be the connected
component of $\sigma(\xi)$ in $q^{-1}(K_i)$.  Then
\begin{equation}
\label{eq:p031-continuous-section-component-surjectivity}
       \sigma(K_i)\subset C_i,
       \qquad q(C_i)=K_i
       \qquad(i\in I),
\end{equation}
\begin{equation}
\label{eq:p031-continuous-section-component-nesting}
                    C_j\subset C_i\qquad(i\leq j),
\end{equation}
and the whole family has the explicit common anchor
\begin{equation}
\label{eq:p031-continuous-section-persistent-anchor}
                       \sigma(x)\in\bigcap_{i\in I}C_i.
\end{equation}
Thus neither properness, global openness, compactness, nor a restriction on
the cardinality of a fiber is needed for the component selected by the
section.
\end{theorem}

\begin{corollary}
\label{cor:p031-canonical-scalar-section-anchor}
Let $K$ be an algebraically closed complete non-Archimedean field, let
$K'/K$ be an analytic extension field, and let $X$ be a $K$-analytic space.
Scalar projection
\[
                    \pi_{K'/K}:X_{K'}\longrightarrow X
\]
has a canonical continuous section
\[
 \sigma_{K'/K}:X\longrightarrow X_{K'},
 \qquad
 \sigma_{K'/K}(x)
  \in\mathcal M\!\left(\mathcal H(x)\widehat\otimes_KK'\right),
\]
whose value is the multiplicative tensor-norm point
\cite[Definition~3.2, Corollary~3.7, Remark~3.13, and
Corollary~3.14]{Poi13Angelic}.
Consequently $\pi_{K'/K}$ is open at every point of the section image and,
for every locally bounded-above
$a:X\to\mathbf R\cup\{-\infty\}$,
\begin{equation}
\label{eq:p031-canonical-scalar-section-regularization}
 (a\circ\pi_{K'/K})^*(\sigma_{K'/K}(x))=a^*(x).
\end{equation}
Moreover, every decreasing family in
\Cref{thm:p031-continuous-section-persistent-component-anchor} has the
surjectivity, nesting, and common-anchor conclusions there after taking
$q=\pi_{K'/K}$ and $\sigma=\sigma_{K'/K}$.
\end{corollary}

\begin{corollary}
\label{cor:p031-canonical-scalar-persistent-terminal-anchor}
Keep the setting of
\Cref{def:p031-rational-defect-rooted-terminal-trace}, assume that $X$ is
a $K$-analytic space for an algebraically closed complete
non-Archimedean field $K$, and let $K'/K$ be an analytic extension field.
Suppose that $g_v(x)>0$, choose
$\xi\in\mathfrak T_{W,v}(x)$, and, for
$0<\alpha<g_v(x)$, put
\[
        K_\alpha:=\overline{\Gamma_{\alpha,W}(x)}^{\,X}.
\]
Let $C_\alpha$ be the connected component of
$\sigma_{K'/K}(\xi)$ in $\pi_{K'/K}^{-1}(K_\alpha)$.  Then
\begin{equation}
\label{eq:p031-canonical-scalar-terminal-component}
 \sigma_{K'/K}(K_\alpha)\subset C_\alpha,
 \qquad
 \pi_{K'/K}(C_\alpha)=K_\alpha,
\end{equation}
the family $(C_\alpha)_{0<\alpha<g_v(x)}$ decreases with $\alpha$, and
\begin{equation}
\label{eq:p031-canonical-scalar-persistent-terminal-anchor}
          \sigma_{K'/K}(x)
             \in\bigcap_{0<\alpha<g_v(x)}C_\alpha.
\end{equation}
Hence the canonical lift of every terminal endpoint has a canonical root
anchor at every strict defect level, although scalar projection need not be
globally open.  This is a topological anchor statement; it asserts neither
fiberwise constancy of a regularized envelope nor compatibility of the
analytic defect with scalar extension.
\end{corollary}

% P03.1 canonical-section persistent anchors proof, live integration.

\begin{proof}[\Cref{lem:p031-continuous-section-pointwise-openness}]
Let $N$ be a neighborhood of $\sigma(z)$.  Continuity of $\sigma$ makes
$\sigma^{-1}(N)$ a neighborhood of $z$, and the section identity gives
\[
                    \sigma^{-1}(N)\subset q(N).
\]
Thus $q(N)$ is a neighborhood of $z$.

For a neighborhood $U$ of $z$, the set $q^{-1}(U)$ is a neighborhood of
$\sigma(z)$, whence
\[
 (a\circ q)^*(\sigma(z))
   \leq \sup_{q^{-1}(U)}a\circ q
   =\sup_U a.
\]
Taking the infimum over $U$ gives the upper bound in
\eqref{eq:p031-continuous-section-usc-regularization}.  Conversely, for
every neighborhood $N$ of $\sigma(z)$, the set $q(N)$ is a neighborhood of
$z$ and
\[
                  \sup_N(a\circ q)=\sup_{q(N)}a\geq a^*(z).
\]
Taking the infimum over $N$ proves the reverse bound and hence the equality.
\end{proof}

\begin{proof}[\Cref{thm:p031-continuous-section-persistent-component-anchor}]
For each $i$, the set $\sigma(K_i)$ is connected, contains
$\sigma(\xi)$, and lies in $q^{-1}(K_i)$.  It is therefore contained in the
component $C_i$.  Applying $q$ gives
\[
        K_i=q(\sigma(K_i))\subset q(C_i)\subset K_i,
\]
which proves
\eqref{eq:p031-continuous-section-component-surjectivity}.

If $i\leq j$, then $C_j$ is a connected subset of $q^{-1}(K_i)$ containing
$\sigma(\xi)$, so $C_j\subset C_i$.  Finally $x\in K_i$ for every $i$, and
the first conclusion gives $\sigma(x)\in\sigma(K_i)\subset C_i$.  This proves
\eqref{eq:p031-continuous-section-component-nesting} and
\eqref{eq:p031-continuous-section-persistent-anchor}.
\end{proof}

\begin{proof}[\Cref{cor:p031-canonical-scalar-section-anchor}]
Poineau's Corollary~3.14 makes every point of $X$ universal.  By
Definition~3.2, its canonical lift to $X_{K'}$ is the point defined by the
multiplicative tensor norm on
$\mathcal H(x)\widehat\otimes_KK'$.  Corollary~3.7 makes the lift continuous
on every affinoid domain, and Remark~3.13 globalizes that conclusion to
analytic spaces.  The lift lies over $x$, so the resulting map
$\sigma_{K'/K}$ is a continuous section of $\pi_{K'/K}$.

Apply
\Cref{lem:p031-continuous-section-pointwise-openness} to obtain pointwise
openness and
\eqref{eq:p031-canonical-scalar-section-regularization}.  Apply
\Cref{thm:p031-continuous-section-persistent-component-anchor} to obtain the
remaining conclusions.
\end{proof}

\begin{proof}[\Cref{cor:p031-canonical-scalar-persistent-terminal-anchor}]
By the proof of
\Cref{thm:p031-rational-defect-terminal-boundary-core}, every $K_\alpha$ is
nonempty and connected, contains $x$ and $\xi$, and satisfies
\[
             K_\beta\subset K_\alpha
             \qquad(0<\alpha<\beta<g_v(x)).
\]
Apply
\Cref{thm:p031-continuous-section-persistent-component-anchor} to scalar
projection and its canonical continuous section from
\Cref{cor:p031-canonical-scalar-section-anchor}.  The theorem gives
\eqref{eq:p031-canonical-scalar-terminal-component}, the stated nesting, and
the explicit common anchor in
\eqref{eq:p031-canonical-scalar-persistent-terminal-anchor}.
\end{proof}


The canonical scalar section also identifies the maximum in the compact-fiber
regularization formula.  Composing scalar projection with the canonical
section gives a continuous idempotent retraction.  The downstairs usc
regularization is exactly the upstairs regularization after this retraction,
and the strict scalar defect is precisely its failure to be retraction
invariant.  For the affine Perron family, this turns classical psh-ness into
an exact subharmonicity test for the retracted curve pullbacks.

\medskip
% P03.1 canonical-retraction normal form, live integration.

\begin{theorem}
\label{thm:p031-canonical-scalar-retraction-normal-form}
Let $K$ be an algebraically closed complete non-Archimedean field, let
$K'/K$ be an analytic extension field, and let $X$ be a $K$-analytic space.
Write
\[
 \pi:=\pi_{K'/K}:X_{K'}\longrightarrow X,
 \qquad
 \sigma:=\sigma_{K'/K}:X\longrightarrow X_{K'}
\]
for scalar projection and its canonical continuous section from
\Cref{cor:p031-canonical-scalar-section-anchor}, and put
\[
                         r:=\sigma\circ\pi.
\]
Then $r$ is a continuous idempotent retraction onto $\sigma(X)$.  If
$a:X\to\mathbf R\cup\{-\infty\}$ is locally bounded above and
\[
                         b:=(a\circ\pi)^*,
\]
then, for every $y\in X_{K'}$,
\begin{equation}
\label{eq:p031-canonical-retraction-fiber-maximum}
 b(r(y))
   =a^*(\pi(y))
   =\max_{z\in\pi^{-1}(\pi(y))}b(z).
\end{equation}
In particular,
\begin{equation}
\label{eq:p031-canonical-retraction-regularization}
                    a^*\circ\pi=b\circ r,
                    \qquad b\leq b\circ r.
\end{equation}
The strict regularization-defect locus is therefore
\begin{equation}
\label{eq:p031-canonical-retraction-defect-locus}
 \mathcal D_a
   :=\{y\in X_{K'}:b(y)<a^*(\pi(y))\}
    =\{y\in X_{K'}:b(y)<b(r(y))\},
\end{equation}
and $\mathcal D_a\cap\sigma(X)=\varnothing$.  The scalar-compatibility
identity $a^*\circ\pi=(a\circ\pi)^*$ holds if and only if
$\mathcal D_a=\varnothing$, if and only if $b=b\circ r$, if and only if
$b$ is constant on every fiber of $\pi$.
\end{theorem}

\begin{corollary}
\label{cor:p031-affine-canonical-retraction-reduction}
Let $K$ be algebraically closed and keep the hypotheses and notation of
\Cref{thm:p031-affine-ground-field-perron}:
\[
                  s:=\sup_{i\in I}u_i,
                  \qquad u:=s^*.
\]
For every analytic extension $K'/K$, let $\pi:=\pi_{K'/K}:V_{K'}\to V$,
let $r:=\sigma_{K'/K}\circ\pi$, and put
\[
                         v_{K'}:=(s\circ\pi)^*.
\]
Then $v_{K'}$ is $K'$-curvewise plurisubharmonic and
\begin{equation}
\label{eq:p031-affine-canonical-retraction-identity}
                         u\circ\pi=v_{K'}\circ r.
\end{equation}
Consequently, $u$ is classically plurisubharmonic if and only if, for every
analytic extension $K'/K$, every smooth connected strict $K'$-analytic
curve $C$, and every $K'$-analytic morphism $g:C\to V_{K'}$, the function
\begin{equation}
\label{eq:p031-affine-retracted-curve-test}
                         v_{K'}\circ r\circ g
\end{equation}
is subharmonic.  For a fixed extension, scalar compatibility holds exactly
when the strict locus
\[
       \mathcal D_{K'}
          :=\{y\in V_{K'}:v_{K'}(y)<v_{K'}(r(y))\}
\]
is empty, equivalently when $v_{K'}$ is fiberwise constant.

This is an exact normal form for the obstruction, not a
psh-preservation theorem: the retraction $r$ is continuous, but is not
asserted to be analytic, and $v_{K'}\circ r$ is not proved subharmonic on
the displayed curves.
\end{corollary}

% P03.1 canonical-retraction normal form proof, live integration.

\begin{proof}[\Cref{thm:p031-canonical-scalar-retraction-normal-form}]
Continuity of $r$ follows from continuity of $\pi$ and $\sigma$.  Since
$\pi\circ\sigma$ is the identity on $X$, one has
\[
 r^2=\sigma\circ\pi\circ\sigma\circ\pi=r,
 \qquad
 r\circ\sigma=\sigma.
\]
Thus $r$ is an idempotent retraction onto $\sigma(X)$.

Fix $y\in X_{K'}$ and put $x:=\pi(y)$.  The section regularization formula
\eqref{eq:p031-canonical-scalar-section-regularization} gives
\[
                     b(r(y))=b(\sigma(x))=a^*(x).
\]
On the other hand, the compact-fiber formula
\eqref{eq:p031-scalar-fiber-maximum} gives
\[
                  a^*(x)=\max_{z\in\pi^{-1}(x)}b(z).
\]
These two equalities prove
\eqref{eq:p031-canonical-retraction-fiber-maximum} and
\eqref{eq:p031-canonical-retraction-regularization}.  They also identify
the two descriptions of $\mathcal D_a$ and show that it is disjoint from
$\sigma(X)$, because $r$ fixes the section image.

The scalar-compatibility identity is equivalent to $b=b\circ r$, hence to
$\mathcal D_a=\varnothing$.  If $b=b\circ r$, then points $y,z$ in one
fiber satisfy $r(y)=r(z)$ and therefore $b(y)=b(z)$.  Conversely, if $b$
is constant on each fiber, then $y$ and $r(y)$ lie in the same fiber, so
$b(y)=b(r(y))$.  This proves all the stated equivalences.
\end{proof}

\begin{proof}[\Cref{cor:p031-affine-canonical-retraction-reduction}]
For a fixed analytic extension $K'/K$, pointwise supremum commutes with
pullback:
\[
                    s\circ\pi=\sup_{i\in I}(u_i\circ\pi).
\]
The scalar-pullback input used in the proof of
\Cref{cor:p031-scalar-compatibility-upgrade} shows that the functions
$u_i\circ\pi$ are classically plurisubharmonic; the family remains nonempty
and locally upper bounded.  Applying
\Cref{thm:p031-affine-ground-field-perron} over $K'$ shows that
$v_{K'}$ is $K'$-curvewise plurisubharmonic.

Apply
\eqref{eq:p031-canonical-retraction-regularization} to $a=s$.  Since
$u=s^*$ and $v_{K'}=(s\circ\pi)^*$, it gives
\eqref{eq:p031-affine-canonical-retraction-identity}.  By the all-extensions
smooth strict curve test used in the proof of
\Cref{cor:p031-scalar-compatibility-upgrade}, the function $u$ is classically
plurisubharmonic exactly when
$u\circ\pi\circ g$ is subharmonic for every extension, curve, and morphism
in the statement.  Substitution from
\eqref{eq:p031-affine-canonical-retraction-identity} proves the asserted
equivalence.

For a fixed extension, the final equivalences are the specialization of
\Cref{thm:p031-canonical-scalar-retraction-normal-form} to $a=s$.  The
construction of $r$ uses only a continuous section.  No analyticity of $r$
has been proved, so analytic pullback functoriality for psh functions does
not apply to $v_{K'}\circ r$.
\end{proof}


Thus finite flatness removes the local pullback obstruction without any
analytic section.  Exact full P03 compatibility is now equivalent to
preservation of the total global gain from rational LFS to semipositive
real PL.  The collar formula identifies pointwise rational LFS
local-to-global cofinality exactly, and the escape theorem shows that every
failure has one non-relatively-compact connected carrier reaching every
sufficiently small boundary, while the positive part of every defect profile
obeys the boundary maximum principle.  On a relatively compact local domain,
the defect is exactly the maximum of its upper boundary trace, and its
pointwise value is the supremal height of an end-reaching connected carrier.
Finite flat pullback preserves this trace pointwise and cannot split a
compact truncated carrier into more branches than the cardinality of one
fiber.  In fact, all strict levels through a fixed root share one terminal
boundary point, and every lift of that point has a common central-fiber
anchor persisting through the entire nested family.  More generally, the
same persistent anchor exists for every continuous proper open surjection
between locally compact Hausdorff spaces, independently of the fiber
cardinality.  In the algebraically closed scalar-extension range, the
canonical section gives the canonical endpoint lift an even stronger anchor:
every pulled-back carrier component contains the entire section of the base
carrier and hence the fixed point $\sigma_{K'/K}(x)$, without properness or
global openness.  The same section yields the canonical-retraction normal
form
\[
       s^*\circ\pi_{K'/K}=v_{K'}\circ r_{K'/K},
       \qquad v_{K'}\leq v_{K'}\circ r_{K'/K},
\]
and the strict scalar defect is exactly the failure of $v_{K'}$ to be
invariant under $r_{K'/K}$.  This does not establish subharmonicity after
composition with the continuous retraction, and hence does not prove the
fiberwise constancy or analytic defect compatibility needed for P03.1.
Rational boundary saturation is a checkable sufficient vanishing condition
and propagates through every finite flat tower.  No theorem traps all such
carriers, or proves boundary saturation, on every normal or geometrically
unibranch higher-dimensional space, and no suitable finite-flat presentation
is constructed for every scalar extension occurring in P03.1.

Thus the full P03.1 obstruction on a smooth affine chart is narrower than
the original ambient curve-selection problem.  After every scalar
extension, the graph-tube proof makes the regularization of the
\emph{pulled-back} raw supremum curvewise psh over the extended field, and
\Cref{cor:p031-fiber-constancy-upgrade} identifies the pullback of the
ground-field regularization with its fiberwise maximum.  Full scalar
compatibility is therefore equivalent to constancy on each scalar fiber.
The type--$4$ example in
\Cref{prop:p031-type-four-topological-defect} shows that non-openness can
produce a strict defect for arbitrary bounded raw data, but it supplies no
psh family and hence no P03.1 counterexample.  The residual-curve lifting
theorem used in
\Cref{thm:gr-pl-comparison} detects degrees of residual line bundles; it
does not supply the missing compatibility for arbitrary psh suprema.  No
normal higher-dimensional counterexample, CLD-v2 conclusion, or continuity
conclusion is inferred.  By
\Cref{thm:p031-completed-algebraic-closure-reduction}, the stronger affine
Perron route may now take the ground field to be complete and algebraically
closed; its remaining obstruction begins with further analytic extensions.
Independently, \Cref{thm:p031-usc-defect-localization} shows that any
failure is supported on the closure of the raw-supremum regularization
defect.  No removability property of that closed set is asserted.

% Author-supplied 24 September 2026 manuscript; provenance and local review are recorded separately.
\section{First-stage approximation on good analytic spaces}
\label{sec:supplied-p071a}

The following argument is adapted from the supplied manuscript
\cite{SuppliedP071Good20260924}.

\label{sec:supplied-p071a-introduction}

Let $K$ be a complete non-Archimedean field and let $X$ be a topologically separated $K$-analytic space. Chambert--Loir and Ducros give $X$ an affine $G$-tropical structure and define a sheaf $\operatorname{Psh}_X$ by starting with ordinary smooth real functions whose smooth $(1,1)$-form $\mathrm{d}'\mathrm{d}''u$ is positive, then repeatedly taking local decreasing sequential limits \cite[\S12.4.6, Definition~8.1.1 and Lemma~8.1.4]{CLD25}. A section obtained after one such limit is called locally approachable \cite[Definition~8.1.13]{CLD25}. The distinction between a single local approximation and membership in the sheaf generated by transfinite repetition matters when the local smooth rows live on different domains.

The first theorem identifies these two classes when every point has an affinoid analytic domain that is an ordinary neighborhood. Further applications treat finite and arbitrary branching-star products at their node tuples. No single tropical chart is fixed for the approximating sequence in either result.

\begin{theorem}
\label{thm:supplied-p071a-analytic-first-stage}
Let $X$ be a good, topologically separated Berkovich $K$-analytic space with its Chambert--Loir--Ducros affine structure. For every ordinal $\alpha\geq1$,
\[
\operatorname{Psh}_{X,\alpha}=\operatorname{Psh}_{X,1}=\operatorname{Psh}_X.
\]
Equivalently, every germ of a section of $\operatorname{Psh}_X$ is the pointwise limit of a decreasing sequence of ordinary smooth functions $u_j$ with $\mathrm{d}'\mathrm{d}''u_j\geq0$ on one fixed ordinary open neighborhood of the germ. The limit may take the value $-\infty$.
\end{theorem}

The first theorem concerns the analytic affine structure of \cite[\S12.4.6]{CLD25}. The product theorem uses the affine structure of \cite[\S9.1]{CLD25}. The finite barrier constructions below use features of those structures and do not assert first-stage collapse for every affine $G$-tropical space.

There are two compactness issues in the proof. First, local decreasing rows with a discontinuous common limit need not become uniformly close on overlaps at the same index. A common finite convex average of indices does become uniformly close; this is a weak-convergence argument in a finite product of $C(K)$ spaces. Second, smooth cutoffs are not themselves plurisubharmonic. On an affinoid spectrum, finitely many cutoffs can be controlled by one convex function of truncated logarithms of global analytic functions. The resulting adjusted cutoffs are smooth plurisubharmonic. A regularized maximum then glues the averaged rows without changing their limit. These constructions use the form-level convex composition rule of \cite[Proposition~8.2.1]{CLD25} and the vanishing of $\mathrm{d}'\mathrm{d}''\log|a|$ on an analytic unit locus \cite[Lemma~16.7.4]{CLD25}.

\subsection{Convex synchronization and diagonalization}
\label{sec:supplied-p071a-synchronization}

We first record two compactness arguments for sheaves of continuous functions. Neither uses analytic geometry. In the first lemma, all coefficients in the convex average are shared by the different rows.

\begin{lemma}
\label{lem:supplied-p071a-convex-synchronization}
Let $K_1,\ldots,K_m$ be compact Hausdorff subspaces of a Hausdorff space. For each $i$, let $(f_{i,k})_{k\geq1}$ be a decreasing sequence in $C(K_i,\mathbb{R})$. Suppose these sequences converge pointwise to the restrictions of one function $v$ on $\bigcup_i K_i$, and suppose $v$ is bounded below. For every $\varepsilon>0$ and $k_0\geq1$, there is a finite set of indices $k\geq k_0$ and nonnegative numbers $\lambda_k$ with $\sum_k\lambda_k=1$ such that, with $F_i=\sum_k\lambda_k f_{i,k}$,
\[
\lVert F_i-F_j\rVert_{K_i\cap K_j}<\varepsilon
\qquad\text{for all }i,j.
\]
Moreover $v\leq F_i\leq f_{i,k_0}$ on $K_i$.
\end{lemma}

\begin{proof}
Empty overlaps impose no condition. On the remaining overlaps form the finite product Banach space
\[
B=\prod_{i<j,\ K_i\cap K_j\ne\varnothing} C(K_i\cap K_j,\mathbb{R})
\]
with the maximum norm, and let $b_k$ have $(i,j)$ component $(f_{i,k}-f_{j,k})|_{K_i\cap K_j}$. The sequence $(b_k)$ is uniformly bounded: each $f_{i,k}$ lies between the common lower bound for $v$ and the bounded function $f_{i,1}$. Every component tends pointwise to zero. By the Riesz--Markov representation of the dual of $C(H,\mathbb{R})$ for compact Hausdorff $H$ \cite[Theorem~18.1]{SchenkerFA}, bounded convergence against each finite signed regular measure gives $b_k\rightharpoonup0$ in $B$.

Fix $k_0$. If zero were outside the norm closure of the convex hull of $\{b_k:k\geq k_0\}$, Hahn--Banach separation \cite[Theorem~2.1]{ButtenschoenHB} would give a bounded real linear functional that is uniformly positive on that hull. Its values on the weakly null tail $b_k$ tend to zero, a contradiction. A finite convex combination of tail vectors therefore has norm below $\varepsilon$. Its coefficients give the asserted $F_i$. Monotonicity of each row gives $v\leq F_i\leq f_{i,k_0}$.
\end{proof}

The next lemma replaces a decreasing array of approximants, already defined on one domain, by one decreasing sequence. It uses upper semicontinuity rather than uniform convergence to the intermediate limits.

\begin{lemma}
\label{lem:supplied-p071a-fixed-domain-diagonal}
Let $K$ be a compact subset of a topological space, and let $\mathscr{S}$ be a class of continuous real functions stable under restriction and addition of constants. For every $j$, suppose $f_{j,k}\in\mathscr{S}(W_j)$ decreases in $k$ to $u_j$ on an open neighborhood $W_j$ of the same $K$. Suppose the functions $u_j|_K$ decrease in $j$ to $u$ on $K$, with values in $\mathbb{R}\cup\{-\infty\}$. Then $u|_{\operatorname{int}K}$ is the decreasing limit of a sequence in $\mathscr{S}(\operatorname{int}K)$.
\end{lemma}

\begin{proof}
On $K$ set $h_n=\min\{f_{j,k}:1\leq j,k\leq n\}$. These functions are continuous, though they need not belong to $\mathscr{S}$, and $h_n\downarrow u$. We use the following compactness fact: if upper semicontinuous $a_r$ decreases to $a$ on $K$ and $b$ is continuous with $a<b$ pointwise, then $a_r<b$ on all of $K$ for some $r$. The open sets $\{a_r<b\}$ increase and cover $K$, so one of them covers $K$.

Set $j_1=k_1=1$ and $F_1=f_{1,1}|_K$. Given $F_{n-1}$ on $K$, put $b_n=\min(h_n,F_{n-1})$. Since $u\leq b_n$, the compactness fact first supplies $j_n$ with $u_{j_n}<b_n+2^{-n-1}$ on $K$, and then supplies $k_n$ with $F_n:=f_{j_n,k_n}|_K<b_n+2^{-n}$ there. Thus $g_n=(f_{j_n,k_n}+2^{-n})|_{\operatorname{int}K}$ belongs to $\mathscr{S}$ and, for $n\geq2$,
\[
g_n\leq F_{n-1}+2^{1-n}=g_{n-1}.
\]
Also $u\leq g_n\leq h_n+2^{1-n}$, so $g_n\downarrow u$ pointwise, including on the set where $u=-\infty$.
\end{proof}

\subsection{Finite logarithmic barriers}
\label{sec:supplied-p071a-barriers}

Let $Y=\mathcal{M}(A)$ be a $K$-affinoid spectrum. Write $|a|(y)$ for the seminorm of $a\in A$ at $y$. A logarithmic expression below is evaluated only where the relevant seminorms are positive. When a scalar function is constant for sufficiently negative arguments, its composition with $\log|a|$ has the specified constant extension across $\{a=0\}$.

\begin{lemma}
\label{lem:supplied-p071a-affinoid-barriers}
Let $K_0\subset Y$ be compact and let $O_1,\ldots,O_m$ be ordinary open subsets covering $K_0$. There are finitely many smooth functions $\theta_1,\ldots,\theta_N:Y\to[0,1]$, each assigned to one of the $O_i$, and a nonnegative smooth function $q:Y\to\mathbb{R}$ such that
\begin{enumerate}
\item the support of each $\theta_j$ lies in its assigned $O_i$, and the ordinary open regions on which some $\theta_j$ is identically $1$ cover $K_0$;
\item $q$, $q+2\theta_j$, and $q-2\theta_j$ are ordinary smooth CLD functions with nonnegative smooth $\mathrm{d}'\mathrm{d}''$ on $Y$ for every $j$.
\end{enumerate}
The functions are built from finitely many elements of $A$, although the number and choice of those elements depend on the cover.
\end{lemma}

\begin{proof}
The evaluation functions $|a|$, $a\in A$, generate the compact Hausdorff topology of $Y$. For each $y\in K_0$ choose a finite system of evaluation inequalities defining an ordinary neighborhood inside one $O_i$, and choose a smaller closed evaluation neighborhood still containing an ordinary neighborhood of $y$. Smooth cutoffs in the finitely many logarithmic variables give a product $\theta$ equal to $1$ near $y$ and supported in the chosen $O_i$. At a coordinate with value zero, use a cutoff equal to $1$ for sufficiently negative logarithmic values and equal to $0$ above a finite threshold. At a positive coordinate, use an ordinary compactly supported cutoff in the logarithmic variable. Each factor is constant at both tails, so the product extends continuously and smoothly across all zero loci. Compactness of $K_0$ yields finitely many products $\theta_i$ whose plateau regions cover it.

List all logarithmic coordinate slots used in these finitely many products. Write $\Theta_i(t_1,\ldots,t_N)$ for the resulting smooth functions on $\mathbb{R}^N$. For each slot $j$ choose a nonnegative smooth compactly supported function $\psi_j$ that is at least $1$ wherever a first or second derivative of the corresponding cutoff factor is nonzero. Put
\[
\rho_j(t)=\int_{-\infty}^{t}(t-s)\psi_j(s)\,\mathrm{d}s.
\]
Then $\rho_j$ is smooth, nonnegative, convex, nondecreasing, and zero on a negative tail. A row and column of $D^2\Theta_i$ vanish whenever the corresponding cutoff factor has zero first and second derivatives. Since the finitely many Hessians $D^2\Theta_i$ are uniformly bounded, one constant $D>0$ can be chosen such that
\[
D\operatorname{diag}(\psi_1,\ldots,\psi_N)\pm 2D^2\Theta_i\geq0
\qquad\text{for every }i\text{ and }t\in\mathbb{R}^N.
\]
Consequently $D\sum_j\rho_j$ and $D\sum_j\rho_j\pm2\Theta_i$ are smooth convex functions on $\mathbb{R}^N$. Compose them with the logarithmic evaluations to define $q$ and $q\pm2\theta_i$.

At a point where some underlying $a_j$ vanishes, its $\rho_j$ summand is identically zero on an ordinary neighborhood and every cutoff factor involving that slot is constant there. Omit those slots. All remaining $a_j$ are analytic units on a smaller ordinary neighborhood: a Laurent domain $|a_j|\geq\varepsilon<|a_j(y)|$ supplies the inverse. Each displayed function is therefore locally a smooth convex function of logarithms of analytic units. The identity $\mathrm{d}'\mathrm{d}''\log|a_j|=0$ as a smooth form on the unit locus \cite[Lemma~16.7.4]{CLD25} and the chain rule show that its smooth Hessian is positive. This includes points on the analytic boundary of $Y$ because the argument concerns smooth forms, not integration by parts.
\end{proof}

\subsection{Gluing descending rows}
\label{sec:supplied-p071a-gluing}

For a sheaf $\mathscr{S}$ of continuous real functions, let $\mathscr{A}_0=\mathscr{S}$, let $\mathscr{A}_{\alpha+1}$ be the sheaf of local decreasing limits of $\mathscr{A}_\alpha$ sections, and at a limit ordinal let $\mathscr{A}_\lambda$ be the local union of the earlier stages. Values $-\infty$ are allowed in the latter stages. We prove the gluing statement in the form needed for the CLD generator.

\begin{proposition}
\label{prop:supplied-p071a-affinoid-gluing}
Let $Y=\mathcal{M}(A)$, and let $\mathscr{S}$ be a sheaf of continuous real functions on $Y$ containing constants and closed under smooth convex coordinatewise nondecreasing composition of finitely many sections. Suppose the functions $q$ and $q\pm2\theta_i$ furnished by \cref{lem:supplied-p071a-affinoid-barriers} belong to $\mathscr{S}(Y)$ for every finite cover. If $v\in\mathscr{A}_1(U)$ and $K_0\subset U$ is compact, there is an ordinary open neighborhood $W$ of $K_0$ inside $U$ on which $v$ is a decreasing limit of functions in $\mathscr{S}(W)$. Consequently $\mathscr{A}_2=\mathscr{A}_1$ and every later stage agrees with them.
\end{proposition}

\begin{proof}
First suppose $v$ is bounded below on $U$. Choose finitely many open sets $O_i\subset U$ covering $K_0$ and decreasing rows $f_{i,k}\in\mathscr{S}(O_i)$ with limit $v|_{O_i}$. Choose cutoffs and $q$ by \cref{lem:supplied-p071a-affinoid-barriers}, then reindex the finite cutoff family and repeat its assigned local row for each cutoff. Denote the resulting data again by $O_i,f_{i,k},\theta_i$. By compact Hausdorff regularity, choose compact $K_i$ with
\[
\operatorname{supp}\theta_i\subset\operatorname{int}K_i\subset K_i\subset O_i,
\]
and set $U_i=\operatorname{int}K_i$. Let $V$ be the union of the open plateau regions of the $\theta_i$. Choose a compact neighborhood $L$ of $K_0$ inside $V$. The sets $K_i$, $V$, $L$, and $q$ will remain fixed as the rows are advanced.

Choose an even smooth probability density $\rho_a$ supported in $[-a,a]$. For $r$ real arguments define
\[
R_a^r(t)=\int_{\mathbb{R}^r}\max_{1\leq i\leq r}(t_i+s_i)\prod_{i=1}^r\rho_a(s_i)\,\mathrm{d}s_i.
\]
This function is smooth, convex, nondecreasing in each argument, and translation equivariant. Symmetry of $\rho_a$ and Jensen's inequality give $\max_i t_i\leq R_a^r(t)\leq\max_i t_i+a$. An argument more than $2a$ below another one can be deleted without changing $R_a^r$.

Fix $\eta>0$. Apply \cref{lem:supplied-p071a-convex-synchronization} to the rows on the compact $K_i$, using a chosen common tail, to obtain finite convex averages $F_i\in\mathscr{S}(O_i)$ such that $|F_i-F_j|<\eta/4$ on $K_i\cap K_j$. Put $p_i=F_i+\eta(2\theta_i-1)$ on $U_i$. At $z\in V$, apply $R_{\eta/16}$ to the arguments $p_i(z)$ for which $z\in U_i$, and call the result $P(z)$. Set $G=P+\eta q$.

We check that $G\in\mathscr{S}(V)$ although the set of available arguments varies. Near each $z\in V$ there is an index $j$ with $\theta_j=1$. If $\theta_i\leq1/4$ at a point of $U_i\cap U_j$, the overlap estimate gives $p_j-p_i>2\eta/16$, so the $i$th argument is inactive. If the boundary of $U_i$ meets $z$, compact support of $\theta_i$ inside $U_i$ makes $\theta_i=0$ near $z$; this argument is inactive wherever it is defined. All other relevant arguments are defined on one neighborhood of $z$. Translation equivariance rewrites $G$ there as the regularized maximum of
\[
F_i+\eta(q+2\theta_i-1).
\]
Each of these is in $\mathscr{S}$ by the hypotheses, and the regularized maximum is an allowed composition. The local formulas agree with the defining formula for $G$, proving membership in the sheaf. The plateau argument and the upper bound for $R_a^r$ also give
\begin{equation}
\label{eq:supplied-p071a-gluing-bounds}
v+\eta\leq G\leq\max_{i:z\in U_i}F_i+B\eta
\qquad(z\in V),\qquad B=\frac{17}{16}+\sup_Yq.
\end{equation}

Construct $G_n$ by taking $\eta_n=2^{-n}/(B+1)$ and convex averages from row indices at least $n$. For $n\geq2$, choose the tail still farther so that $f_{i,k}<G_{n-1}+\eta_n$ on $K_i\cap L$ for every selected index $k$. Such a common tail exists: the continuous function $G_{n-1}$ strictly exceeds $v$ on $L$ by \eqref{eq:supplied-p071a-gluing-bounds}, and the decreasing row $f_{i,k}$ converges pointwise to $v$; the increasing open sets $\{f_{i,k}<G_{n-1}+\eta_n\}$ cover the compact $K_i\cap L$. Synchronization can be applied to any prescribed tail. The upper bound in \eqref{eq:supplied-p071a-gluing-bounds} now yields
\[
G_n\leq G_{n-1}+2^{-n}\quad\text{on }L.
\]
For every $z\in V$, the same bound and $F_i\leq f_{i,n}$ give $G_n(z)\to v(z)$, since only finitely many indices occur. Hence $(G_n+2^{-n})|_{\operatorname{int}L}$ decreases to $v$ and consists of $\mathscr{S}$ sections on one neighborhood of $K_0$.

For arbitrary $v$, put $v^{(m)}=\max(v,-m)$. On $O_i$ the functions
\[
f_{i,k}^{(m)}=\frac{f_{i,k}-m+\sqrt{(f_{i,k}+m)^2+k^{-2}}}{2}
\]
belong to $\mathscr{S}$ by convex monotone composition and decrease to $v^{(m)}$. They are bounded below. Apply the preceding construction to these rows. Its cutoffs, compact sets, and neighborhood $\operatorname{int}L$ depend on the $O_i$ but not on $m$; therefore every $v^{(m)}$ has a decreasing $\mathscr{S}$ row on the same open $\operatorname{int}L$. Since $v^{(m)}\downarrow v$, \cref{lem:supplied-p071a-fixed-domain-diagonal} on a smaller compact neighborhood of $K_0$ inside $\operatorname{int}L$ gives one decreasing $\mathscr{S}$ sequence converging to $v$ near $K_0$.

For the stage assertion, let $u_j\in\mathscr{A}_1(U)$ decrease to $u$, and fix a compact neighborhood $K_0\subset U$ of a point. Apply the just-proved property to each $u_j$ with this same $K_0$. Each resulting row is defined on a possibly different neighborhood of $K_0$, exactly as allowed in \cref{lem:supplied-p071a-fixed-domain-diagonal}. The lemma gives a decreasing $\mathscr{S}$ sequence converging to $u$ on $\operatorname{int}K_0$. Thus $\mathscr{A}_2\subseteq\mathscr{A}_1$. The reverse inclusion is immediate, and transfinite induction gives equality at all later stages.
\end{proof}

\subsection{The analytic CLD sheaf}
\label{sec:supplied-p071a-cld-sheaf}

\begin{proof}[of \cref{thm:supplied-p071a-analytic-first-stage}]
On an affinoid analytic space $Y=\mathcal{M}(A)$, let $\mathscr{S}$ be the actual sheaf of ordinary smooth CLD functions with positive smooth $\mathrm{d}'\mathrm{d}''$. Smooth CLD functions form a sheaf of continuous real functions on the ordinary topology \cite[\S6.2.1--6.2.2]{CLD25}. The stage-zero case of the convex composition theorem \cite[Proposition~8.2.1]{CLD25} gives the closure property in \cref{prop:supplied-p071a-affinoid-gluing}. The functions in \cref{lem:supplied-p071a-affinoid-barriers} supply its other hypothesis. That proposition shows that every locally approachable function on $Y$ has a decreasing smooth psh approximation on a fixed ordinary neighborhood of any prescribed compact subset, and that the hierarchy stabilizes at stage one. This hierarchy is exactly the one in \cite[Definition~8.1.1 and Lemma~8.1.4]{CLD25}.

For $x\in X$, goodness gives an affinoid analytic domain $Y\subset X$ containing an ordinary open neighborhood $V$ of $x$. The induced analytic structures on $V$ agree, so the ordinary smooth functions, forms, and form positivity on $V$ agree. Their stage-zero sheaves agree, and the sheafified successor and limit operations commute with restriction to $V$. Transfinite induction therefore identifies the corresponding stages on $V$. The affinoid conclusion supplies one ordinary open $W\subset V$ containing $x$ and one decreasing sequence of ordinary smooth psh functions on $W$ converging to the prescribed generated germ. Equality of the stage-one and generated sheaves follows stalkwise.
\end{proof}

\begin{remark}
\label{rem:supplied-p071a-scope}
The barrier in \cref{lem:supplied-p071a-affinoid-barriers} uses finitely many global elements of one affinoid algebra to control finitely many changing local presentations. An arbitrary affine $G$-tropical space need not have an affinoid neighborhood or an analogous family of global evaluation functions. The separate criterion in \cref{sec:supplied-p071a-double-star} uses a common family of smooth affine coordinates.
\end{remark}

\subsection{A finite double-star product}
\label{sec:supplied-p071a-double-star}

The gluing argument also applies when ordinary topology is generated on one neighborhood by smooth affine functions. This criterion does not require an analytic realization.

\begin{proposition}
\label{prop:supplied-p071a-affine-coordinate-criterion}
Let $B$ be a locally compact Hausdorff ordinary open subset of an affine $G$-tropical CLD space. Suppose a family $\mathscr{H}$ of globally defined smooth affine real functions on $B$ generates its ordinary topology. Then the CLD decreasing-limit hierarchy on $B$ stabilizes at stage one. Every generated germ has a decreasing sequence of ordinary smooth psh functions on one fixed ordinary neighborhood.
\end{proposition}

\begin{proof}
Let $K\subset B$ be compact and let finitely many ordinary open row domains cover it. By local compactness, choose for each $x\in K$ a relatively compact ordinary neighborhood whose closure lies in one row domain. Finite-coordinate cylinders in $\mathscr{H}$ form a basis for the ordinary topology, so a smaller such cylinder lies inside this neighborhood. Products of smooth one-variable cutoffs in its finitely many affine coordinates give a smooth function $\theta_x$ equal to $1$ near $x$, with compact support in the assigned row domain. Finitely many of these plateau regions cover $K$.

Collect the finitely many affine coordinates used by these cutoffs. Apply the diagonal Hessian construction in the proof of \cref{lem:supplied-p071a-affinoid-barriers} to their Euclidean cutoff functions. It gives a nonnegative smooth $q$ such that $q$ and $q\pm2\theta_x$ are smooth convex functions of these affine coordinates. Since each coordinate has zero smooth $\mathrm{d}'\mathrm{d}''$, the convex chain rule makes these functions ordinary smooth psh. The union of the plateau regions lies in the finite union of the compact cutoff supports, so $q$ is bounded on the region used for gluing.

These are precisely the barrier and compactness properties used in \cref{prop:supplied-p071a-affinoid-gluing}. Its convex synchronization, localizing smooth maximum, truncation, and diagonal argument therefore apply to $B$ with the same estimates. They give a decreasing smooth psh row on one neighborhood of $K$ for every stage-one section and imply equality of all later stages with stage one.
\end{proof}

For a finite set $B_m$ of cardinality $m\geq3$, let $X_{B_m}$ denote the CLD star with branches $[-\infty,0]$ indexed by $B_m$ and with their endpoints at $0$ identified to a node $\eta_m$ \cite[\S6.7.3]{CLD25}. Give products of these stars the induced CLD product affine structure \cite[\S\S6.1.3 and~9.1]{CLD25}.

\begin{theorem}
\label{thm:supplied-p071a-double-star}
Let $m,n\geq3$. Every germ at $(\eta_m,\eta_n)$ in the generated CLD psh sheaf on $X_{B_m}\times X_{B_n}$ belongs to the first decreasing-limit stage. In particular it has a decreasing sequence of ordinary smooth psh approximants on one fixed ordinary neighborhood of the double node.
\end{theorem}

\begin{proof}
Choose $R,S>0$. In $X_{B_m}$ let $U_R$ be the union of the branch intervals $(-R,0]$, and define $V_S\subset X_{B_n}$ similarly. Because the branch sets are finite, $U_R$ and $V_S$ are ordinary open neighborhoods of their nodes. Enumerate $B_m$ and put
\[
v_1=e_1,\ldots,v_{m-1}=e_{m-1},\qquad v_m=-\sum_{j=1}^{m-1}e_j
\quad\text{in }\mathbb{R}^{m-1}.
\]
On branch $b$, parametrized by $t\in(-R,0]$, define $H_j(b,t)=t(v_b)_j$. Each $H_j$ is smooth affine on $U_R$: it is affine on each branch, and its node slopes sum to zero, as required in the star affine structure \cite[\S6.7.3]{CLD25}. Thus $\mathrm{d}'\mathrm{d}''H_j=0$ as a smooth form by the affine form calculus \cite[\S3.3.11]{CLD25}. The map $H=(H_1,\ldots,H_{m-1})$ embeds $U_R$ into $\mathbb{R}^{m-1}$, since its distinct branch rays meet only at the origin. Indeed, it is the restriction of the continuous injection from the compact truncated star with intervals $[-R,0]$, hence a topological embedding. Construct an analogous embedding $G$ on $V_S$.

The functions $H_j\circ\mathrm{pr}_1$ and $G_\ell\circ\mathrm{pr}_2$ are smooth affine on $U_R\times V_S$ by the CLD product affine rule \cite[\S9.1.4]{CLD25}. They have zero smooth Hessian and together generate the ordinary product topology because $(H,G)$ is a topological embedding. Apply \cref{prop:supplied-p071a-affine-coordinate-criterion} on this ordinary open neighborhood of $(\eta_m,\eta_n)$.
\end{proof}

The finite-coordinate construction in \cref{thm:supplied-p071a-double-star} uses a neighborhood from which the $-\infty$ endpoints have been removed. It gives no assertion for arbitrary nonambient branching systems whose ordinary neighborhoods lack a common family of topology-generating affine functions.

\subsection{Products of arbitrarily branching stars}
\label{sec:supplied-p071a-arbitrary-star-products}

\begin{proposition}
\label{prop:supplied-p071a-positive-coordinate-criterion}
Let $B$ be a locally compact Hausdorff ordinary open subset of an affine $G$-tropical space. Suppose its ordinary topology is generated by a family of globally defined ordinary smooth psh functions. Then the decreasing-limit hierarchy on $B$ stabilizes at stage one.
\end{proposition}

\begin{proof}
Fix a compact $K\subset B$ and a finite ordinary open cover on whose members a stage-one section has decreasing smooth psh rows. Choose finitely many smooth cutoffs $\theta_i$ that equal $1$ on ordinary open sets covering $K$ and have compact support in their assigned row domains. They can be taken as smooth functions $\Theta_i(h_1,\ldots,h_N)$ of finitely many topology-generating smooth psh functions: take a finite-coordinate cylinder around each point with closure in an assigned row domain, apply compactly supported Euclidean cutoff functions, and collect the finitely many coordinates used. After collecting them, multiply each cutoff by extra factors in the other coordinates, chosen equal to $1$ near its plateau point. This keeps its support in the assigned domain and makes each $\Theta_i$ compactly supported on $\mathbb{R}^N$.

The first and second derivatives of all $\Theta_i$ are supported in a common compact subset of $\mathbb{R}^N$. Choose $D>0$ such that the functions
\[
Q(t)=D\sum_{j=1}^N\mathrm{e}^{t_j},\qquad Q(t)\pm2\Theta_i(t)
\]
are convex and nondecreasing in each coordinate on $\mathbb{R}^N$. On the compact derivative support, this follows by making $D\mathrm{e}^{t_j}$ dominate the finitely many first derivatives and the operator norms of the Hessians; outside it the derivatives of $\Theta_i$ vanish. Hence $q=Q(h_1,\ldots,h_N)$ and $q\pm2\theta_i$ are smooth psh by the stage-zero convex monotone composition rule \cite[Proposition~8.2.1]{CLD25}. The function $q$ is bounded on the compact region used for gluing. The row synchronization and gluing proof of \cref{prop:supplied-p071a-affinoid-gluing} now applies on that region, including the truncation and diagonal arguments. Every stage-two germ is stage one, and transfinite induction gives the assertion.
\end{proof}

For an arbitrary set $B$ with at least two elements, let $X_B$ be the CLD star of \cite[\S6.7.3]{CLD25} and let $o$ be its node. Fix a branch $b_*$ and $R>0$, and remove the closed tail $\{b_*\}\times[-\infty,-R]$. The complement $U$ is an ordinary open node neighborhood: all other branches remain complete. For each $b\neq b_*$ and $a<0$, choose a smooth convex nondecreasing function $\rho_a$ that vanishes on $(-\infty,a]$, is strictly increasing on $(a,0]$, and is affine of slope $1$ near $0$. Put $c_a=\rho_a(0)$ and define $h_{b,a}$ on $U$ by
\[
h_{b,a}(b,t)=\rho_a(t),\qquad h_{b,a}(b_*,t)=c_a-t,
\qquad h_{b,a}=c_a\quad\text{on all other branches}.
\]
It is ordinary smooth: it is constant near every included $-\infty$ endpoint, changes on only two branches, and its node slopes $1$ and $-1$ sum to zero. Its smooth Hessian vanishes near the node and on the reservoir branch; on the transition part of branch $b$ it equals $\rho_a''(t)\,\mathrm{d}'t\wedge\mathrm{d}''t\geq0$. Thus $h_{b,a}$ is stage zero.

These functions generate the ordinary topology of $U$. On a nonreservoir branch, $h_{b,a}<c_a$ isolates that branch and varying $a$ separates its points and endpoint neighborhoods. On the reservoir branch, $h_{b,a}>c_a$ isolates it and gives a strictly monotone coordinate. At the node, a basic ordinary neighborhood excludes tails in only finitely many branches. For each such nonreservoir branch one inequality in $h_{b,a}$ removes its tail; one further inequality removes any prescribed reservoir tail. Their finite intersection lies inside the basic neighborhood. Continuity gives the reverse inclusion of the generated topology.

\begin{theorem}
\label{thm:supplied-p071a-arbitrary-star-products}
Let $X_{B_1},\ldots,X_{B_r}$ be CLD stars with arbitrary branch sets $B_i$ of cardinality at least two, where $r$ is finite. At the tuple of central nodes in their product, every germ of the generated CLD psh sheaf belongs to the first decreasing-limit stage.
\end{theorem}

\begin{proof}
In each factor choose the ordinary open neighborhood $U_i$ and smooth psh family $\{h_{b,a}\}$ above. Pullbacks of these functions by the product projections remain smooth psh by \cite[Proposition~8.1.5 and \S9.1.4]{CLD25}. Their union generates the ordinary product topology on $\prod_i U_i$. Apply \cref{prop:supplied-p071a-positive-coordinate-criterion} there and restrict to the node tuple.
\end{proof}


% Author-supplied 24 September 2026 manuscript; provenance and local review are recorded separately.
\section{A second-stage affine tropical obstruction}
\label{sec:supplied-p071b}

The following argument is adapted from the supplied manuscript
\cite{SuppliedP071Star20260924}.


The transfinite plurisubharmonic sheaf of Chambert--Loir and Ducros starts with ordinary smooth functions with positive smooth Hessian, then repeatedly takes local decreasing sequential limits \cite[Definition~8.1.1 and Lemma~8.1.4]{CLD25}. Its first stage consists of functions locally approachable by one decreasing sequence of smooth plurisubharmonic functions on a fixed ordinary neighborhood \cite[Definition~8.1.13]{CLD25}. We show that first-stage collapse fails for an affine structure obtained by restricting the permitted slopes on the compact infinite star of \cite[\S6.7.3]{CLD25}.

\begin{theorem}\label{thm:supplied-p071b-second-stage-star}
There is a compact connected affine $G$-tropical space $X$ and a bounded continuous function $u\in\operatorname{Psh}_{X,2}(X)$ whose germ at a point $o$ is not in $\operatorname{Psh}_{X,1}$. More strongly, no sequence of ordinary smooth CLD plurisubharmonic functions on any fixed ordinary neighborhood of $o$ converges pointwise to $u$ there.
\end{theorem}

This settles the unrestricted first-stage assertion for affine $G$-tropical spaces negatively. The construction does not supply an analytic space with this affinage. Hence it does not decide the version requiring a Berkovich analytic counterexample, for which compatibility with analytic logarithmic coordinates and scalar extension would also be necessary.

\subsection{A restricted affinage on the compact star}\label{sec:supplied-p071b-affinage}

We use the rational parlinear convention of \cite[\S\S1.8.2 and~1.8.8]{CLD25}. Let $D$ be the countable branch set
\[
 D=\{r_j,b_j:j\geq1\}\cup\{a_{j,n}:j,n\geq1\}.
\]
Let $X$ be the compact star of \cite[\S6.7.3]{CLD25} with branch set $D$: each branch has coordinate $t\in[-\infty,0]$, all points $t=0$ form the node $o$, and every ordinary neighborhood of $o$ contains all but finitely many complete branches. The domain category and the standard parlinear sheaf are those of that example. Write $e_d$ for the coordinate vector of branch $d$, and put
\begin{equation}\label{eq:supplied-p071b-slope-space}
 L_{\mathbb Q}=\operatorname{span}_{\mathbb Q}\{\ell_{j,n}:j,n\geq1\},\qquad
 \ell_{j,n}=e_{r_j}-n e_{b_j}-e_{a_{j,n}}.
\end{equation}
Every vector in this span has finite support. Set $L_{\mathbb R}=L_{\mathbb Q}\otimes_{\mathbb Q}\mathbb R$.

On a domain $W$ of the standard star, let $A(W)$ consist of continuous functions that are rational affine on each connected branch portion, constant near every included $-\infty$ endpoint, and whose slopes on terminal branch intervals present at $o\in W$ extend to some vector in $L_{\mathbb Q}$. There is no node-slope condition when $o\notin W$; separated branch components away from the node have their ordinary rational affine structure. Let $P$ be the subsheaf of the standard parlinear sheaf whose sections are $G$-locally sections of $A$.

\begin{lemma}\label{lem:supplied-p071b-site-inheritance}
The inherited domain category with parlinear sheaf $P$ is a rational $G$-tropical space, and $A$ is an affinage. On an ordinary open neighborhood of $o$, every $A$-coordinate has a unique full terminal slope vector in $L_{\mathbb Q}$.
\end{lemma}

\begin{proof}
The slope-extension condition is stable under restriction: discarding a terminal branch condition preserves the same extension vector. It is also ordinary local on each domain. A relative ordinary neighborhood of $o$ in $W$ retains a terminal subinterval of every branch present at $o$; on each connected terminal portion, a rational affine function has one slope. Thus one local extension vector already gives the required slopes throughout those portions. Other components impose no node condition. Constants belong to $A$, and $A$ is a rational vector subpresheaf of the standard parlinear sheaf.

The $G$-local definition makes $P$ a $G$-sheaf. Every finite $P$-map is a map for the standard star, so the inherited domain, inverse-image, and compact-image axioms continue to hold. To check composition with a parlinear target map $\phi$, take its $G$-cover by rational affine cells. Their preimages under the given map are inherited domains and form a $G$-cover: for each source point, finitely many target cells through its image cover an ordinary target neighborhood, whose inverse image is an ordinary source neighborhood inside their finite preimage union. On each preimage, $\phi$ is a rational affine combination of the components and hence lies in $P$; the $G$-sheaf property glues these expressions. These are the axioms and site-morphism condition of \cite[\S\S3.1.2--3.1.4]{CLD25}. The ordinary locality of $A$ and its $G$-generation of $P$ are precisely the affinage conditions \cite[\S\S6.1.1--6.1.2]{CLD25}.

Finally, an ordinary open node neighborhood has a terminal interval in every branch. Its actual slope vector is therefore uniquely determined, and the defining extension condition puts it in $L_{\mathbb Q}$.
\end{proof}

Ordinary smoothness at $o$ means a smooth real function of finitely many $A$-coordinates on one ordinary neighborhood \cite[Definition~6.2.1]{CLD25}. Differentiating such a presentation at $o$ shows that every ordinary smooth function has terminal slope vector in $L_{\mathbb R}$. On a finite branch interior the affine structure is the usual one-dimensional rational structure, so positivity of the smooth Hessian is ordinary convexity. At an included $-\infty$ endpoint every $A$-coordinate is locally constant. Consequently each ordinary smooth psh function is convex along branches and constant near their included endpoints. We will use this necessary condition, not its converse at arbitrary node germs.

\subsection{The second-stage function}\label{sec:supplied-p071b-construction}

Choose a smooth convex nondecreasing $\rho:\mathbb R\to[0,\infty)$ that is zero on $(-\infty,-1]$ and has derivative $1$ on $[-1/2,\infty)$. Put $c=\rho(0)>0$. For each $j$ define $h_j$ on $X$ by $h_j=\rho(t)-c$ on $b_j$ and $h_j=0$ elsewhere. Thus $-c\leq h_j\leq0$ and $h_j(o)=0$.

\begin{lemma}\label{lem:supplied-p071b-first-stage-blocks}
Each $h_j$ belongs to $\operatorname{Psh}_{X,1}(X)$.
\end{lemma}

\begin{proof}
Remove the closed tail $\{r_j\}\times[-\infty,-1]$ and call the resulting ordinary open node neighborhood $B_j$. For $n\geq1$ define $f_{j,n}$ on $B_j$, with value $2c/n$ at $o$, by
\[
 f_{j,n}=\begin{cases}
 \rho(t)-c+2c/n,&\text{on }b_j,\\
 -t/n+2c/n,&\text{on }r_j,\\
 (\rho(t)-c)/n+2c/n,&\text{on }a_{j,n},\\
 2c/n,&\text{on all other branches.}
 \end{cases}
\]
Near $o$ it is affine with slope vector $-\ell_{j,n}/n\in L_{\mathbb Q}$, so it is itself an $A$-coordinate there and has zero smooth Hessian. On branch interiors the formulas are smooth and convex; at included $-\infty$ endpoints they are locally constant. The only transition Hessians are $\rho''(t)$ and $\rho''(t)/n$, both nonnegative. Positivity is $G$-local \cite[Definition~3.5.1]{CLD25}, so $f_{j,n}$ is ordinary smooth CLD psh on $B_j$.

The four branch formulas give, everywhere on the fixed $B_j$,
\begin{equation}\label{eq:supplied-p071b-block-error}
 \frac cn\leq f_{j,n}-h_j\leq\frac{2c+1}{n}.
\end{equation}
Choose increasing integers $n_k$ with $n_{k+1}\geq((2c+1)/c)n_k$. Then \eqref{eq:supplied-p071b-block-error} implies $f_{j,n_{k+1}}\leq f_{j,n_k}$, and the decreasing limit is $h_j$ on $B_j$. At every point away from $o$, the function $h_j$ is locally ordinary smooth psh. Ordinary locality in the first stage therefore gives $h_j\in\operatorname{Psh}_{X,1}(X)$.
\end{proof}

Put $u_m=\sum_{j=1}^{m}2^{-j}h_j$. Finite sums and positive scalings preserve the first stage: at each point intersect finitely many approximation neighborhoods and add their decreasing smooth psh rows termwise. Since $h_j\leq0$, the sequence $u_m$ decreases. It converges uniformly to
\begin{equation}\label{eq:supplied-p071b-witness}
 u=\sum_{j\geq1}2^{-j}h_j,
\end{equation}
because $\|2^{-j}h_j\|_\infty\leq c2^{-j}$. Hence $u$ is bounded and continuous and belongs to $\operatorname{Psh}_{X,2}(X)$ by the successor-stage definition \cite[Definition~8.1.1]{CLD25}.

\begin{proof}[of \Cref{thm:supplied-p071b-second-stage-star}]
It remains to exclude first-stage approximation at $o$. Let $V$ be any ordinary neighborhood of $o$. The cofinite node topology gives an index $j$ for which all branches $r_j,b_j,a_{j,n}$, $n\geq1$, lie wholly in $V$. Let $f$ be any ordinary smooth CLD psh function on $V$. Its terminal slope vector $s$ lies in $L_{\mathbb R}$ by the smooth-presentation rule in \cref{sec:supplied-p071b-affinage}. Along every full branch of block $j$, convexity and constancy near $-\infty$ imply that its terminal slope is nonnegative.

The block-$j$ component of $s$ is a finite sum $\sum_n\alpha_n\ell_{j,n}$. On each $a_{j,n}$ its slope is $-\alpha_n\geq0$, while on $r_j$ the slope is $\sum_n\alpha_n\geq0$. Since only finitely many $\alpha_n$ are nonzero, all vanish. In particular the terminal slope on $b_j$ is zero. A convex function whose derivative is zero near $-\infty$ and at $0$ is constant throughout that branch; hence $f|_{b_j}=f(o)$. The same equality holds for every pointwise limit of ordinary smooth psh functions on $V$.

But \eqref{eq:supplied-p071b-witness} gives $u(o)=0$ and $u(b_j,-\infty)=-2^{-j}c$. Thus $u$ is no pointwise limit of such functions on $V$. As $V$ was arbitrary, its germ is not in the first stage.
\end{proof}


% Author-supplied 24 September 2026 manuscript; provenance and local review are recorded separately.
\section{A codimension-two envelope obstruction}
\label{sec:supplied-p031}

The following argument is adapted from the supplied manuscript
\cite{SuppliedP03120260924}.


Let $K$ be a complete nontrivially valued non-Archimedean field. For a finite upper semicontinuous function $H$ on a good strict $K$-analytic space $X$, consider the upper semicontinuous regularization of the supremum of all continuous locally Fubini--Study functions bounded above by $H$. The classical envelope question asks for geometric hypotheses under which this regularization is classically plurisubharmonic, including its analytic-curve tests after extension of the ground field. One proposed weakening of normality or geometric unibranchness everywhere is to impose either condition only in codimension one. The following construction rules out that weakening, even for a geometrically integral surface with a seminormal singularity.

Write $B=K\langle u,v\rangle$, and set
\[
 A=\{f\in B:f(0,0)=f(1,0)\},\qquad X=\mathcal M(A).
\]
The normalization of $X$ is the closed bidisc $\mathcal M(B)$; its map $\nu$ identifies the two rational points $p=(0,0)$ and $q=(1,0)$ to one point $m\in X$.

\begin{theorem}\label{thm:supplied-p031-codimension-two-envelope}
For every complete nontrivially valued non-Archimedean field $K$, the space $X$ is a good strict geometrically integral and geometrically seminormal analytic surface, smooth away from the codimension-two point $m$. There is a finite upper semicontinuous obstacle $H$ on $X$ such that the upper semicontinuous regularization of its full continuous locally Fubini--Study minorant envelope equals $H$, but $H$ is not classically plurisubharmonic. The same envelope identity and the same failure of an analytic-curve test hold after every complete extension $L/K$.
\end{theorem}

The result concerns the sufficiency of codimension-one hypotheses. It leaves open the envelope question under normality or geometric unibranchness at every point. The obstruction is the compatibility of the two branches over $m$: continuity forces every minorant to take the smaller branch value at $m$, whereas upper regularization restores the larger value and creates a positive spike on the other branch.

\subsection{The pinched bidisc}\label{sec:supplied-p031-pinched-bidisc}

Put $s=u(u-1)$, $t=us$, and $w=uv$, and regard $C=K\langle s,v\rangle$ as a subalgebra of $B$. The following explicit decomposition controls the topology as well as the algebra.

\begin{lemma}\label{lem:supplied-p031-banach-decomposition}
The inclusion $C\subset B$ is isometric, and $B=C\oplus uC$ with the maximum norm. Moreover,
\[
 A=C\oplus u(s,v)C=K\langle s,v,t,w\rangle,
\]
where the last equality is an equality of strict affinoid algebras with the induced norm.
\end{lemma}

\begin{proof}
The relation $u^2=u+s$ reduces every polynomial in $u,v$ to a unique expression $F(s,v)+uG(s,v)$. The reduction does not increase the Gauss norm, because the recursion has integral coefficients. Conversely, let $F$ and $G$ be Tate series and suppose their maximum Gauss norm is nonzero. Scale by a coefficient of that norm and reduce the proposed image modulo the maximal ideal of the valuation ring. The homomorphism
\[
 \widetilde K[S,V]\oplus U\widetilde K[S,V]\longrightarrow\widetilde K[U,V],\qquad S\longmapsto U^2-U,
\]
is injective: the leading $U$-degrees of nonzero terms from its two summands have opposite parity. The image therefore has the original maximum norm. Applying this argument to the polynomial truncations and using completeness gives the asserted isometric decomposition for Tate series, including the isometry of $C$.

The equality $f(p)=f(q)$ for $f=F(s,v)+uG(s,v)$ is equivalent to $G(0,0)=0$. Every such $G$ has a decomposition $G=sG_1+vG_2$ with $\max(\|G_1\|,\|G_2\|)\leq\|G\|$, obtained by splitting off its monomials divisible by $s$ and the remaining monomials divisible by $v$. Thus $f=F+tG_1+wG_2$. This expression supplies a bounded quotient map from $K\langle S,V,T,W\rangle$ to $A$ whose quotient norm is the induced norm.
\end{proof}

\begin{proposition}\label{prop:supplied-p031-normalization-base-change}
The inclusion $A\subset B$ is a finite birational normalization. The associated analytic map $\nu:\mathcal M(B)\to X$ is an isomorphism away from $m=\{s=v=t=w=0\}$ and has precisely the two points $p,q$ above $m$. The point $m$ is $K$-rational of codimension two. For every complete extension $L/K$,
\[
 A\widehat\otimes_K L=\{f\in L\langle u,v\rangle:f(0,0)=f(1,0)\},
\]
and all these assertions hold over $L$. In particular, $X$ is geometrically integral and geometrically seminormal, and is smooth outside $m$.
\end{proposition}

\begin{proof}
By \cref{lem:supplied-p031-banach-decomposition}, $B=A[u]$ and $u^2-u=s$. Also $u=t/s$ in the fraction field of $A$, so $A\subset B$ is finite and birational. The Tate algebra $B$ is factorial and hence normal \cite[Theorem~5.4]{JonssonHorawa20}; consequently it is the integral closure of $A$.

On the locus $s\neq0$, the identity $u=t/s$ defines an analytic inverse to $\nu$. On $v\neq0$ the inverse is $u=w/v$, and the two expressions agree on their overlap. The relations
\[
 t^2-st-s^3=0,\qquad w^2-vw-sv^2=0,\qquad tv=sw
\]
show that a point with $s=v=0$ also has $t=w=0$. Thus the complement of these two inverse charts is the single rational point $m$. Since $B/(s,v)=K[u]/(u(u-1))\simeq K\times K$, the fiber of $m$ consists of $p$ and $q$. The affinoid spectra are compact Hausdorff, so the continuous surjection $\nu$ is a quotient map.

The finite extension $C\subset A$ and the dimension of $C$ give $\dim A=2$. The ideal $(v,w)=vB$ in $A$ is prime: its quotient is $\{f\in K\langle u\rangle:f(0)=f(1)\}$, a domain. It lies strictly between $(0)$ and the maximal ideal $m$, so $m$ has height two. The isomorphism away from $m$ identifies that locus with an open subset of the smooth bidisc.

The decomposition $A=C\oplus u(s,v)C$ is a split Banach-space decomposition. Completed extension of scalars therefore gives the same equalizer description over $L$, and the preceding normalization and fiber arguments apply verbatim to $L$. In particular $A\widehat\otimes_K L$ is a domain. For seminormality, let $b$ in its fraction field satisfy $b^2,b^3\in A\widehat\otimes_K L$. It is integral and therefore belongs to the normalizing algebra $L\langle u,v\rangle$. If $\alpha=b(p)$ and $\beta=b(q)$, then $\alpha^2=\beta^2$ and $\alpha^3=\beta^3$, forcing $\alpha=\beta$ in every characteristic. Hence $b$ lies in the equalizer. This is the square-and-cube criterion \cite[Definition~29.48.1, \href{https://stacks.math.columbia.edu/tag/0EUL}{Tag~0EUL}]{StacksProject}, and it proves seminormality after every complete extension. The two distinct points $p,q$ above $m$ also show that this singularity is not geometrically unibranch.
\end{proof}

\subsection{The full minorant envelope}\label{sec:supplied-p031-envelope}

Choose $\pi\in K$ with $0<|\pi|<1$ and put $c=-\log|\pi|>0$. On $\mathcal M(B)$ consider the continuous function
\[
 h=\max\{0,c+\log|u-1|\},
\]
with the value $0$ at $u=1$. It takes values in $[0,c]$, equals $c$ where $|u|<1$, and equals $0$ where $|u-1|\leq|\pi|$. In particular $h(p)=c$ and $h(q)=0$. Define $H(m)=c$ and $H(x)=h(\nu^{-1}(x))$ for $x\neq m$. The function $H$ is upper semicontinuous: it is continuous off $m$ and its value at $m$ is its global upper bound.

Let $\mathcal F_H$ be the full family of continuous locally Fubini--Study functions $\varphi$ on $X$ with $\varphi\leq H$, and let $g=\sup_{\varphi\in\mathcal F_H}\varphi$. The constant zero function belongs to $\mathcal F_H$.

\begin{proposition}\label{prop:supplied-p031-envelope-identity}
For this obstacle, $g(m)=0$, $g=H$ on $X\setminus\{m\}$, and $g^*=H$ on $X$. The same identities hold for the full family of continuous semipositive real piecewise-linear minorants whenever that class contains global Fubini--Study functions.
\end{proposition}

\begin{proof}
For $n\geq1$ define $f_n=(u-1)^n s$ and $k_n=(u-1)^n v$. Both lie in $A$ because they vanish at $p$ and $q$. The function
\begin{equation}\label{eq:supplied-p031-fs-minorants}
 \psi_n=\frac1n\log\max\bigl\{1,|\pi|^{-n}|f_n|,|\pi|^{-n}|k_n|\bigr\}
\end{equation}
is globally Fubini--Study: its three sections of the trivial bundle are $1,\pi^{-n}f_n,\pi^{-n}k_n$, and the first section removes any base locus. Set $R=\max\{|s|,|v|\}\leq1$ on $\mathcal M(B)$. Then
\[
 \psi_n\circ\nu=\max\left\{0,c+\log|u-1|+\frac1n\log R\right\},
\]
interpreted as zero at $R=0$ or $u=1$. Thus $0\leq\psi_n\leq H$, including at $m$ where $\psi_n(m)=0$. Away from $p,q$ one has $R>0$, so $\psi_n(x)\to H(x)$ for each $x\neq m$. It follows that $g=H$ away from $m$.

Every continuous minorant $\varphi\leq H$ satisfies $\varphi(m)\leq0$: the points $\nu(1,\pi^j)$ tend to $m$, avoid it, and have obstacle value zero. Since the zero function is a minorant, $g(m)=0$. Conversely, $\nu(0,\pi^j)\to m$ and $g(\nu(0,\pi^j))=H(\nu(0,\pi^j))=c$. Therefore $g^*(m)\geq c$. The inequality $g\leq H$ and upper semicontinuity of $H$ give $g^*\leq H$, proving $g^*=H$. This argument only uses continuity and the presence of the functions $0,\psi_n$ in the chosen full minorant class, so it applies to the stated semipositive real piecewise-linear class as well.
\end{proof}

\begin{proof}[of \Cref{thm:supplied-p031-codimension-two-envelope}]
For the geometry apply \cref{prop:supplied-p031-normalization-base-change}; for the envelope apply \cref{prop:supplied-p031-envelope-identity}. To test positivity, let $\gamma:\mathcal M(K\langle z\rangle)\to X$ be the analytic map induced by $u=1$, $v=z$ on the normalization. Its pullback is
\[
 (g^*\circ\gamma)(z)=
 \begin{cases}
 c,&z=0,\\
 0,&z\neq0.
 \end{cases}
\]
On the closed subdisc $|z|\leq|\pi|$, the origin is a relative-interior point, and the function is zero at its Gauss boundary point. Comparison with the harmonic function zero would bound its value at the origin by zero, contrary to $c>0$. Thus $g^*$ fails the analytic-curve test for classical plurisubharmonicity at an interior point. The description of $A$ in \cref{prop:supplied-p031-normalization-base-change} shows that the same formulas define the scalar extensions of $H$, $\psi_n$, and $\gamma$. The two sequences and the disc comparison then prove the claims over every complete extension $L/K$.
\end{proof}


The supplied results above distinguish the analytic and general affine
tropical categories: first-stage approximation holds on good separated
analytic spaces, while the restricted-slope star is a second-stage
counterexample. The pinched bidisc shows that codimension-one conditions
alone do not ensure the envelope property at a singular point.

\section{Dictionary and checklist}
\label{sec:ch15-dictionary}

\subsection{The correspondence table}

The following correspondences are used without further comment from
\Cref{chap:cld-bedford-taylor} onwards.
\[
\begin{array}{ccl}
 \text{regular meromorphic }f
 &\longleftrightarrow&
 \text{section of }\cM_X^\times,\ \divis(f)\text{ Cartier},\\[1mm]
 \text{invertibility locus of }f
 &\longleftrightarrow&
 \text{locus where }\log|f|\text{ is smooth},\\[1mm]
 \dprimeddprime[\log|f|]=\delta_{\divis(f)}
 &\longleftrightarrow&
 \text{slope formula on a skeleton},\\[1mm]
 \text{multiplicity of a component}
 &\longleftrightarrow&
 \text{calibration of }\Trop(\{f=0\}),\\[1mm]
 \text{regular sequence }f_1,\dots,f_d
 &\longleftrightarrow&
 \text{codimension }d,\ (S_d),\ \text{extension by zero},\\[1mm]
 \log\max_i|f_i|
 &\longleftrightarrow&
 \text{potential of }\delta_Y\text{ (Bochner--Martinelli)},\\[1mm]
 \text{smooth }u,\ \dprimeddprime u\ge0
 &\longleftrightarrow&
 \operatorname{Psh}_{S,0},\\[1mm]
 \text{decreasing limit of smooth psh}
 &\longleftrightarrow&
 \operatorname{Psh}_{S,1}=\text{locally approachable},\\[1mm]
 \text{generated sheaf}
 &\longleftrightarrow&
 \operatorname{Psh}_S=\operatorname{Psh}_{S,\alpha},\ \alpha\gg0,\\[1mm]
 \text{convex function on }\mathbf R^n
 &\longleftrightarrow&
 \operatorname{Psh}_{\mathbf R^n}=\operatorname{Psh}_{\mathbf R^n,1},\\[1mm]
 \{u=-\infty\}
 &\longleftrightarrow&
 \text{polar locus},\ =\{f=0\}\text{ for }u=\log|f|,\\[1mm]
 \text{balancing at a vertex}
 &\longleftrightarrow&
 \text{no }u(p)\text{-term in the vertex atom},\\[1mm]
 \text{curve test over all }K'/K
 &\longleftrightarrow&
 \text{classically psh}.
\end{array}
\]
The third line is \Cref{thm:cld-poincare-lelong} together with
\Cref{thm:ch08-slope-formula}; the fourth is
\Cref{thm:trop-multiplicity-formula}; the sixth is
\Cref{thm:cld-local-bochner-martinelli}; the eighth is
\Cref{def:ch15-approachable}; the ninth is
\Cref{prop:ch15-hierarchy-stabilizes}; the tenth is
\Cref{thm:ch15-convexity-model}; the eleventh is
\Cref{prop:ch15-log-approachable}; the twelfth is
\Cref{lem:ch15-vertex-atom}.

\subsection{Four separations}

Each of the following is a place where two objects that share a name or a
notation are not the same, and each has been the source of an error in the
literature or in a first draft of this book.

\emph{First, three different canonical extensions.}  The current
$[\log|f|]$ extends by local integrability; the currents
$u(\dprimeddprime u)^k$ with $k<d$ extend by zero; and the Levine current
$\Lambda$ extends as a pushforward from a blowup.  Only the first two are
characterized by a uniqueness statement, and the second only in bidegree
$<d$.  None of the three follows from a general
uniqueness-by-codimension principle.

\emph{Second, smooth psh and positive current.}  On a boundaryless space
these agree for smooth functions (\Cref{prop:cld-psh-positive-current}).
Without balancing they do not (\Cref{ex:ch15-unbalanced-max}).  And for
nonsmooth functions the implication runs one way only, except on
$\mathbf R^n$ (\Cref{thm:ch15-convexity-model}) and for continuous functions
on smooth proper curves
(\Cref{thm:p063-current-psh-ordinary-regularization}).

\emph{Third, the two hierarchy questions.}  Eventual stabilization is a
theorem (\Cref{prop:ch15-hierarchy-stabilizes}); collapse at stage one is open
in general (\Cref{prob:frontier-cld-hierarchy}).  It does hold on real affine
spaces, tropical toric partial compactifications, finite ambient-induced
real-affine stars, the two constant-tail interval models, and the exact CLD
stars $X_B$ for arbitrary nonempty branch sets, in the scopes of
\Cref{thm:ch15-convexity-model,prop:ch16-collapse-tropical-toric,%
thm:ch16-collapse-finite-tropical-curves,%
thm:p071-cld-constant-collar-endpoint,thm:p071-cld-infinite-collar-endpoint,%
thm:p071-infinite-star-collapse-candidate}.  The class named by CLD is the
stage-one class, and the adjective is \emph{locally approachable}; a psh
function that is not known to be approachable is not thereby known not to be.
Thuillier's \Cref{thm:ch15-thuillier-approximation} is a stage-one theorem
on curves for \emph{his} subharmonic class and \emph{his} smooth functions;
it is not by itself a collapse theorem for the CLD hierarchy, for the three
reasons listed in \Cref{sec:ch15-thuillier-curves}.  The additional tree and
residual-divisor argument needed for continuous smooth proper curves is
\Cref{thm:p063-current-psh-ordinary-regularization}.

\emph{Fourth, the two psh classes.}  CLD-psh implies classically psh
(\Cref{thm:cld-to-classical-psh}); the converse is not asserted.  A maximum
principle is available for the second class, hence, under its hypotheses, for
the first (\Cref{cor:ch15-cld-curve-max}); it is not available for the CLD
class in general.

\subsection{The sign checklist}

Before using any formula of this chapter in a mixed computation, run the
following four checks.

\begin{enumerate}
\item \emph{Log or valuation.}  Is the potential $\log|f|$ or
  $v(f)=-\log|f|$?  The two differ by a sign in
  \Cref{thm:cld-poincare-lelong}; see \eqref{eq:appJ-cld-pl} and
  \eqref{eq:appJ-valuation-pl}.
\item \emph{Norm or squared norm.}  A complex comparison uses
  $-\log|s|_h^2$; the non-Archimedean weight is $-\log\|s\|$.  The factor $2$
  is \eqref{eq:appJ-complex-weight}, and the safest comparison is
  \Cref{cor:ch15-pl-section} against
  \eqref{eq:appJ-complex-section-pl}.
\item \emph{Which end of the skeleton.}  On a segment, $\theta=v(T)$
  increases towards the zero of $T$ and $\log|T|$ decreases towards it; an
  outgoing slope at an endpoint has one tangent direction, not two.  See
  \Cref{ex:ch15-disc-coordinate}.
\item \emph{Balanced or not.}  A vertex atom is the weighted sum of outgoing
  slopes only when the defect vanishes
  (\Cref{lem:ch15-vertex-atom}).  If a computation produces a negative atom
  for a function that ought to be psh, the space has a boundary there.
\end{enumerate}

\section*{Source notes}

\paragraph{Attributions.}
Poincar\'e--Lelong and the higher-codimension formulae are from
\cite[\S16.7 and \S18.7]{CLD25}.  The approximation definition of psh is
from \cite[Chapter~8]{CLD25}.  Classically psh and pluriharmonic functions,
their maximum principles, and the comparison with semipositive PL functions
are from \cite{GR25Classical,GR25PL}.  The potential theory of Berkovich
curves that both of those sources presuppose --- subharmonic functions,
their identification with the degree-zero currents of positive Laplacian,
the approximation theorem, the maximum and Harnack principles, and the
capacity theory --- is Thuillier's thesis \cite{Thu05}; the statements used
here are \Cref{thm:ch15-thuillier-currents},
\Cref{thm:ch15-thuillier-approximation} and the three facts recorded after
them in \Cref{sec:ch15-thuillier-curves}, all quoted and none reproved.
These two psh languages meet in the
directions stated above; no unproved converse is used.

In detail: \Cref{lem:ch15-divisor-basic} and
\Cref{prop:ch15-log-current} assemble \cite[\S\S16.7.1--16.7.2 and
18.1.5]{CLD25} with the extension lemma \cite[Lemma~7.1.13]{CLD25} and the
analytic Hauptidealsatz \cite[Lemma~11.5.2]{CLD25}.
\Cref{thm:cld-poincare-lelong} and \Cref{lem:pl-formal-reductions} are
\cite[Theorem~16.7.3 and Lemmas~16.7.4--16.7.5]{CLD25}, and the six-step
presentation of \Cref{sec:ch15-pl-proof-detail} follows that proof, with the
cutoff, the Green formula \cite[Proposition~5.11.1]{CLD25} and the tropical
product \Cref{thm:trop-local-germ} isolated as separate steps.
\Cref{cor:ch15-pl-section} is \cite[Theorem~18.1.6]{CLD25} and
\Cref{cor:ch15-degree-zero} is the trivially metrized case of
\cite[Corollary~18.1.7]{CLD25}.
\Cref{thm:cld-local-bochner-martinelli} is Mihatsch's theorem, quoted as
\cite[Theorem~18.7.2]{CLD25}, with \Cref{prop:ch15-bm-fibre} from
\cite[Corollary~18.6.2]{CLD25} and
\Cref{prop:ch15-bm-fibre-stabilizes} from
\cite[Proposition~18.7.3]{CLD25}; \Cref{thm:cld-levine} is
\cite[Theorem~18.7.5]{CLD25}.
\Cref{prop:ch15-psh-cone}, \Cref{prop:ch15-hierarchy-stabilizes},
\Cref{prop:ch15-uniform-limit}, \Cref{lem:ch15-shift},
\Cref{prop:ch15-L1-hierarchy}, \Cref{def:ch15-approachable},
\Cref{prop:ch15-composition}, \Cref{prop:ch15-max-stability} and
\Cref{prop:ch15-pullback-stagewise} are, respectively,
\cite[Lemma~8.1.2, Lemma~8.1.4, Proposition~8.1.9, Lemma~8.1.10,
Proposition~8.1.12, Definition~8.1.13, Proposition~8.2.1, Corollary~8.2.3 and
Proposition~8.1.5]{CLD25}.  \Cref{thm:ch15-convexity-model} is
\cite[Proposition~8.3.2]{CLD25}, and its proof uses the radial mollification
of \Cref{prop:ch09-radial-mollification} and the distributional
characterization \Cref{thm:ch09-distributional-convexity}, both proved in
this book from \cite[\S10]{Roc70} and \cite{EvansGariepy}.
\Cref{prop:ch15-exhaustion} is the corrected compact-positive-sublevel form
of \cite[Proposition~16.9.1]{CLD25}; its regularized maximum is
\cite[Lemma~5.18]{DemBook}.  The distinction between the negative value at
the chosen point and the positive compact sublevel is retained in
\Cref{thm:ch16-barrier}.
\Cref{thm:ch15-negligibility-recalled} recalls
\cite[Theorem~16.9.2 and Corollaries~16.9.3--16.9.4]{CLD25}, proved in
\Cref{chap:cld-bedford-taylor}.

\paragraph{What is new in the presentation.}
Four things in this chapter are arranged differently from the source, and one
correction is made.

The proof of \Cref{thm:cld-poincare-lelong} is presented in six steps so that
the role of each hypothesis is localized; in particular the two places where
boundarylessness is used, and the single place where purity of the relative
dimension is used, are marked.  The derivation of the constant in
\Cref{lem:ch15-pl-constant} is done by a support argument --- the current is a
measure carried by a point, so one test function determines it --- rather than
by an explicit boundary computation, which shortens it considerably.
\Cref{sec:ch15-quartic} carries a four-root example to explicit numbers and
checks it twice, once on the tree and once after retraction; the second
computation reproduces \Cref{ex:ch08-newton-cubic} in a case with an extra
Newton segment.  \Cref{prop:ch15-log-approachable}, that $\log|f|$ is locally
approachable, is proved directly with the truncated regularized maximum
rather than quoted; the truncation property
\Cref{lem:ch15-regmax}\,\cref{itm:ch15-regmax-trunc} is what makes the
approximants smooth across the zero locus, and no other regularized maximum
has it.  \Cref{thm:ch15-max-onedim} and
\Cref{ex:ch15-unbalanced-max} isolate the one-dimensional maximum principle
and its failure without balancing; the computation of the vertex atom in
\Cref{lem:ch15-vertex-atom} makes explicit which term balancing removes.

\paragraph{A citation correction.}
\index{CLD hierarchy!source correction}
The sentence in \cite[Remark~8.1.6]{CLD25} stating that stationarity of the
transfinite psh hierarchy is not known is superseded by the same authors'
\cite[Lemma~8.1.4]{CLD25}, which proves eventual stationarity by a
cardinality argument and which they use repeatedly.  We have therefore stated
eventual stabilization as a theorem
(\Cref{prop:ch15-hierarchy-stabilizes}) and identified the live question as
collapse at stage one, that is, whether every psh function is
\emph{locally approachable} in the sense of
\cite[Definition~8.1.13]{CLD25}.  That question is
\Cref{prob:frontier-cld-hierarchy}.  Readers comparing this chapter with the
source should be aware of the discrepancy; the second sentence of
\cite[Remark~8.1.6]{CLD25}, that all known examples lie at the first stage,
records the intended content and is consistent with
\Cref{thm:ch15-approachable-stability}.

\paragraph{Historical note.}
The complex Poincar\'e--Lelong formula, in the form
$\ddc\log|f|^2=[\divis f]$, goes back to Lelong's work of the 1940s and
1950s on plurisubharmonic functions and the currents they define
\cite{Lel45,Lel68}; the systematic theory of currents of integration over
analytic sets, including the local structure of their Lelong numbers, is due
to Lelong and to King \cite{Kin71}, with the analyticity of Lelong sublevel
sets proved by Siu \cite{Siu74}.  The Chern--Levine--Nirenberg inequalities
\cite{CLN69} and the Bedford--Taylor theory \cite{BT76,BT82} supplied the
tools that make the higher-codimension formulae possible; a modern account of
all of this is \cite{DemBook}.  On the non-Archimedean side, the
smoothing theorem of Fornaess and Narasimhan \cite{FN80} is the complex
statement that CLD's approximation definition of psh is designed to imitate.
Zhang's semipositive metrics \cite{Zhang1995} are the closest antecedent of
the local psh notion used here, and the relation between the two is a
recurring theme of \Cref{chap:cld-metrized-bundles,chap:compare-ma}.  The
non-Archimedean Bochner--Martinelli formula is due to Mihatsch, in the
context of $\delta$-forms on Lubin--Tate spaces, and reached the present
setting through the Gubler--K\"unnemann calculus \cite{GK17,GK19}; the
$\delta$-form version of Poincar\'e--Lelong, in which the right-hand side can
itself be multiplied, is \Cref{thm:gk-poincare-lelong}.  The curve-test
class and its maximum principle are the recent contribution of Gubler and
Rabinoff \cite{GR25Classical,GR25PL}, building on Thuillier's potential
theory on Berkovich curves \cite{Thu05} and on the comparison of the two
notions of subharmonicity established by Wanner \cite{Wanner19}; Jell's Poincar\'e lemma
\cite{Jell16} and the harmonic-tropicalization framework \cite{GJR26} belong
to the same circle of ideas and are taken up in \Cref{chap:delta-forms}.
